%%%%%%%%%%%%%%%%%%%%%%%%%%%%%%%%%%%%%%%%%%%%%%%%%%%%%%%%%%%%%%%%%
\chapter{Unit 2: Data}
\label{chap:unit2}
%%%%%%%%%%%%%%%%%%%%%%%%%%%%%%%%%%%%%%%%%%%%%%%%%%%%%%%%%%%%%%%%%

%%===============================================================
\section{Data}
\label{sec:data}
%%===============================================================

\paragraph{Data}



Some of the different types of data are shown below.



% [Image: image]



\paragraph{Overview}



Data are individual facts, statistics, or items of information, often numeric, collected through observation. Technically, data are a set of values of qualitative or quantitative variables about one or more persons or objects. A datum (singular of data) is a single value of a single variable.



Although "data" and "information" are often used interchangeably, they have distinct meanings. Data are considered transformed into information when viewed in context or post-analysis. In academic contexts, data are simply units of information. Data are used in scientific research, business management (e.g., sales data, revenue, profits, stock price), finance, governance (e.g., crime rates, unemployment rates, literacy rates), and various other forms of human organizational activity (e.g., censuses of homeless people by non-profit organizations).



Data are measured, collected, reported, and analyzed, and used to create visualizations such as graphs, tables, or images. Raw data ("unprocessed data") is a collection of numbers or characters before it has been "cleaned" and corrected by researchers. Raw data needs correction to remove outliers or obvious errors (e.g., a thermometer reading from an Arctic location recording a tropical temperature). Data processing occurs in stages, and the "processed data" from one stage may be considered "raw data" for the next stage. Field data is raw data collected in an uncontrolled "in situ" environment. Experimental data is data generated within a scientific investigation by observation and recording.



Data has been described as the new oil of the digital economy.





\paragraph{Etymology and Terminology}



The first English use of the word "data" is from the 1640s. The term "data" was first used to mean "transmissible and storable computer information" in 1946. The expression "data processing" was first used in 1954.



The Latin word \textit{data} is the plural of \textit{datum}, "(thing) given," neuter past participle of \textit{dare} "to give". In English, "data" may be used as a plural noun with some writers using \textit{datum} for singular and "data" for plural, especially in the natural sciences, life sciences, and social sciences. In everyday language and much of software development and computer science, "data" is most commonly used in the singular as a mass noun (like "sand" or "rain"). The term "big data" takes the singular.



\paragraph{Meaning}



Data, information, knowledge, and wisdom are closely related concepts, but each has its role. Data are collected and analyzed; it becomes information suitable for decisions once analyzed. The extent to which data is informative depends on its unexpectedness to the user. The amount of information in a data stream is characterized by its Shannon entropy.



Knowledge is understanding based on extensive experience dealing with information. For example, the height of Mount Everest is data, which can be measured and entered into a database. This data, included in a book with other data on Mount Everest, describes the mountain for those planning a climb. Understanding based on climbing experience may be seen as "knowledge," and practical climbing based on this knowledge as "wisdom." Wisdom refers to the practical application of knowledge.



Data are often assumed to be the least abstract, with information being less abstract, and knowledge being the most abstract. Data becomes information by interpretation; e.g., the height of Mount Everest is "data," a book on its geological characteristics is "information," and a climber''s guidebook is "knowledge." Information ranges from everyday usage to technical use. The concept of information is related to constraint, communication, control, data, form, instruction, knowledge, meaning, mental stimulus, pattern, perception, and representation. Data are a series of symbols, while information occurs when symbols refer to something.



Before computing devices, people manually collected and imposed patterns on data. Now, computers collect, sort, and process data, creating patterns interpreted as "truth" and authorized as aesthetic or ethical criteria. Mechanical computing devices represent data as physical quantities or sequences of symbols from a fixed alphabet. Computer programs are collections of data that can be instructions. Metadata, such as library catalogs, describe other data.



\paragraph{Data Documents}



Data exist in the form of data documents when registered. Types include:



\begin{itemize}
    \item Data repository
    \item Data study
    \item Data set
    \item Software
    \item Data paper
    \item Database
    \item Data handbook
    \item Data journal
\end{itemize}

Some of these documents (repositories, studies, sets, and software) are indexed in Data Citation Indexes, while data papers are indexed in traditional bibliographic databases (e.g., Science Citation Index).



\paragraph{Data Collection}



Data can be gathered from primary sources (first-hand) or secondary sources (previously collected). Data analysis includes methodologies such as triangulation and percolation, offering methods to collect, classify, and analyze data using various angles to maximize objectivity and understanding. Data are "percolated" through predetermined steps to extract relevant information.



\paragraph{In Other Fields}



In other fields, data are increasingly used, but their interpretive nature might conflict with the concept of data as "given." Peter Checkland''s term \textit{capta} (from Latin \textit{capere}, “to take”) distinguishes between possible data and a sub-set of data focused on. Johanna Drucker argues that using data in the humanities might introduce assumptions that are counterproductive, such as discrete phenomena or observer-independence. The term \textit{capta} emphasizes the act of observation as constitutive, offering an alternative to data for visual representations in the humanities.



\paragraph{Licensing}



Content obtained and/or adapted from:



\begin{itemize}
    \item \href{https://en.wikipedia.org/wiki/Data}{Data, Wikipedia} under a CC BY-SA license
\end{itemize}

\subsection{Learning Objectives}

\begin{enumerate}
    \item 1.\textbf{Identify and differentiate between data and information:}
    \item Understand the distinct meanings of data and information.
    \item Recognize how data transforms into information through context and analysis.
    \item Give examples of data and information in various fields such as science, business, and governance. 2.\textbf{Explain the types and processing of data:}
    \item Describe different types of data including raw data, field data, and experimental data.
    \item Understand the stages of data processing and the transformation from raw data to processed data.
    \item Identify methods of data collection and the importance of correcting raw data. 3.\textbf{Analyze the role and evolution of data:}
    \item Discuss the historical context and etymology of the term "data".
    \item Explain the relationship between data, information, knowledge, and wisdom.
    \item Evaluate the impact of data in the digital economy and its use in creating data visualizations.
\end{enumerate}

\subsection{Assessment Instrument}

\begin{enumerate}
    \item 1.\textbf{Multiple Choice:} Which of the following best describes the difference between data and information?
    \item a) Data is always numeric; information is always textual.
    \item b) Data is raw and unprocessed; information is processed and contextualized.
    \item c) Data and information are interchangeable terms.
    \item d) Information is collected through observation; data is derived from information. 2.\textbf{Short Answer:} Provide an example of data and explain how it can be transformed into information. 3.\textbf{True or False:} In academic contexts, data and information are considered the same. Explain your answer. 4.\textbf{Multiple Choice:} Which of the following is considered raw data?
    \item a) A cleaned and corrected dataset.
    \item b) A collection of unprocessed numbers or characters.
    \item c) A graph created from data analysis.
    \item d) An infographic summarizing research findings. 5.\textbf{Short Answer:} Describe the difference between field data and experimental data. 6.\textbf{Essay:} Explain the stages of data processing and why correcting raw data is important. Provide an example of an error that might need to be corrected. 7.\textbf{Multiple Choice:} When was the term "data" first used to mean "transmissible and storable computer information"?
    \item a) 1640
    \item b) 1946
    \item c) 1954
    \item d) 2000 8.\textbf{Short Answer:} Discuss the relationship between data, information, knowledge, and wisdom. Provide an example to illustrate each concept. 9.\textbf{Essay:} Evaluate the statement "Data has been described as the new oil of the digital economy." Discuss the significance of this analogy and the role of data visualizations in modern business.
\end{enumerate}

%%===============================================================
\section{Learning Community}
\label{sec:learning-community}
%%===============================================================

where the matrices

`{{math|''R''1, ..., ''R''''k''}}`{=mediawiki} are

2 × 2 rotation matrices, and with the remaining entries zero.

Exceptionally, a rotation block may be diagonal,

`{{math|±''I''}}`{=mediawiki}. Thus, negating one column if necessary,

and noting that a 2 × 2 reflection diagonalizes to a +1 and -1, any

orthogonal matrix can be brought to the form

$P^\mathrm{T}QP = \begin{bmatrix}

\begin{matrix}R\_1 \& \& \\\\ \& \ddots \& \\\\ \& \& R\_k\end{matrix} \& 0 \\\0 \& \begin{matrix}\pm 1 \& \& \\\\ \& \ddots \& \\\\ \& \& \pm 1\end{matrix} \\\\end{bmatrix},$

%%===============================================================
\section{First-degree equation involving percentages}
\label{sec:first-degree-equation-involving-percentages}
%%===============================================================

\paragraph{First-degree equation involving percentages}



In mathematics, a percentage (from Latin \textit{per centum} "by a hundred") is a number or ratio expressed as a fraction of 100. It is often denoted using the percent sign, "\%", although the abbreviations "pct.", "pct" and sometimes "pc" are also used. A percentage is a dimensionless number (pure number); it has no unit of measurement.



\paragraph{Calculations}



The percent value is computed by multiplying the numeric value of the ratio by 100. For example, to find 50 apples as a percentage of 1250 apples, one first computes the ratio $\frac{50}{1250} = 0.04$ and then multiplies by 100 to obtain 4\%. The percent value can also be found by multiplying first instead of later, so in this example, the 50 would be multiplied by 100 to give 5,000, and this result would be divided by 1250 to give 4\%.



To calculate a percentage of a percentage, convert both percentages to fractions of 100, or to decimals, and multiply them. For example, 50\% of 40\% is: $\frac{50}{100} \times \frac{40}{100} = 0.50 \times 0.40 = 0.20 = \frac{20}{100} = 20\%$



It is not correct to divide by 100 and use the percent sign at the same time; it would literally imply division by 10,000. For example, 25\% = $\frac{25}{100} = 0.25$ not $\frac{25\%}{100}$ which actually is $\frac{25/100}{100} = 0.0025$



A term such as $\frac{100}{100}\%$ would also be incorrect, since it would be read as 1 percent, even if the intent was to say 100\%.



Whenever communicating about a percentage, it is important to specify what it is relative to (i.e., what is the total that corresponds to 100\%). The following problem illustrates this point.



In a certain college 60\% of all students are female, and 10\% of all students are computer science majors. If 5\% of female students are computer science majors, what percentage of computer science majors are female?

We are asked to compute the ratio of female computer science majors to all computer science majors. We know that 60\% of all students are female, and among these 5\% are computer science majors, so we conclude that $\frac{60}{100} \times \frac{5}{100} = \frac{3}{100}$ or 3\% of all students are female computer science majors. Dividing this by the 10\% of all students that are computer science majors, we arrive at the answer: $\frac{3\%}{10\%} = \frac{30}{100}$ or 30\% of all computer science majors are female.



This example is closely related to the concept of conditional probability.



\paragraph{Percentage increase and decrease}



Due to inconsistent usage, it is not always clear from the context what a percentage is relative to. When speaking of a "10\% rise" or a "10\% fall" in a quantity, the usual interpretation is that this is relative to the initial value of that quantity. For example, if an item is initially priced at `$`200 and the price rises 10% (an increase of `$`20), the new price will be `$`220. Note that this final price is 110\% of the initial price (100\% + 10\% = 110\%).



Some other examples of percent changes:



\begin{itemize}
    \item An increase of 100\% in a quantity means that the final amount is 200\% of the initial amount (100\% of initial + 100\% of increase = 200\% of initial). In other words, the quantity has doubled.
    \item An increase of 800\% means the final amount is 9 times the original (100\% + 800\% = 900\% = 9 times as large).
    \item A decrease of 60\% means the final amount is 40\% of the original (100\% – 60\% = 40\%).
    \item A decrease of 100\% means the final amount is zero (100\% – 100\% = 0\%).
\end{itemize}

In general, a change of $x$ percent in a quantity results in a final amount that is $100 + x$ percent of the original amount (equivalently, $(1 + 0.01x)$ times the original amount).



\paragraph{Resources}

\begin{itemize}
    \item \href{https://www.khanacademy.org/math/cc-sixth-grade-math/cc-6th-ratios-prop-topic/cc-6th-percent-problems/v/solving-percent-problems}{Solving Percent Problems, Khan Academy}
    \item \href{https://www.youtube.com/watch?v=_SpE4hQ8D_o}{Taking Percentages, Khan Academy}
    \item \href{https://www.youtube.com/watch?v=Uo8HgcyfRFI}{Part, Whole, \& Percent Proportion Word Problems, The Organic Chemistry Tutor}
\end{itemize}

\paragraph{Licensing}

Content obtained and/or adapted from:



\begin{itemize}
    \item \href{https://en.wikipedia.org/wiki/Percentage}{Percentage, Wikipedia} under a CC BY-SA license
\end{itemize}

\subsection{Learning Objectives}

\begin{enumerate}
    \item 1.\textbf{Calculate Percentages}: Students will be able to calculate a percentage by converting a ratio to a fraction of 100 and multiplying the result by 100. 2.\textbf{Determine Percentage of a Percentage}: Students will understand how to find a percentage of a percentage by converting both percentages to fractions or decimals and multiplying them. 3.\textbf{Apply Percentage Change}: Students will be able to calculate and interpret percentage increases and decreases, understanding the impact on the original quantity.
\end{enumerate}

\subsection{Assessment Instrument}

\begin{enumerate}
    \item 1.\textbf{Question 1}: What is the percentage of 30 out of 200? 2.\textbf{Question 2}: If a class has 45 students and 12 of them are left-handed, what percentage of the class is left-handed? 3.\textbf{Question 3}: What is 25\% of 80\%? 4.\textbf{Question 4}: Find 40\% of 50\%. 5.\textbf{Question 5}: The price of a laptop increased from `$`800 to `$`960. What is the percentage increase? 6.\textbf{Question 6}: A population of 5,000 decreased by 15\%. What is the new population?
\end{enumerate}

%%===============================================================
\section{Ranking methods}
\label{sec:ranking-methods}
%%===============================================================

Ranking methods

%%===============================================================
\section{Graphical Display}
\label{sec:graphical-display}
%%===============================================================

\paragraph{Graphical Display}

A \textbf{chart} is a graphical representation for data visualization, in which "the data is represented by symbols, such as bars in a bar chart, lines in a line chart, or slices in a pie chart." A chart can represent tabular numeric data, functions, or some kinds of qualitative structure and provides different information.

The term "chart" as a graphical representation of data has multiple meanings:

\begin{itemize}
    \item A data chart is a type of diagram or graph that organizes and represents a set of numerical or qualitative data.
    \item Maps adorned with extra information (map surround) for a specific purpose are often known as charts, such as a nautical chart or aeronautical chart, typically spread over several map sheets.
    \item Other domain-specific constructs are sometimes called charts, such as the chord chart in music notation or a record chart for album popularity.
\end{itemize}

Charts are often used to ease the understanding of large quantities of data and the relationships between parts of the data. Charts can usually be read more quickly than the raw data. They are used in a wide variety of fields and can be created by hand (often on graph paper) or by computer using a charting application. Certain types of charts are more useful for presenting a given data set than others. For example, data that presents percentages in different groups (such as "satisfied, not satisfied, unsure") are often displayed in a pie chart but may be more easily understood when presented in a horizontal bar chart. On the other hand, data that represents numbers that change over a period of time (such as "annual revenue from 1990 to 2000") might be best shown as a line chart.

\paragraph{Features}

A chart can take a large variety of forms. However, there are common features that provide the chart with its ability to extract meaning from data.

Typically, the data in a chart is represented graphically since humans can infer meaning from pictures more quickly than from text. Thus, text is generally used only to annotate the data.

One of the most important uses of text in a graph is the title. A graph''''''''s title usually appears above the main graphic and provides a succinct description of what the data in the graph refers to.

Dimensions in the data are often displayed on axes. If a horizontal and a vertical axis are used, they are usually referred to as the x-axis and y-axis. Each axis will have a scale, denoted by periodic graduations and usually accompanied by numerical or categorical indications. Each axis will typically also have a label displayed outside or beside it, briefly describing the dimension represented. If the scale is numerical, the label will often be suffixed with the unit of that scale in parentheses. For example, "Distance traveled (m)" is a typical x-axis label and would mean that the distance traveled, in units of meters, is related to the horizontal position of the data within the chart.

Within the graph, a grid of lines may appear to aid in the visual alignment of data. The grid can be enhanced by visually emphasizing the lines at regular or significant graduations. The emphasized lines are then called major gridlines, and the remainder is minor gridlines.

A chart''''''''s data can appear in all manner of formats and may include individual textual labels describing the datum associated with the indicated position in the chart. The data may appear as dots or shapes, connected or unconnected, and in any combination of colors and patterns. In addition, inferences or points of interest can be overlaid directly on the graph to further aid information extraction.

When the data appearing in a chart contains multiple variables, the chart may include a legend (also known as a key). A legend contains a list of the variables appearing in the chart and an example of their appearance. This information allows the data from each variable to be identified in the chart.

\paragraph{Types}

\paragraph{Common Charts}

Four of the most common charts are:

$\phantom{--}$\textbf{Histogram}$\phantom{-----}$\textbf{Bar Chart}$\phantom{----}$\textbf{Pie Chart}$\phantom{-----}$\textbf{Line Chart}

% [Image: histogram]$\phantom{--}$% [Image: bar]$\phantom{--}$% [Image: Pie]$\phantom{--}$% [Image: Line]

_______________________________________________________

Other common charts are:

$\phantom{-}$\textbf{Timeline Chart}$\phantom{--}$\textbf{Organizational Chart}$\phantom{---}$\textbf{Tree Chart}$\phantom{---}$\textbf{Flow Chart}

% [Image: timeline]$\phantom{--}$% [Image: org chart]$\phantom{--}$% [Image: tree]$\phantom{--}$% [Image: Flow]

_______________________________________________________

$\phantom{--}$\textbf{Area Chart}$\phantom{-----}$\textbf{Cartogram} $\phantom{---}$ \textbf{Pedigree Chart} $\phantom{--}$ \textbf{Sunburst Chart}

% [Image: Area]$\phantom{--}$% [Image: cart]  $\phantom{--}$% [Image: pedigree]$\phantom{--}$% [Image: sunburst]

\paragraph{Mean}

The mean, or more precisely the arithmetic mean, is simply the arithmetic average of a group of numbers (or data set) and is shown using the $\bar{x}$ symbol. So the mean of the variable $x$ is $\bar{x}$, pronounced "x-bar". It is calculated by adding up all of the values in a data set and dividing by the number of values in that data set:
$$
\bar{x} = \frac{\sum x}{n}
$$

For example, take the following set of data: $\{1, 2, 3, 4, 5\}$. The mean of this data would be:
$$
\bar{x} = \frac{\sum x}{n} = \frac{1 + 2 + 3 + 4 + 5}{5} = \frac{15}{5} = 3
$$

Here is a more complicated data set: $\{10, 14, 86, 2, 68, 99, 1\}$. The mean would be calculated like this:
$$
\bar{x} = \frac{\sum x}{n} = \frac{10 + 14 + 86 + 2 + 68 + 99 + 1}{7} = \frac{280}{7} = 40
$$

\paragraph{Dot Plots}

A dot plot of 50 random values from 0 to 9:

% [Image: dotplot]

The dot plot as a representation of a distribution consists of a group of data points plotted on a simple scale. Dot plots are used for continuous, quantitative, univariate data. Data points may be labeled if there are few of them.

Dot plots are one of the simplest statistical plots, and are suitable for small to moderate-sized data sets. They are useful for highlighting clusters and gaps, as well as outliers. Their other advantage is the conservation of numerical information. When dealing with larger data sets (around 20–30 or more data points) the related stem plot, box plot, or histogram may be more efficient, as dot plots may become too cluttered after this point. Dot plots may be distinguished from histograms in that dots are not spaced uniformly along the horizontal axis.

Although the plot appears to be simple, its computation and the statistical theory underlying it are not simple. The algorithm for computing a dot plot is closely related to kernel density estimation. The size chosen for the dots affects the appearance of the plot. Choice of dot size is equivalent to choosing the bandwidth for a kernel density estimate.

In the R programming language, this type of plot is also referred to as a stripchart.

\paragraph{Histogram}

A histogram is an approximate representation of the distribution of numerical data. It was first introduced by Karl Pearson. To construct a histogram, the first step is to "bin" (or "bucket") the range of values—that is, divide the entire range of values into a series of intervals—and then count how many values fall into each interval. The bins are usually specified as consecutive, non-overlapping intervals of a variable. The bins (intervals) must be adjacent and are often (but not required to be) of equal size.

If the bins are of equal size, a rectangle is erected over the bin with height proportional to the frequency—the number of cases in each bin. A histogram may also be normalized to display "relative" frequencies. It then shows the proportion of cases that fall into each of several categories, with the sum of the heights equaling 1.

However, bins need not be of equal width; in that case, the erected rectangle is defined to have its area proportional to the frequency of cases in the bin. The vertical axis is then not the frequency but frequency density—the number of cases per unit of the variable on the horizontal axis. Examples of variable bin width are displayed on Census bureau data below.

As the adjacent bins leave no gaps, the rectangles of a histogram touch each other to indicate that the original variable is continuous.

Histograms give a rough sense of the density of the underlying distribution of the data, and often for density estimation: estimating the probability density function of the underlying variable. The total area of a histogram used for probability density is always normalized to 1. If the length of the intervals on the x-axis are all 1, then a histogram is identical to a relative frequency plot.

A histogram can be thought of as a simplistic kernel density estimation, which uses a kernel to smooth frequencies over the bins. This yields a smoother probability density function, which will in general more accurately reflect the distribution of the underlying variable. The density estimate could be plotted as an alternative to the histogram and is usually drawn as a curve rather than a set of boxes. Histograms are nevertheless preferred in applications, when their statistical properties need to be modeled. The correlated variation of a kernel density estimate is very difficult to describe mathematically, while it is simple for a histogram where each bin varies independently.

An alternative to kernel density estimation is the average shifted histogram, which is fast to compute and gives a smooth curve estimate of the density without using kernels.

The histogram is one of the seven basic tools of quality control.

Histograms are sometimes confused with bar charts. A histogram is used for continuous data, where the bins represent ranges of data, while a bar chart is a plot of categorical variables. Some authors recommend that bar charts have gaps between the rectangles to clarify the distinction.

\paragraph{Examples}

% [Image: hist]

This is the data for the histogram above, using 500 items:

$$
\begin{array}{|r|r|} \hline
   \textbf{Bin/Interval} & \textbf{Count/Frequency} \\\\hline
   -3.5 \text{ to } -2.51 & 9 \\\\ \hline
   -2.5 \text{ to } -1.51 & 32 \\\\ \hline
   -1.5 \text{ to } -0.51 & 109 \\\\ \hline
   -0.5 \text{ to } 0.49 & 180 \\\\ \hline
   0.5 \text{ to } 1.49 & 132 \\\\ \hline
   1.5 \text{ to } 2.49 & 34 \\\\ \hline
   2.5 \text{ to } 3.49 & 4 \\\\ \hline
\end{array}
$$

The words used to describe the patterns in a histogram are: "symmetric", "skewed left" or "right", "unimodal", "bimodal" or "multimodal".

\textbf{Symmetric, Unimodal}$\phantom{---}$\textbf{Skewed Right}$\phantom{---}$\textbf{Skewed Left}

$\phantom{-}$% [Image: sym uni]$\phantom{--}$% [Image: skew r]$\phantom{--}$% [Image: skew l]

_______________________________________________________________

$\phantom{---}$\textbf{Bimodal}$\phantom{----}$\textbf{Multimodal}$\phantom{----}$\textbf{Symmetric}

$\phantom{-}$% [Image: bimodal]$\phantom{--}$% [Image: multi]$\phantom{--}$% [Image: sym]

It is a good idea to plot the data using several different bin widths to learn more about it. Here is an example on tips given in a restaurant.

Tips using a `$`1 bin width, skewed right, unimodal:

% [Image: tips 1]

Tips using a 10c bin width, still skewed right, multimodal with modes at `$2` and `$`3 amounts, indicates rounding, also some outliers

% [Image: tips 2]

The U.S. Census Bureau found that there were 124 million people who work outside of their homes. Using their data on the time occupied by travel to work, the table below shows the absolute number of people who responded with travel times "at least 30 but less than 35 minutes" is higher than the numbers for the categories above and below it. This is likely due to people rounding their reported journey time. The problem of reporting values as somewhat arbitrarily rounded numbers is a common phenomenon when collecting data from people.

\textbf{Data by Absolute Numbers}
$$
\begin{array}{|c|c|c|c|}
\hline
\textbf{Interval} & \textbf{Width} & \textbf{Quantity} & \textbf{Quantity/Width} \\\\ \hline
0 & 5 & 4180 & 836 \\\\ \hline
5 & 5 & 13687 & 2737 \\\\ \hline
10 & 5 & 18618 & 3723 \\\\ \hline
15 & 5 & 19634 & 3926 \\\\ \hline
20 & 5 & 17981 & 3596 \\\\ \hline
25 & 5 & 7190 & 1438 \\\\ \hline
30 & 5 & 16369 & 3273 \\\\ \hline
35 & 5 & 3212 & 642 \\\\ \hline
40 & 5 & 4122 & 824 \\\\ \hline
45 & 15 & 9200 & 613 \\\\ \hline
60 & 30 & 6461 & 215 \\\\ \hline
90 & 60 & 3435 & 57 \\\\ \hline
\end{array}
$$

This histogram shows the number of cases per unit interval as the height of each block, so that the area of each block is equal to the number of people in the survey who fall into its category. The area under the curve represents the total number of cases (124 million). This type of histogram shows absolute numbers, with Q in thousands.

% [Image: hist]

\paragraph{Licensing}

Content obtained and/or adapted from:

\begin{itemize}
    \item \href{https://en.wikipedia.org/wiki/Chart}{Chart, Wikipedia} under a CC BY-SA license
\end{itemize}

\subsection{Learning Objectives}

\begin{enumerate}
    \item \textbf{Interpret Data in Common Charts} Students will be able to interpret and analyze data presented in common chart types, including histograms, bar charts, pie charts, and line charts. This involves understanding the information conveyed through each chart type and describing data trends and patterns.
    \item \textbf{Compare and Choose Appropriate Chart Types} Students will be able to compare different types of charts (e.g., histograms, line charts, pie charts) and select the most appropriate chart type to effectively represent a given data set. This includes understanding the advantages and limitations of each chart type for visualizing specific data characteristics.
    \item \textbf{Construct and Label Charts} Students will be able to construct various types of charts (e.g., histograms, bar charts, line charts) from provided data, ensuring that they accurately label axes, include titles, and use legends where necessary to enhance clarity and understanding.
\end{enumerate}

\subsection{Assessment Instrument}

\begin{enumerate}
    \item Part 1: Identification and Description
    \item \textbf{Multiple Choice} \textbf{Which type of chart is best suited for displaying percentages of different categories?}
    \item a) Histogram
    \item b) Line Chart
    \item c) Pie Chart
    \item d) Dot Plot
    \item \textbf{True/False} \textbf{A bar chart and a histogram are identical in function and representation.}
    \item True
    \item False
    \item \textbf{Fill in the Blank} \textbf{In a histogram, the vertical axis typically represents the __________ of cases in each bin.}
    \item Frequency
    \item Mean
    \item Median
    \item Mode Part 2: Chart Features and Analysis
    \item \textbf{Short Answer} \textbf{Describe the purpose of gridlines in a chart.}
    \item \textbf{Short Answer} \textbf{Explain the role of the legend in a chart.}
    \item \textbf{Matching} \textbf{Match the chart type with its primary use.}
    \item \textbf{Chart Type}
    \item a) Line Chart
    \item b) Dot Plot
    \item c) Pie Chart
    \item d) Histogram
    \item \textbf{Use}
    \item 1) Showing distribution of data across intervals
    \item 2) Displaying data changes over time
    \item 3) Visualizing percentage shares of a whole
    \item 4) Representing individual data points Part 3: Data Interpretation
    \item \textbf{Data Analysis} \textbf{Given the following histogram data, calculate the mean value.} $$ \begin{array}{|c|c|} \hline \textbf{Bin/Interval} & \textbf{Count/Frequency} \\\\ \hline 0-1 & 5 \\\\ \hline 1-2 & 15 \\\\ \hline 2-3 & 25 \\\\ \hline 3-4 & 10 \\\\ \hline 4-5 & 5 \\\\ \hline \end{array} $$
    \item \textbf{Chart Creation} \textbf{Create a histogram using the following data set. Assume equal-width bins of size 5.} Data Set: $ \{1, 3, 7, 8, 10, 15, 18, 20, 25, 30\} $
    \item \textbf{Interpretation of Dot Plot} \textbf{Analyze the dot plot provided below and describe the shape of the data} !\href{https://upload.wikimedia.org/wikipedia/commons/thumb/b/bc/Dotplot_of_random_values_2.png/300px-Dotplot_of_random_values_2.png}{Dot Plot} Part 4: Advanced Analysis
    \item \textbf{Short Answer} \textbf{Discuss how varying bin widths in a histogram can affect the interpretation of the data.}
    \item \textbf{Calculation} \textbf{Using the following frequency data, plot a histogram and calculate the mean of the data.} $$ \begin{array}{|c|c|} \hline \textbf{Bin/Interval} & \textbf{Count/Frequency} \\\\ \hline 0-10 & 50 \\\\ \hline 10-20 & 75 \\\\ \hline 20-30 & 100 \\\\ \hline 30-40 & 80 \\\\ \hline \end{array} $$
    \item \textbf{Chart Identification} \textbf{Given a chart with a title "Annual Revenue from 2000 to 2020," identify which type of chart it most likely is and justify your answer.}
\end{enumerate}

%%===============================================================
\section{Analyzing graphical displays}
\label{sec:analyzing-graphical-displays}
%%===============================================================

Analyzing graphical displays

%%===============================================================
\section{Sampling}
\label{sec:sampling}
%%===============================================================

\paragraph{Sampling}



A visual representation of the sampling process:



% [Image: Sample]



In statistics, quality assurance, and survey methodology, sampling is the selection of a subset (a statistical sample) of individuals from within a statistical population to estimate characteristics of the whole population. Statisticians attempt for the samples to represent the population in question. Two advantages of sampling are lower cost and faster data collection than measuring the entire population.



Each observation measures one or more properties (such as weight, location, color) of observable bodies distinguished as independent objects or individuals. In survey sampling, weights can be applied to the data to adjust for the sample design, particularly in stratified sampling. Results from probability theory and statistical theory are employed to guide the practice. In business and medical research, sampling is widely used for gathering information about a population. Acceptance sampling is used to determine if a production lot of material meets the governing specifications.





\paragraph{Population Definition}



Successful statistical practice is based on focused problem definition. In sampling, this includes defining the "population" from which our sample is drawn. A population can be defined as including all people or items with the characteristic one wishes to understand. Because there is very rarely enough time or money to gather information from everyone or everything in a population, the goal becomes finding a representative sample (or subset) of that population.



Sometimes what defines a population is obvious. For example, a manufacturer needs to decide whether a batch of material from production is of high enough quality to be released to the customer, or should be sentenced for scrap or rework due to poor quality. In this case, the batch is the population.



Although the population of interest often consists of physical objects, sometimes it is necessary to sample over time, space, or some combination of these dimensions. For instance, an investigation of supermarket staffing could examine checkout line length at various times, or a study on endangered penguins might aim to understand their usage of various hunting grounds over time. For the time dimension, the focus may be on periods or discrete occasions.



In other cases, the examined ''''population'''' may be even less tangible. For example, Joseph Jagger studied the behavior of roulette wheels at a casino in Monte Carlo, and used this to identify a biased wheel. In this case, the ''''population'''' Jagger wanted to investigate was the overall behavior of the wheel (i.e. the probability distribution of its results over infinitely many trials), while his ''''sample'''' was formed from observed results from that wheel. Similar considerations arise when taking repeated measurements of some physical characteristic such as the electrical conductivity of copper.



This situation often arises when seeking knowledge about the cause system of which the observed population is an outcome. In such cases, sampling theory may treat the observed population as a sample from a larger ''''superpopulation''''. For example, a researcher might study the success rate of a new ''''quit smoking'''' program on a test group of 100 patients, in order to predict the effects of the program if it were made available nationwide. Here the superpopulation is "everybody in the country, given access to this treatment" – a group which does not yet exist, since the program isn''''t yet available to all.



The population from which the sample is drawn may not be the same as the population about which information is desired. Often there is large but not complete overlap between these two groups due to frame issues etc. (see below). Sometimes they may be entirely separate – for instance, one might study rats in order to get a better understanding of human health, or one might study records from people born in 2008 in order to make predictions about people born in 2009.



Time spent in making the sampled population and population of concern precise is often well spent, because it raises many issues, ambiguities and questions that would otherwise have been overlooked at this stage.



\paragraph{Sampling Frame}



In the most straightforward case, such as the sampling of a batch of material from production (acceptance sampling by lots), it would be most desirable to identify and measure every single item in the population and to include any one of them in our sample. However, in the more general case this is not usually possible or practical. There is no way to identify all rats in the set of all rats. Where voting is not compulsory, there is no way to identify which people will vote at a forthcoming election (in advance of the election). These imprecise populations are not amenable to sampling in any of the ways below and to which we could apply statistical theory.



As a remedy, we seek a sampling frame which has the property that we can identify every single element and include any in our sample. The most straightforward type of frame is a list of elements of the population (preferably the entire population) with appropriate contact information. For example, in an opinion poll, possible sampling frames include an electoral register and a telephone directory.



A probability sample is a sample in which every unit in the population has a chance (greater than zero) of being selected in the sample, and this probability can be accurately determined. The combination of these traits makes it possible to produce unbiased estimates of population totals, by weighting sampled units according to their probability of selection.



\textbf{Example:} We want to estimate the total income of adults living in a given street. We visit each household in that street, identify all adults living there, and randomly select one adult from each household. (For example, we can allocate each person a random number, generated from a uniform distribution between 0 and 1, and select the person with the highest number in each household). We then interview the selected person and find their income.



People living on their own are certain to be selected, so we simply add their income to our estimate of the total. But a person living in a household of two adults has only a one-in-two chance of selection. To reflect this, when we come to such a household, we would count the selected person''''s income twice towards the total. (The person who is selected from that household can be loosely viewed as also representing the person who isn''''t selected.)



In the above example, not everybody has the same probability of selection; what makes it a probability sample is the fact that each person''''s probability is known. When every element in the population does have the same probability of selection, this is known as an ''''equal probability of selection'''' (EPS) design. Such designs are also referred to as ''''self-weighting'''' because all sampled units are given the same weight.



Probability sampling includes: Simple Random Sampling, Systematic Sampling, Stratified Sampling, Probability Proportional to Size Sampling, and Cluster or Multistage Sampling. These various ways of probability sampling have two things in common:



1. Every element has a known nonzero probability of being sampled and

2. Involves random selection at some point.



\paragraph{Nonprobability Sampling}



Nonprobability sampling is any sampling method where some elements of the population have no chance of selection (these are sometimes referred to as ''''out of coverage''''/''''undercovered''''), or where the probability of selection can''''t be accurately determined. It involves the selection of elements based on assumptions regarding the population of interest, which forms the criteria for selection. Hence, because the selection of elements is nonrandom, nonprobability sampling does not allow the estimation of sampling errors. These conditions give rise to exclusion bias, placing limits on how much information a sample can provide about the population. Information about the relationship between sample and population is limited, making it difficult to extrapolate from the sample to the population.



\textbf{Example:} We visit every household in a given street, and interview the first person to answer the door. In any household with more than one occupant, this is a nonprobability sample, because some people are more likely to answer the door (e.g. an unemployed person who spends most of their time at home is more likely to answer than an employed housemate who might be at work when the interviewer calls) and it''''s not practical to calculate these probabilities.



Nonprobability sampling methods include convenience sampling, quota sampling, and purposive sampling. In addition, nonresponse effects may turn any probability design into a nonprobability design if the characteristics of nonresponse are not well understood, since nonresponse effectively modifies each element''''s probability of being sampled.



\paragraph{Sampling Methods}



Within any of the types of frames identified above, a variety of sampling methods can be employed, individually or in combination. Factors commonly influencing the choice between these designs include:



\begin{itemize}
    \item Nature and quality of the frame
    \item Availability of auxiliary information about units on the frame
    \item Accuracy requirements, and the need to measure accuracy
    \item Whether detailed analysis of the sample is expected
    \item Cost/operational concerns
\end{itemize}

\paragraph{Simple Random Sampling}



A visual representation of selecting a simple random sample:



% [Image: SRS]



In a simple random sample (SRS) of a given size, all subsets of a sampling frame have an equal probability of being selected. Each element of the frame thus has an equal probability of selection: the frame is not subdivided or partitioned. Furthermore, any given pair of elements has the same chance of selection as any other such pair (and similarly for triples, and so on). This minimizes bias and simplifies analysis of results. In particular, the variance between individual results within the sample is a good indicator of variance in the overall population, which makes it relatively easy to estimate the accuracy of results.



Simple random sampling can be vulnerable to sampling error because the randomness of the selection may result in a sample that doesn''''t reflect the makeup of the population. For instance, a simple random sample of ten people from a given country will on average produce five men and five women, but any given trial is likely to overrepresent one sex and underrepresent the other. Systematic and stratified techniques attempt to overcome this problem by "using information about the population" to choose a more "representative" sample.



Also, simple random sampling can be cumbersome and tedious when sampling from a large target population. In some cases, investigators are interested in research questions specific to subgroups of the population. For example, researchers might be interested in examining whether cognitive ability as a predictor of job performance is equally applicable across racial groups. Simple random sampling cannot accommodate the needs of researchers in this situation, because it does not provide subsamples of the population, and other sampling strategies, such as stratified sampling, can be used instead.





\paragraph{Systematic Sampling}



A visual representation of selecting a random sample using the systematic sampling technique:



% [Image: sys]



Systematic sampling (also known as interval sampling) relies on arranging the study population according to some ordering scheme and then selecting elements at regular intervals through that ordered list. Systematic sampling involves a random start and then proceeds with the selection of every $k$-th element from then onwards. In this case, $k = \frac{\text{population size}}{\text{sample size}}$. It is important that the starting point is not automatically the first in the list, but is instead randomly chosen from within the first to the $k$-th element in the list. A simple example would be to select every 10th name from the telephone directory (an ''''every 10th'''' sample, also referred to as ''''sampling with a skip of 10'''').



As long as the starting point is randomized, systematic sampling is a type of probability sampling. It is easy to implement and the stratification induced can make it efficient, if the variable by which the list is ordered is correlated with the variable of interest. ''''Every 10th'''' sampling is especially useful for efficient sampling from databases.



\textbf{Example:} Suppose we wish to sample people from a long street that starts in a poor area (house No. 1) and ends in an expensive district (house No. 1000). A simple random selection of addresses from this street could easily end up with too many from the high end and too few from the low end (or vice versa), leading to an unrepresentative sample. Selecting (e.g.) every 10th street number along the street ensures that the sample is spread evenly along the length of the street, representing all of these districts. (Note that if we always start at house #1 and end at #991, the sample is slightly biased towards the low end; by randomly selecting the start between #1 and #10, this bias is eliminated.)



However, systematic sampling is especially vulnerable to periodicities in the list. If periodicity is present and the period is a multiple or factor of the interval used, the sample is especially likely to be unrepresentative of the overall population, making the scheme less accurate than simple random sampling.



\textbf{Example:} Consider a street where the odd-numbered houses are all on the north (expensive) side of the road, and the even-numbered houses are all on the south (cheap) side. Under the sampling scheme given above, it is impossible to get a representative sample; either the houses sampled will all be from the odd-numbered, expensive side, or they will all be from the even-numbered, cheap side, unless the researcher has previous knowledge of this bias and avoids it by using a skip which ensures jumping between the two sides (any odd-numbered skip).



Another drawback of systematic sampling is that even in scenarios where it is more accurate than SRS, its theoretical properties make it difficult to quantify that accuracy. (In the two examples of systematic sampling that are given above, much of the potential sampling error is due to variation between neighboring houses – but because this method never selects two neighboring houses, the sample will not give us any information on that variation.)



As described above, systematic sampling is an EPS method, because all elements have the same probability of selection (in the example given, one in ten). It is not ''''simple random sampling'''' because different subsets of the same size have different selection probabilities – e.g. the set {$4, 14, 24, \ldots, 994$} has a one-in-ten probability of selection, but the set {$4, 13, 24, 34, \ldots$} has zero probability of selection.



Systematic sampling can also be adapted to a non-EPS approach; for an example, see discussion of PPS samples below.



\paragraph{Stratified Sampling}



A visual representation of selecting a random sample using the stratified sampling technique:



% [Image: strat]



When the population embraces a number of distinct categories, the frame can be organized by these categories into separate "strata." Each stratum is then sampled as an independent sub-population, out of which individual elements can be randomly selected. The ratio of the size of this random selection (or sample) to the size of the population is called a sampling fraction. There are several potential benefits to stratified sampling.



First, dividing the population into distinct, independent strata can enable researchers to draw inferences about specific subgroups that may be lost in a more generalized random sample.



Second, utilizing a stratified sampling method can lead to more efficient statistical estimates (provided that strata are selected based upon relevance to the criterion in question, instead of availability of the samples). Even if a stratified sampling approach does not lead to increased statistical efficiency, such a tactic will not result in less efficiency than would simple random sampling, provided that each stratum is proportional to the group''''s size in the population.



Third, it is sometimes the case that data are more readily available for individual, pre-existing strata within a population than for the overall population; in such cases, using a stratified sampling approach may be more convenient than aggregating data across groups (though this may potentially be at odds with the previously noted importance of utilizing criterion-relevant strata).



Finally, since each stratum is treated as an independent population, different sampling approaches can be applied to different strata, potentially enabling researchers to use the approach best suited (or most cost-effective) for each identified subgroup within the population.



There are, however, some potential drawbacks to using stratified sampling. First, identifying strata and implementing such an approach can increase the cost and complexity of sample selection, as well as leading to increased complexity of population estimates. Second, when examining multiple criteria, stratifying variables may be related to some, but not to others, further complicating the design, and potentially reducing the utility of the strata. Finally, in some cases (such as designs with a large number of strata, or those with a specified minimum sample size per group), stratified sampling can potentially require a larger sample than would other methods (although in most cases, the required sample size would be no larger than would be required for simple random sampling).



A stratified sampling approach is most effective when three conditions are met:

1. Variability within strata is minimized

2. Variability between strata is maximized

3. The variables upon which the population is stratified are strongly correlated with the desired dependent variable.



\paragraph{Advantages over other sampling methods}



\begin{itemize}
    \item Focuses on important subpopulations and ignores irrelevant ones.
    \item Allows use of different sampling techniques for different subpopulations.
    \item Improves the accuracy/efficiency of estimation.
    \item Permits greater balancing of statistical power of tests of differences between strata by sampling equal numbers from strata varying widely in size.
\end{itemize}

\paragraph{Disadvantages}



\begin{itemize}
    \item Requires selection of relevant stratification variables which can be difficult.
    \item Is not useful when there are no homogeneous subgroups.
    \item Can be expensive to implement.
\end{itemize}

\paragraph{Poststratification}



Stratification is sometimes introduced after the sampling phase in a process called "poststratification." This approach is typically implemented due to a lack of prior knowledge of an appropriate stratifying variable or when the experimenter lacks the necessary information to create a stratifying variable during the sampling phase. Although the method is susceptible to the pitfalls of post hoc approaches, it can provide several benefits in the right situation. Implementation usually follows a simple random sample. In addition to allowing for stratification on an ancillary variable, poststratification can be used to implement weighting, which can improve the precision of a sample''''s estimates.



\paragraph{Oversampling}



Choice-based sampling is one of the stratified sampling strategies. In choice-based sampling, the data are stratified on the target and a sample is taken from each stratum so that the rare target class will be more represented in the sample. The model is then built on this biased sample. The effects of the input variables on the target are often estimated with more precision with the choice-based sample even when a smaller overall sample size is taken, compared to a random sample. The results usually must be adjusted to correct for the oversampling.



\paragraph{Probability-Proportional-to-Size Sampling}



In some cases the sample designer has access to an "auxiliary variable" or "size measure", believed to be correlated to the variable of interest, for each element in the population. These data can be used to improve accuracy in sample design. One option is to use the auxiliary variable as a basis for stratification, as discussed above.



Another option is probability proportional to size (''''PPS'''') sampling, in which the selection probability for each element is set to be proportional to its size measure, up to a maximum of 1. In a simple PPS design, these selection probabilities can then be used as the basis for Poisson sampling. However, this has the drawback of variable sample size, and different portions of the population may still be over- or under-represented due to chance variation in selections.



Systematic sampling theory can be used to create a probability proportionate to size sample. This is done by treating each count within the size variable as a single sampling unit. Samples are then identified by selecting at even intervals among these counts within the size variable. This method is sometimes called PPS-sequential or monetary unit sampling in the case of audits or forensic sampling.



\textbf{Example:} Suppose we have six schools with populations of 150, 180, 200, 220, 260, and 490 students respectively (total 1500 students), and we want to use student population as the basis for a PPS sample of size three. To do this, we could allocate the first school numbers 1 to 150, the second school 151 to 330 (= 150 + 180), the third school 331 to 530, and so on to the last school (1011 to 1500). We then generate a random start between 1 and 500 (equal to 1500/3) and count through the school populations by multiples of 500. If our random start was 137, we would select the schools which have been allocated numbers 137, 637, and 1137, i.e. the first, fourth, and sixth schools.



The PPS approach can improve accuracy for a given sample size by concentrating sample on large elements that have the greatest impact on population estimates. PPS sampling is commonly used for surveys of businesses, where element size varies greatly and auxiliary information is often available – for instance, a survey attempting to measure the number of guest-nights spent in hotels might use each hotel''''s number of rooms as an auxiliary variable. In some cases, an older measurement of the variable of interest can be used as an auxiliary variable when attempting to produce more current estimates.



\paragraph{Cluster Sampling}



A visual representation of selecting a random sample using the cluster sampling technique:



% [Image: cluster]



Sometimes it is more cost-effective to select respondents in groups (''''clusters''''). Sampling is often clustered by geography, or by time periods. (Nearly all samples are in some sense ''''clustered'''' in time – although this is rarely taken into account in the analysis.) For instance, if surveying households within a city, we might choose to select 100 city blocks and then interview every household within the selected blocks.



Clustering can reduce travel and administrative costs. In the example above, an interviewer can make a single trip to visit several households in one block, rather than having to drive to a different block for each household.



It also means that one does not need a sampling frame listing all elements in the target population. Instead, clusters can be chosen from a cluster-level frame, with an element-level frame created only for the selected clusters. In the example above, the sample only requires a block-level city map for initial selections, and then a household-level map of the 100 selected blocks, rather than a household-level map of the whole city.



Cluster sampling (also known as clustered sampling) generally increases the variability of sample estimates above that of simple random sampling, depending on how the clusters differ between one another as compared to the within-cluster variation. For this reason, cluster sampling requires a larger sample than SRS to achieve the same level of accuracy – but cost savings from clustering might still make this a cheaper option.



Cluster sampling is commonly implemented as multistage sampling. This is a complex form of cluster sampling in which two or more levels of units are embedded one in the other. The first stage consists of constructing the clusters that will be used to sample from. In the second stage, a sample of primary units is randomly selected from each cluster (rather than using all units contained in all selected clusters). In following stages, in each of those selected clusters, additional samples of units are selected, and so on. All ultimate units (individuals, for instance) selected at the last step of this procedure are then surveyed. This technique, thus, is essentially the process of taking random subsamples of preceding random samples.



Multistage sampling can substantially reduce sampling costs, where the complete population list would need to be constructed (before other sampling methods could be applied). By eliminating the work involved in describing clusters that are not selected, multistage sampling can reduce the large costs associated with traditional cluster sampling. However, each sample may not be a full representative of the whole population.



\paragraph{Quota Sampling}



In quota sampling, the population is first segmented into mutually exclusive sub-groups, just as in stratified sampling. Then judgement is used to select the subjects or units from each segment based on a specified proportion. For example, an interviewer may be told to sample 200 females and 300 males between the age of 45 and 60.



It is this second step which makes the technique one of non-probability sampling. In quota sampling the selection of the sample is non-random. For example, interviewers might be tempted to interview those who look most helpful. The problem is that these samples may be biased because not everyone gets a chance of selection. This random element is its greatest weakness and quota versus probability has been a matter of controversy for several years.



\paragraph{Minimax Sampling}



In imbalanced datasets, where the sampling ratio does not follow the population statistics, one can resample the dataset in a conservative manner called minimax sampling. The minimax sampling has its origin in Anderson minimax ratio whose value is proved to be 0.5: in a binary classification, the class-sample sizes should be chosen equally. This ratio can be proved to be minimax ratio only under the assumption of LDA classifier with Gaussian distributions. The notion of minimax sampling is recently developed for a general class of classification rules, called class-wise smart classifiers. In this case, the sampling ratio of classes is selected so that the worst case classifier error over all the possible population statistics for class prior probabilities, would be the best.



\paragraph{Accidental Sampling}



Accidental sampling (sometimes known as grab, convenience or opportunity sampling) is a type of nonprobability sampling which involves the sample being drawn from that part of the population which is close to hand. That is, a population is selected because it is readily available and convenient. It may be through meeting the person or including a person in the sample when one meets them or chosen by finding them through technological means such as the internet or through phone. The researcher using such a sample cannot scientifically make generalizations about the total population from this sample because it would not be representative enough. For example, if the interviewer were to conduct such a survey at a shopping center early in the morning on a given day, the people that he/she could interview would be limited to those given there at that given time, which would not represent the views of other members of society in such an area, if the survey were to be conducted at different times of day and several times per week. This type of sampling is most useful for pilot testing. Several important considerations for researchers using convenience samples include:



\begin{itemize}
    \item Are there controls within the research design or experiment which can serve to lessen the impact of a non-random convenience sample, thereby ensuring the results will be more representative of the population?
    \item Is there good reason to believe that a particular convenience sample would or should respond or behave differently than a random sample from the same population?
    \item Is the question being asked by the research one that can adequately be answered using a convenience sample?
\end{itemize}

In social science research, snowball sampling is a similar technique, where existing study subjects are used to recruit more subjects into the sample. Some variants of snowball sampling, such as respondent driven sampling, allow calculation of selection probabilities and are probability sampling methods under certain conditions.



\paragraph{Voluntary Sampling}



The voluntary sampling method is a type of non-probability sampling. Volunteers choose to complete a survey.



Volunteers may be invited through advertisements in social media. The target population for advertisements can be selected by characteristics like location, age, sex, income, occupation, education or interests using tools provided by the social medium. The advertisement may include a message about the research and link to a survey. After following the link and completing the survey the volunteer submits the data to be included in the sample population. This method can reach a global population but is limited by the campaign budget. Volunteers outside the invited population may also be included in the sample.



It is difficult to make generalizations from this sample because it may not represent the total population. Often, volunteers have a strong interest in the main topic of the survey.



\paragraph{Line-Intercept Sampling}



Line-intercept sampling is a method of sampling elements in a region whereby an element is sampled if a chosen line segment, called a "transect", intersects the element.



\paragraph{Panel Sampling}



Panel sampling is the method of first selecting a group of participants through a random sampling method and then asking that group for (potentially the same) information several times over a period of time. Therefore, each participant is interviewed at two or more time points; each period of data collection is called a "wave". The method was developed by sociologist Paul Lazarsfeld in 1938 as a means of studying political campaigns. This longitudinal sampling-method allows estimates of changes in the population, for example with regard to chronic illness to job stress to weekly food expenditures. Panel sampling can also be used to inform researchers about within-person health changes due to age or to help explain changes in continuous dependent variables such as spousal interaction. There have been several proposed methods of analyzing panel data, including MANOVA, growth curves, and structural equation modeling with lagged effects.



\paragraph{Snowball Sampling}



Snowball sampling involves finding a small group of initial respondents and using them to recruit more respondents. It is particularly useful in cases where the population is hidden or difficult to enumerate.



\paragraph{Theoretical Sampling}



Theoretical sampling occurs when samples are selected on the basis of the results of the data collected so far with a goal of developing a deeper understanding of the area or develop theories. Extreme or very specific cases might be selected in order to maximize the likelihood a phenomenon will actually be observable.



\paragraph{Replacement of Selected Units}



Sampling schemes may be without replacement (''''WOR'''' – no element can be selected more than once in the same sample) or with replacement (''''WR'''' – an element may appear multiple times in the one sample). For example, if we catch fish, measure them, and immediately return them to the water before continuing with the sample, this is a WR design, because we might end up catching and measuring the same fish more than once. However, if we do not return the fish to the water or tag and release each fish after catching it, this becomes a WOR design.



\paragraph{Sample Size Determination}



Formulas, tables, and power function charts are well known approaches to determine sample size.



\textbf{Steps for using sample size tables:}



1. Postulate the effect size of interest, a, and ß.

2. Check sample size table

3. Select the table corresponding to the selected a

4. Locate the row corresponding to the desired power

5. Locate the column corresponding to the estimated effect size.

6. The intersection of the column and row is the minimum sample size required.



\paragraph{Sampling and Data Collection}



Good data collection involves:



\begin{itemize}
    \item Following the defined sampling process
    \item Keeping the data in time order
    \item Noting comments and other contextual events
    \item Recording non-responses
\end{itemize}

\paragraph{Applications of Sampling}



Sampling enables the selection of right data points from within the larger data set to estimate the characteristics of the whole population. For example, there are about 600 million tweets produced every day. It is not necessary to look at all of them to determine the topics that are discussed during the day, nor is it necessary to look at all the tweets to determine the sentiment on each of the topics. A theoretical formulation for sampling Twitter data has been developed.



In manufacturing different types of sensory data such as acoustics, vibration, pressure, current, voltage and controller data are available at short time intervals. To predict down-time it may not be necessary to look at all the data but a sample may be sufficient.



\paragraph{Errors in Sample Surveys}



Survey results are typically subject to some error. Total errors can be classified into sampling errors and non-sampling errors. The term "error" here includes systematic biases as well as random errors.



\paragraph{Sampling Errors and Biases}



Sampling errors and biases are induced by the sample design. They include:



\begin{itemize}
    \item \textbf{Selection Bias:} When the true selection probabilities differ from those assumed in calculating the results.
    \item \textbf{Random Sampling Error:} Random variation in the results due to the elements in the sample being selected at random.
\end{itemize}

\paragraph{Non-Sampling Error}



Non-sampling errors are other errors which can impact final survey estimates, caused by problems in data collection, processing, or sample design. Such errors may include:



\begin{itemize}
    \item \textbf{Over-coverage:} Inclusion of data from outside of the population
    \item \textbf{Under-coverage:} Sampling frame does not include elements in the population.
    \item \textbf{Measurement Error:} e.g., when respondents misunderstand a question, or find it difficult to answer
    \item \textbf{Processing Error:} Mistakes in data coding
    \item \textbf{Non-response or Participation Bias:} Failure to obtain complete data from all selected individuals
\end{itemize}

After sampling, a review should be held of the exact process followed in sampling, rather than that intended, in order to study any effects that any divergences might have on subsequent analysis.



A particular problem involves non-response. Two major types of non-response exist:



\begin{itemize}
    \item \textbf{Unit Nonresponse:} Lack of completion of any part of the survey
    \item \textbf{Item Non-response:} Submission or participation in survey but failing to complete one or more components/questions of the survey
\end{itemize}

In survey sampling, many of the individuals identified as part of the sample may be unwilling to participate, not have the time to participate (opportunity cost), or survey administrators may not have been able to contact them. In this case, there is a risk of differences between respondents and nonrespondents, leading to biased estimates of population parameters. This is often addressed by improving survey design, offering incentives, and conducting follow-up studies which make a repeated attempt to contact the unresponsive and to characterize their similarities and differences with the rest of the frame. The effects can also be mitigated by weighting the data (when population benchmarks are available) or by imputing data based on answers to other questions. Nonresponse is particularly a problem in internet sampling. Reasons for this problem may include improperly designed surveys, over-surveying (or survey fatigue), and the fact that potential participants may have multiple e-mail addresses, which they don''''t use anymore or don''''t check regularly.



\paragraph{Survey Weights}



In many situations the sample fraction may be varied by stratum and data will have to be weighted to correctly represent the population. Thus for example, a simple random sample of individuals in the United Kingdom might not include some in remote Scottish islands who would be inordinately expensive to sample. A cheaper method would be to use a stratified sample with urban and rural strata. The rural sample could be under-represented in the sample, but weighted up appropriately in the analysis to compensate.



More generally, data should usually be weighted if the sample design does not give each individual an equal chance of being selected. For instance, when households have equal selection probabilities but one person is interviewed from within each household, this gives people from large households a smaller chance of being interviewed. This can be accounted for using survey weights. Similarly, households with more than one telephone line have a greater chance of being selected in a random digit dialing sample, and weights can adjust for this.



Weights can also serve other purposes, such as helping to correct for non-response.



\paragraph{Methods of Producing Random Samples}



\begin{itemize}
    \item Random number table
    \item Mathematical algorithms for pseudo-random number generators
    \item Physical randomization devices such as coins, playing cards or sophisticated devices such as ERNIE
\end{itemize}

\paragraph{Resources}



\begin{itemize}
    \item \href{https://courses.lumenlearning.com/math4libarts/chapter/sampling-methods/}{Sampling Methods, Lumen Learning}
    \item \href{https://www.khanacademy.org/math/statistics-probability/designing-studies/sampling-methods-stats/a/sampling-methods-review}{Types of Sampling Methods, Khan Academy}
\end{itemize}

\paragraph{Licensing}



Content obtained and/or adapted from:



\begin{itemize}
    \item \href{https://en.wikipedia.org/wiki/Sampling_(statistics}{Sampling (statistics), Wikipedia}) under a CC BY-SA license
\end{itemize}

\subsection{Learning Objectives}

\begin{enumerate}
    \item \textbf{SLO 1: Understanding Sampling Techniques}
    \item \textbf{Description:} Students will be able to identify and describe various sampling techniques such as simple random sampling, stratified sampling, cluster sampling, and systematic sampling. \textbf{SLO 2: Application of Sampling Methods}
    \item \textbf{Description:} Students will demonstrate the ability to apply appropriate sampling methods to different types of research problems. \textbf{SLO 3: Addressing Sampling Errors and Biases}
    \item \textbf{Description:} Students will be able to recognize potential sampling errors and biases, and propose strategies to mitigate these issues.
\end{enumerate}

\subsection{Assessment Instrument}

\begin{enumerate}
    \item \textbf{Part 1: Multiple-Choice Quiz}
    \item Which of the following is an example of systematic sampling?
    \item A) Every 10th person in a list is selected.
    \item B) The population is divided into strata and random samples are taken from each stratum.
    \item C) A random sample of clusters is selected, and then all individuals within the selected clusters are sampled.
    \item D) A random number table is used to select individuals.
    \item In stratified sampling, the population is divided into:
    \item A) Random groups
    \item B) Homogeneous subgroups
    \item C) Clusters
    \item D) Systematic intervals
    \item Which sampling method is best suited for a population that is geographically dispersed?
    \item A) Simple random sampling
    \item B) Cluster sampling
    \item C) Stratified sampling
    \item D) Systematic sampling \textbf{Part 2: Scenario-Based Application} \textbf{Scenario:} A researcher wants to study the dietary habits of high school students across a large city. \textbf{Instructions:}
    \item Select the most appropriate sampling method for this study. Justify your choice.
    \item Outline the steps you would take to implement this sampling method.
    \item Describe how you would address potential sampling errors and biases. \textbf{Rubric:}
    \item \textbf{Selection of Sampling Method (30 points):} Appropriateness and justification of the selected method.
    \item \textbf{Implementation Steps (40 points):} Clear and detailed steps for conducting the sampling.
    \item \textbf{Addressing Errors and Biases (30 points):} Identification of potential issues and proposed mitigation strategies.
\end{enumerate}

%%===============================================================
\section{Mean and Central Limit Theorem}
\label{sec:mean-and-central-limit-theorem}
%%===============================================================

\paragraph{Mean and Central Limit Theorem}





\paragraph{Mean}



The mean, or more precisely the arithmetic mean, is simply the arithmetic average of a group of numbers (or data set) and is shown using the  $\phantom{-} \bar{} \phantom{-}$  symbol. So the mean of the variable $x$ is $\bar{x}$, pronounced "x-bar". It is calculated by adding up all of the values in a data set and dividing by the number of values in that data set:



$$

\bar{x} = \frac{\sum x}{n}

$$



For example, take the following set of data: {$1, 2, 3, 4, 5$}. The mean of this data would be:



$$

\bar{x} = \frac{1 + 2 + 3 + 4 + 5}{5} = \frac{15}{5} = 3

$$



Here is a more complicated data set: {$10, 14, 86, 2, 68, 99, 1$}. The mean would be calculated like this:



$$

\bar{x} = \frac{10 + 14 + 86 + 2 + 68 + 99 + 1}{7} = \frac{280}{7} = 40

$$



\paragraph{Central Limit Theorem}



In probability theory, the central limit theorem (CLT) establishes that, in many situations, when independent random variables are summed up, their properly normalized sum tends toward a normal distribution (informally a bell curve) even if the original variables themselves are not normally distributed. The theorem is a key concept in probability theory because it implies that probabilistic and statistical methods that work for normal distributions can be applicable to many problems involving other types of distributions. This theorem has seen many changes during the formal development of probability theory. Previous versions of the theorem date back to 1811, but in its modern general form, this fundamental result in probability theory was precisely stated as late as 1920, thereby serving as a bridge between classical and modern probability theory.



If $(X\_1, X\_2, \ldots, X\_n)$ are $n$ random samples drawn from a population with overall mean $\mu$ and finite variance $\sigma^2$, and if $\bar{X}\_n$ is the sample mean, then the limiting form of the distribution:



$$

Z = \lim\_{n \to \infty} \sqrt{n} \left(\frac{\bar{X}\_n - \mu}{\sigma}\right)

$$



is a standard normal distribution.



For example, suppose that a sample is obtained containing many observations, each observation being randomly generated in a way that does not depend on the values of the other observations, and that the arithmetic mean of the observed values is computed. If this procedure is performed many times, the central limit theorem says that the probability distribution of the average will closely approximate a normal distribution. A simple example of this is that if one flips a coin many times, the probability of getting a given number of heads will approach a normal distribution, with the mean equal to half the total number of flips. At the limit of an infinite number of flips, it will equal a normal distribution.



The central limit theorem has several variants. In its common form, the random variables must be identically distributed. In variants, convergence of the mean to the normal distribution also occurs for non-identical distributions or for non-independent observations, if they comply with certain conditions.



The earliest version of this theorem, that the normal distribution may be used as an approximation to the binomial distribution, is the de Moivre–Laplace theorem.



\paragraph{Resources}



\begin{itemize}
    \item \href{https://www.khanacademy.org/math/ap-statistics/sampling-distribution-ap/what-is-sampling-distribution/v/central-limit-theorem}{Central Limit Theorem, Khan Academy}
    \item \href{https://www.youtube.com/watch?v=4YLtvNeRIrg}{Central Limit Theorem - Sampling Distribution of Sample Means - Stats \& Probability, The Organic Chemistry Tutor}
\end{itemize}

\paragraph{Licensing}



Content obtained and/or adapted from:



\begin{itemize}
    \item \href{https://en.wikibooks.org/wiki/Statistics/Summary/Averages/mean}{Statistics/Summary/Averages/Mean, Wikibooks} under a CC BY-SA license
    \item \href{https://en.wikipedia.org/wiki/Central_limit_theorem}{Central limit theorem, Wikipedia} under a CC BY-SA license
\end{itemize}

\subsection{Learning Objectives}

\begin{enumerate}
    \item \textbf{Calculate the Mean:}
    \item Students will be able to calculate the mean of a given data set using the formula: $$ \bar{x} = \frac{\sum x}{n} $$
    \item \textbf{Apply the Central Limit Theorem:}
    \item Students will be able to explain and apply the Central Limit Theorem to determine how the distribution of sample means approaches a normal distribution as the sample size increases.
    \item \textbf{Interpret Data with CLT:}
    \item Students will be able to interpret the implications of the Central Limit Theorem in practical scenarios, such as predicting the normality of sample means from random variables.
\end{enumerate}

\subsection{Assessment Instrument}

\begin{enumerate}
    \item 1. Calculation of Mean \textbf{Instructions:} Compute the mean for the following data sets. a. Data set: {5, 15, 25, 35, 45} b. Data set: {7, 14, 21, 28, 35, 42} 2. Application of Central Limit Theorem \textbf{Instructions:} Explain how the Central Limit Theorem applies to the following scenario.
    \item \textbf{Scenario:} A researcher flips a fair coin 100 times and records the number of heads obtained. If this experiment is repeated many times, describe the expected distribution of the number of heads. 3. Interpretation with CLT \textbf{Instructions:} Answer the following question based on the Central Limit Theorem.
    \item \textbf{Question:} A sample of 30 measurements from a population with mean 100 and variance 25 is taken. If we repeatedly take samples of size 30, what distribution will the sample means approximate?
\end{enumerate}

%%===============================================================
\section{Standard Deviation}
\label{sec:standard-deviation}
%%===============================================================

\paragraph{Standard Deviation}



\textbf{A plot of normal distribution (or bell-shaped curve) where each band has a width of 1 standard deviation}



% [Image: bell]



\textbf{Cumulative probability of a normal distribution with expected value 0 and standard deviation 1}



% [Image: cumulative]



In statistics, the standard deviation is a measure of the amount of variation or dispersion of a set of values. A low standard deviation indicates that the values tend to be close to the mean (also called the expected value) of the set, while a high standard deviation indicates that the values are spread out over a wider range.



Standard deviation may be abbreviated SD, and is most commonly represented in mathematical texts and equations by the lower case Greek letter $\sigma$ (for the population standard deviation) or the Latin letter $s$ (for the sample standard deviation).



The standard deviation of a random variable, sample, statistical population, data set, or probability distribution is the square root of its variance. It is algebraically simpler, though in practice, less robust than the average absolute deviation. A useful property of the standard deviation is that unlike the variance, it is expressed in the same unit as the data.



The standard deviation of a population or sample and the standard error of a statistic (e.g., of the sample mean) are quite different, but related. The sample mean''''s standard error is the standard deviation of the set of means that would be found by drawing an infinite number of repeated samples from the population and computing a mean for each sample. The mean''''s standard error turns out to equal the population standard deviation divided by the square root of the sample size, and is estimated by using the sample standard deviation divided by the square root of the sample size. For example, a poll''''s standard error (what is reported as the margin of error of the poll) is the expected standard deviation of the estimated mean if the same poll were to be conducted multiple times. Thus, the standard error estimates the standard deviation of an estimate, which itself measures how much the estimate depends on the particular sample that was taken from the population.



In science, it is common to report both the standard deviation of the data (as a summary statistic) and the standard error of the estimate (as a measure of potential error in the findings). By convention, only effects more than two standard errors away from a null expectation are considered "statistically significant", a safeguard against spurious conclusions that are really due to random sampling error.



When only a sample of data from a population is available, the term standard deviation of the sample or sample standard deviation can refer to either the above-mentioned quantity as applied to those data, or to a modified quantity that is an unbiased estimate of the population standard deviation (the standard deviation of the entire population).





\paragraph{Basic Examples}



\paragraph{Population Standard Deviation of Grades of Eight Students}



Suppose that the entire population of interest is eight students in a particular class. For a finite set of numbers, the population standard deviation is found by taking the square root of the average of the squared deviations of the values subtracted from their average value. The marks of a class of eight students (that is, a statistical population) are the following eight values:



$$

2, 4, 4, 4, 5, 5, 7, 9

$$



These eight data points have the mean (average) of $5$:



$$

\mu = \frac{2 + 4 + 4 + 4 + 5 + 5 + 7 + 9}{8} = \frac{40}{8} = 5

$$



First, calculate the deviations of each data point from the mean, and square the result of each:



$$

(2 - 5)^2 = (-3)^2 = 9

$$



$$

(5 - 5)^2 = 0^2 = 0

$$



$$

(4 - 5)^2 = (-1)^2 = 1

$$



$$

(5 - 5)^2 = 0^2 = 0

$$



$$

(4 - 5)^2 = (-1)^2 = 1

$$



$$

(7 - 5)^2 = 2^2 = 4

$$



$$

(4 - 5)^2 = (-1)^2 = 1

$$



$$

(9 - 5)^2 = 4^2 = 16

$$



The variance is the mean of these values:



$$

\sigma^2 = \frac{9 + 0 + 1 + 0 + 1 + 4 + 1 + 16}{8} = \frac{32}{8} = 4

$$



and the population standard deviation is equal to the square root of the variance:



$$

\sigma = \sqrt{4} = 2

$$



This formula is valid only if the eight values with which we began form the complete population. If the values instead were a random sample drawn from some large parent population (for example, they were 8 students randomly and independently chosen from a class of 2 million), then one divides by $7$ (which is $n - 1$) instead of $8$ (which is $n$) in the denominator of the last formula, and the result is: $s = \sqrt{\frac{32}{7}} \approx 2.1$.  In that case, the result of the original formula would be called the sample standard deviation and denoted by $s$ instead of $\sigma$. Dividing by $n - 1$ rather than by $n$ gives an unbiased estimate of the variance of the larger parent population. This is known as Bessel''''s correction. Roughly, the reason for it is that the formula for the sample variance relies on computing differences of observations from the sample mean, and the sample mean itself was constructed to be as close as possible to the observations, so just dividing by $n$ would underestimate the variability.



\paragraph{Standard Deviation of Average Height for Adult Men}



If the population of interest is approximately normally distributed, the standard deviation provides information on the proportion of observations above or below certain values. For example, the average height for adult men in the United States is about $70$ inches ($177.8$ cm), with a standard deviation of around $3$ inches ($7.62$ cm). This means that most men (about $68\%$, assuming a normal distribution) have a height within $3$ inches ($7.62$ cm) of the mean ($67$–$73$ inches ($170.18$–$185.42$ cm)) - one standard deviation - and almost all men (about $95\%$) have a height within $6$ inches ($15.24$ cm) of the mean ($64$–$76$ inches ($162.56$–$193.04$ cm)) - two standard deviations. If the standard deviation were zero, then all men would be exactly $70$ inches ($177.8$ cm) tall. If the standard deviation were $20$ inches ($50.8$ cm), then men would have much more variable heights, with a typical range of about $50$–$90$ inches ($127$–$228.6$ cm). Three standard deviations account for $99.7\%$ of the sample population being studied, assuming the distribution is normal or bell-shaped (see the $68$–$95$–$99.7$ rule, or the empirical rule, for more information).





\paragraph{Definition of Population Values}



Let $\mu$ be the expected value (the average) of random variable $X$ with density $f(x)$:



$$

\mu \equiv \operatorname{E}[X] = \int_{-\infty}^{+\infty} x f(x) \, dx

$$



The standard deviation $\sigma$ of $X$ is defined as



$$

\sigma \equiv \sqrt{\operatorname{E}[(X - \mu)^2]} = \sqrt{\int_{-\infty}^{+\infty} (x - \mu)^2 f(x) \, dx}

$$



which can be shown to equal $\sqrt{\operatorname{E}[X^2] - (\operatorname{E}[X])^2}$



Using words, the standard deviation is the square root of the variance of $X$.



The standard deviation of a probability distribution is the same as that of a random variable having that distribution.



Not all random variables have a standard deviation. If the distribution has fat tails going out to infinity, the standard deviation might not exist, because the integral might not converge. The normal distribution has tails going out to infinity, but its mean and standard deviation do exist, because the tails diminish quickly enough. The Pareto distribution with parameter $\alpha \in (1, 2]$ has a mean, but not a standard deviation (loosely speaking, the standard deviation is infinite). The Cauchy distribution has neither a mean nor a standard deviation.



\paragraph{Discrete Random Variable}



In the case where $X$ takes random values from a finite data set $x_1, x_2, \ldots, x_N$, with each value having the same probability, the standard deviation is



$$

\sigma = \sqrt{\frac{1}{N}\[(x\_1 - \mu\)^2 + (x\_1 - \mu)^2 + (x\_2 - \mu)^2 + \cdots + (x\_N - \mu)^2]}

$$



where



$$

\mu = \frac{1}{N} (x_1 + \cdots + x\_N)

$$



or, by using summation notation,



$$

\sigma = \sqrt{\frac{1}{N} \sum\_{i=1}^{N} (x\_i - \mu)^2}

$$



where



$$

\mu = \frac{1}{N} \sum\_{i=1}^{N} x_i

$$



If, instead of having equal probabilities, the values have different probabilities, let $x\_1$ have probability $p\_1$, $x\_2$ have probability $p\_2$, $\ldots$, $x\_N$ have probability $p\_N$. In this case, the standard deviation will be



$$

\sigma = \sqrt{\sum\_{i=1}^{N} p\_i (x\_i - \mu)^2}

$$



where



$$

\mu = \sum\_{i=1}^{N} p_i x_i

$$



\paragraph{Continuous Random Variable}



The standard deviation of a continuous real-valued random variable $X$ with probability density function $p(x)$ is



$$

\sigma = \sqrt{\int\_{X} (x - \mu)^2 p(x) \, dx}

$$



where



$$

\mu = \int\_{X} x p(x) \, dx

$$



and where the integrals are definite integrals taken for $x$ ranging over the set of possible values of the random variable $X$.



In the case of a parametric family of distributions, the standard deviation can be expressed in terms of the parameters. For example, in the case of the log-normal distribution with parameters $\mu$ and $\sigma^2$, the standard deviation is



$$

\sqrt{\left(e^{\sigma^2} - 1\right) e^{2\mu + \sigma^2}}

$$



\paragraph{Estimation}



One can find the standard deviation of an entire population in cases (such as standardized testing) where every member of a population is sampled. In cases where that cannot be done, the standard deviation $\sigma$ is estimated by examining a random sample taken from the population and computing a statistic of the sample, which is used as an estimate of the population standard deviation. Such a statistic is called an estimator, and the estimator (or the value of the estimator, namely the estimate) is called a sample standard deviation, and is denoted by $s$ (possibly with modifiers).



Unlike in the case of estimating the population mean, for which the sample mean is a simple estimator with many desirable properties (unbiased, efficient, maximum likelihood), there is no single estimator for the standard deviation with all these properties, and unbiased estimation of standard deviation is a very technically involved problem. Most often, the standard deviation is estimated using the corrected sample standard deviation (using $N - 1$), defined below, and this is often referred to as the "sample standard deviation", without qualifiers. However, other estimators are better in other respects: the uncorrected estimator (using $N$) yields lower mean squared error, while using $N - 1.5$ (for the normal distribution) almost completely eliminates bias.



\paragraph{Uncorrected Sample Standard Deviation}



The formula for the population standard deviation (of a finite population) can be applied to the sample, using the size of the sample as the size of the population (though the actual population size from which the sample is drawn may be much larger). This estimator, denoted by $s\_N$, is known as the uncorrected sample standard deviation, or sometimes the standard deviation of the sample (considered as the entire population), and is defined as follows:



$$

s\_{N} = \sqrt{\frac{1}{N} \sum\_{i=1}^{N} (x\_{i} - \bar{x})^{2}}

$$



where $\{x\_{1}, x\_{2}, \ldots, x\_{N}\}$ are the observed values of the sample items, and $\bar{x}$ is the mean value of these observations, while the denominator $N$ stands for the size of the sample: this is the square root of the sample variance, which is the average of the squared deviations about the sample mean.



This is a consistent estimator (it converges in probability to the population value as the number of samples goes to infinity), and is the maximum-likelihood estimate when the population is normally distributed. However, this is a biased estimator, as the estimates are generally too low. The bias decreases as sample size grows, dropping off as $1/N$, and thus is most significant for small or moderate sample sizes; for $N > 75$ the bias is below 1\%. Thus for very large sample sizes, the uncorrected sample standard deviation is generally acceptable. This estimator also has a uniformly smaller mean squared error than the corrected sample standard deviation.



\paragraph{Corrected Sample Standard Deviation}



If the biased sample variance (the second central moment of the sample, which is a downward-biased estimate of the population variance) is used to compute an estimate of the population''''s standard deviation, the result is



$$

s\_{N} = \sqrt{\frac{1}{N} \sum\_{i=1}^{N} (x\_{i} - \bar{x})^{2}}

$$



Here taking the square root introduces further downward bias, by Jensen''''s inequality, due to the square root''''s being a concave function. The bias in the variance is easily corrected, but the bias from the square root is more difficult to correct, and depends on the distribution in question.



An unbiased estimator for the variance is given by applying Bessel''''s correction, using $N - 1$ instead of $N$ to yield the unbiased sample variance, denoted $s^2$:



$$

s^2 = \frac{1}{N - 1} \sum\_{i=1}^{N} (x\_{i} - \bar{x})^{2}

$$



This estimator is unbiased if the variance exists and the sample values are drawn independently with replacement. $N - 1$ corresponds to the number of degrees of freedom in the vector of deviations from the mean, $(x\_{1} - \bar{x}, \ldots, x\_{N} - \bar{x})$.



Taking square roots reintroduces bias (because the square root is a nonlinear function, which does not commute with the expectation), yielding the corrected sample standard deviation, denoted by $s$:



$$

s = \sqrt{\frac{1}{N - 1} \sum\_{i=1}^{N} (x\_{i} - \bar{x})^{2}}

$$



As explained above, while $s^2$ is an unbiased estimator for the population variance, $s$ is still a biased estimator for the population standard deviation, though markedly less biased than the uncorrected sample standard deviation. This estimator is commonly used and generally known simply as the "sample standard deviation". The bias may still be large for small samples ($N$ less than 10). As sample size increases, the amount of bias decreases. We obtain more information and the difference between $\frac{1}{N}$ and $\frac{1}{N - 1}$ becomes smaller.



\paragraph{Unbiased Sample Standard Deviation}



For unbiased estimation of standard deviation, there is no formula that works across all distributions, unlike for mean and variance. Instead, $s$ is used as a basis, and is scaled by a correction factor to produce an unbiased estimate. For the normal distribution, an unbiased estimator is given by $s / c\_4$, where the correction factor (which depends on $N$) is given in terms of the Gamma function, and equals:



$$

c\_{4}(N) = \sqrt{\frac{2}{N - 1}} \frac{\Gamma \left(\frac{N}{2}\right)}{\Gamma \left(\frac{N - 1}{2}\right)}

$$



This arises because the sampling distribution of the sample standard deviation follows a (scaled) chi distribution, and the correction factor is the mean of the chi distribution.



An approximation can be given by replacing $N - 1$ with $N - 1.5$, yielding:



$$

\hat{\sigma} = \sqrt{\frac{1}{N - 1.5} \sum\_{i=1}^{N} (x\_{i} - \bar{x})^{2}}

$$



The error in this approximation decays quadratically (as $1/N^2$), and it is suited for all but the smallest samples or highest precision: for $N = 3$ the bias is equal to 1.3\%, and for $N = 9$ the bias is already less than 0.1\%.



A more accurate approximation is to replace $N - 1.5$ above with $N - 1.5 + \frac{1}{8 (N - 1)}$.



For other distributions, the correct formula depends on the distribution, but a rule of thumb is to use the further refinement of the approximation:



$$

\hat{\sigma} = \sqrt{\frac{1}{N - 1.5 - \frac{1}{4} \gamma\_{2}} \sum\_{i=1}^{N} (x\_{i} - \bar{x})^{2}}

$$



where $\gamma\_2$ denotes the population excess kurtosis. The excess kurtosis may be either known beforehand for certain distributions, or estimated from the data.



\paragraph{Confidence Interval of a Sampled Standard Deviation}



The standard deviation we obtain by sampling a distribution is itself not absolutely accurate, both for mathematical reasons (explained here by the confidence interval) and for practical reasons of measurement (measurement error). The mathematical effect can be described by the confidence interval or CI.



To show how a larger sample will make the confidence interval narrower, consider the following examples: A small population of $N = 2$ has only 1 degree of freedom for estimating the standard deviation. The result is that a 95\% CI of the SD runs from $0.45 \times \text{SD}$ to $31.9 \times \text{SD}$; the factors here are as follows:



$$

\Pr \left(q\_{\frac{\alpha}{2}} < k \frac{s^{2}}{\sigma^{2}} < q\_{1 - \frac{\alpha}{2}} \right) = 1 - \alpha

$$



where $q\_p$ is the $p$-th quantile of the chi-square distribution with $k$ degrees of freedom, and $1 - \alpha$ is the confidence level. This is equivalent to the following:



$$

\Pr \left( k \frac{s^{2}}{q\_{1 - \frac{\alpha}{2}}} < \sigma^{2} < k \frac{s^{2}}{q\_{\frac{\alpha}{2}}} \right) = 1 - \alpha

$$



With $k = 1$, $q\_{0.025} = 0.000982$ and $q\_{0.975} = 5.024$. The reciprocals of the square roots of these two numbers give us the factors 0.45 and 31.9 given above.



A larger population of $N = 10$ has 9 degrees of freedom for estimating the standard deviation. The same computations as above give us in this case a 95\% CI running from $0.69 \times \text{SD}$ to $1.83 \times \text{SD}$. So even with a sample population of 10, the actual SD can still be almost a factor 2 higher than the sampled SD. For a sample population $N=100$, this is down to $0.88 \times \text{SD}$ to $1.16 \times \text{SD}$. To be more certain that the sampled SD is close to the actual SD we need to sample a large number of points.



These same formulae can be used to obtain confidence intervals on the variance of residuals from a least squares fit under standard normal theory, where $k$ is now the number of degrees of freedom for error.



\paragraph{Bounds on Standard Deviation}



For a set of $N > 4$ data spanning a range of values $R$, an upper bound on the standard deviation $s$ is given by $s = 0.6R$.



An estimate of the standard deviation for $N > 100$ data taken to be approximately normal follows from the heuristic that 95\% of the area under the normal curve lies roughly two standard deviations to either side of the mean, so that, with 95\% probability the total range of values $R$ represents four standard deviations so that $s \approx R/4$. This so-called range rule is useful in sample size estimation, as the range of possible values is easier to estimate than the standard deviation. Other divisors $K(N)$ of the range such that $s \approx R / K(N)$ are available for other values of $N$ and for non-normal distributions.



\paragraph{Identities and Mathematical Properties}



The standard deviation is invariant under changes in location, and scales directly with the scale of the random variable. Thus, for a constant $c$ and random variables $X$ and $Y$:



$$

\sigma(c) = 0 \\sigma(X + c) = \sigma(X) \\sigma(cX) = |c| \sigma(X)

$$



The standard deviation of the sum of two random variables can be related to their individual standard deviations and the covariance between them:



$$

\sigma(X + Y) = \sqrt{\operatorname{var}(X) + \operatorname{var}(Y) + 2 \operatorname{cov}(X, Y)}

$$



where $\operatorname{var} = \sigma^2$ and $\operatorname{cov}$ stand for variance and covariance, respectively.



The calculation of the sum of squared deviations can be related to moments calculated directly from the data. In the following formula, the letter $E$ is interpreted to mean expected value, i.e., mean.



$$

\sigma(X) = \sqrt{\operatorname{E} [(X - \operatorname{E}[X])^{2}]} = \sqrt{\operatorname{E}[X^{2}] - (\operatorname{E}[X])^{2} ]}

$$



The sample standard deviation can be computed as:



$$

s(X) = \sqrt{\frac{N}{N - 1} \operatorname{E} \left[(X - \operatorname{E}[X])^{2}\right]}

$$



For a finite population with equal probabilities at all points, we have:



$$

\sqrt{\frac{1}{N} \sum\_{i=1}^{N} (x\_{i} - \bar{x})^{2}} = \sqrt{\frac{1}{N} \left(\sum\_{i=1}^{N} x\_{i}^{2}\right) - \bar{x}^{2}} = \sqrt{\left(\frac{1}{N} \sum\_{i=1}^{N} x\_{i}^{2}\right) - \left(\frac{1}{N} \sum\_{i=1}^{N} x\_{i}\right)^{2}}

$$



which means that the standard deviation is equal to the square root of the difference between the average of the squares of the values and the square of the average value.



See computational formula for the variance for proof, and for an analogous result for the sample standard deviation.



\paragraph{Interpretation and Application}



\paragraph{Example}



Consider samples from two populations with the same mean but different standard deviations. 



\begin{itemize}
    \item \textbf{Red population}: Mean = 100, Standard Deviation (SD) = 10
    \item \textbf{Blue population}: Mean = 100, SD = 50
\end{itemize}

% [Image: sd graph]



A large standard deviation indicates that the data points can spread far from the mean, while a small standard deviation indicates that they are clustered closely around the mean.



For example, each of the following three populations {$0, 0, 14, 14$}, {$0, 6, 8, 14$}, and {$6, 6, 8, 8$} has a mean of 7. Their standard deviations are 7, 5, and 1, respectively. The third population has a much smaller standard deviation than the other two because its values are all close to 7. These standard deviations have the same units as the data points themselves. 



\begin{itemize}
    \item If the data set {$0, 6, 8, 14$} represents the ages of a population of four siblings in years, the standard deviation is 5 years.
    \item If the population {$1000, 1006, 1008, 1014$} represents the distances traveled by four athletes, measured in meters, it has a mean of 1007 meters and a standard deviation of 5 meters.
\end{itemize}

\paragraph{Measure of Uncertainty}



Standard deviation may serve as a measure of uncertainty. In physical science, for example, the reported standard deviation of a group of repeated measurements gives the precision of those measurements. When deciding whether measurements agree with a theoretical prediction, the standard deviation of those measurements is crucial: if the mean of the measurements is too far away from the prediction (with the distance measured in standard deviations), then the theory being tested probably needs to be revised. This is because they fall outside the range of values that could reasonably be expected to occur, if the prediction were correct and the standard deviation appropriately quantified.



While the standard deviation does measure how far typical values tend to be from the mean, other measures are available. An example is the mean absolute deviation, which might be considered a more direct measure of average distance compared to the root mean square distance inherent in the standard deviation.



\paragraph{Relationship between Standard Deviation and Mean}



The mean and the standard deviation of a set of data are descriptive statistics usually reported together. In a certain sense, the standard deviation is a "natural" measure of statistical dispersion if the center of the data is measured about the mean. This is because the standard deviation from the mean is smaller than from any other point. The precise statement is the following: suppose $x\_1, ..., x\_n$ are real numbers and define the function:



$$

\sigma (r) = \sqrt{\frac{1}{N-1} \sum\_{i=1}^{N} (x\_i - r)^2}

$$



Using calculus or by completing the square, it is possible to show that $\sigma(r)$ has a unique minimum at the mean:



$$

r = \bar{x}

$$



Variability can also be measured by the coefficient of variation, which is the ratio of the standard deviation to the mean. It is a dimensionless number.



\paragraph{Standard Deviation of the Mean}



Often, we want some information about the precision of the mean we obtained. We can obtain this by determining the standard deviation of the sampled mean. Assuming statistical independence of the values in the sample, the standard deviation of the mean is related to the standard deviation of the distribution by:



$$

\sigma\_{\text{mean}} = \frac{\sigma}{\sqrt{N}}

$$



where $N$ is the number of observations in the sample used to estimate the mean. This can easily be proven with (see basic properties of the variance):



$$

\operatorname{var}(X) \equiv \sigma\_X^2

$$



$$

\operatorname{var}(X\_1 + X\_2) \equiv \operatorname{var}(X\_1) + \operatorname{var}(X\_2)

$$



(Statistical independence is assumed.)



$$

\operatorname{var}(cX\_1) \equiv c^2 \operatorname{var}(X\_1)

$$



hence



$$

\operatorname{var}(\text{mean}) = \operatorname{var}\left(\frac{1}{N} \sum\_{i=1}^{N} X\_i\right) = \frac{1}{N^2} \operatorname{var}\left(\sum\_{i=1}^{N} X\_i\right) = \frac{1}{N^2} \sum\_{i=1}^{N} \operatorname{var}(X\_i) = \frac{N}{N^2} \operatorname{var}(X) = \frac{1}{N} \operatorname{var}(X)

$$



Resulting in:



$$

\sigma\_{\text{mean}} = \frac{\sigma}{\sqrt{N}}

$$



In order to estimate the standard deviation of the mean $\sigma\_{\text{mean}}$, it is necessary to know the standard deviation of the entire population $\sigma$ beforehand. However, in most applications, this parameter is unknown. For example, if a series of 10 measurements of a previously unknown quantity is performed in a laboratory, it is possible to calculate the resulting sample mean and sample standard deviation, but it is impossible to calculate the standard deviation of the mean.



\paragraph{Rapid Calculation Methods}



The following two formulas can represent a running (repeatedly updated) standard deviation. A set of two power sums $s\_1$ and $s\_2$ are computed over a set of $N$ values of $x$, denoted as $x\_1, ..., x\_N$:



$$

s\_j = \sum\_{k=1}^{N} x\_k^j

$$



Given the results of these running summations, the values $N$, $s\_1$, $s\_2$ can be used at any time to compute the current value of the running standard deviation:



$$

\sigma = \sqrt{\frac{N s\_2 - s\_1^2}{N}}

$$



Where $N$, as mentioned above, is the size of the set of values (or can also be regarded as $s\_0$).



Similarly for sample standard deviation,



$$

s = \sqrt{\frac{N s\_2 - s\_1^2}{N (N - 1)}}

$$



In a computer implementation, as the two $s\_j$ sums become large, we need to consider round-off error, arithmetic overflow, and arithmetic underflow. The method below calculates the running sums method with reduced rounding errors. This is a "one pass" algorithm for calculating variance of $n$ samples without the need to store prior data during the calculation. Applying this method to a time series will result in successive values of standard deviation corresponding to $n$ data points as $n$ grows larger with each new sample, rather than a constant-width sliding window calculation.



For $k = 1, ..., n$:



$$

A\_0 = 0

$$



$$

A\_k = A\_{k-1} + \frac{x\_k - A\_{k-1}}{k}

$$



where $A$ is the mean value.



$$

Q\_0 = 0

$$



$$

Q\_k = Q\_{k-1} + \frac{k-1}{k} (x\_k - A\_{k-1})^2 = Q\_{k-1} + (x\_k - A\_{k-1}) (x\_k - A\_k)

$$



Note: $Q\_1 = 0$ since $k-1 = 0$ or $x\_1 = A\_1$



Sample variance:



$$

s\_n^2 = \frac{Q\_n}{n-1}

$$



Population variance:



$$

\sigma\_n^2 = \frac{Q\_n}{n}

$$



\paragraph{Weighted Calculation}



When the values $x\_i$ are weighted with unequal weights $w\_i$, the power sums $s\_0$, $s\_1$, $s\_2$ are each computed as:



$$

s\_j = \sum\_{k=1}^{N} w\_k x\_k^j

$$



And the standard deviation equations remain unchanged. $s\_0$ is now the sum of the weights and not the number of samples $N$.



The incremental method with reduced rounding errors can also be applied, with some additional complexity.



A running sum of weights must be computed for each $k$ from 1 to $n$:



$$

W\_0 = 0

$$



$$

W\_k = W\_{k-1} + w\_k

$$



and places where $1/n$ is used above must be replaced by $w\_i/W\_n$:



$$

A\_0 = 0

$$



$$

A\_k = A\_{k-1} + \frac{w\_k}{W\_k} (x\_k - A\_{k-1})

$$



$$

Q\_0 = 0

$$



$$

Q\_k = Q\_{k-1} + \frac{w\_k W\_{k-1}}{W\_k} (x\_k - A\_{k-1})^2 = Q\_{k-1} + w\_k (x\_k - A\_{k-1}) (x\_k - A\_k)

$$



In the final division,



$$

\sigma\_n^2 = \frac{Q\_n}{W\_n}

$$



and



$$

s\_n^2 = \frac{Q\_n}{W\_n - 1}

$$



or



$$

s\_n^2 = \frac{n''}{n'' - 1} \sigma\_n^2

$$



where $n$ is the total number of elements, and $n''$ is the number of elements with non-zero weights.



The above formulas become equal to the simpler formulas given above if weights are taken as equal to one.



\paragraph{Resources}



\begin{itemize}
    \item \href{https://www.youtube.com/watch?v=IaTFpp-uzp0}{How To Calculate The Standard Deviation, The Organic Chemistry Tutor}
    \item \href{https://www.youtube.com/watch?v=Uk98hiMQgN0}{Standard Deviation Formula, Statistics, Variance, Sample and Population Mean, The Organic Chemistry Tutor}
\end{itemize}

\paragraph{Licensing}



Content obtained and/or adapted from:



\href{https://en.wikipedia.org/wiki/Standard_deviation}{Standard deviation, Wikipedia} under a CC BY-SA license

\subsection{Learning Objectives}

\begin{enumerate}
    \item 1.\textbf{Understand and Calculate Population Standard Deviation}
    \item SLO: Students will be able to define and calculate the population standard deviation for a given data set.
    \item Example: Given the data set $ \{2, 4, 4, 4, 5, 5, 7, 9\} $, students will calculate the mean $ \mu $, variance $ \sigma^2 $, and standard deviation $ \sigma $. 2.\textbf{Interpret Standard Deviation in the Context of Normal Distribution}
    \item SLO: Students will be able to interpret the standard deviation in the context of a normal distribution and explain its significance.
    \item Example: Given the average height for adult men in the United States is about $ 70 $ inches with a standard deviation of $ 3 $ inches, students will explain what it means for a data point to be within one, two, or three standard deviations from the mean. 3.\textbf{Calculate and Differentiate Between Population and Sample Standard Deviation}
    \item SLO: Students will understand the difference between population and sample standard deviation and be able to calculate both.
    \item Example: Given a data set, students will calculate both the population standard deviation $ \sigma $ and the sample standard deviation $ s $, applying Bessel's correction for the sample standard deviation.
\end{enumerate}

\subsection{Assessment Instrument}

\begin{enumerate}
    \item Question 1: Calculate the population standard deviation for the following data set: $ \{10, 12, 23, 23, 16, 23, 21, 16\} $. a. Calculate the mean $ \mu $. b. Calculate the variance $ \sigma^2 $. c. Calculate the standard deviation $ \sigma $. Question 2: Explain what it means for a data point to be within one, two, or three standard deviations from the mean in the context of a normal distribution. Question 3: Given the data set $ \{15, 22, 19, 30, 18, 25\} $, calculate both the population standard deviation $ \sigma $ and the sample standard deviation $ s $.
\end{enumerate}

%%===============================================================
\section{Probability}
\label{sec:probability}
%%===============================================================

\paragraph{Probability}



The probabilities of rolling several numbers using two dice:



% [Image: dice]



Probability is the branch of mathematics concerning numerical descriptions of how likely an event is to occur, or how likely it is that a proposition is true. The probability of an event is a number between $0$ and $1$, where, roughly speaking, $0$ indicates impossibility of the event and $1$ indicates certainty. The higher the probability of an event, the more likely it is that the event will occur. A simple example is the tossing of a fair (unbiased) coin. Since the coin is fair, the two outcomes ("heads" and "tails") are both equally probable; the probability of "heads" equals the probability of "tails"; and since no other outcomes are possible, the probability of either "heads" or "tails" is $\frac{1}{2}$ (which could also be written as $0.5$ or $50\%$).



These concepts have been given an axiomatic mathematical formalization in probability theory, which is used widely in areas of study such as statistics, mathematics, science, finance, gambling, artificial intelligence, machine learning, computer science, game theory, and philosophy to, for example, draw inferences about the expected frequency of events. Probability theory is also used to describe the underlying mechanics and regularities of complex systems.





\paragraph{Terminology of the probability theory}



\textbf{Experiment}: An operation which can produce some well-defined outcomes, is called an Experiment.



\textit{Example}: When we toss a coin, we know that either head or tail shows up. So, the operation of tossing a coin may be said to have two well-defined outcomes, namely, (a) heads showing up; and (b) tails showing up.



\textbf{Random Experiment}: When we roll a die we are well aware of the fact that any of the numerals $1, 2, 3, 4, 5$, or $6$ may appear on the upper face but we cannot say that which exact number will show up.



Such an experiment in which all possible outcomes are known and the exact outcome cannot be predicted in advance, is called a Random Experiment.



\textbf{Sample Space}: All the possible outcomes of an experiment as a whole form the Sample Space.



\textit{Example}: When we roll a die we can get any outcome from $1$ to $6$. All the possible numbers which can appear on the upper face form the Sample Space (denoted by $S$). Hence, the Sample Space of a dice roll is $S=\{1,2,3,4,5,6$}.



\textbf{Outcome}: Any possible result out of the Sample Space $S$ for a Random Experiment is called an Outcome.



\textit{Example}: When we roll a die, we might obtain $3$ or when we toss a coin, we might obtain heads.



\textbf{Event}: Any subset of the Sample Space $S$ is called an Event (denoted by $E$). When an outcome which belongs to the subset $E$ takes place, it is said that an Event has occurred. Whereas, when an outcome which does not belong to subset $E$ takes place, the Event has not occurred.



\textit{Example}: Consider the experiment of throwing a die. Over here the Sample Space $S=${$1,2,3,4,5,6$}. Let $E$ denote the event of ''a number appearing less than $4$.'' Thus the Event $E=${$1,2,3$}. If the number $1$ appears, we say that Event $E$ has occurred. Similarly, if the outcomes are $2$ or $3$, we can say Event $E$ has occurred since these outcomes belong to subset $E$.



\textbf{Trial}: By a trial, we mean performing a random experiment.



\textit{Example}: (i) Tossing a fair coin, (ii) rolling an unbiased die.



\paragraph{Theory}



Like other theories, the theory of probability is a representation of its concepts in formal terms—that is, in terms that can be considered separately from their meaning. These formal terms are manipulated by the rules of mathematics and logic, and any results are interpreted or translated back into the problem domain.



There have been at least two successful attempts to formalize probability, namely the Kolmogorov formulation and the Cox formulation. In Kolmogorov''s formulation (see also probability space), sets are interpreted as events and probability as a measure on a class of sets. In Cox''s theorem, probability is taken as a primitive (i.e., not further analyzed), and the emphasis is on constructing a consistent assignment of probability values to propositions. In both cases, the laws of probability are the same, except for technical details.



There are other methods for quantifying uncertainty, such as the Dempster–Shafer theory or possibility theory, but those are essentially different and not compatible with the usually-understood laws of probability.



\paragraph{Applications}



Probability theory is applied in everyday life in risk assessment and modeling. The insurance industry and markets use actuarial science to determine pricing and make trading decisions. Governments apply probabilistic methods in environmental regulation, entitlement analysis, and financial regulation.



An example of the use of probability theory in equity trading is the effect of the perceived probability of any widespread Middle East conflict on oil prices, which have ripple effects in the economy as a whole. An assessment by a commodity trader that a war is more likely can send that commodity''s prices up or down, and signals other traders of that opinion. Accordingly, the probabilities are neither assessed independently nor necessarily rationally. The theory of behavioral finance emerged to describe the effect of such groupthink on pricing, on policy, and on peace and conflict.



In addition to financial assessment, probability can be used to analyze trends in biology (e.g., disease spread) as well as ecology (e.g., biological Punnett squares). As with finance, risk assessment can be used as a statistical tool to calculate the likelihood of undesirable events occurring, and can assist with implementing protocols to avoid encountering such circumstances. Probability is used to design games of chance so that casinos can make a guaranteed profit, yet provide payouts to players that are frequent enough to encourage continued play.



Another significant application of probability theory in everyday life is reliability. Many consumer products, such as automobiles and consumer electronics, use reliability theory in product design to reduce the probability of failure. Failure probability may influence a manufacturer''s decisions on a product''s warranty.



The cache language model and other statistical language models that are used in natural language processing are also examples of applications of probability theory.



\paragraph{Mathematical treatment}



Consider an experiment that can produce a number of results. The collection of all possible results is called the sample space of the experiment, sometimes denoted as $\Omega$. The power set of the sample space is formed by considering all different collections of possible results. For example, rolling a die can produce six possible results. One collection of possible results gives an odd number on the die. Thus, the subset {$1,3,5$} is an element of the power set of the sample space of dice rolls. These collections are called "events". In this case, {$1,3,5$} is the event that the die falls on some odd number. If the results that actually occur fall in a given event, the event is said to have occurred.



A probability is a way of assigning every event a value between zero and one, with the requirement that the event made up of all possible results (in our example, the event {$1,2,3,4,5,6$}) is assigned a value of one. To qualify as a probability, the assignment of values must satisfy the requirement that for any collection of mutually exclusive events (events with no common results, such as the events {$1,6$}, {$3$}, and {$2,4$}), the probability that at least one of the events will occur is given by the sum of the probabilities of all the individual events.



The probability of an event $A$ is written as $P(A)$, $p(A)$, or $\text{Pr}(A)$. This mathematical definition of probability can extend to infinite sample spaces, and even uncountable sample spaces, using the concept of a measure.



The opposite or complement of an event $A$ is the event [not $A$] (that is, the event of $A$ not occurring), often denoted as $A''$, $A^c$, $\overline{A}$, $A^{\complement}$, $\neg A$, or $\sim A$; its probability is given by $P(\text{not } A) = 1 - P(A)$. As an example, the chance of not rolling a six on a six-sided die is $1 - (\text{chance of rolling a six}) = 1 - \frac{1}{6} = \frac{5}{6}$.



If two events $A$ and $B$ occur on a single performance of an experiment, this is called the intersection or joint probability of $A$ and $B$, denoted as $P(A \cap B)$.



\paragraph{Independent events}



If two events, $A$ and $B$ are independent then the joint probability is



$$P(A \text{ and } B) = P(A \cap B) = P(A)P(B).$$



For example, if two coins are flipped, then the chance of both being heads is $\frac{1}{2} \times \frac{1}{2} = \frac{1}{4}$. In this case, $P(A) = \frac{1}{2}$, $P(B) = \frac{1}{2}$, and $P(A \cap B) = \frac{1}{4}$.



\paragraph{Mutually exclusive events}



If either event $A$ or event $B$ but never both occurs on a single performance of an experiment, then they are called mutually exclusive events. If two events are mutually exclusive then the probability of both occurring is denoted as



$$P(A \text{ and } B) = P(A \cap B) = 0.$$



For example, if a single card is drawn from a standard deck of cards, the chance of getting both a heart and a spade is zero, because a single card cannot be both a heart and a spade.



\paragraph{Not mutually exclusive events}



If events $A$ and $B$ are not mutually exclusive then



$$P(A \text{ or } B) = P(A \cup B) = P(A) + P(B) - P(A \cap B).$$



For example, if a single card is drawn from a standard deck of cards, the chance of getting either a heart or a face card (jack, queen, or king) or both is given by



$$P(\text{heart}) = \frac{13}{52},$$

$$P(\text{face card}) = \frac{12}{52},$$

$$P(\text{heart} \cap \text{face card}) = \frac{3}{52}.$$



So,



$$P(\text{heart} \cup \text{face card}) = P(\text{heart}) + P(\text{face card}) - P(\text{heart} \cap \text{face card}) = \frac{13}{52} + \frac{12}{52} - \frac{3}{52} = \frac{22}{52} = \frac{11}{26}.$$



\paragraph{Conditional Probability}



Conditional probability is the probability of some event $A$, given the occurrence of some other event $B$. Conditional probability is written $P(A \mid B)$, and is read "the probability of $A$, given $B$". It is defined by



$$

P(A \mid B) = \frac{P(A \cap B)}{P(B)}.

$$



If $P(B) = 0$, then $P(A \mid B)$ is formally undefined by this expression. In this case, $A$ and $B$ are independent, since $P(A \cap B) = P(A)P(B) = 0$. However, it is possible to define a conditional probability for some zero-probability events using a $\sigma$-algebra of such events (such as those arising from a continuous random variable).



For example, in a bag of 2 red balls and 2 blue balls (4 balls in total), the probability of taking a red ball is $\frac{1}{2}$; however, when taking a second ball, the probability of it being either a red ball or a blue ball depends on the ball previously taken. For example, if a red ball was taken, then the probability of picking a red ball again would be $\frac{1}{3}$, since only 1 red and 2 blue balls would have been remaining. And if a blue ball was taken previously, the probability of taking a red ball will be $\frac{2}{3}$.



\paragraph{Inverse Probability}



In probability theory and applications, Bayes'' rule relates the odds of event $A\_1$ to event $A\_2$, before (prior to) and after (posterior to) conditioning on another event $B$. The odds on $A\_1$ to event $A\_2$ is simply the ratio of the probabilities of the two events. When arbitrarily many events $A$ are of interest, not just two, the rule can be rephrased as posterior is proportional to prior times likelihood,



$$

P(A \mid B) \propto P(A)P(B \mid A)

$$



where the proportionality symbol means that the left-hand side is proportional to (i.e., equals a constant times) the right-hand side as $A$ varies, for fixed or given $B$ (Lee, 2012; Bertsch McGrayne, 2012). In this form, it goes back to Laplace (1774) and to Cournot (1843); see Fienberg (2005).



\paragraph{Summary of Probabilities}



$$

\begin{array}{|c|c|}

\hline

\text{Event} & \text{Probability} \\\\ \hline

A & P(A) \in [0, 1] \\\\hline

\text{not } A & P(A^{\complement}) = 1 - P(A) \\\\ \hline

A \text{ or } B & \begin{aligned} 

P(A \cup B) &= P(A) + P(B) - P(A \cap B) \\\P(A \cup B) &= P(A) + P(B) \quad \text{if } A \text{ and } B \text{ are mutually exclusive}

\end{aligned} \\\\ \hline

A \text{ and } B & \begin{aligned}

P(A \cap B) &= P(A \mid B)P(B) = P(B \mid A)P(A) \\\P(A \cap B) &= P(A)P(B) \quad \text{if } A \text{ and } B \text{ are independent}

\end{aligned} \\\\ \hline

A \text{ given } B & P(A \mid B) = \frac{P(A \cap B)}{P(B)} = \frac{P(B \mid A)P(A)}{P(B)} \\\\ \hline

\end{array}

$$





\paragraph{Resources}



\begin{itemize}
    \item \href{https://www.khanacademy.org/math/cc-seventh-grade-math/cc-7th-probability-statistics/cc-7th-basic-prob/v/basic-probability}{Introduction to Theoretical Probability, Khan Academy}
\end{itemize}

\paragraph{Licensing}



Content obtained and/or adapted from:



\begin{itemize}
    \item \href{https://en.wikipedia.org/wiki/Probability}{Probability, Wikipedia} under a CC BY-SA license
\end{itemize}

\subsection{Learning Objectives}

\begin{enumerate}
    \item \textbf{Identify and Describe}: Students will be able to identify and describe the sample space and events of a simple random experiment, such as rolling a die or tossing a coin. They should be able to list all possible outcomes and determine the event space.
    \item \textbf{Calculate Basic Probabilities}: Students will be able to calculate the probability of a single event occurring in a random experiment. This includes determining the probability of rolling a specific number on a die or obtaining a particular outcome in other simple experiments.
    \item \textbf{Apply Probability Rules}: Students will be able to apply basic probability rules, including calculating the probability of the union and intersection of events. They should be able to use formulas for mutually exclusive and non-mutually exclusive events and understand conditional probability.
\end{enumerate}

\subsection{Assessment Instrument}

\begin{enumerate}
    \item Multiple-Choice Questions 1.\textbf{Sample Space Identification}
    \item What is the sample space when rolling a single six-sided die?
    \item A) {1, 2, 3, 4, 5, 6}
    \item B) {1, 2, 3, 4, 5}
    \item C) {1, 2, 3}
    \item D) {1, 2, 3, 4, 5, 6, 7} 2.\textbf{Basic Probability Calculation}
    \item What is the probability of rolling a 4 on a fair six-sided die?
    \item A) $\frac{1}{6}$
    \item B) $\frac{1}{4}$
    \item C) $\frac{1}{2}$
    \item D) $\frac{1}{3}$ 3.\textbf{Union of Events}
    \item If event A is rolling an even number on a six-sided die and event B is rolling a number greater than 4, what is the probability of either A or B occurring?
    \item A) $\frac{5}{6}$
    \item B) $\frac{4}{6}$
    \item C) $\frac{3}{6}$
    \item D) $\frac{2}{6}$ Short Answer Questions 4.\textbf{Event Description}
    \item Define what is meant by an "event" in the context of probability and provide an example using the experiment of rolling a die. 5.\textbf{Conditional Probability Calculation}
    \item In a bag with 3 red balls and 2 blue balls, if one ball is randomly drawn and it is red, what is the probability that the next ball drawn will also be red? (Assume no replacement.) True/False Questions 6.\textbf{Probability Rules}
    \item The probability of rolling a number less than 7 on a fair six-sided die is 1. (True/False) 7.\textbf{Mutually Exclusive Events}
    \item The events "rolling an odd number" and "rolling a 4" on a fair six-sided die are mutually exclusive. (True/False)
\end{enumerate}

%%===============================================================
\section{Conditional Probability}
\label{sec:conditional-probability}
%%===============================================================

\paragraph{Conditional Probability}

In probability theory, conditional probability is a measure of the probability of an event occurring, given that another event (by assumption, presumption, assertion, or evidence) has already occurred. If the event of interest is $A$ and the event $B$ is known or assumed to have occurred, "the conditional probability of $A$ given $B$," or "the probability of $A$ under the condition $B$," is usually written as $P(A|B)$ or occasionally $P\_{B}(A)$. This can also be understood as the fraction of probability $B$ that intersects with $A$:
$$
P(A \mid B) = \frac{P(A \cap B)}{P(B)}
$$

For example, the probability that any given person has a cough on any given day may be only 5\%. But if we know or assume that the person is sick, then they are much more likely to be coughing. For example, the conditional probability that someone unwell is coughing might be 75\%, in which case we would have that $P(\text{Cough}) = 5\%$ and $P(\text{Cough} \mid \text{Sick}) = 75\%$. Although, there doesn''t have to be a relationship or dependence between $A$ and $B$, and they don''t have to occur simultaneously.

$P(A|B)$ may or may not be equal to $P(A)$ (the unconditional probability of $A$). If $P(A|B) = P(A)$, then events $A$ and $B$ are said to be independent: in such a case, knowledge about either event does not alter the likelihood of each other. $P(A|B)$ (the conditional probability of $A$ given $B$) typically differs from $P(B|A)$. For example, if a person has dengue fever, they might have a 90\% chance of testing positive for the disease. In this case, what is being measured is that if event $B$ (having dengue) has occurred, the probability of $A$ (testing positive) given that $B$ occurred is 90\%: $P(A|B) = 90\%$. Alternatively, if a person tests positive for dengue fever, they may have only a 15\% chance of actually having this rare disease due to high false positive rates. In this case, the probability of the event $B$ (having dengue) given that the event $A$ (testing positive) has occurred is 15\%: $P(B|A) = 15\%$. It should be apparent now that falsely equating the two probabilities can lead to various errors of reasoning, which is commonly seen through base rate fallacies.

While conditional probabilities can provide extremely useful information, limited information is often supplied or at hand. Therefore, it can be useful to reverse or convert a conditional probability using Bayes'' theorem:
$$
P(A \mid B) = \frac{P(B \mid A) \cdot P(A)}{P(B)}
$$
Another option is to display conditional probabilities in a conditional probability table to illuminate the relationship between events.

\paragraph{Definition}

Illustration of conditional probabilities with an Euler diagram. The unconditional probability $P(A) = 0.30 + 0.10 + 0.12 = 0.52$. However, the conditional probability $P(A \mid B\_1) = 1$, $P(A \mid B\_2) = \frac{0.12}{0.12 + 0.04} = 0.75$, and $P(A \mid B\_3) = 0$.

% [Image: Euler]

On a tree diagram, branch probabilities are conditional on the event associated with the parent node. (Here, the overbars indicate that the event does not occur.)

% [Image: tree]

Venn Pie Chart describing conditional probabilities

% [Image: Venn]

\paragraph{Conditioning on an Event}

\paragraph{Kolmogorov Definition}

Given two events $A$ and $B$ from the sigma-field of a probability space, with the unconditional probability of $B$ being greater than zero (i.e., $P(B)>0$), the conditional probability of $A$ given $B$ is defined to be the quotient of the probability of the joint of events $A$ and $B$, and the probability of $B$:
$$
P(A \mid B) = \frac{P(A \cap B)}{P(B)}
$$
where $P(A \cap B)$ is the probability that both events $A$ and $B$ occur. This may be visualized as restricting the sample space to situations in which $B$ occurs. The logic behind this equation is that if the possible outcomes for $A$ and $B$ are restricted to those in which $B$ occurs, this set serves as the new sample space.

Note that the above equation is a definition—not a theoretical result. We just denote the quantity $\frac{P(A \cap B)}{P(B)}$ as $P(A \mid B)$, and call it the conditional probability of $A$ given $B$.

\paragraph{As an Axiom of Probability}

Some authors, such as de Finetti, prefer to introduce conditional probability as an axiom of probability:
$$
P(A \cap B) = P(A \mid B) \cdot P(B)
$$
Although mathematically equivalent, this may be preferred philosophically; under major probability interpretations, such as the subjective theory, conditional probability is considered a primitive entity. Moreover, this "multiplication rule" can be practically useful in computing the probability of $A \cap B$ and introduces a symmetry with the summation axiom for mutually exclusive events:
$$
P(A \cup B) = P(A) + P(B) - P(A \cap B)
$$

\paragraph{As the Probability of a Conditional Event}

Conditional probability can be defined as the probability of a conditional event $A\_B$. The Goodman–Nguyen–Van Fraassen conditional event can be defined as:
$$
A\_B = \bigcup\_{i \geq 1} \left( \bigcap\_{j < i} \overline{B}\_j, A\_i B\_i \right)
$$
It can be shown that:
$$
P(A\_B) = \frac{P(A \cap B)}{P(B)}
$$
which meets the Kolmogorov definition of conditional probability.

\paragraph{Conditioning on an Event of Probability Zero}

If $P(B) = 0$, then according to the definition, $P(A \mid B)$ is undefined.

The case of greatest interest is that of a random variable $Y$, conditioned on a continuous random variable $X$ resulting in a particular outcome $x$. The event $B = \{X = x\}$ has probability zero and, as such, cannot be conditioned on.

Instead of conditioning on $X$ being exactly $x$, we could condition on it being closer than distance $\epsilon$ away from $x$. The event $B = \{x - \epsilon < X < x + \epsilon\}$ will generally have nonzero probability and hence, can be conditioned on. We can then take the limit:
$$
\lim\_{\epsilon \to 0} P(A \mid x - \epsilon < X < x + \epsilon)
$$
For example, if two continuous random variables $X$ and $Y$ have a joint density $f\_{X,Y}(x, y)$, then by L''Hôpital''s rule:
$$
\lim\_{\epsilon \to 0} P(Y \in U \mid x\_0 - \epsilon < X < x\_0 + \epsilon) = \lim\_{\epsilon \to 0} \frac{\int\_{x\_0 - \epsilon}^{x\_0 + \epsilon} \int\_{U} f\_{X,Y}(x, y) \, dy \, dx}{\int\_{x\_0 - \epsilon}^{x\_0 + \epsilon} \int\_{\mathbb{R}} f\_{X,Y}(x, y) \, dy \, dx} = \frac{\int\_{U} f\_{X,Y}(x\_0, y) \, dy}{\int\_{\mathbb{R}} f\_{X,Y}(x\_0, y) \, dy}
$$
The resulting limit is the conditional probability distribution of $Y$ given $X$ and exists when the denominator, the probability density $f\_X(x\_0)$, is strictly positive.

It is tempting to define the undefined probability $P(A \mid X = x)$ using this limit, but this cannot be done in a consistent manner. In particular, it is possible to find random variables $X$ and $W$ and values $x$, $w$ such that the events {$X = x$} and {$W = w$} are identical but the resulting limits are not:

$$
\lim\_{\epsilon \to 0} P(A \mid x - \epsilon \leq X \leq x + \epsilon) \neq \lim\_{\epsilon \to 0} P(A \mid w - \epsilon \leq W \leq w + \epsilon).
$$

\paragraph{Conditioning on a Discrete Random Variable}

Let $X$ be a discrete random variable and its possible outcomes denoted $V$. For example, if $X$ represents the value of a rolled die, then $V$ is the set {$1, 2, 3, 4, 5, 6$}.

Let us assume for the sake of presentation that $X$ is a discrete random variable, so that each value in $V$ has a nonzero probability.

For a value $x$ in $V$ and an event $A$, the conditional probability is given by
$$
P(A \mid X = x).
$$
Writing
$$
c(x, A) = P(A \mid X = x)
$$
for short, we see that it is a function of two variables, $x$ and $A$.

For a fixed $A$, we can form the random variable
$$
Y = c(X, A).
$$
It represents an outcome of
$$
P(A \mid X = x)
$$
whenever a value $x$ of $X$ is observed.

The conditional probability of $A$ given $X$ can thus be treated as a random variable $Y$ with outcomes in the interval
$$
[0, 1].
$$
From the law of total probability, its expected value is equal to the unconditional probability of $A$.

\paragraph{Partial Conditional Probability}

The partial conditional probability
$$
P(A \mid B\_{1} \equiv b\_{1}, \ldots, B\_{m} \equiv b\_{m})
$$
is about the probability of event $A$ given that each of the condition events $B\_i$ has occurred to a degree $b\_i$ (degree of belief, degree of experience) that might be different from 100\%. Frequentistically, partial conditional probability makes sense if the conditions are tested in experiment repetitions of appropriate length $n$. Such $n$-bounded partial conditional probability can be defined as the conditionally expected average occurrence of event $A$ in testbeds of length $n$ that adhere to all of the probability specifications $B\_i \equiv b\_i$, i.e.:
$$
P^n(A \mid B\_{1} \equiv b\_{1}, \ldots, B\_{m} \equiv b\_{m}) = \operatorname{E}(\overline{A}^n \mid \overline{B}\_{1}^n = b\_{1}, \ldots, \overline{B}\_{m}^n = b\_{m}).
$$
Based on that, partial conditional probability can be defined as
$$
P(A \mid B\_{1} \equiv b\_{1}, \ldots, B\_{m} \equiv b\_{m}) = \lim\_{n \to \infty} P^n(A \mid B\_{1} \equiv b\_{1}, \ldots, B\_{m} \equiv b\_{m}),
$$
where $b\_i n \in \mathbb{N}$.

Jeffrey conditionalization is a special case of partial conditional probability, in which the condition events must form a partition:
$$
P(A \mid B\_{1} \equiv b\_{1}, \ldots, B\_{m} \equiv b\_{m}) = \sum\_{i=1}^{m} b\_{i} P(A \mid B\_{i}).
$$

\paragraph{Example}

Suppose that somebody secretly rolls two fair six-sided dice, and we wish to compute the probability that the face-up value of the first one is 2, given the information that their sum is no greater than 5.

Let $D\_1$ be the value rolled on dice 1.  
Let $D\_2$ be the value rolled on dice 2.

\textbf{Probability that $D\_1 = 2$}

Table 1 shows the sample space of 36 combinations of rolled values of the two dice, each of which occurs with probability $\frac{1}{36}$, with the numbers displayed in the lower right quadrant being $D\_1 + D\_2$.  The numbers in red are considered a "successful outcome."

$D\_1 = 2$ in exactly 6 of the 36 outcomes; thus $P(D\_1 = 2) = \frac{6}{36} = \frac{1}{6}$.

$$
\textbf{Table 1}
\begin{array}{c|cccccc}
D _1 \backslash D _2 & 1 & 2 & 3 & 4 & 5 & 6 \\\\ \hline
1 & 2 & 3 & 4 & 5 & 6 & 7 \\\2 & \textcolor{red}{3} & \textcolor{red}{4} & \textcolor{red}{5} & \textcolor{red}{6} & \textcolor{red}{7} & \textcolor{red}{8} \\\3 & 4 & 5 & 6 & 7 & 8 & 9 \\\4 & 5 & 6 & 7 & 8 & 9 & 10 \\\5 & 6 & 7 & 8 & 9 & 10 & 11 \\\6 & 7 & 8 & 9 & 10 & 11 & 12 \\\\ 
\end{array}
$$

\textbf{Probability that $D\_1 + D\_2 \leq 5$}

Table 2 shows that $D\_1 + D\_2 \leq 5$ for exactly 10 of the 36 outcomes, thus $P(D\_1 + D\_2 \leq 5) = \frac{10}{36}$:

$$
\textbf{Table 2}
\begin{array}{c|cccccc}
D\_1 \backslash D\_2 & 1 & 2 & 3 & 4 & 5 & 6 \\\\ \hline
1 & \textcolor{red}{2} & \textcolor{red}{3} & \textcolor{red}{4} & \textcolor{red}{5} & 6 & 7 \\\2 & \textcolor{red}{3} & \textcolor{red}{4} & \textcolor{red}{5} & 6 & 7 & 8 \\\3 & \textcolor{red}{4} & \textcolor{red}{5} & 6 & 7 & 8 & 9 \\\4 & \textcolor{red}{5} & 6 & 7 & 8 & 9 & 10 \\\5 & 6 & 7 & 8 & 9 & 10 & 11 \\\6 & 7 & 8 & 9 & 10 & 11 & 12 \\\\ \end{array}
$$

\textbf{Probability that $D\_1 = 2$ given that $D\_1 + D\_2 \leq 5$}

Table 3 shows that for 3 of these 10 outcomes, $D\_1 = 2$.  The numbers in blue represent the numbers that are less than or equal to 5, while the red numbers are the solution set.

Thus, the conditional probability $P(D\_1 = 2 \mid D\_1 + D\_2 \leq 5) = \frac{3}{10} = 0.3$:

$$\textbf{Table 3}
\begin{array}{c|cccccc}
D\_1 \backslash D\_2 & 1 & 2 & 3 & 4 & 5 & 6 \\\\ \hline
1 & \textcolor{blue}{2} & \textcolor{blue}{3} & \textcolor{blue}{4} & \textcolor{blue}{5} & 6 & 7 \\\2 & \textcolor{red}{3} & \textcolor{red}{4} & \textcolor{red}{5} & 6 & 7 & 8 \\\3 & \textcolor{blue}{4} & \textcolor{blue}{5} & 6 & 7 & 8 & 9 \\\4 & \textcolor{blue}{5} & 6 & 7 & 8 & 9 & 10 \\\5 & 6 & 7 & 8 & 9 & 10 & 11 \\\6 & 7 & 8 & 9 & 10 & 11 & 12 
\\\\ \end{array}
$$

Here, in the earlier notation for the definition of conditional probability, the conditioning event $B$ is that $D\_1 + D\_2 \leq 5$, and the event $A$ is $D\_1 = 2$. We have $P(A \mid B) = \frac{P(A \cap B)}{P(B)} = \frac{\frac{3}{36}}{\frac{10}{36}} = \frac{3}{10},$ as seen in the table.

\paragraph{Use in Inference}

In statistical inference, the conditional probability is an update of the probability of an event based on new information. The new information can be incorporated as follows:

1. Let $A$, the event of interest, be in the sample space, say $(X,P)$.
2. The occurrence of the event $A$ knowing that event $B$ has or will have occurred, means the occurrence of $A$ as it is restricted to $B$, i.e. $A \cap B$.
3. Without the knowledge of the occurrence of $B$, the information about the occurrence of $A$ would simply be $P(A)$.
4. The probability of $A$ knowing that event $B$ has or will have occurred, will be the probability of $A \cap B$ relative to $P(B)$, the probability that $B$ has occurred.

This results in

$$
P(A \mid B) = \frac{P(A \cap B)}{P(B)}
$$

whenever $P(B) > 0$ and $0$ otherwise. This approach results in a probability measure that is consistent with the original probability measure and satisfies all the Kolmogorov axioms. This conditional probability measure also could have resulted by assuming that the relative magnitude of the probability of $A$ with respect to $X$ will be preserved with respect to $B$ (cf. a Formal Derivation below).

The wording "evidence" or "information" is generally used in the Bayesian interpretation of probability. The conditioning event is interpreted as evidence for the conditioned event. That is, $P(A)$ is the probability of $A$ before accounting for evidence $E$, and $P(A \mid E)$ is the probability of $A$ after having accounted for evidence $E$ or after having updated $P(A)$. This is consistent with the frequentist interpretation, which is the first definition given above.

\paragraph{Statistical Independence}

Events $A$ and $B$ are defined to be statistically independent if

$$
P(A \cap B) = P(A)P(B).
$$

If $P(B)$ is not zero, then this is equivalent to the statement that

$$
P(A \mid B) = P(A).
$$

Similarly, if $P(A)$ is not zero, then

$$
P(B \mid A) = P(B)
$$

is also equivalent. Although the derived forms may seem more intuitive, they are not the preferred definition as the conditional probabilities may be undefined, and the preferred definition is symmetrical in $A$ and $B$.

\paragraph{Independent Events vs. Mutually Exclusive Events}

The concepts of mutually independent events and mutually exclusive events are separate and distinct. The following table contrasts results for the two cases (provided that the probability of the conditioning event is not zero).

$$
\begin{array}{|c|c|c|} \hline
\textbf{If statistically independent} & \textbf{If mutually exclusive} \\\\ \hline
P(A \mid B) = P(A) & P(A \mid B) = 0 \\\\ \hline
P(B \mid A) = P(B) & P(B \mid A) = 0 \\\\ \hline
P(A \cap B) = P(A)P(B) & P(A \cap B) = 0 \\\\ \hline
\end{array}
$$

In fact, mutually exclusive events cannot be statistically independent (unless both of them are impossible), since knowing that one occurs gives information about the other (in particular, that the latter will certainly not occur).

\paragraph{Common Fallacies}

These fallacies should not be confused with Robert K. Shope''s 1978 "conditional fallacy," which deals with counterfactual examples that beg the question.

\paragraph{Assuming Conditional Probability is of Similar Size to its Inverse}

A geometric visualization of Bayes'' theorem. In the table, the values 2, 3, 6, and 9 give the relative weights of each corresponding condition and case. The figures denote the cells of the table involved in each metric, the probability being the fraction of each figure that is shaded. This shows that $P(A|B) P(B) = P(B|A) P(A)$, i.e.,

$$
P(A|B) = \frac{P(B|A) P(A)}{P(B)}
$$

Similar reasoning can be used to show that $P(\bar{A}|B) = \frac{P(\bar{B}|A) P(\bar{A})}{P(B)}$, etc. 

In general, it cannot be assumed that $P(A|B) \approx P(B|A)$. This can be an insidious error, even for those who are highly conversant with statistics. The relationship between $P(A|B)$ and $P(B|A)$ is given by Bayes'' theorem:

$$
\begin{aligned}
P(B|A) &= \frac{P(A|B) P(B)}{P(A)} \\Leftrightarrow \frac{P(B|A)}{P(A|B)} &= \frac{P(B)}{P(A)}
\end{aligned}
$$

That is, $P(A|B) \approx P(B|A)$ only if $\frac{P(B)}{P(A)} \approx 1$, or equivalently, $P(A) \approx P(B)$.

\paragraph{Assuming Marginal and Conditional Probabilities are of Similar Size}

In general, it cannot be assumed that $P(A) \approx P(A|B)$. These probabilities are linked through the law of total probability:

$$
P(A) = \sum\_{n} P(A \cap B\_n) = \sum\_{n} P(A|B\_n) P(B\_n)
$$

where the events $(B\_n)$ form a countable partition of $\Omega$.

This fallacy may arise through selection bias. For example, in the context of a medical claim, let $S\_C$ be the event that a sequela (chronic disease) $S$ occurs as a consequence of circumstance (acute condition) $C$. Let $H$ be the event that an individual seeks medical help. Suppose that in most cases, $C$ does not cause $S$ (so that $P(S\_C)$ is low). Suppose also that medical attention is only sought if $S$ has occurred due to $C$. From experience of patients, a doctor may therefore erroneously conclude that $P(S\_C)$ is high. The actual probability observed by the doctor is $P(S\_C|H)$.

\paragraph{Over- or Under-weighting Priors}

Not taking prior probability into account partially or completely is called base rate neglect. The reverse, insufficient adjustment from the prior probability is conservatism.

\paragraph{Formal Derivation}

Formally, $P(A | B)$ is defined as the probability of $A$ according to a new probability function on the sample space, such that outcomes not in $B$ have probability 0 and that it is consistent with all original probability measures.

Let $\Omega$ be a sample space with elementary events $\{\omega\}$, and let $P$ be the probability measure with respect to the $\sigma$-algebra of $\Omega$. Suppose we are told that the event $B \subseteq \Omega$ has occurred. A new probability distribution (denoted by the conditional notation) is to be assigned on $\{\omega\}$ to reflect this. All events that are not in $B$ will have null probability in the new distribution. For events in $B$, two conditions must be met: the probability of $B$ is one and the relative magnitudes of the probabilities must be preserved. The former is required by the axioms of probability, and the latter stems from the fact that the new probability measure has to be the analog of $P$ in which the probability of $B$ is one - and every event that is not in $B$, therefore, has a null probability. Hence, for some scale factor $\alpha$, the new distribution must satisfy:

$$
\omega \in B : P(\omega \mid B) = \alpha P(\omega)
$$

$$
\omega \notin B : P(\omega \mid B) = 0
$$

$$
\sum\_{\omega \in \Omega} P(\omega \mid B) = 1.
$$

Substituting 1 and 2 into 3 to select $\alpha$:

% [Image: pic]

So the new probability distribution is:

% [Image: pic]

Now for a general event $A$:

% [Image: pic]

\paragraph{Licensing}

Content obtained and/or adapted from:

\href{https://en.wikipedia.org/wiki/Conditional_probability}{Conditional probability, Wikipedia} under a CC BY-SA license

\subsection{Learning Objectives}

\begin{enumerate}
    \item \textbf{Identify and Calculate Conditional Probabilities}
    \item Students will be able to define conditional probability and calculate the conditional probability $P(A \mid B)$ using the formula $P(A \mid B) = \frac{P(A \cap B)}{P(B)}$.
    \item \textbf{Interpret Conditional Probabilities in Real-World Contexts}
    \item Students will be able to interpret and explain the meaning of conditional probabilities in real-world contexts, such as medical testing or other scenarios involving dependent events.
    \item \textbf{Apply Bayes'' Theorem}
    \item Students will be able to apply Bayes'' theorem to find the conditional probability of an event, given prior probabilities and the conditional probabilities of related events, using the formula $P(A \mid B) = \frac{P(B \mid A) \cdot P(A)}{P(B)}$.
\end{enumerate}

\subsection{Assessment Instrument}

\begin{enumerate}
    \item Instructions Answer the following questions based on the provided lesson on conditional probability. Show all your work and calculations where applicable. Questions
    \item \textbf{Definition and Formula} Given two events $A$ and $B$, if $P(A) = 0.4$, $P(B) = 0.5$, and $P(A \cap B) = 0.2$, find the conditional probability $P(A \mid B)$. Use the formula: $$ P(A \mid B) = \frac{P(A \cap B)}{P(B)} $$
    \item \textbf{Independence} Consider two events $A$ and $B$ such that $P(A \mid B) = P(A)$. Explain what this implies about the relationship between $A$ and $B$.
    \item \textbf{Bayes'' Theorem} If $P(A \mid B) = 0.8$, $P(B \mid A) = 0.6$, and $P(A) = 0.5$, find $P(B)$ using Bayes'' theorem. Recall that: $$ P(A \mid B) = \frac{P(B \mid A) \cdot P(A)}{P(B)} $$
    \item \textbf{Euler Diagram Analysis} Create a Euler diagram where:
    \item $P(A) = 0.52$
    \item $P(A \mid B\_1) = 1$
    \item $P(A \mid B\_2) = 0.75$
    \item $P(A \mid B\_3) = 0$ Determine the following:
    \item What is the probability of $A \cap B\_2$ if $P(B\_2) = 0.16$?
    \item What does $P(A \mid B\_1) = 1$ imply about $A$ and $B\_1$?
    \item \textbf{Tree Diagram Probability} Create a tree diagram represents probabilities for the following events:
    \item $P(B\_1) = 0.3$
    \item $P(B\_2) = 0.4$
    \item $P(B\_3) = 0.3$
    \item $P(A \mid B\_1) = 0.7$
    \item $P(A \mid B\_2) = 0.5$
    \item $P(A \mid B\_3) = 0.2$ Find $P(A)$ using the law of total probability: $$ P(A) = P(A \mid B\_1)P(B\_1) + P(A \mid B\_2)P(B\_2) + P(A \mid B\_3)P(B\_3) $$
    \item \textbf{Venn Diagram Interpretation} Create a Venn diagram and answer the following:
    \item If $P(A) = 0.3$, $P(B) = 0.4$, and $P(A \cap B) = 0.1$, find $P(A \cup B)$.
    \item Explain how the Venn diagram can be used to visualize $P(A \mid B)$.
    \item \textbf{Kolmogorov Definition} Suppose $P(A \cap B) = 0.18$ and $P(B) = 0.6$. Calculate $P(A \mid B)$ using the Kolmogorov definition of conditional probability: $$ P(A \mid B) = \frac{P(A \cap B)}{P(B)} $$
    \item \textbf{Continuous Random Variables} For two continuous random variables $X$ and $Y$ with joint density $f\_{X,Y}(x, y)$, calculate:
    \item $P(Y \in U \mid X = x\_0)$ given: $$ \lim\_{\epsilon \to 0} \frac{\int\_{x\_0 - \epsilon}^{x\_0 + \epsilon} \int\_U f\_{X,Y}(x, y) \, dy \, dx}{\int\_{x\_0 - \epsilon}^{x\_0 + \epsilon} \int\_{\mathbb{R}} f\_{X,Y}(x, y) \, dy \, dx} $$
    \item \textbf{Discrete Random Variables} If $X$ is a discrete random variable with possible outcomes $1, 2, 3$ and $P(X = 1) = 0.2$, $P(X = 2) = 0.5$, and $P(X = 3) = 0.3$, find the conditional probability $P(A \mid X = 2)$ where $P(A) = 0.4$ and $P(A \cap X = 2) = 0.3$.
    \item \textbf{Partial Conditional Probability} Given that $P(A \mid B\_1 \equiv b\_1, B\_2 \equiv b\_2) = 0.7$, where $b\_1$ and $b\_2$ are degrees of belief, and $P(B\_1) = 0.5$, $P(B\_2) = 0.4$, calculate the partial conditional probability $P(A \mid B\_1 \equiv b\_1, B\_2 \equiv b\_2)$ if the events are independent.
\end{enumerate}

%%===============================================================
\section{Conversions}
\label{sec:conversions}
%%===============================================================

\paragraph{Process Overview}



The process of conversion depends on the specific situation and the intended purpose. This may be governed by regulation, contract, technical specifications, or other published standards. Engineering judgment may include such factors as:



\begin{itemize}
    \item The precision and accuracy of measurement and the associated uncertainty of measurement.
    \item The statistical confidence interval or tolerance interval of the initial measurement.
    \item The number of significant figures of the measurement.
    \item The intended use of the measurement including the engineering tolerances.
    \item Historical definitions of the units and their derivatives used in old measurements; e.g., international foot vs. US survey foot.
\end{itemize}

Some conversions from one system of units to another need to be exact, without increasing or decreasing the precision of the first measurement. This is sometimes called soft conversion. It does not involve changing the physical configuration of the item being measured.



By contrast, a hard conversion or an adaptive conversion may not be exactly equivalent. It changes the measurement to convenient and workable numbers and units in the new system. It sometimes involves a slightly different configuration or size substitution of the item. Nominal values are sometimes allowed and used.



\paragraph{Conversion Factors}



A conversion factor is used to change the units of a measured quantity without changing its value. The unity bracket method of unit conversion consists of a fraction in which the denominator is equal to the numerator, but they are in different units. Because of the identity property of multiplication, the value of a quantity will not change as long as it is multiplied by one. Also, if the numerator and denominator of a fraction are equal to each other, then the fraction is equal to one. So as long as the numerator and denominator of the fraction are equivalent, they will not affect the value of the measured quantity.



The following example demonstrates how the unity bracket method is used to convert the rate $5 \text{ km/s}$ to $ \text{ m/s}$. The symbols $\text{km}$, $\text{m}$, and $\text{s}$ represent kilometer, meter, and second, respectively.



% [Image: pic1] % [Image: pic2] % [Image: pic3] % [Image: pic4]



Thus, it is found that $5 \text{ km/s}$ is equal to $5000 \text{ m/s}$.



\paragraph{Software Tools}



There are many conversion tools. They are found in the function libraries of applications such as spreadsheets, databases, calculators, and in macro packages and plugins for many other applications such as mathematical, scientific, and technical applications.



There are many standalone applications that offer thousands of various units with conversions. For example, the free software movement offers a command line utility GNU units for Linux and Windows.



\paragraph{Calculation Involving Non-SI Units}



In cases where non-SI units are used, the numerical calculation of a formula can be done by first working out the pre-factor, and then plugging in the numerical values of the given/known quantities.



For example, in the study of Bose–Einstein condensate, atomic mass $m$ is usually given in daltons, instead of kilograms, and chemical potential $\mu$ is often given in Boltzmann constant times nanokelvin. The condensate''s healing length is given by:



$$

\xi = \frac{\hbar}{\sqrt{2m\mu}}

$$



For a $^{23}\text{Na}$ condensate with a chemical potential of (Boltzmann constant times) $128 \text{ nK}$, the calculation of healing length (in micrometres) can be done in two steps:



1. Calculate the pre-factor



   Assume that $m = 1 \text{ Da}$, $\mu = k_B \cdot 1 \text{ nK}$, this gives



   $$

   \xi = \frac{\hbar}{\sqrt{2m\mu}} = 15.574 \text{ µm}

   $$



   which is our pre-factor.



2. Calculate the numbers



   Now, make use of the fact that $\xi \propto \frac{1}{\sqrt{m\mu}}$. With $m = 23 \text{ Da}$, $\mu = 128 \, k_B \cdot \text{nK}$,



   $$

   \xi = \frac{15.574}{\sqrt{23 \cdot 128}} \text{ µm} = 0.287 \text{ µm}

   $$



This method is especially useful for programming and/or making a worksheet, where input quantities are taking multiple different values; for example, with the pre-factor calculated above, it''s very easy to see that the healing length of $^{174}\text{Yb}$ with a chemical potential of $20.3 \text{ nK}$ is



$$

\xi = \frac{15.574}{\sqrt{174 \cdot 20.3}} \text{ µm} = 0.262 \text{ µm}

$$



\paragraph{Resources}



\begin{itemize}
    \item \href{https://www.khanacademy.org/math/geometry-home/geometry-volume-surface-area/geometry-volume-rect-prism/v/conversion-between-metric-units}{Conversion Between Metric Units, Khan Academy}
    \item \href{https://www.youtube.com/watch?v=HZ9weUkSdoY}{Understanding Conversion Factors, Tyler DeWitt}
    \item \href{https://www.youtube.com/watch?v=eK8gXP3pImU}{Converting Units With Conversion Factors - Metric System Review \& Dimensional Analysis, The Organic Chemistry Tutor}
\end{itemize}

\paragraph{Licensing}



Content obtained and/or adapted from:



\href{https://en.wikipedia.org/wiki/Conversion_of_units}{Conversion of units, Wikipedia} under a CC BY-SA license

\subsection{Learning Objectives}

\begin{enumerate}
    \item SLO 1 \textbf{Identify and describe} the factors influencing the process of conversion, including precision, accuracy, and historical definitions of units. SLO 2 \textbf{Apply} the unity bracket method to convert units using conversion factors. SLO 3 \textbf{Calculate} the healing length of a Bose-Einstein condensate given specific non-SI units and the formula provided.
\end{enumerate}

\subsection{Assessment Instrument}

\begin{enumerate}
    \item \textbf{Question 1:} Describe the role of precision and accuracy in the conversion process. How do these factors impact the measurement''s reliability? Provide examples of soft and hard conversions. \textbf{Question 2:} Use the unity bracket method to convert a rate of $7 \text{ miles per hour}$ to $ \text{ meters per second}$. Show the fraction setup, conversion factors, and the final result. \textbf{Question 3:} Given the formula for healing length: $$ \xi = \frac{\hbar}{\sqrt{2m\mu}} $$ and a $^{87}\text{Rb}$ condensate with a chemical potential of $150 \text{ nK}$ and an atomic mass of $85 \text{ Da}$, calculate $\xi$ in micrometers. Show your calculations and the final answer. \textbf{Question 4:} Compare and contrast soft and hard conversions. Provide an example where each type of conversion would be appropriate and explain why.
\end{enumerate}

%%===============================================================
\section{Index numbers}
\label{sec:index-numbers}
%%===============================================================

\paragraph{Index Numbers}



\paragraph{Ratio}



In mathematics, a ratio indicates how many times one number contains another. For example, if there are eight oranges and six lemons in a bowl of fruit, then the ratio of oranges to lemons is $8:6$ (that is, $8:6$, which is equivalent to the ratio $4:3$). Similarly, the ratio of lemons to oranges is $6:8$ (or $3:4$) and the ratio of oranges to the total amount of fruit is $8:14$ (or $4:7$).



The numbers in a ratio may be quantities of any kind, such as counts of people or objects, or such as measurements of lengths, weights, time, etc. In most contexts, both numbers are restricted to be positive.



A ratio may be specified either by giving both constituting numbers, written as "a to b" or "a:b", or by giving just the value of their quotient $\frac{a}{b}$. Equal quotients correspond to equal ratios.



Consequently, a ratio may be considered as an ordered pair of numbers, a fraction with the first number in the numerator and the second in the denominator, or as the value denoted by this fraction. Ratios of counts, given by (non-zero) natural numbers, are rational numbers, and may sometimes be natural numbers. When two quantities are measured with the same unit, as is often the case, their ratio is a dimensionless number. A quotient of two quantities that are measured with different units is called a rate.



\paragraph{Converting Between Decimals and Fractions}



To change a common fraction to a decimal, do a long division of the decimal representations of the numerator by the denominator (this is idiomatically also phrased as "divide the denominator into the numerator"), and round the answer to the desired accuracy. For example, to change $\frac{1}{4}$ to a decimal, divide $1.00$ by $4$ ("$4$ into $1.00$"), to obtain $0.25$. To change $\frac{1}{3}$ to a decimal, divide $1.000...$ by $3$ ("$3$ into $1.0000...$"), and stop when the desired accuracy is obtained, e.g., at $4$ decimals with $0.3333$. The fraction $\frac{1}{4}$ can be written exactly with two decimal digits, while the fraction $\frac{1}{3}$ cannot be written exactly as a decimal with a finite number of digits. To change a decimal to a fraction, write in the denominator a $1$ followed by as many zeroes as there are digits to the right of the decimal point, and write in the numerator all the digits of the original decimal, just omitting the decimal point. Thus $12.3456 = \frac{123456}{10000}$.



\paragraph{Converting Between Ratios and Percents}



If a mixture contains substances A, B, C, and D in the ratio $5:9:4:2$ then there are 5 parts of A for every 9 parts of B, 4 parts of C, and 2 parts of D. As $5+9+4+2=20$, the total mixture contains $\frac{5}{20}$ of A (5 parts out of 20), $\frac{9}{20}$ of B, $\frac{4}{20}$ of C, and $\frac{2}{20}$ of D. If we divide all numbers by the total and multiply by 100, we have converted to percentages: $25\%$ A, $45\%$ B, $20\%$ C, and $10\%$ D (equivalent to writing the ratio as $25:45:20:10$).



If the two or more ratio quantities encompass all of the quantities in a particular situation, it is said that "the whole" contains the sum of the parts: for example, a fruit basket containing two apples and three oranges and no other fruit is made up of two parts apples and three parts oranges. In this case, $\frac{2}{5}$, or $40\%$ of the whole is apples and $\frac{3}{5}$, or $60\%$ of the whole is oranges. This comparison of a specific quantity to "the whole" is called a proportion.



The percent value is computed by multiplying the numeric value of the ratio by 100. For example, to find $50$ apples as a percentage of $1250$ apples, one first computes the ratio $\frac{50}{1250}=0.04$, and then multiplies by 100 to obtain $4\%$. The percent value can also be found by multiplying first instead of later, so in this example, the $50$ would be multiplied by 100 to give $5000$, and this result would be divided by $1250$ to give $4\%$.



To calculate a percentage of a percentage, convert both percentages to fractions of $100$, or to decimals, and multiply them. For example, $50\%$ of $40\%$ is: $\frac{50}{100} \times \frac{40}{100} = 0.50 \times 0.40 = 0.20 = \frac{20}{100} = 20\%$.



\paragraph{Averages}



An average is simply a number that is representative of data. More particularly, it is a measure of central tendency. There are several types of average. Averages are useful for comparing data, especially when sets of different size are being compared. It acts as a representative figure of the whole set of data.



Perhaps the simplest and commonly used average is the arithmetic mean or more simply mean which is explained in the next section.



Other common types of average are the median, the mode, the geometric mean, and the harmonic mean, each of which may be the most appropriate one to use under different circumstances.



\paragraph{Mean}



The mean, or more precisely the arithmetic mean, is simply the arithmetic average of a group of numbers (or data set) and is shown using -bar symbol $\bar{x}$. So the mean of the variable $x$ is $\bar{x}$, pronounced "x-bar". It is calculated by adding up all of the values in a data set and dividing by the number of values in that data set:



$$

\bar{x} = \frac{\sum x}{n}

$$



For example, take the following set of data: {$1,2,3,4,5$}. The mean of this data would be:



$$

\bar{x} = \frac{\sum x}{n} = \frac{1 + 2 + 3 + 4 + 5}{5} = \frac{15}{5} = 3

$$



Here is a more complicated data set: {$10,14,86,2,68,99,1$}. The mean would be calculated like this:



$$

\bar{x} = \frac{\sum x}{n} = \frac{10 + 14 + 86 + 2 + 68 + 99 + 1}{7} = \frac{280}{7} = 40

$$



\paragraph{Median}



The median is the "middle value" in a set. That is, the median is the number in the center of a data set that has been ordered sequentially.



For example, let''s look at the data in our second data set from above: {$10,14,86,2,68,99,1$}. What is its median?



First, we sort our data set sequentially: {$1,2,10,14,68,85,99$}.



Next, we determine the total number of points in our data set (in this case, 7).



Finally, we determine the central position of our data set (in this case, the 4th position), and the number in the central position is our median - {$1,2,10,14,68,85,99$}, making $14$ our median.



Because our data set had an odd number of points, determining the central position was easy - it will have the same number of points before it as after it. But what if our data set has an even number of points?



Let''s take the same data set, but add a new number to it: {$1,2,10,14,68,85,99,100$}. What is the median of this set?



When you have an even number of points, you must determine the two central positions of the data set. (See side box for instructions.) So for a set of $8$ numbers, we get $(8 + 1) / 2 = 9 / 2 = 4 \frac{1}{2}$, which has $4$ and $5$ on either side.



Looking at our dataset, we see that the $4$th and $5$th numbers are $14$ and $68$. From there, we return to our trusty friend the mean to determine the median. $(14 + 68) / 2 = 82 / 2 = 41$. 



For the set {$2, 4, 6, 8$}, we count the numbers to determine if it''s odd or even. As we see it is even, so we can write: $M = (4 + 6) / 2 = 10 / 2 = 5$. 



$5$ is the median of the above sequential numbers.



\paragraph{Mode}



The mode is the most common or "most frequent" value in a data set. Example: the mode of the following data set {$1, 2, 5, 5, 6, 3$} is $5$ since it appears twice. This is the most common value of the data set. Data sets having one mode are said to be unimodal, with two are said to be bimodal and with more than two are said to be multimodal. An example of a unimodal dataset is {$1, 2, 3, 4, 4, 4, 5, 6, 7, 8, 8, 9$}. The mode for this data set is $4$. An example of a bimodal data set is {$1, 2, 2, 3, 3$}. This is because both $2$ and $3$ are modes. Please note: If all points in a data set occur with equal frequency, it is equally accurate to describe the data set as having many modes or no mode.



\paragraph{Midrange}



The midrange is the arithmetic mean strictly between the minimum and the maximum value in a data set.



\paragraph{Relationship of the Mean, Median, and Mode}



The relationship of the mean, median, and mode to each other can provide some information about the relative shape of the data distribution. If the mean, median, and mode are approximately equal to each other, the distribution can be assumed to be approximately symmetrical. If the mean > median > mode, the distribution will be skewed to the right. If the mean < median < mode, the distribution will be skewed to the left.



\paragraph{Licensing}



Content obtained and/or adapted from:



\begin{itemize}
    \item \href{https://en.wikipedia.org/wiki/Ratio}{Ratio, Wikipedia} under a CC BY-SA license
    \item \href{https://en.wikipedia.org/wiki/Fraction}{Fraction, Wikipedia} under a CC BY-SA license
    \item \href{https://en.wikibooks.org/wiki/Statistics/Summary/Averages/mean}{Statistics/Summary/Averages/Mean, Wikibooks} under a CC BY-SA license
\end{itemize}

\subsection{Learning Objectives}

\begin{enumerate}
    \item \textbf{SLO 1:} Convert a given ratio to its simplest form and express it as a percentage.
    \item \textbf{SLO 2:} Convert between fractions and decimals, and apply this knowledge to calculate the decimal representation of a given fraction.
    \item \textbf{SLO 3:} Calculate and interpret the mean, median, and mode of a given data set, and explain their significance in terms of data analysis.
\end{enumerate}

\subsection{Assessment Instrument}

\begin{enumerate}
    \item Question 1: Ratio and Percentage Conversion Convert the following ratio to its simplest form and express it as a percentage: $$12:15$$ Question 2: Fraction to Decimal Conversion Convert the following fraction to a decimal and round to three decimal places: $$\frac{7}{16}$$ Question 3: Mean, Median, and Mode Calculation Given the data set {4, 8, 6, 5, 7, 8, 5}:
    \item Calculate the mean.
    \item Determine the median.
    \item Identify the mode.
    \item Explain what each of these statistics tells you about the data set.
\end{enumerate}

%%===============================================================
\section{Weighted averages}
\label{sec:weighted-averages}
%%===============================================================

\paragraph{Weighted Averages}



The weighted arithmetic mean is similar to an ordinary arithmetic mean (the most common type of average), except that instead of each of the data points contributing equally to the final average, some data points contribute more than others. The notion of weighted mean plays a role in descriptive statistics and also occurs in a more general form in several other areas of mathematics.



If all the weights are equal, then the weighted mean is the same as the arithmetic mean. While weighted means generally behave in a similar fashion to arithmetic means, they do have a few counterintuitive properties, as captured for instance in Simpson's paradox.



\paragraph{Examples}



\paragraph{Basic Example}



Given two school classes, one with 20 students, and one with 30 students, the grades in each class on a test were:



\begin{itemize}
    \item Morning class: $62, 67, 71, 74, 76, 77, 78, 79, 79, 80, 80, 81, 81, 82, 83, 84, 86, 89, 93, 98$
    \item Afternoon class: $81, 82, 83, 84, 85, 86, 87, 87, 88, 88, 89, 89, 89, 90, 90, 90, 90, 91, 91, 91, 92, 92, 93, 93, 94, 95, 96, 97, 98, 99$
\end{itemize}

The mean for the morning class is $80$ and the mean of the afternoon class is $90$. The unweighted mean of the two means is $85$. However, this does not account for the difference in the number of students in each class (20 versus 30); hence the value of $85$ does not reflect the average student grade (independent of class). The average student grade can be obtained by averaging all the grades, without regard to classes (add all the grades up and divide by the total number of students):



$$

\bar{x} = \frac{4300}{50} = 86

$$



Or, this can be accomplished by weighting the class means by the number of students in each class. The larger class is given more "weight":



$$

\bar{x} = \frac{(20 \times 80) + (30 \times 90)}{20 + 30} = 86

$$



Thus, the weighted mean makes it possible to find the mean average student grade without knowing each student's score. Only the class means and the number of students in each class are needed.



\paragraph{Convex Combination Example}



Since only the relative weights are relevant, any weighted mean can be expressed using coefficients that sum to one. Such a linear combination is called a convex combination.



Using the previous example, we would get the following weights:



$$

\frac{20}{20 + 30} = 0.4

$$



$$

\frac{30}{20 + 30} = 0.6

$$



Then, apply the weights like this:



$$

\bar{x} = (0.4 \times 80) + (0.6 \times 90) = 86

$$



\paragraph{Mathematical Definition}



Formally, the weighted mean of a non-empty finite multiset of data $\{x\_1, x\_2, \dots, x\_n\}$, with corresponding non-negative weights $\{w\_1, w\_2, \dots, w\_n\}$ is



$$

\bar{x} = \frac{\sum\_{i=1}^{n} w\_i x\_i}{\sum\_{i=1}^{n} w\_i}

$$



which expands to:



$$

\bar{x} = \frac{w\_1 x\_1 + w\_2 x\_2 + \cdots + w\_n x\_n}{w\_1 + w\_2 + \cdots + w\_n}

$$



Therefore, data elements with a high weight contribute more to the weighted mean than do elements with a low weight. The weights cannot be negative. Some may be zero, but not all of them (since division by zero is not allowed).



The formulas are simplified when the weights are normalized such that they sum up to $1$, i.e.:



$$

\sum\_{i=1}^{n} w\_i' = 1

$$



For such normalized weights, the weighted mean is then:



$$

\bar{x} = \sum\_{i=1}^{n} w\_i' x\_i

$$



Note that one can always normalize the weights by making the following transformation on the original weights:



$$

w\_i' = \frac{w\_i}{\sum\_{j=1}^{n} w\_j}

$$



Using the normalized weight yields the same results as when using the original weights:



$$

\bar{x} = \sum\_{i=1}^{n} w\_i' x\_i = \sum\_{i=1}^{n} \frac{w\_i}{\sum\_{j=1}^{n} w\_j} x\_i = \frac{\sum\_{i=1}^{n} w\_i x\_i}{\sum\_{j=1}^{n} w\_j} = \frac{\sum\_{i=1}^{n} w\_i x\_i}{\sum\_{i=1}^{n} w\_i}

$$



The ordinary mean $\frac{1}{n} \sum\_{i=1}^{n} x\_i$ is a special case of the weighted mean where all data have equal weights.



If the data elements are independent and identically distributed random variables with variance $\sigma$, the standard error of the weighted mean, $\sigma\_{\bar{x}}$, can be shown via uncertainty propagation to be:



$$

\sigma\_{\bar{x}} = \sigma \sqrt{\sum\_{i=1}^{n} w\_i'^2}

$$



\paragraph{Resources}



\begin{itemize}
    \item \href{https://www.youtube.com/watch?v=LdrBNhWw9AM}{How To Find The Weighted Mean and Weighted Average In Statistics, The Organic Chemistry Tutor}
\end{itemize}

\paragraph{Licensing}



Content obtained and/or adapted from:



\begin{itemize}
    \item \href{https://en.wikipedia.org/wiki/Weighted_arithmetic_mean}{Weighted arithmetic mean, Wikipedia} under a CC BY-SA license
\end{itemize}

\subsection{Learning Objectives}

\begin{enumerate}
    \item Understand and explain the concept of weighted arithmetic mean and its differences from the ordinary arithmetic mean.
    \item Calculate the weighted mean for a given dataset with varying weights.
    \item Apply the concept of weighted mean to real-life scenarios, such as combining class averages with different student counts.
\end{enumerate}

\subsection{Assessment Instrument}

\begin{enumerate}
    \item Question 1 \textbf{Explain the difference between the weighted arithmetic mean and the ordinary arithmetic mean. Provide an example to illustrate your explanation.} Question 2 \textbf{Given the following datasets and weights, calculate the weighted mean:}
    \item Data: $4, 7, 10$
    \item Weights: $1, 2, 3$ Question 3 \textbf{A school has two classes with the following average scores and number of students:}
    \item Class A: Average score = 75, Students = 25
    \item Class B: Average score = 85, Students = 35 \textbf{Calculate the overall average score for the school using the weighted mean.}
\end{enumerate}

%%===============================================================
\section{Expected Value}
\label{sec:expected-value}
%%===============================================================

\paragraph{Expected Value}

Let $X$ be a (discrete) random variable with a finite number of finite outcomes $x\_1, x\_2, \ldots, x\_k$ occurring with probabilities $p\_1, p\_2, \ldots, p\_k,$ respectively. The expectation of $X$ is defined as

$$
\operatorname{E}[X] = \sum\_{i=1}^{k} x\_i p\_i = x\_1 p\_1 + x\_2 p\_2 + \cdots + x\_k p\_k.
$$

Since $p\_1 + p\_2 + \cdots + p\_k = 1,$ the expected value is the weighted sum of the $x\_i$ values, with the probabilities $p\_i$ as the weights.

If all outcomes $x\_i$ are equiprobable (that is, $p\_1 = p\_2 = \cdots = p\_k$), then the weighted average turns into the simple average. On the other hand, if the outcomes $x\_i$ are not equiprobable, then the simple average must be replaced with the weighted average, which takes into account the fact that some outcomes are more likely than others.

\paragraph{Examples}

Let $X$ represent the outcome of a roll of a fair six-sided dice. More specifically, $X$ will be the number of pips showing on the top face of the dice after the toss. The possible values for $X$ are 1, 2, 3, 4, 5, and 6, all of which are equally likely with a probability of $\frac{1}{6}$. The expectation of $X$ is

$$
\operatorname{E}[X] = 1 \cdot \frac{1}{6} + 2 \cdot \frac{1}{6} + 3 \cdot \frac{1}{6} + 4 \cdot \frac{1}{6} + 5 \cdot \frac{1}{6} + 6 \cdot \frac{1}{6} = 3.5.
$$

If one rolls the dice $n$ times and computes the average (arithmetic mean) of the results, then as $n$ grows, the average will almost surely converge to the expected value, a fact known as the strong law of large numbers.

The roulette game consists of a small ball and a wheel with 38 numbered pockets around the edge. As the wheel is spun, the ball bounces around randomly until it settles down in one of the pockets. Suppose random variable $X$ represents the (monetary) outcome of a `$`1 bet on a single number ("straight up" bet). If the bet wins (which happens with probability $\frac{1}{38}$ in American roulette), the payoff is `$`35; otherwise the player loses the bet. The expected profit from such a bet will be

$$
\operatorname{E}[\text{gain from \$1 bet}] = -\$1 \cdot \frac{37}{38} + \$35 \cdot \frac{1}{38} = -\frac{\$1}{19}.
$$

That is, the bet of `$`1 stands to lose $-\frac{\$1}{19}$, so its expected value is $-\frac{\$1}{19}$.

\paragraph{Resources}

\begin{itemize}
    \item \href{https://mathresearch.utsa.edu/wikiFiles/MAT1053/Expected_Value/MAT1053_M12.2Expected_Value.pdf}{Expected Value, Book Chapter}
    \item \href{https://mathresearch.utsa.edu/wikiFiles/MAT1053/Expected_Value/MAT1053_M12.2Expected_Value.pdf}{Guided Notes}
\end{itemize}

\paragraph{Licensing}

Content obtained and/or adapted from:

\begin{itemize}
    \item \href{https://en.wikipedia.org/wiki/Expected_value}{Expected Value, Wikipedia} under a CC BY-SA license
\end{itemize}

\subsection{Learning Objectives}

\begin{enumerate}
    \item Define and calculate the expected value of a discrete random variable.
    \item Differentiate between equiprobable and non-equiprobable outcomes when calculating expected value.
    \item Apply the concept of expected value to solve real-world problems, including games of chance.
\end{enumerate}

\subsection{Assessment Instrument}

\begin{enumerate}
    \item Question 1 Let $X$ be a discrete random variable with outcomes $x\_1, x\_2, \ldots, x\_k$ occurring with probabilities $p\_1, p\_2, \ldots, p\_k,$ respectively. Calculate the expected value of $X$ given the following data:
    \item $x\_1 = 2, p\_1 = 0.1$
    \item $x\_2 = 4, p\_2 = 0.2$
    \item $x\_3 = 6, p\_3 = 0.3$
    \item $x\_4 = 8, p\_4 = 0.4$ Question 2 A fair six-sided dice is rolled. Calculate the expected value of the number of pips showing on the top face after the roll. Show your calculations. Question 3 In a game of American roulette, a player bets $\$1$ on a single number ("straight up" bet). The probability of winning the bet is $\frac{1}{38}$, and the payoff is $\$35$ if the bet wins. Calculate the expected profit for the player from such a bet. Explain whether this is a favorable bet for the player based on the expected value.
\end{enumerate}

%%===============================================================
\section{Weighted moving average graphs}
\label{sec:weighted-moving-average-graphs}
%%===============================================================

\paragraph{Weighted Moving Average Graphs}



Smoothing of a noisy sine (blue curve) with a moving average (red curve). In statistics, a moving average (rolling average or running average) is a calculation to analyze data points by creating a series of averages of different subsets of the full data set. It is also called a moving mean (MM) or rolling mean and is a type of finite impulse response filter. Variations include: simple, cumulative, or weighted forms (described below).



% [Image: smoothing]



Given a series of numbers and a fixed subset size, the first element of the moving average is obtained by taking the average of the initial fixed subset of the number series. Then the subset is modified by "shifting forward"; that is, excluding the first number of the series and including the next value in the subset.



A moving average is commonly used with time series data to smooth out short-term fluctuations and highlight longer-term trends or cycles. The threshold between short-term and long-term depends on the application, and the parameters of the moving average will be set accordingly. For example, it is often used in technical analysis of financial data, like stock prices, returns or trading volumes. It is also used in economics to examine gross domestic product, employment or other macroeconomic time series. Mathematically, a moving average is a type of convolution and so it can be viewed as an example of a low-pass filter used in signal processing. When used with non-time series data, a moving average filters higher frequency components without any specific connection to time, although typically some kind of ordering is implied. Viewed simplistically it can be regarded as smoothing the data.





\paragraph{Simple Moving Average}



In financial applications, a simple moving average (SMA) is the unweighted mean of the previous $k$ data points. However, in science and engineering, the mean is normally taken from an equal number of data on either side of a central value. This ensures that variations in the mean are aligned with the variations in the data rather than being shifted in time. An example of a simple equally weighted running mean is the mean over the last $k$ entries of a data set containing $n$ entries. Let those data points be $p\_1, p\_2, \ldots, p\_n$. This could be closing prices of a stock. The mean over the last $k$ data points (days in this example) is denoted as $SMA\_k$ and calculated as:



$$

SMA\_k = \frac{p\_{n-k+1} + p\_{n-k+2} \cdots + p\_{n}}{k} = \frac{1}{k} \sum\_{i=n-k+1}^{n} p\_i

$$



When calculating the next mean $SMA\_{k,\text{next}}$ with the same sampling width $k$, the range from $n-k+2$ to $n+1$ is considered. A new value $p\_{n+1}$ comes into the sum and the oldest value $p\_{n-k+1}$ drops out. This simplifies the calculations by reusing the previous mean $SMA\_{k,\text{prev}}$.



$$

SMA\_{k,\text{next}} = \frac{1}{k} \sum\_{i=n-k+2}^{n+1} p\_i = \frac{1}{k} \left( \underbrace{p\_{n-k+2} + p\_{n-k+3} + \ldots + p\_{n} + p\_{n+1}}\_{\sum\_{i=n-k+2}^{n+1} p\_i} + \underbrace{p\_{n-k+1} - p\_{n-k+1}}\_{=0} \right) 

$$



$$

= \underbrace{\frac{1}{k} \left( p\_{n-k+1} + p\_{n-k+2} + \ldots + p\_{n} \right)}\_{=SMA\_{k,\text{prev}}} - \frac{p\_{n-k+1}}{k} + \frac{p\_{n+1}}{k} 

$$



$$

= SMA\_{k,\text{prev}} + \frac{1}{k} \left( p\_{n+1} - p\_{n-k+1} \right)

$$





This means that the moving average filter can be computed quite cheaply on real-time data with a FIFO / circular buffer and only 3 arithmetic steps.



During the initial filling of the FIFO / circular buffer, the sampling window is equal to the data set size thus $k=n$ and the average calculation is performed as a cumulative moving average.



The period selected ($k$) depends on the type of movement of interest, such as short, intermediate, or long-term. In financial terms, moving-average levels can be interpreted as support in a falling market or resistance in a rising market.



If the data used are not centered around the mean, a simple moving average lags behind the latest datum by half the sample width. An SMA can also be disproportionately influenced by old data dropping out or new data coming in. One characteristic of the SMA is that if the data has a periodic fluctuation, then applying an SMA of that period will eliminate that variation (the average always containing one complete cycle). But a perfectly regular cycle is rarely encountered.



For a number of applications, it is advantageous to avoid the shifting induced by using only "past" data. Hence a central moving average can be computed, using data equally spaced on either side of the point in the series where the mean is calculated. This requires using an odd number of points in the sample window.



A major drawback of the SMA is that it lets through a significant amount of the signal shorter than the window length. Worse, it actually inverts it. This can lead to unexpected artifacts, such as peaks in the smoothed result appearing where there were troughs in the data. It also leads to the result being less smooth than expected since some of the higher frequencies are not properly removed.



\paragraph{Cumulative Moving Average}



In a cumulative moving average (CMA), the data arrive in an ordered datum stream, and the user would like to get the average of all the data up until the current datum. For example, an investor may want the average price of all the stock transactions for a particular stock up until the current time. As each new transaction occurs, the average price at the time of the transaction can be calculated for all the transactions up to that point using the cumulative average, typically an equally weighted average of the sequence of $n$ values $x\_1, \ldots, x\_n$ up to the current time:



$$

CMA\_n = \frac{x\_1 + \cdots + x\_n}{n}

$$



The brute-force method to calculate this would be to store all the data and calculate the sum and divide by the number of points every time a new datum arrived. However, it is possible to simply update the cumulative average as a new value, $x\_{n+1}$, becomes available, using the formula



$$

CMA\_{n+1} = \frac{x\_{n+1} + n \cdot CMA\_n}{n+1}

$$



Thus, the current cumulative average for a new datum is equal to the previous cumulative average, times $n$, plus the latest datum, all divided by the number of points received so far, $n+1$. When all the data arrive $(n = N)$, the cumulative average will equal the final average. It is also possible to store a running total of the data as well as the number of points and dividing the total by the number of points to get the CMA each time a new datum arrives.



The derivation of the cumulative average formula is straightforward. Using



$$

x\_1 + \cdots + x\_n = n \cdot CMA\_n

$$



and similarly for $n+1$, it is seen that



$$

x\_{n+1} = (x\_1 + \cdots + x\_{n+1}) - (x\_1 + \cdots + x\_n)

$$



Solving this equation for $CMA\_{n+1}$ results in



$$

CMA\_{n+1} = \frac{x\_{n+1} + n \cdot CMA\_n}{n+1} = \frac{x\_{n+1} + (n+1-1) \cdot CMA\_n}{n+1} = \frac{(n+1) \cdot CMA\_n + x\_{n+1} - CMA\_n}{n+1} = CMA\_n + \frac{x\_{n+1} - CMA\_n}{n+1}

$$



\paragraph{Weighted Moving Average}



A weighted average is an average that has multiplying factors to give different weights to data at different positions in the sample window. Mathematically, the weighted moving average is the convolution of the data with a fixed weighting function. One application is removing pixelization from a digital graphical image.



In technical analysis of financial data, a weighted moving average (WMA) has the specific meaning of weights that decrease in arithmetical progression. In an $n$-day WMA, the latest day has weight $n$, the second latest $n-1$, etc., down to one.



$$

WMA\_M = \frac{n p\_M + (n-1) p\_{M-1} + \cdots + 2 p\_{(M-n)+2} + p\_{(M-n)+1}}{n + (n-1) + \cdots + 2 + 1}

$$



The denominator is a triangle number equal to $\frac{n(n+1)}{2}$. In the more general case, the denominator will always be the sum of the individual weights.



When calculating the WMA across successive values, the difference between the numerators of $WMA\_{M+1}$ and $WMA\_M$ is $np\_{M+1} - p\_M - \ldots - p\_{M-n+1}$. If we denote the sum $p\_M + \ldots + p\_{M-n+1}$ by $Total\_M$, then



$$

Total\_{M+1} = Total\_M + p\_{M+1} - p\_{M-n+1}

$$



$$

Numerator\_{M+1} = Numerator\_M + n p\_{M+1} - Total\_M

$$



$$

WMA\_{M+1} = \frac{Numerator\_{M+1}}{n + (n-1) + \cdots + 2 + 1}

$$



The graph below shows how the weights decrease, from the highest weight for the most recent data, down to zero. It can be compared to the weights in the exponential moving average, which follows.



% [Image: decrease]



\paragraph{Exponential Moving Average}





An exponential moving average (EMA), also known as an exponentially weighted moving average (EWMA), is a first-order infinite impulse response filter that applies weighting factors which decrease exponentially. The weighting for each older datum decreases exponentially, never reaching zero. The graph below shows an example of the weight decrease.



% [Image: weight]



The EMA for a series $ Y $ may be calculated recursively:



$$

S\_t = 

\begin{cases} 

Y\_1, & t = 1 \\\\alpha Y\_t + (1 - \alpha) \cdot S\_{t-1}, & t > 1 

\end{cases}

$$



Where:



\begin{itemize}
    \item The coefficient $ \alpha $ represents the degree of weighting decrease, a constant smoothing factor between 0 and 1. A higher $ \alpha $ discounts older observations faster.
    \item $ Y\_t $ is the value at a time period $ t $.
    \item $ S\_t $ is the value of the EMA at any time period $ t $.
\end{itemize}

$ S\_1 $ may be initialized in a number of different ways, most commonly by setting $ S\_1 $ to $ Y\_1 $ as shown above, though other techniques exist, such as setting $ S\_1 $ to an average of the first 4 or 5 observations. The importance of the $ S\_1 $ initialisation''s effect on the resultant moving average depends on $ \alpha $; smaller $ \alpha $ values make the choice of $ S\_1 $ relatively more important than larger $ \alpha $ values, since a higher $ \alpha $ discounts older observations faster.



Whatever is done for $ S\_1 $ assumes something about values prior to the available data and is necessarily in error. In view of this, the early results should be regarded as unreliable until the iterations have had time to converge. This is sometimes called a ''spin-up'' interval. One way to assess when it can be regarded as reliable is to consider the required accuracy of the result. For example, if 3\% accuracy is required, initializing with $ Y\_1 $ and taking data after five time constants (defined above) will ensure that the calculation has converged to within 3\% (only <3\% of $ Y\_1 $ will remain in the result). Sometimes with very small $ \alpha $, this can mean little of the result is useful. This is analogous to the problem of using a convolution filter (such as a weighted average) with a very long window.



This formulation is according to Hunter (1986). By repeated application of this formula for different times, we can eventually write $ S\_t $ as a weighted sum of the datum points $ Y\_t $:



$$

S\_t = \alpha \left[Y\_t + (1 - \alpha)Y\_{t-1} + (1 - \alpha)^2 Y\_{t-2} + \cdots + (1 - \alpha)^k Y\_{t-k} \right] + (1 - \alpha)^{k+1} S\_{t-(k+1)}

$$



for any suitable $ k \in$ {$0, 1, 2, \ldots$}. The weight of the general datum $ Y\_{t-i} $ is $ \alpha (1 - \alpha)^i $.



This formula can also be expressed in technical analysis terms as follows, showing how the EMA steps towards the latest datum, but only by a proportion of the difference (each time):



$$

\text{EMA}\_{\text{today}} = \text{EMA}\_{\text{yesterday}} + \alpha \left[ \text{price}\_{\text{today}} - \text{EMA}\_{\text{yesterday}} \right]

$$



Expanding out $ \text{EMA}\_{\text{yesterday}} $ each time results in the following power series, showing how the weighting factor on each datum $ p\_1, p\_2, \ldots $ decreases exponentially:



$$

\text{EMA}\_{\text{today}} = \alpha \left[ p\_1 + (1 - \alpha)p\_2 + (1 - \alpha)^2 p\_3 + (1 - \alpha)^3 p\_4 + \cdots \right]

$$



where



\begin{itemize}
    \item $ p\_1 $ is $ \text{price}\_{\text{today}} $
    \item $ p\_2 $ is $ \text{price}\_{\text{yesterday}} $
\end{itemize}

and so on



$$

\text{EMA}\_{\text{today}} = \frac{p\_1 + (1 - \alpha)p\_2 + (1 - \alpha)^2 p\_3 + (1 - \alpha)^3 p\_4 + \cdots }{1 + (1 - \alpha) + (1 - \alpha)^2 + (1 - \alpha)^3 + \cdots }

$$



since



$$

\frac{1}{\alpha} = 1 + (1 - \alpha) + (1 - \alpha)^2 + \cdots

$$



It can also be calculated recursively without introducing the error when initializing the first estimate (n starts from 1):



$$

\text{EMA}\_n = \frac{\text{WeightedSum}\_n}{\text{WeightedCount}\_n}

$$



$$

\text{WeightedSum}\_n = p\_n + (1 - \alpha)\text{WeightedSum}\_{n-1}

$$



$$

\text{WeightedCount}\_n = 1 + (1 - \alpha)\text{WeightedCount}\_{n-1} = \frac{1 - (1 - \alpha)^n}{1 - (1 - \alpha)} = \frac{1 - (1 - \alpha)^n}{\alpha}

$$



Assume



$$

\text{WeightedSum}\_0 = \text{WeightedCount}\_0 = 0

$$



This is an infinite sum with decreasing terms.



\paragraph{Approximating the EMA with a Limited Number of Terms}



The question of how far back to go for an initial value depends, in the worst case, on the data. Large price values in old data will affect the total even if their weighting is very small. If prices have small variations then just the weighting can be considered. The power formula above gives a starting value for a particular day, after which the successive days formula shown first can be applied. The weight omitted by stopping after $ k $ terms is



$$

\alpha \left[(1 - \alpha)^k + (1 - \alpha)^{k+1} + (1 - \alpha)^{k+2} + \cdots \right],

$$



which is



$$

\alpha (1 - \alpha)^k \left[1 + (1 - \alpha) + (1 - \alpha)^2 + \cdots \right],

$$



i.e. a fraction



$$

\frac{\text{weight omitted by stopping after } k \text{ terms}}{\text{total weight}} = \frac{\alpha \left[(1 - \alpha)^k + (1 - \alpha)^{k+1} + (1 - \alpha)^{k+2} + \cdots \right]}{\alpha \left[1 + (1 - \alpha) + (1 - \alpha)^2 + \cdots \right]} = (1 - \alpha)^k

$$



out of the total weight.



For example, to have 99.9\% of the weight, set the above ratio equal to 0.1\% and solve for $ k $:



$$

k = \frac{\log(0.001)}{\log(1 - \alpha)}

$$



to determine how many terms should be used. Since $ \alpha \to 0 $ as $ N \to \infty $, we know $ \log(1 - \alpha) $ approaches $ -\alpha $ as $ N $ increases. This gives:



$$

k \approx \frac{\log(0.001)}{-\alpha}

$$



When $ \alpha $ is related to $ N $ as $ \alpha = \frac{2}{N + 1} $, this simplifies to approximately



$$

k \approx 3.45(N + 1)

$$



for this example (99.9\% weight).



\paragraph{Relationship Between SMA and EMA}



Note that there is no "accepted" value that should be chosen for $ \alpha $, although there are some recommended values based on the application. A commonly used value for $ \alpha $ is $ \alpha = \frac{2}{N + 1} $. This is because the weights of an SMA and EMA have the same "center of mass" when $ \alpha\_{\mathrm{EMA}} = \frac{2}{N\_{\mathrm{SMA}} + 1} $.



\textbf{Proof}



The weights of an $N$-day SMA have a "center of mass" on the $R^{\text{th}}$ day, where



$$

R = \frac{N + 1}{2}

$$



(or 



$$

R = \frac{N - 1}{2}

$$



if we use zero-based indexing).



For the remainder of this proof we will use one-based indexing.



Now, meanwhile, the weights of an EMA have center of mass



$$

R\_{\text{EMA}} = \alpha \left[ 1 + 2(1 - \alpha) + 3(1 - \alpha)^2 + \cdots + k(1 - \alpha)^{k-1} \right]

$$



That is,



$$

R\_{\text{EMA}} = \alpha \sum\_{k=1}^{\infty} k(1 - \alpha)^{k-1}

$$



We also know the Maclaurin Series



$$

\frac{1}{1 - x} = \sum\_{k=0}^{\infty} x^k

$$



Taking derivatives of both sides with respect to $x$ gives:



$$

(x - 1)^{-2} = \sum\_{k=0}^{\infty} kx^{k-1}

$$



or



$$

(x - 1)^{-2} = 0 + \sum\_{k=1}^{\infty} kx^{k-1}

$$



Substituting $x = 1 - \alpha$, we get



$$

R\_{\text{EMA}} = \alpha \left( \alpha \right)^{-2}

$$



or



$$

R\_{\text{EMA}} = \left( \alpha \right)^{-1}

$$



So the value of $\alpha$ that sets $R\_{\text{SMA}} = R\_{\text{EMA}}$ is, in fact:



$$

\frac{N\_{\text{SMA}} + 1}{2} = \left( \alpha\_{\text{EMA}} \right)^{-1}

$$



or



$$

\frac{2}{N\_{\text{SMA}} + 1} = \alpha\_{\text{EMA}}

$$



q.e.d.



And so $\frac{2}{N + 1}$ is the value of $\alpha$ that creates an EMA whose weights have the same center of gravity as would the equivalent $N$-day SMA. This is also why sometimes an EMA is referred to as an $N$-day EMA. Despite the name suggesting there are $N$ periods, the terminology only specifies the $\alpha$ factor. $N$ is not a stopping point for the calculation in the way it is in an SMA or WMA. For sufficiently large $N$, the first $N$ datum points in an EMA represent about 86\% of the total weight in the calculation when $\alpha = \frac{2}{N + 1}$



\textbf{Proof}



The sum of the weights of all the terms (i.e., infinite number of terms) in an exponential moving average is 1. The sum of the weights of $N$ terms is $1 - (1 - \alpha)^{N + 1}$.  Both of these sums can be derived by using the formula for the sum of a geometric series. The weight omitted after $N$ terms is given by subtracting this from 1, and you get $1 - \left[1 - (1 - \alpha)^{N + 1}\right] = (1 - \alpha)^{N + 1}$ (this is essentially the formula given previously for the weight omitted).



We now substitute the commonly used value for $\alpha = \frac{2}{N + 1}$ in the formula for the weight of $N$ terms. If you make this substitution, and you make use of $\lim\_{n \to \infty} \left(1 + \frac{a}{n}\right)^n = e^a$ then you get:



$$

\frac{\alpha \left[1 + (1 - \alpha) + (1 - \alpha)^2 + \cdots + (1 - \alpha)^N\right]}{\alpha \left[1 + (1 - \alpha) + (1 - \alpha)^2 + \cdots \right]} = 1 - \left(1 - \frac{2}{N + 1}\right)^N

$$



i.e.



$$

\lim\_{N \to \infty} \left[1 - \left(1 - \frac{2}{N + 1}\right)^{N + 1}\right]

$$



simplified, tends to 



$$

1 - e^{-2} \approx 0.8647

$$



the 0.8647 approximation. Intuitively, what this is telling us is that the weight after $N$ terms of an "N-period" exponential moving average converges to 0.8647.



q.e.d.



The designation of $\alpha = \frac{2}{N + 1}$ is not a requirement. (For example, a similar proof could be used to just as easily determine that the EMA with a half-life of $N$-days is $\alpha = 1 - 0.5^{\frac{1}{N}}$ or that the EMA with the same median as an $N$-day SMA is $\alpha = 1 - 0.5^{\frac{1}{0.5N}}$). In fact, $\frac{2}{N + 1}$ is merely a common convention to form an intuitive understanding of the relationship between EMAs and SMAs, for industries where both are commonly used together on the same datasets. In reality, an EMA with any value of $\alpha$ can be used, and can be named either by stating the value of $\alpha$, or with the more familiar $N$-day EMA terminology letting $N = \left(\frac{2}{\alpha}\right) - 1$



\paragraph{Exponentially Weighted Moving Variance and Standard Deviation}



In addition to the mean, we may also be interested in the variance and in the standard deviation to evaluate the statistical significance of a deviation from the mean.



EWMVar can be computed easily along with the moving average. The starting values are $\text{EMA}\_1 = x\_1$ and $\text{EMVar}\_1 = 0$ and we then compute the subsequent values using:



$$

\begin{aligned}

\delta\_i &= x\_i - \text{EMA}\_{i-1} \\text{EMA}\_i &= \text{EMA}\_{i-1} + \alpha \cdot \delta\_i \\text{EMVar}\_i &= (1 - \alpha) \left(\text{EMVar}\_{i-1} + \alpha \cdot \delta\_i^2\right)

\end{aligned}

$$



From this, the exponentially weighted moving standard deviation can be computed as $\text{EMSD}\_i = \sqrt{\text{EMVar}\_i}$.  We can then use the standard score to normalize data with respect to the moving average and variance. This algorithm is based on Welford''s algorithm for computing the variance.



\paragraph{Modified Moving Average}



A \textbf{modified moving average} (MMA), \textbf{running moving average} (RMA), or \textbf{smoothed moving average} (SMMA) is defined as:



$$

\overline{p}\_{MM, \text{today}} = \frac{(N - 1) \overline{p}\_{MM, \text{yesterday}} + p\_{\text{today}}}{N}

$$



In short, this is an exponential moving average, with $\alpha = \frac{1}{N}$.  The only difference between EMA and SMMA/RMA/MMA is how $\alpha$ is computed from $N$. For EMA the customary choice is $\alpha = \frac{2}{N + 1}$



\paragraph{Application to Measuring Computer Performance}



Some computer performance metrics, e.g., the average process queue length, or the average CPU utilization, use a form of exponential moving average.



$$

S\_n = \alpha (t\_n - t\_{n-1}) Y\_n + \left[1 - \alpha (t\_n - t\_{n-1})\right] S\_{n-1}

$$



Here $\alpha$ is defined as a function of time between two readings. An example of a coefficient giving bigger weight to the current reading, and smaller weight to the older readings is



$$

\alpha (t\_n - t\_{n-1}) = 1 - \exp \left(-\frac{t\_n - t\_{n-1}}{W \cdot 60}\right)

$$



where $\exp()$ is the exponential function, time for readings $t\_n$ is expressed in seconds, and $W$ is the period of time in minutes over which the reading is said to be averaged (the mean lifetime of each reading in the average). Given the above definition of $\alpha$, the moving average can be expressed as



$$

S\_n = \left[1 - \exp \left(-\frac{t\_n - t\_{n-1}}{W \cdot 60}\right)\right] Y\_n + \exp \left(-\frac{t\_n - t\_{n-1}}{W \cdot 60}\right) S\_{n-1}

$$



For example, a 15-minute average $L$ of a process queue length $Q$, measured every 5 seconds (time difference is 5 seconds), is computed as



$$

\begin{aligned}

L\_n &= \left[1 - \exp \left(-\frac{5}{15 \cdot 60}\right)\right] Q\_n + e^{-\frac{5}{15 \cdot 60}} L\_{n-1} \&= \left[1 - \exp \left(-\frac{1}{180}\right)\right] Q\_n + e^{-\frac{1}{180}} L\_{n-1} \&= Q\_n + e^{-\frac{1}{180}} \left(L\_{n-1} - Q\_n\right)

\end{aligned}

$$



\paragraph{Other Weightings}



Other weighting systems are used occasionally – for example, in share trading a \textbf{volume weighting} will weight each time period in proportion to its trading volume.



A further weighting, used by actuaries, is Spencer''s 15-Point Moving Average (a central moving average). Its symmetric weight coefficients are $[-3, -6, -5, 3, 21, 46, 67, 74, 67, 46, 21, 3, -5, -6, -3]$, which factors as $\frac{[1, 1, 1, 1] \times [1, 1, 1, 1] \times [1, 1, 1, 1, 1] \times [-3, 3, 4, 3, -3]}{320}$ and leaves samples of any cubic polynomial unchanged.



Outside the World of Finance, weighted running means have many forms and applications. Each weighting function or "kernel" has its own characteristics. In engineering and science, the frequency and phase response of the filter is often of primary importance in understanding the desired and undesired distortions that a particular filter will apply to the data.



A mean does not just "smooth" the data. A mean is a form of low-pass filter. The effects of the particular filter used should be understood in order to make an appropriate choice. On this point, the French version of this article discusses the spectral effects of 3 kinds of means (cumulative, exponential, Gaussian).



\paragraph{Moving Median}



From a statistical point of view, the moving average, when used to estimate the underlying trend in a time series, is susceptible to rare events such as rapid shocks or other anomalies. A more robust estimate of the trend is the \textbf{simple moving median} over $n$ time points:



$$

\tilde{p}\_{\text{SM}} = \text{Median}(p\_M, p\_{M-1}, \ldots, p\_{M-n+1})

$$



where the median is found by, for example, sorting the values inside the brackets and finding the value in the middle. For larger values of $n$, the median can be efficiently computed by updating an indexable skiplist.



Statistically, the moving average is optimal for recovering the underlying trend of the time series when the fluctuations about the trend are normally distributed. However, the normal distribution does not place high probability on very large deviations from the trend which explains why such deviations will have a disproportionately large effect on the trend estimate. It can be shown that if the fluctuations are instead assumed to be Laplace distributed, then the moving median is statistically optimal. For a given variance, the Laplace distribution places higher probability on rare events than does the normal, which explains why the moving median tolerates shocks better than the moving mean.



When the simple moving median above is central, the smoothing is identical to the median filter which has applications in, for example, image signal processing.



\paragraph{Moving Average Regression Model}



In a moving average regression model, a variable of interest is assumed to be a weighted moving average of unobserved independent error terms; the weights in the moving average are parameters to be estimated.



Those two concepts are often confused due to their name, but while they share many similarities, they represent distinct methods and are used in very different contexts.



\paragraph{Licensing}



Content obtained and/or adapted from:



\href{https://en.wikipedia.org/wiki/Moving_average}{Moving average, Wikipedia} under a CC BY-SA license

\subsection{Learning Objectives}

\begin{enumerate}
    \item \textbf{Identify and interpret the components of weighted moving average graphs.}
    \item Understand the concept and applications of moving averages.
    \item Recognize the different types of moving averages: simple, cumulative, and weighted.
    \item \textbf{Calculate simple moving averages (SMA), cumulative moving averages (CMA), and weighted moving averages (WMA).}
    \item Apply formulas to compute SMA, CMA, and WMA for given data sets.
    \item Understand the differences between these types of moving averages and their respective calculation methods.
    \item \textbf{Graph and analyze time series data using moving averages.}
    \item Plot moving averages on a graph to smooth out data and identify trends.
    \item Interpret the implications of the moving averages on the overall data trend.
\end{enumerate}

\subsection{Assessment Instrument}

\begin{enumerate}
    \item Question 1: Identifying Components Given the following description, identify and interpret the components of a weighted moving average graph.
    \item A weighted moving average (WMA) gives more importance to recent data points.
    \item The formula for WMA with weights decreasing arithmetically is: $$ WMA\_M = \frac{n p\_M + (n-1) p\_{M-1} + \cdots + 2 p\_{(M-n)+2} + p\_{(M-n)+1}}{n + (n-1) + \cdots + 2 + 1} $$ a. What does the variable $n$ represent in the formula? b. How does the weighted moving average differ from the simple moving average? Question 2: Calculating Moving Averages Given the following data set representing stock prices over 5 days: [10, 12, 14, 16, 18] a. Calculate the 3-day simple moving average (SMA). b. Calculate the cumulative moving average (CMA) for each day. c. Calculate the 3-day weighted moving average (WMA) with weights 3, 2, and 1. Question 3: Graphing Moving Averages Using the data set from Question 2, plot the original data, 3-day SMA, and 3-day WMA on a graph. a. Describe the trends shown by the original data and the smoothed data. b. Explain the benefits of using moving averages to analyze time series data.
\end{enumerate}

%%===============================================================
\section{Weighted moving average graphs continued}
\label{sec:weighted-moving-average-graphs-continued}
%%===============================================================

\paragraph{Weighted Moving Average Graphs Continued}







\paragraph{Calculating the Mean}

The mean, or more precisely the arithmetic mean, is simply the arithmetic average of a group of numbers (or data set) and is shown using the \(\bar{}\) symbol. So the mean of the variable \(x\) is \(\bar{x}\), pronounced "x-bar". It is calculated by adding up all of the values in a data set and dividing by the number of values in that data set:  $ \bar{x} = \frac{\sum x}{n} $.  For example, take the following set of data: \{1,2,3,4,5\}. The mean of this data would be:



$$ \bar{x} = \frac{1 + 2 + 3 + 4 + 5}{5} = \frac{15}{5} = 3 $$



Here is a more complicated data set: \{10,14,86,2,68,99,1\}. The mean would be calculated like this:



$$ \bar{x} = \frac{10 + 14 + 86 + 2 + 68 + 99 + 1}{7} = \frac{280}{7} = 40 $$



\paragraph{Spreadsheets}



\paragraph{What is a Spreadsheet}



% [Image: spreadsheet]



A spreadsheet is a group of values and other data organized into rows and columns similar to the ruled paper worksheets traditionally used by bookkeepers and accountants. Spreadsheet software is mandatory to create computerized spreadsheets. Microsoft Excel is a form of a spreadsheet. Spreadsheets can support keeping track of data, support in quickly formulating subtotals, populating visual graphs and charts, and essentially is a working tool that can easily be shared. A worksheet is a single spreadsheet document. A workbook allows multiple worksheets to be saved together in a single spreadsheet file. Worksheets are divided into rows and columns. The intersection of a row and a column is called a cell. One must enter content into the active cell, or current cell; it has a border around it to make it easily identified. Data is entered directly into worksheet cells by clicking a cell to make it the active cell. Labels, constant values, formulas, and functions are the data that is entered into a cell. Before one enters a formula or function into a cell, one must begin with some type of mathematical symbol, usually the equal sign (=). Spreadsheets are used to organize and calculate data. There is a maximum number of rows and columns in a spreadsheet, which varies depending on the version of software you have. It is essential to know how to use spreadsheets for school, work, sports, or anything that requires data!





\paragraph{Tables, Graphics, and Templates}



% [Image: table]



Tables, graphics, and templates are all available to a user with application software, such as Microsoft Word, Microsoft Excel, and PowerPoint. Tables are ways a user can organize data and information at their convenience. According to Microsoft Word, there are now many different available options for users who are looking for various kinds of tables. These different options include the following: the Graphic Grid, Insert Table, Draw Table, insert a new or existing Excel Spreadsheet table, and Quick Tables. The concept of using tables for data input is relatively simple. In order for a user to insert a table, the user must first open Microsoft Word. Once they have done this, they must click the "table" button to customize the table to achieve their needs. The overall format for a table consists of a large (or small) grid that can be altered by the amount of information the user has, e.g., four columns five rows. Next, the user must insert the table into the Word document by selecting "insert table" from the dropdown menu. Microsoft Excel contains pivot tables that are tables that include data from a spreadsheet with columns and rows that can be specifically selected. Graphics in Microsoft Word are pictures, or clip art that are able to be inserted into a Microsoft Word document, Excel Spreadsheet, PowerPoint slide, or any other Office application. The most common graphic used in Excel is graphs. You can create graphs based on data taken from your spreadsheet. Graphics are inserted into these Office Applications to enhance the information presented in a Word Document, Excel worksheet, or PowerPoint slide. A user can insert their own picture through their Office documents; add clip art, shapes, SmartArt, screenshot, or Word Art. Templates are pre-constructed document layouts whose primary use is to assist a user in creating a specific type of document in a convenient amount of time. The different options of templates vary, but a few of the following are common ones used every day: agendas, brochures, calendars, flyers, fax covers, and many more. Templates are used to save a user time and confusion in creating their document.



\paragraph{How to use a Spreadsheet}



When using a spreadsheet application, the user can use various concepts to compute the data entered into the cells in the spreadsheet. These different concepts are provided within the program. Some very common concepts that are utilized are charts, functions, formulas, and cell referencing.



\paragraph{Charts}



% [Image: chart]



A chart can be created as its own object or embedded within the sheet itself. This is helpful when a user needs to analyze data or represent changing data. Some forms of charts are: line graphs, scatter plot charts, bar charts, Venn diagram charts, and the list goes on.





\paragraph{Functions}



% [Image: function]



A function is a pre-programmed mathematical formula to allow the user to make calculations based on the data input. The functions under spreadsheets are there to perform a simple calculation by using a certain value, called arguments. If the user wishes to create their own formula, Visual Basic can be used to write a formula and then the spreadsheet program can input the values into the newly written formula, reporting the data back into the sheet. There are many different reasons to have functions on spreadsheets. One would be for arithmetic functions to process numerical data. The next would be statistical functions that use analysis tools and averaging tools. This would be useful for finding the average of the numbers in a certain row/column on a spreadsheet. The next function is date that processes and converts dates. This function could be used to put the sequential dates in order on the spreadsheet. The next function is logic functions that process logic data. An example of logic data would be an AND/OR function. If there was something that needed to be marked yes if it is above 5 and marked no if it is below 5 then that would be a logic function. The last one is financial functions that process monetary data. They all must start with an equal sign, the name of the function, parenthesis opening and closing. In the function a comma or semicolon is used as the delimiter, depending on what settings are there in the spreadsheet would depend on which one to use. An example would be \textbf{=SUM(A1:A4)}, this function would find the sum in those cells. Some of the most common functions are \textbf{SUM}, \textbf{AVERAGE}, \textbf{IF}, \textbf{COUNT}, \textbf{MAX} and \textbf{MIN}.



\paragraph{Formulas}



% [Image: formula]





A formula identifies the calculation needed to place the result in the cell it is contained within. This means a cell has two display components; the formula itself and the resulting value. Typically, a formula consists of five expressions: value, references, arithmetic operations, relation operations, and functions. By using these expressions, formulas can help to make tables, solve math problems, calculate a mortgage, figure out accounting tasks, and many other business-related tasks that used to be done tediously on paper.



A formula always starts with an equals sign ($=$), followed by a constant, a function or a reference, then followed by an operator, and then followed by another constant, function or reference. A constant is a value that never changes; this includes numbers, dates, titles and other text input. References represent a certain cell, such as $A2$. An operator is usually a math symbol, such as $+$ or $\*$ which tells the computer how to compute (add or multiply, respectively) the given constants or functions given in the formula. It is good to be careful that one knows the difference between a constant and a reference. If the constant $30$ is input into cell $A3$, and the formula says $=30+2$, then if $A3$’s value changes, the expression of the formula will not change unless the formula itself changes. If one wishes to have a formula that returns the value of a cell, then the formula should read $=A3+2$. Another thing to note is that the operators will follow the basic “rules” of calculation. For example, the formula $=3+2\*4$ will add $3$ to the product of $2$ and $4$, rather than add $3$ and $2$, then multiply the sum by $4$. (Parentheses can be used to change the order: $(3+2)*4$ would add first, then multiply.) Operators are not always arithmetic; they can also be comparison, text concatenation, and reference operators. Comparison includes greater than, less than, greater than or equal to, and less than or equal to. To connect two values into one value, a text concatenation (the “and” sign i.e. \&) is used. The signs used as reference operators are the following: a colon is used to reference two cells and all the cells between them (i.e. $B1:B10$); a comma is used to combine multiple references into one reference (i.e. $B1:B10,C1:C10$); and a space is used as an intersection operator.



% [Image: graphics]



\paragraph{Cell Referencing}

Cell referencing refers to the ability to utilize a cell or range of cells in a spreadsheet and is commonly used to create formulas to calculate data. Formulas can retrieve data from one cell in the worksheet, different areas of the worksheet, or different cells throughout an entire workbook. There are two ways of doing this: relative and absolute cell referencing. A relative cell reference will adjust as the formula is copied from another cell while an absolute will not adjust. An example of this would be $=D2+F2$ which, when copied from row 2 to row 3, will become $D3+F3$. It is also important to note that a user can reference both the same sheet and other sheets in a workbook using this concept.



\paragraph{Pivot Tables}

One of the most powerful features available in the Microsoft Office spreadsheet program Excel is pivot tables. Pivot tables allow you to manipulate large amounts of raw data. They make it easy to analyze the data in different ways, with a simple click and drag. Vast quantities of data can be summarized in a variety of ways. Calculations can be performed by row or column. Data can be filtered or sorted automatically by any or all of the fields. Excel can even recommend a basic layout of a pivot table based on the type of data selected. A wizard is available to assist in the creation of the table. An important thing to remember when using pivot tables is that any time the original data source is modified, the data must be refreshed in the pivot table.



% [Image: pivot]



Once the pivot table has been created and the data has been analyzed in a meaningful way, it can then be represented graphically using pivot charts. All the basic chart types available in Excel are available in the pivot chart menu. Much like the pivot tables they are built on, they can also be manipulated with ease. They can be filtered to display only the relevant information from the main data source. They can also be added to and refreshed very easily.



\paragraph{Licensing}

Content obtained and/or adapted from:



\href{https://en.wikibooks.org/wiki/Statistics/Summary/Averages/mean}{Mean, Wikibooks} under a CC BY-SA license  

\href{https://en.wikibooks.org/wiki/Introduction_to_Computer_Information_Systems/Application_Software}{Application Software, Wikibooks} under a CC BY-SA license

\subsection{Learning Objectives}

\begin{enumerate}
    \item Calculate the mean (arithmetic average) of a given data set.
    \item Utilize spreadsheet software to input data and perform mean calculations.
    \item Create and interpret weighted moving average graphs using spreadsheet functions and pivot tables.
\end{enumerate}

\subsection{Assessment Instrument}

\begin{enumerate}
    \item Part 1: Calculating the Mean
    \item Calculate the mean of the following data set: \{5, 8, 12, 20, 25\}.
    \item Given the data set {$7, 15, 23, 30, 50$}, calculate the mean. Part 2: Using Spreadsheet Software
    \item Input the following data set into a spreadsheet: \{3, 6, 9, 12, 15\}. Use the spreadsheet functions to calculate the mean. Provide a screenshot of your spreadsheet showing the data and the calculated mean.
    \item Create a table in a spreadsheet with the following data: $$ \begin{array}{|c|c|c|} \hline \text{A} & \text{B} & \text{C} \\\   \hline 5 & 10 & 15 \\\\ \hline 20 & 25 & 30 \\\\ \hline 35 & 40 & 45 \\\   \hline \end{array} $$ Calculate the mean of each column using the appropriate spreadsheet functions. Provide screenshots of your results. Part 3: Creating Weighted Moving Average Graphs
    \item Given the following data set for sales over 7 weeks: {$100, 150, 200, 250, 300, 350, 400$}, create a weighted moving average graph in a spreadsheet with weights {$0.1, 0.2, 0.3, 0.4$}. Provide a screenshot of your graph and the steps you followed to create it.
    \item Using the data from the previous question, create a pivot table in the spreadsheet to summarize the sales data. Then, create a pivot chart to visually represent the data. Provide screenshots of both the pivot table and the pivot chart.
\end{enumerate}

%%===============================================================
\section{Ratios and percentages}
\label{sec:ratios-and-percentages}
%%===============================================================

\paragraph{Ratios and Percentages}





\paragraph{Ratio}

In mathematics, a ratio indicates how many times one number contains another. For example, if there are eight oranges and six lemons in a bowl of fruit, then the ratio of oranges to lemons is eight to six (that is, $8:6$, which is equivalent to the ratio $4:3$). Similarly, the ratio of lemons to oranges is $6:8$ (or $3:4$) and the ratio of oranges to the total amount of fruit is $8:14$ (or $4:7$).



The numbers in a ratio may be quantities of any kind, such as counts of people or objects, or such as measurements of lengths, weights, time, etc. In most contexts, both numbers are restricted to be positive.



A ratio may be specified either by giving both constituting numbers, written as "a to b" or "a:b", or by giving just the value of their quotient $\frac{a}{b}$. Equal quotients correspond to equal ratios.



Consequently, a ratio may be considered as an ordered pair of numbers, a fraction with the first number in the numerator and the second in the denominator, or as the value denoted by this fraction. Ratios of counts, given by (non-zero) natural numbers, are rational numbers, and may sometimes be natural numbers. When two quantities are measured with the same unit, as is often the case, their ratio is a dimensionless number. A quotient of two quantities that are measured with different units is called a rate.



\paragraph{Simplify Fractions}

If both the numerator and the denominator of a fraction are multiplied or divided by the same number, then the fraction does not change its value.



This means we can make some fractions simpler, by making the numbers involved smaller. When the numbers are as small as possible, the fraction is said to be expressed in its lowest terms, or reduced. Let''s look at an example: $ \frac{4}{6} = \frac{2}{3} $.



Remember what that $=$ sign means? It means that $\frac{4}{6}$ and $\frac{2}{3}$, although they look different, have the same value. Same value, different appearance.



This is because

$$

\frac{4}{6} = \frac{4/2}{6/2} = \frac{2}{3}.

$$

Try it out yourself with a calculator or with long division. From here, there is nothing we can divide both numbers by, so the fraction is reduced, which is our goal. To find out what to divide them by, you need to find a common divisor, which is a number they are both divisible by.



\paragraph{Percentages}

Percentages are another way of representing rational numbers. Rational numbers can be represented as decimals, fractions, or percentages.



Percent, when you actually break the word apart, consists of two words: per and cent. Per is a small but extremely powerful word. It means "to divide". Cent is a Latin word for 100. Percent, then, means "divide by 100".



To convert from a fraction to a percentage, we simply multiply the fraction by 100. For example,

$$

\frac{1}{4} = \left( \frac{1}{4} \times 100 \right)\% = 25\%.

$$



To convert from a decimal to a percentage, we first convert the decimal into a fraction, and then proceed with the approach outlined above. For example,

$$

0.15 = \frac{15}{100} = \left( \frac{15}{100} \times 100 \right)\% = 15\%.

$$



\paragraph{Pie Charts}





% [Image: pie]



A Pie-Chart/Diagram is a graphical device - a circular shape broken into sub-divisions. The sub-divisions are called "sectors", whose areas are proportional to the various parts into which the whole quantity is divided. The sectors may be colored differently to show the relationship of parts to the whole. A pie diagram is an alternative of the sub-divided bar diagram.



To construct a pie-chart, first we draw a circle of any suitable radius then the whole quantity which is to be divided is equated to 360 degrees. The different parts of the circle in terms of angles are calculated by the following formula.



$$

\frac{\text{Component Value}}{\text{Whole Quantity}} \times 360

$$



The component parts, i.e., sectors, have been cut beginning from the top in clockwise order.



Note that the percentages in a list may not add up to exactly 100\% due to rounding. For example, if a person spends a third of their time on each of three activities: 33\%, 33\%, and 33\%, the sums to 99\%.



Warning: Pie charts are a poor way of communicating information. The eye is good at judging linear measures and bad at judging relative areas. A bar chart or dot chart is a preferable way of displaying this type of data.



Cleveland (1985), page 264: "Data that can be shown by pie charts always can be shown by a dot chart. This means that judgments of position along a common scale can be made instead of the less accurate angle judgments." This statement is based on the empirical investigations of Cleveland and McGill as well as investigations by perceptual psychologists.



Three-dimensional (3D) pie charts compound perceptual misinterpretation of statistical information by altering the relative angle of pie slices to create the impression of depth into a vanishing point. Angles and areas at the bottom of the chart must be exaggerated and the angles and areas at the top of the chart reduced in order to create the dimensional effect; a specifically false depiction of the data.



\paragraph{Licensing}

Content obtained and/or adapted from:



\href{https://en.wikipedia.org/wiki/Ratio}{Ratio, Wikipedia} under a CC BY-SA license  

\href{https://en.wikibooks.org/wiki/Fractions}{Fractions, Wikibooks} under a CC BY-SA license  

\href{https://en.wikibooks.org/wiki/Statistics}{Statistics, Wikibooks} under a CC BY-SA license  

\href{https://en.wikibooks.org/wiki/Statistics/Pie_Charts}{Pie Charts, Wikibooks} under a CC BY-SA license

\subsection{Learning Objectives}

\begin{enumerate}
    \item \textbf{Identify and Interpret Ratios}
    \item Students will be able to identify and interpret ratios from various contexts, including counts and measurements.
    \item Example: Determine the ratio of apples to oranges in a basket containing 12 apples and 8 oranges.
    \item \textbf{Simplify Fractions to Lowest Terms}
    \item Students will be able to simplify fractions by dividing both the numerator and the denominator by their greatest common divisor (GCD).
    \item Example: Simplify the fraction $\frac{18}{24}$ to its lowest terms.
    \item \textbf{Convert Between Fractions, Decimals, and Percentages}
    \item Students will be able to convert between fractions, decimals, and percentages.
    \item Example: Convert the fraction $\frac{3}{5}$ to a decimal and a percentage.
\end{enumerate}

\subsection{Assessment Instrument}

\begin{enumerate}
    \item Question 1: Identify and Interpret Ratios
    \item \textbf{Problem}: In a class of 30 students, 18 are girls and 12 are boys.
    \item What is the ratio of girls to boys?
    \item What is the ratio of boys to the total number of students? Question 2: Simplify Fractions
    \item \textbf{Problem}: Simplify the fraction $\frac{28}{42}$ to its lowest terms. Question 3: Convert Between Fractions, Decimals, and Percentages
    \item \textbf{Problem}: Convert the decimal $0.75$ to a fraction and a percentage.
\end{enumerate}

%%===============================================================
\section{Part-to-part ratios \& Part-to-whole ratios}
\label{sec:part-to-part-ratios-part-to-whole-ratios}
%%===============================================================

\paragraph{Part-to-part ratios \& Part-to-whole ratios}



\textbf{Pie chart of populations of English native speakers}



% [Image: pie]



A pie chart (or a circle chart) is a circular statistical graphic, which is divided into slices to illustrate numerical proportion. In a pie chart, the arc length of each slice (and consequently its central angle and area), is proportional to the quantity it represents. While it is named for its resemblance to a pie which has been sliced, there are variations on the way it can be presented. The earliest known pie chart is generally credited to William Playfair''s \textit{Statistical Breviary} of 1801.



Pie charts are very widely used in the business world and the mass media. However, they have been criticized, and many experts recommend avoiding them, as research has shown it is difficult to compare different sections of a given pie chart, or to compare data across different pie charts. Pie charts can be replaced in most cases by other plots such as the bar chart, box plot, dot plot, etc.



\paragraph{History}



The earliest known pie chart is generally credited to William Playfair''s \textit{Statistical Breviary} of 1801, in which two such graphs are used. Playfair presented an illustration, which contained a series of pie charts. One of those charts depicted the proportions of the Turkish Empire located in Asia, Europe, and Africa before 1789. This invention was not widely used at first.



Playfair thought that pie charts were in need of a third dimension to add additional information.



Florence Nightingale may not have invented the pie chart, but she adapted it to make it more readable, which fostered its wide use, still today. Indeed, Nightingale reconfigured the pie chart making the length of the wedges variable instead of their width. The graph, then, resembled a cock''s comb. She was later assumed to have created it due to the obscurity and lack of practicality of Playfair''s creation. Nightingale''s polar area diagram, or occasionally the Nightingale rose diagram, equivalent to a modern circular histogram, to illustrate seasonal sources of patient mortality in the military field hospital she managed, was published in \textit{Notes on Matters Affecting the Health, Efficiency, and Hospital Administration of the British Army} and sent to Queen Victoria in 1858. According to the historian Hugh Small, "she may have been the first to use [pie charts] for persuading people of the need for change."



The French engineer Charles Joseph Minard also used pie charts, in 1858. A map of his from 1858 used pie charts to represent the cattle sent from all around France for consumption in Paris.



\paragraph{Early types of pie charts in the 19th century}



\begin{itemize}
    \item Pie charts from William Playfair''s \textit{Statistical Breviary}, 1801
\end{itemize}

% [Image: pie]



\begin{itemize}
    \item One of the pie charts, 1801
\end{itemize}

% [Image: pie]



\begin{itemize}
    \item Minard''s map, 1858
\end{itemize}

% [Image: map]



\begin{itemize}
    \item Polar chart by Florence Nightingale, 1858
\end{itemize}

% [Image: polar]



\paragraph{Variants and similar charts}





\paragraph{3D pie chart and perspective pie cake}





\textbf{3D pie chart showing Atmospheric air components percentage}

% [Image: 3D]



A 3D pie chart, or perspective pie chart, is used to give the chart a 3D look. Often used for aesthetic reasons, the third dimension does not improve the reading of the data; on the contrary, these plots are difficult to interpret because of the distorted effect of perspective associated with the third dimension. The use of superfluous dimensions not used to display the data of interest is discouraged for charts in general, not only for pie charts.



\paragraph{Doughnut Chart}



\textbf{Information on the data as a hole in the center of a doughnut chart}



% [Image: doughnut]



A doughnut chart (also spelled donut) is a variant of the pie chart, with a blank center allowing for additional information about the data as a whole to be included. Doughnut charts are similar to pie charts in that their aim is to illustrate proportions. This type of circular graph can support multiple statistics at once and it provides a better data intensity ratio to standard pie charts. It does not have to contain information in the center.



\paragraph{Exploded Pie Chart}



\textit{Exploded pie chart for the example data (see below), with the largest party group exploded.}



% [Image: exploded]



A chart with one or more sectors separated from the rest of the disk is known as an exploded pie chart. This effect is used to either highlight a sector, or to highlight smaller segments of the chart with small proportions.



\paragraph{Polar Area Diagram}



\textit{"Diagram of the causes of mortality in the army in the East" by Florence Nightingale.}



% [Image: polar]



The polar area diagram is similar to a usual pie chart, except sectors have equal angles and differ rather in how far each sector extends from the center of the circle. The polar area diagram is used to plot cyclic phenomena (e.g., counts of deaths by month). For example, if the counts of deaths in each month for a year are to be plotted then there will be 12 sectors (one per month) all with the same angle of 30 degrees each. The radius of each sector would be proportional to the square root of the death count for the month, so the area of a sector represents the number of deaths in a month. If the death count in each month is subdivided by cause of death, it is possible to make multiple comparisons on one diagram, as is seen in the polar area diagram famously developed by Florence Nightingale.



The first known use of polar area diagrams was by André-Michel Guerry, which he called \textit{courbes circulaires} (circular curves), in an 1829 paper showing seasonal and daily variation in wind direction over the year and births and deaths by hour of the day. Léon Lalanne later used a polar diagram to show the frequency of wind directions around compass points in 1843. The wind rose is still used by meteorologists. Nightingale published her rose diagram in 1858. Although the name "coxcomb" has come to be associated with this type of diagram, Nightingale originally used the term to refer to the publication in which this diagram first appeared—an attention-getting book of charts and tables—rather than to this specific type of diagram.



\paragraph{Ring Chart, Sunburst Chart, and Multilevel Pie Chart}



\textbf{Multi-level pie chart representing disk usage in a Linux file system}



% [Image: multi]



A ring chart, also known as a sunburst chart or a multilevel pie chart, is used to visualize hierarchical data, depicted by concentric circles. The circle in the center represents the root node, with the hierarchy moving outward from the center. A segment of the inner circle bears a hierarchical relationship to those segments of the outer circle which lie within the angular sweep of the parent segment.



\paragraph{Spie Chart}



\textbf{A spie chart comparing number of students with student costs across four different schools}



% [Image: spie]



A spie chart comparing number of students with student costs across four different schools

A variant of the polar area chart is the spie chart, designed by Dror Feitelson. The design superimposes a normal pie chart with a modified polar area chart to permit the comparison of two sets of related data. The base pie chart represents the first data set in the usual way, with different slice sizes. The second set is represented by the superimposed polar area chart, using the same angles as the base, and adjusting the radii to fit the data. For example, the base pie chart could show the distribution of age and gender groups in a population, and the overlay their representation among road casualties. Age and gender groups that are especially susceptible to being involved in accidents then stand out as slices that extend beyond the original pie chart.



\paragraph{Square Chart / Waffle Chart}



**Square pie chart (waffle chart), showing how smaller percentages are more easily shown than on circular charts. On the 10x10 grid, each cell represents 1\%.

**



% [Image: square]



Square pie chart (waffle chart), showing how smaller percentages are more easily shown than on circular charts. On the 10x10 grid, each cell represents 1\%.

Square charts, also called waffle charts, are a form of pie charts that use squares instead of circles to represent percentages. Similar to basic circular pie charts, square pie charts take each percentage out of a total 100\%. They are often 10 by 10 grids, where each cell represents 1\%. Despite the name, circles, pictograms (such as of people), and other shapes may be used instead of squares. One major benefit to square charts is that smaller percentages, difficult to see on traditional pie charts, can be easily depicted.



\paragraph{Example}



A pie chart for the example data

The following example chart is based on preliminary results of the election for the European Parliament in 2004. The table lists the number of seats allocated to each party group, along with the derived percentage of the total that they each make up. The values in the last column, the derived central angle of each sector, is found by multiplying the percentage by 360°.



$$

\begin{array}{|c|c|c|c|}

\hline

\text{Group} & \text{Seats} & \text{Percent (\%)} & \text{Central Angle (Degrees)} \\\\ \hline

\text{EUL}   & 39    & 5.3 & 19.2 \\\\ \hline

\text{PES}   & 200   & 27.3 & 98.4 \\\\ \hline

\text{EFA}   & 42    & 5.7 & 20.7 \\\\ \hline

\text{EDD}   & 15    & 2.0 & 7.4 \\\\ \hline

\text{ELDR}  & 67    & 9.2 & 33.0 \\\\ \hline

\text{EPP}   & 276   & 37.7 & 135.7 \\\\ \hline

\text{UEN}   & 27    & 3.7 & 13.3 \\\\ \hline

\text{Other} & 66    & 9.0 & 32.5 \\\\ \hline

\text{Total} & 732   & 99.9* & 360.2* \\\\ \hline

\end{array}

$$



\textbf{Because of rounding, these totals do not add up to 100 and 360.}



% [Image: pie]



\paragraph{Use and Effectiveness}





A flaw exhibited by pie charts is that they cannot show more than a few values without separating the visual encoding (the “slices”) from the data they represent (typically percentages). When slices become too small, pie charts have to rely on colors, textures or arrows so the reader can understand them. This makes them unsuitable for use with larger amounts of data. Pie charts also take up a larger amount of space on the page compared to the more flexible bar charts, which do not need to have separate legends, and can display other values such as averages or targets at the same time.



Statisticians generally regard pie charts as a poor method of displaying information, and they are uncommon in scientific literature. One reason is that it is more difficult for comparisons to be made between the size of items in a chart when area is used instead of length and when different items are shown as different shapes.



\textbf{Three sets of percentages, plotted as both pie charts and bar charts}



% [Image: pie bar]



Comparing the data on bar charts is generally easier. Further, in research performed at AT\&T Bell Laboratories, it was shown that comparison by angle was less accurate than comparison by length. Most subjects have difficulty ordering the slices in the pie chart by size; when an equivalent bar chart is used the comparison is much easier. Similarly, comparisons between data sets are easier using the bar chart. However, if the goal is to compare a given category (a slice of the pie) with the total (the whole pie) in a single chart and the multiple is close to 25 or 50 percent, then a pie chart can often be more effective than a bar graph.



\textbf{An example of a pie chart with 18 values}



% [Image: pie]



In a pie chart with many sections, several values may be represented with the same or similar colors, making interpretation difficult.



\textbf{An example of a doughnut shape pie chart}



% [Image: donut]



Several studies presented at the European Visualization Conference analyzed the relative accuracy of several pie chart formats, reaching the conclusion that pie charts and doughnut charts produce similar error levels when reading them, and square pie charts provide the most accurate reading.



\paragraph{Licensing}



Content obtained and/or adapted from:



\href{https://en.wikipedia.org/wiki/Pie_chart}{Pie chart, Wikipedia} under a CC BY-SA license

\subsection{Learning Objectives}

\begin{enumerate}
    \item SLO 1: Define and differentiate between part-to-part ratios and part-to-whole ratios. Identify these ratios in pie charts and explain their significance in representing data. SLO 2: Calculate the central angle of each sector in a pie chart given the percentage of the whole represented by each sector. Convert between percentages and central angles to interpret pie chart data. SLO 3: Compare the effectiveness of pie charts, doughnut charts, and bar charts in representing part-to-whole and part-to-part ratios. Discuss scenarios where one type of chart may be preferred over others based on data presentation needs.
\end{enumerate}

\subsection{Assessment Instrument}

\begin{enumerate}
    \item Multiple Choice Questions
    \item \textbf{Which of the following correctly describes a part-to-whole ratio in a pie chart?}
    \item A) The ratio of one sector to another sector in the chart.
    \item B) The ratio of the central angle of one sector to the central angle of another sector.
    \item C) The ratio of the central angle of a sector to the total central angle of the pie chart.
    \item D) The ratio of the number of slices to the total number of slices in the chart.
    \item \textbf{Given that a sector in a pie chart represents 25\% of the whole, what is the central angle of this sector?}
    \item A) 25°
    \item B) 75°
    \item C) 90°
    \item D) 360°
    \item \textbf{Which chart type is recommended when dealing with a large number of data categories to avoid clutter and improve readability?}
    \item A) Pie Chart
    \item B) Doughnut Chart
    \item C) Bar Chart
    \item D) Polar Area Diagram Short Answer Questions
    \item \textbf{Explain how a doughnut chart differs from a pie chart in terms of data representation. Provide an example of when a doughnut chart might be more useful than a pie chart.}
    \item \textbf{For a pie chart representing the distribution of four categories with percentages 20\%, 30\%, 25\%, and 25\%, calculate the central angle for each sector. Draw the pie chart showing these sectors.} Data Interpretation
    \item \textbf{Refer to the provided pie chart of English native speakers. Identify the part-to-whole ratio for the segment representing North American English speakers. Explain how this ratio is represented in the chart.} !\href{https://upload.wikimedia.org/wikipedia/commons/thumb/d/db/English_dialects1997.svg/450px-English_dialects1997.svg.png}{pie}
    \item \textbf{Analyze the pie chart for the 2004 European Parliament election. Discuss the benefits of using an exploded pie chart for this data and how it helps in interpreting the information compared to a standard pie chart.} !\href{https://upload.wikimedia.org/wikipedia/commons/thumb/2/2e/Pie_chart_EP_election_2004.svg/450px-Pie_chart_EP_election_2004.svg.png}{exploded}
\end{enumerate}

%%===============================================================
\section{Absolute change (additive reasoning) \& Relative change (multiplicative reasoning)}
\label{sec:absolute-change-additive-reasoning-relative-change-multiplicative-reasoning}
%%===============================================================

\paragraph{Absolute Change (Additive Reasoning) \& Relative Change (Multiplicative Reasoning)}



\paragraph{Measuring Errors}



Learning objectives



\begin{itemize}
    \item identify true and relative true errors
    \item identify approximate and relative approximate errors
    \item explain the relationship between the absolute relative approximate error and the number of significant digits
    \item identify significant digits
\end{itemize}



\paragraph{True and Relative True Errors}

A true error ($E\_t$) is defined as the difference between the true (exact) value and an approximate value. This type of error is only measurable when the true value is available. For example, a true error ($E\_t$) can be represented as:



$ \text{true error} (E\_t) = \text{true value} - \text{approximate value} $



A true error doesn''t signify the importance of an error. For instance, a 0.1 pound error is small when measuring a person''s weight, but the same error can be disastrous when measuring a medicine dosage. Relative true error ($\epsilon\_t$) is defined as the ratio between the true error and the true value.



$ \text{relative true error} (\epsilon\_t) = \frac{\text{true error}}{\text{true value}} $



\paragraph{Approximate and Relative Approximate Errors}

Often the true value is unknown, especially in numerical computing. In this case, we quantify errors using approximate values only. When an iterative method is used, we get an approximate value at the end of each iteration. The approximate error ($E\_a$) is defined as the difference between the present approximate value and the previous approximation.



$ \text{approximate error} (E\_a) = \text{present approximation} - \text{previous approximation} $



Similarly, we can calculate the relative approximate error ($\epsilon\_a$) by dividing the approximate error by the present approximate value.



$ \text{relative approximate error} (\epsilon\_a) = \frac{\text{approximate error}}{\text{present approximation}} $



\paragraph{Relative Approximate Error and Significant Digits}

Assuming our iterative method yields a better approximation as the iteration progresses, we often set an acceptable tolerance to stop the iteration when the relative approximate error is small enough. This tolerance is usually set in terms of significant digits - the number of digits that carry meaning contributing to precision. The number of significant digits corresponds to the number of digits in the scientific notation representing a number''s significand or mantissa.



An approximate rule for minimizing the error is: if the absolute relative approximate error is less than or equal to a predefined tolerance (usually in terms of significant digits), then the acceptable error has been reached, and no more iterations are required. Given the absolute relative approximate error, we can derive the least number of significant digits using the same equation.



$ |\epsilon\_a| \leq 0.5 \times 10^{2-m}\% $



\paragraph{Absolute Value}



The absolute value of a number is its distance from zero (0) on a number line. This action ignores the "+" or "–" sign of a number because distance in mathematics is never negative. The symbol $|x|$ represents the absolute value of $x$. It is also called the modulus of $x$.



\paragraph{Lesson}



You identify the absolute value of a number by writing the number between two vertical bars referred to as absolute value brackets: $| \text{number} |$.



A helpful way of thinking about absolute value is relating it to a railroad track. If you were to stand on a railroad track, more specifically on any one of the railroad ties and mark that spot as zero, railroad ties to the left would represent negative numbers and railroad ties to the right would represent positive numbers.



For example:



\begin{itemize}
    \item The number $-7$ is 7 units away from zero on the negative side of the railroad track. So, $|-7| = 7$.
    \item The number $16$ is 16 units away from zero on the positive side of the railroad track. So, $|16| = 16$.
    \item The number $0$ is 0 units from zero on the railroad track. So $|0| = 0$.
\end{itemize}

Therefore, the absolute value of any number is a positive number or zero.



In summary, THE ABSOLUTE VALUE OF A NUMBER:



\begin{itemize}
    \item If $x$ is a positive number, then $|x| = x$. Example: $|5| = 5$
    \item If $x$ is zero, then $|x| = 0$. Example: $|0| = 0$
    \item If $x$ is a negative number, $|x| = -x$. Example: $|-6| = -(-6) = 6$
\end{itemize}

You can find the absolute value of expressions as well. When addressed with this, treat the absolute value brackets as you would parentheses. Simplify everything inside the absolute value brackets by performing all necessary operations following the order of operations. The last step, once you have a single number inside the absolute value brackets, is to take the absolute value. For example:



$ |-5 + 1 \times 3| = |-5 + 3| = |-2| = 2 $



You may use the absolute value to find the distance between two numbers on the number line. Let $a$ and $b$ be variables. Then $|a - b|$ is the distance between $a$ and $b$. For example:



\begin{itemize}
    \item If $a = 3$ and $b = 7$, then $|3 - 7| = |-4| = 4$.
    \item If $a = 7$ and $b = 3$, then $|7 - 3| = |4| = 4$.
\end{itemize}

Two things to watch out for are an opposite sign and/or an operation outside the absolute value brackets. As stated above, simplify everything inside the absolute value brackets by performing all necessary operations following the order of operations. Your last step, once you have a single number inside the absolute value brackets, is to take the absolute value. Once you have taken the absolute value, perform the other necessary operations by following the order of operations from left to right in the expression. For example:



$ -|5| = -5 $

$ 7 + |-5 + 1 \times 3| = 7 + |-5 + 3| = 7 + |-2| = 7 + 2 = 9 $



\paragraph{Example Problems}

$$ |0| = 0 $$

$$ -|-21| = -21 $$

$$ |7 - 1| + 9 = |6| + 9 = 6 + 9 = 15 $$

$$ |0.5| = 0.5 $$

$$ |-5 \times 3| = |-15| = 15 $$

$$ \left| \frac{-2}{3} \right| = \frac{2}{3} $$

$$ |25 - 16| = |9| = 9 $$

$$ |9| = 9 $$

$$ |3.5 - 5.7| = |-2.2| = 2.2 $$

$$ |-11.5| = 11.5 $$

$$ 5 - |3| = 5 - 3 = 2 $$

$$ -|6.7| = -6.7 $$

$$ 8 + |-6 \times 2| = 8 + |-12| = 8 + 12 = 20 $$



\paragraph{Licensing}

Content obtained and/or adapted from:



\href{https://en.wikibooks.org/wiki/Introduction_to_Numerical_Methods/Measuring_Errors}{Measuring Errors, Wikibooks} under a CC BY-SA license  

\href{https://en.wikibooks.org/wiki/Basic_Algebra/Working_with_Numbers/Absolute_Value}{Absolute Value, Wikibooks} under a CC BY-SA license

\subsection{Learning Objectives}

\begin{enumerate}
    \item Identify true and relative true errors.
    \item Identify approximate and relative approximate errors.
    \item Explain the relationship between absolute relative approximate error and the number of significant digits.
\end{enumerate}

\subsection{Assessment Instrument}

\begin{enumerate}
    \item Part 1: True and Relative True Errors
    \item Given the true value of a quantity is $10.5$ and an approximate value is $10.2$, calculate the true error ($E\_t$).
    \item If the true error ($E\_t$) is $0.3$ and the true value is $15$, determine the relative true error ($\epsilon\_t$). Part 2: Approximate and Relative Approximate Errors
    \item During an iterative process, the present approximation is $8.4$ and the previous approximation is $8.1$. Calculate the approximate error ($E\_a$).
    \item Given the approximate error ($E\_a$) is $0.5$ and the present approximation is $9.5$, find the relative approximate error ($\epsilon\_a$). Part 3: Relative Approximate Error and Significant Digits
    \item If the absolute relative approximate error is $0.005$, determine the minimum number of significant digits required for the calculation.
    \item Explain how the relative approximate error impacts the number of significant digits in numerical computations.
\end{enumerate}

%%===============================================================
\section{Adjusting claims and hypothesis}
\label{sec:adjusting-claims-and-hypothesis}
%%===============================================================

\paragraph{Adjusting Claims and Hypothesis}



An example of a comparative line graph.



\textbf{Common Logarithms}

% [Image: line]



Line graphs are a good way to compare anything between two or more variables. These are one of the easier types of graphs to look at because you can clearly see the changes in the line as time goes on. A line graph should never be used to show statistics at a set moment in time. A disadvantage of line graphs is it can be difficult to see exactly what number is being specified at a data point. If it is important that the numbers be easily read, they may be inserted within the graph at their respective data points. Since most line graphs are used to convey a trend over time, however, specific numbers are not always necessary.



Main things to be included in a line graph are as follows:



\begin{itemize}
    \item \textbf{Title}: The most important thing (what is a chart without a title). Must be to the point.
    \item \textbf{X-axis (Horizontal) and Y-axis (Vertical)}: Both axes should be labeled and provide units of measurement. They should extend slightly past the highest graphed value.
    \item \textbf{Tick Marks}: Use tick marks on both axes to show units. Make sure they are not cluttering the graph.
    \item \textbf{Line}: The line shows the actual data being displayed. It can be displayed as a single data point, data points with a line connecting the points, or as a single smooth line with no data points. Different colors can also be used to designate various elements.
    \item \textbf{Source}: Should be included if the source will not be obvious to readers.
\end{itemize}

\textbf{Tips for making line graphs:}



\begin{itemize}
    \item Use different colors to ensure that the reader can distinguish between the lines. If you cannot use different colors, use different types of lines, for example, dots or dashes.
    \item If possible, start axes at zero to avoid misleading readers. If this is not possible, use hash marks to inform your readers that your graph does not start at zero.
\end{itemize}

\paragraph{Example}



In the experimental sciences, data collected from experiments are often visualized by a graph. For example, if one were to collect data on the speed of a body at certain points in time, one could visualize the data by a data table such as the following:



$$

\begin{array}{|c|c|}

\hline

\text{Elapsed Time (s)} & \text{Speed (mps)} \\\\hline

0 & 0 \\\\hline

1 & 3 \\\\hline

2 & 7 \\\\hline

3 & 12 \\\\hline

4 & 18 \\\\hline

5 & 30 \\\\hline

6 & 45.6 \\\\hline

\end{array}

$$



% [Image: line]



The table representation of data is a great way of displaying exact values, but can be a poor way to understand the underlying patterns that those values represent. Because of these qualities, the table display is often erroneously conflated with the data itself; whereas it is just another visualization of the data.



Understanding the process described by the data in the table is aided by producing a graph or line chart of Speed versus Time. Such a visualization appears in the figure above.



Mathematically, if we denote time by the variable $t$, and speed by $v$, then the function plotted in the graph would be denoted $v(t)$ indicating that $v$ (the dependent variable) is a function of $t$.



\paragraph{Best-fit}



A parody line graph (1919) by William Addison Dwiggins.

Charts often include an overlaid mathematical function depicting the best-fit trend of the scattered data. This layer is referred to as a best-fit layer and the graph containing this layer is often referred to as a line graph.



It is simple to construct a "best-fit" layer consisting of a set of line segments connecting adjacent data points; however, such a "best-fit" is usually not an ideal representation of the trend of the underlying scatter data for the following reasons:



\begin{itemize}
    \item It is highly improbable that the discontinuities in the slope of the best-fit would correspond exactly with the positions of the measurement values.
    \item It is highly unlikely that the experimental error in the data is negligible, yet the curve falls exactly through each of the data points.
\end{itemize}

In either case, the best-fit layer can reveal trends in the data. Further, measurements such as the gradient or the area under the curve can be made visually, leading to more conclusions or results from the data table.



A true best-fit layer should depict a continuous mathematical function whose parameters are determined by using a suitable error-minimization scheme, which appropriately weights the error in the data values. Such curve fitting functionality is often found in graphing software or spreadsheets. Best-fit curves may vary from simple linear equations to more complex quadratic, polynomial, exponential, and periodic curves.



\paragraph{Licensing}



Content obtained and/or adapted from:



\begin{itemize}
    \item \href{https://en.wikibooks.org/wiki/Charts}{Charts, Wikibooks} under a CC BY-SA license
    \item \href{https://en.wikipedia.org/wiki/Line_chart}{Line chart, Wikipedia} under a CC BY-SA license
\end{itemize}

\subsection{Learning Objectives}

\begin{enumerate}
    \item Understand the components and structure of a line graph, including titles, axes, tick marks, and data points.
    \item Interpret trends and patterns from line graphs to analyze the relationship between variables over time.
    \item Apply best-fit techniques to create a line graph that accurately represents data trends, including identifying potential errors and misrepresentations.
\end{enumerate}

\subsection{Assessment Instrument}

\begin{enumerate}
    \item Part 1: Multiple Choice Questions
    \item Which component is essential for identifying the subject of a line graph?
    \item a) Tick Marks
    \item b) Title
    \item c) Data Points
    \item d) Source
    \item Why should axes in a line graph extend slightly past the highest graphed value?
    \item a) To ensure clarity of data points
    \item b) To add aesthetic appeal
    \item c) To fit more data points
    \item d) To confuse the reader
    \item What is a primary disadvantage of using a line graph?
    \item a) Difficulty in creating the graph
    \item b) Difficulty in identifying the exact value at a data point
    \item c) Difficulty in labeling the axes
    \item d) Difficulty in drawing the line Part 2: Short Answer Questions
    \item Explain why a line graph should not be used to show statistics at a set moment in time.
    \item Describe two tips for making an effective line graph. Part 3: Practical Application
    \item Given the following data table, plot a line graph showing the relationship between elapsed time and speed. Label all necessary components, including title, axes, and data points. $$ \begin{array}{|c|c|} \hline \text{Elapsed Time (s)} & \text{Speed (mps)} \\\\ \hline 0 & 0 \\\\ \hline 1 & 3 \\\\ \hline 2 & 7 \\\\ \hline 3 & 12 \\\\ \hline 4 & 18 \\\\ \hline 5 & 30 \\\\ \hline 6 & 45.6 \\\\ \hline \end{array} $$
    \item Using the graph from question 6, add a best-fit line and discuss any trends or patterns you observe. Part 4: Error Analysis
    \item Discuss why discontinuities in the slope of a best-fit line are generally improbable and how experimental error can affect the representation of data on a line graph.
\end{enumerate}

%%===============================================================
\section{Debt-to-income (DTI) ratios}
\label{sec:debt-to-income-dti-ratios}
%%===============================================================

\paragraph{Debt-to-Income (DTI) Ratios}

In the consumer mortgage industry, debt-to-income ratio (often abbreviated DTI) is the percentage of a consumer''''s monthly gross income that goes toward paying debts. (Speaking precisely, DTIs often cover more than just debts; they can include principal, taxes, fees, and insurance premiums as well. Nevertheless, the term is a set phrase that serves as a convenient, well-understood shorthand.) There are two main kinds of DTI, as discussed below.

\paragraph{Two Main Kinds of DTI}

The two main kinds of DTI are expressed as a pair using the notation $x/y$ (for example, $28/36$).

\begin{itemize}
    \item \textbf{Front-end DTI}: Indicates the percentage of income that goes toward housing costs, which for renters is the rent amount and for homeowners is PITI (mortgage principal and interest, mortgage insurance premium [when applicable], hazard insurance premium, property taxes, and homeowners'''' association dues [when applicable]).
\end{itemize}

\begin{itemize}
    \item \textbf{Back-end DTI}: Indicates the percentage of income that goes toward paying all recurring debt payments, including those covered by the first DTI, and other debts such as credit card payments, car loan payments, student loan payments, child support payments, alimony payments, and legal judgments.
\end{itemize}

DTI ratios are used to determine a person’s ability to obtain a mortgage. Most lenders examine front-end DTI and back-end DTI to determine a buyer’s suitability for a mortgage loan.

\begin{itemize}
    \item \textbf{Front-end DTI} = $\frac{\text{Housing payments}}{\text{Total income}}$. Housing refers to rent or mortgage payments.
\end{itemize}

\begin{itemize}
    \item \textbf{Back-end DTI} = $\frac{\text{All recurring debt payments}}{\text{Total income}}$. Recurring debts include rent or mortgage, plus car payments, student loans, credit card payments, etc.
\end{itemize}

According to Bankrate, your front-end DTI should not exceed $28\%$ and your back-end DTI should not exceed $36\%$.

\paragraph{Example}

If the lender requires a debt-to-income ratio of $28/36$, then to qualify a borrower for a mortgage, the lender would go through the following process to determine what expense levels they would accept:

\textbf{Using Yearly Figures:}

\begin{itemize}
    \item Gross Income of `$`45,000
    \item `$`45,000 x 0.28 = `$`12,600 allowed for housing expense.
    \item `$`45,000 x 0.36 = `$`16,200 allowed for housing expense plus recurring debt.
\end{itemize}

\textbf{Using Monthly Figures:}

\begin{itemize}
    \item Gross Income of `$`3,750 (`$`45,000/12)
    \item `$`3,750 x 0.28 = `$`1,050 allowed for housing expense.
    \item `$`3,750 x 0.36 = `$`1,350 allowed for housing expense plus recurring debt.
\end{itemize}

\paragraph{DTI Limits Used in Qualifying Borrowers}

\paragraph{United States}

\paragraph{Current Limits}

\textbf{Conforming Loans:}

\begin{itemize}
    \item Conventional financing limits are typically $28/36$ for manually underwritten loans. The maximum can be exceeded up to $45\%$ if the borrower meets additional credit score and reserve requirements.
    \item FHA limits are currently $31/43$. When using the FHA''''s Energy Efficient Mortgage program, the "stretch ratios" of $33/45$ are used.
    \item VA loan limits are only calculated with one DTI of $41$. (This is effectively equal to $41/41$, although VA does not use that notation.)
    \item USDA $29/41$
\end{itemize}

\textbf{Nonconforming Loans:}

\begin{itemize}
    \item Back ratio limits up to $55\%$ became common for nonconforming loans in the 2000s, as the financial industry experimented with looser credit, with innovative terms and mechanisms, fueled by a real estate bubble. The mortgage business underwent a shift as the traditional mortgage banking industry was shadowed by an infusion of lending from the shadow banking system that eventually rivaled the size of the conventional financing sector. The subprime mortgage crisis produced a market correction that revised these limits downward again for many borrowers, reflecting a predictable tightening of credit after the laxness of the credit bubble. Creative financing (involving riskier ratios) still exists, but nowadays is granted with tighter, more sensible qualification of customers.
\end{itemize}

\paragraph{Historical Limits}

\begin{itemize}
    \item The business of lending and borrowing money has evolved qualitatively in the post–World War II era. It was not until that era that the FHA and the VA (through the G.I. Bill) led the creation of a mass market in 30-year, fixed-rate, amortized mortgages. It was not until the 1970s that the average working person carried credit card balances. Thus the typical DTI limit in use in the 1970s was PITI<25\%, with no codified limit for the second DTI ratio (the one including credit cards). In other words, in today''''s notation, it could be expressed as $25/25$, or perhaps more accurately, $25/\text{NA}$, with the NA limit left to the discretion of lenders on a case-by-case basis. In the following decades these limits gradually climbed higher, and the second limit was codified (coinciding with the evolution of modern credit scoring), as lenders determined empirically how much risk was profitable. This empirical process continues today.
\end{itemize}

\paragraph{Canada}

\begin{itemize}
    \item The Vanier Institute of the Family measures debt to income as total family debt to net income. This is a different ratio, because it compares a cashflow number (yearly after-tax income) to a static number (accumulated debt) - rather than to the debt payment as above. The Institute reported on February 17, 2010 that the average Canadian Family owes `$`100,000, therefore having a debt to net income after taxes of $150\\%$.
\end{itemize}

\paragraph{United Kingdom}

\begin{itemize}
    \item The Bank of England (as of June 26, 2014) implemented a debt to income multiplier on mortgages of $4.5$ (A consumer mortgage can be $4.5$ times the size of annual income), in an attempt to cool rapidly rising house prices. Previously internal standards were relied upon in order to assess the risk of defaults; however, in the wake of the 2008 financial crisis, it was decided that the risk of contagion between housing markets was too great to rely solely on voluntary regulation.
\end{itemize}

\paragraph{Licensing}

Content obtained and/or adapted from:

\begin{itemize}
    \item \href{https://en.wikipedia.org/wiki/Debt-to-income_ratio}{Debt-to-income ratio, Wikipedia} under a CC BY-SA license
\end{itemize}

\subsection{Learning Objectives}

\begin{enumerate}
    \item \textbf{Understand DTI Ratios}: Students will be able to explain the concept of debt-to-income (DTI) ratios, including the definitions and uses of front-end and back-end DTI ratios, and their importance in mortgage qualification.
    \item \textbf{Calculate Front-end and Back-end DTI}: Students will be able to compute both front-end and back-end DTI ratios using given income and debt figures, and interpret these ratios in the context of mortgage approval criteria.
    \item \textbf{Analyze DTI Limits Across Countries}: Students will be able to compare DTI limits used in different countries (United States, Canada, and the United Kingdom), and understand historical changes in these limits.
\end{enumerate}

\subsection{Assessment Instrument}

\begin{enumerate}
    \item \textbf{1. Conceptual Understanding:}
    \item Define debt-to-income (DTI) ratio and explain the difference between front-end and back-end DTI ratios.
    \item Explain why DTI ratios are important for mortgage qualification. \textbf{2. Calculation:}
    \item Given the following annual gross income of `$45,000` and monthly recurring debts of `$1,350`, calculate:
    \item The front-end DTI ratio if the annual housing expense is `$12,600`.
    \item The back-end DTI ratio with the given debts. \textbf{3. Comparative Analysis:}
    \item Using the provided DTI limits for the United States, Canada, and the United Kingdom, answer the following:
    \item Compare the current DTI limits for conforming loans in the United States with historical limits. What changes have occurred over time?
    \item Analyze the DTI limits in Canada and the United Kingdom. How do these limits differ from those in the United States, and what might be the reasons for these differences?
\end{enumerate}

%%===============================================================
\section{Proportional reasoning}
\label{sec:proportional-reasoning}
%%===============================================================

\paragraph{Proportional Reasoning}



\paragraph{Vocabulary}



\begin{itemize}
    \item \textbf{Ratio:} The comparison of two numbers.
\end{itemize}

\begin{itemize}
    \item \textbf{Reciprocal:} The multiplicative inverse of a number -- when the numerator and denominator of a fraction are switched around.
\end{itemize}

\begin{itemize}
    \item \textbf{Equivalent Ratios:} Two ratios that reduce to the same ratio.
\end{itemize}

\begin{itemize}
    \item \textbf{Proportion:} An equation stating two ratios are equal.
\end{itemize}

\begin{itemize}
    \item \textbf{To Cross-multiply:} See below.
\end{itemize}

\paragraph{Lesson}



There are several ways to express a "ratio". Let's compare the number of boys with the number of girls in a particular classroom. Let's say that our classroom has 25 students, 10 of whom are boys. That means there are 15 girls. So, the ratio of boys to girls is $10$ to $15$.



\paragraph{Writing a Ratio}



There are three ways of expressing a ratio. You can simply use words like we did above, or you can separate the two numbers using a colon or a fraction:



\begin{itemize}
    \item $10:15$ or $10/15$ or $\frac{10}{15}$
\end{itemize}

In mathematics, we always use fractions to represent ratios. In this classroom example, many other ratios can be made:



\begin{itemize}
    \item $\frac{10}{25}$, the number of boys out of the total number of students
\end{itemize}

\begin{itemize}
    \item $\frac{15}{25}$, the number of girls out of the total number of students
\end{itemize}

\begin{itemize}
    \item $\frac{15}{10}$, the number of girls to the number of boys
\end{itemize}

\paragraph{Simplifying a Ratio}



Since we're using a fraction to represent ratios, the ratios can sometimes be reduced. For example, the ratio of boys to girls in our hypothetical classroom is $\frac{10}{15}$, but can be reduced to $\frac{2}{3}$. So, it would be correct to say that the ratio of boys to girls in our hypothetical classroom is $2:3$, or two boys for every 3 girls.



\paragraph{Proportions}



The following equation is a proportion:



$$

\frac{10}{15} = \frac{2}{3}

$$



Any proportion is simply an equation that states two ratios are equal to one another. Sometimes, a proportion may contain variables:



$$

\frac{2}{3} = \frac{x}{30}

$$



If we wish to solve such an equation, we can use a process called cross-multiplication.



\paragraph{Solving Proportions using Cross-Multiplication}



If $\frac{a}{b} = \frac{c}{d}$, then the products that are formed by diagonals across the equal sign are also equal:



$$

ad = bc

$$



Note that this would be equivalent to saying:



$$

bc = ad

$$



\paragraph{Example Problems}



1. \textbf{Is the ratio} $\frac{4}{6}$ \textbf{equal to the ratio} $\frac{12}{16}$?



   We can test by cross-multiplying. If the cross-products are equal, then so are the original two ratios.



   $$4 \cdot 16 \stackrel{?}{=} 6 \cdot 12$$



   $$64 \stackrel{?}{=} 72$$



   No, $64 \neq 72$, so $\frac{4}{6} \neq \frac{12}{16}$.



2. \textbf{Solve the following proportion for} $x$:



   $$

   \frac{3}{4} = \frac{x}{10}

   $$



   Cross-multiply:



   $$

   3 \cdot 10 = 4x

   $$



   $$

   30 = 4x

   $$



   Solve for $x$ by dividing both sides of the equation by $4$. Your result:



   $$

   x = \frac{30}{4} = 7.5

   $$



3. \textbf{Solve for} $x$:



   $$

   \frac{x+1}{4} = \frac{3}{8}

   $$



   Cross multiply:



   $$

   8 \cdot (x+1) = 3 \cdot 4

   $$



   $$

   8x + 8 = 12

   $$



   Now, solve for $x$:



   Subtract $8$ from both sides:



   $$

   8x = 4

   $$



   Divide both sides by $8$:



   $$

   x = \frac{1}{2}

   $$



4. \textbf{A 9th-grade algebra classroom has a ratio of boys to girls of} $\frac{1}{2}$. \textbf{If there are 14 girls, how many boys are in the class?}



   We can create a proportion that compares the ratios of boys to girls, where $x$ is the number of boys in the class:



   $$

   \frac{1}{2} = \frac{x}{14}

   $$



   Solve the proportion by cross-multiplying:



   $$

   2x = 14

   $$



   $$

   x = 7

   $$



   There are $7$ boys in the class.



\paragraph{Licensing}



Content obtained and/or adapted from:



\begin{itemize}
    \item \href{https://en.wikibooks.org/wiki/Proportions}{Proportions, Wikibooks} under a CC BY-SA license
\end{itemize}

\subsection{Learning Objectives}

\begin{enumerate}
    \item 1.\textbf{Identify and Explain Ratios:}
    \item Students will be able to identify and explain the concept of ratios, including how to express them using words, colons, or fractions. 2.\textbf{Simplify Ratios:}
    \item Students will be able to simplify ratios by reducing fractions to their lowest terms and interpreting equivalent ratios in various contexts. 3.\textbf{Solve Proportions:}
    \item Students will be able to solve proportions by applying cross-multiplication and solving for the unknown variable in a proportion equation.
\end{enumerate}

\subsection{Assessment Instrument}

\begin{enumerate}
    \item Instructions: For each question, solve the proportion or simplify the ratio as directed. Show all work clearly. 1.\textbf{Determine if the ratios are equivalent:}
    \item Are the ratios $\frac{6}{8}$ and $\frac{9}{12}$ equivalent? Use cross-multiplication to check. 2.\textbf{Solve the proportion for $x$:}
    \item Solve the proportion $\frac{5}{7} = \frac{x}{14}$. 3.\textbf{Simplify the ratio:}
    \item Simplify the ratio $\frac{18}{24}$ to its lowest terms and express it in the form $a:b$. 4.\textbf{Word Problem:}
    \item In a recipe, the ratio of flour to sugar is $\frac{3}{4}$. If you use 12 cups of flour, how many cups of sugar are needed? Set up and solve the proportion to find the answer.
\end{enumerate}

%%===============================================================
\section{Mathematical (Linear) relationships}
\label{sec:mathematical-linear-relationships}
%%===============================================================

\paragraph{Mathematical (Linear) Relationships}



We are going to start by looking at simple functions called linear equations. When none of the instances of $x$ and $y$ in the algebraic expression defining the function rule have exponents, then all the instances of $x$ and $y$ can be combined into just two occurrences. The graph of the expression can be represented as a straight line. The equation that expresses the function is considered a linear equation with two variables. The following equation is a simple example of such a linear equation:



$$

y - x = 2

$$



Since $y$ is the dependent variable, it is standing in for the function. We can re-write the expression as $f(x) - x = 2$. If we add an $x$ to both sides the equality property holds and we get the expression $f(x) - x + x = x + 2$. Simplifying, we get $f(x) = x + 2$. In the following table, we''ll pick 3 values for $x$, and then calculate the dependent ($y$) values from $f(x)$.



$$

\begin{array}{|c|c|c|}

\hline

\text{x value} & \text{y value (abscissa)} & \text{Coordinates (x, y)} \\\\ \hline

-1 & 1 & (-1, 1) \\\\ \hline

0 & 2 & (0, 2) \\\\ \hline

1 & 3 & (1, 3) \\\\ \hline

\end{array}

$$





This function is equivalent to the previous example of a linear equation, $y - x = 2$. The arrows at each end of the line indicate that the line extends infinitely in both directions. All linear functions of a single input variable have or can be algebraically arranged to have the general form:



$$

y = f(x) = mx + b

$$



where $x$ and $y$ are variables, $f(x)$ is the function of $x$, $m$ is a constant called the slope of the line, and $b$ is a constant which is the ordinate of the $y$-intercept (i.e., the value of $y$ where the function line crosses the $y$-axis). The slope indicates the steepness of the line. In the previous example where $y = x + 2$, the slope $m = 1$ and the $y$-intercept ordinate $b = 2$. The $y = mx + b$ form of a linear function is called the slope-intercept form.



Unless a domain for $x$ is otherwise stated, the domain for linear functions will be assumed to be all real numbers and so the lines in graphs of all linear functions extend infinitely in both directions. Also, in linear functions with all real number domains, the range of a linear function will cover the entire set of real numbers for $y$, unless the slope $m = 0$ and the function equals a constant. In such cases, the range is simply the constant.



Conversely, if a function has the general form $y = mx + b$ or if it can be arranged to have that form, the function is linear. A linear equation with two variables has or can be algebraically rearranged to have the general form:



$$

Ax + By = C

$$



where $x$ and $y$ represent the linear variables, and the letters $A$, $B$, and $C$ can represent any real constants, either positive or negative. Conversely, an equation with two variables $x$ and $y$ having that general form, or being able to be arranged in that form, would be linear as long as $A$ and $B$ are not both equal to 0. In the preceding equation, capital letters are used to avoid confusion with other constants in this chapter and for consistency.



If one divides the preceding equation by $B$ (when $B \neq 0$) and solves for $y$, the following form can be obtained:



$$

y = \left(-\frac{A}{B}\right)x + \frac{C}{B}

$$



If one equates $-\frac{A}{B}$ to the slope $m$ and $\frac{C}{B}$ to the $y$-intercept ordinate $b$, it can be seen that the general form for a linear equation and the slope-intercept form for a linear function are practically interconvertible except for the fact that, in a linear function, the $B$ constant in the linear equation form cannot equal 0.



\paragraph{General Equation Forms of a Line}



The most general form applicable to all lines on a two-dimensional Cartesian graph is:



$$

Ax + By = C

$$



with three constants, $A$, $B$, and $C$. These constants are not unique to the line because multiplying the whole equation by a constant factor gives a new set of valid constants for the same line. When $B = 0$, the rest of the equation represents a vertical line, which is not a function. If $B \neq 0$, then the line is a function. Such a linear function can be represented by the slope-intercept form which has two constants.



\textbf{Slope-intercept form:}



$$

y = mx + b

$$



The two constants, $m$ and $b$, used together are unique to the line. In other words, a certain line can have only one pair of values for $m$ and $b$ in this form.



\textbf{The point-slope form given here:}



$$

y - y\_1 = m(x - x\_1)

$$



uses three constants; $m$ is unique for a given line; $x_1$ and $y_1$ are not unique and can be from any point on the line. The point-slope form cannot represent a vertical line.



\textbf{The intercept form of a line, given here:}



$$

\frac{x}{a} + \frac{y}{b} = 1 \quad (a \neq 0 \text{ and } b \neq 0)

$$



uses two unique constants which are the $x$ and $y$ intercepts but cannot be made to represent horizontal or vertical lines or lines crossing through $(0,0)$. It is the least applicable of the general forms in this summary.



Of the last three general forms of a linear function, the slope-intercept form is the most useful because it uses only constants unique to a given line and can represent any linear function. All of the problems in this book and in mathematics in general can be solved without using the point-slope form or the intercept form unless they are specifically called for in a problem. Generally, problems involving linear functions can be solved using the slope-intercept form $(y = mx + b)$ and the formula for slope.



\paragraph{Licensing}



Content obtained and/or adapted from:



\begin{itemize}
    \item \href{https://en.wikibooks.org/wiki/Linear_Equations_and_Functions}{Linear Equations and Functions, Wikibooks} under a CC BY-SA license
    \item \href{https://en.wikibooks.org/wiki/Function_Graphing}{Function Graphing, Wikibooks} under a CC BY-SA license
\end{itemize}

\subsection{Learning Objectives}

\begin{enumerate}
    \item \textbf{Identify and represent linear equations:} Students will identify linear equations and represent them in different forms, including $y = mx + b$ and $Ax + By = C$.
    \item \textbf{Calculate slope and intercepts:} Students will calculate the slope $m$ and $y$-intercept $b$ from the equation of a line and use these to graph the line on a Cartesian plane.
    \item \textbf{Convert between forms:} Students will convert linear equations between slope-intercept form ($y = mx + b$) and general form ($Ax + By = C$) and interpret the resulting graph.
\end{enumerate}

\subsection{Assessment Instrument}

\begin{enumerate}
    \item \textbf{1. Multiple Choice Questions}
    \item \textbf{Question 1:} Given the linear equation $3x - 4y = 12$, what is the slope of the line?
    \item A) $3$
    \item B) $-3/4$
    \item C) $4/3$
    \item D) $-3$
    \item \textbf{Question 2:} Which of the following is the slope-intercept form of the equation $2x + 3y = 6$?
    \item A) $y = -2x + 6$
    \item B) $y = \frac{-2}{3}x + 2$
    \item C) $y = \frac{2}{3}x - 2$
    \item D) $y = \frac{-2}{3}x + 2$ \textbf{2. Short Answer Questions}
    \item \textbf{Question 3:} Convert the linear equation $4x - 5y = 20$ to slope-intercept form and identify the slope and $y$-intercept.
    \item \textbf{Question 4:} Graph the equation $y = -x + 4$ on a Cartesian plane and identify the $x$-intercept and $y$-intercept. \textbf{3. Problem Solving}
    \item \textbf{Question 5:} A line passes through the points $(1, 2)$ and $(3, 4)$. Determine the slope of the line and write its equation in slope-intercept form.
    \item \textbf{Question 6:} Convert the following equation from general form to slope-intercept form and describe the graph: $6x + 2y = 8$.
\end{enumerate}

%%===============================================================
\section{Proportionality vs. Linearity}
\label{sec:proportionality-vs-linearity}
%%===============================================================

\paragraph{Proportionality vs. Linearity}





\paragraph{Proportionality}



In mathematics, two varying quantities are said to be in a relation of proportionality if their ratio or product yields a constant. This constant is called the coefficient of proportionality or proportionality constant.



If the ratio $\frac{y}{x}$ of two variables ($x$ and $y$) is equal to a constant ($k = \frac{y}{x}$), then the variable in the numerator of the ratio ($y$) can be expressed as the product of the other variable and the constant ($y = k \cdot x$). In this case, $y$ is said to be directly proportional to $x$ with proportionality constant $k$. Equivalently, one may write $x = \frac{1}{k} \cdot y$; that is, $x$ is directly proportional to $y$ with proportionality constant $\frac{1}{k} = \frac{x}{y}$. If the term proportional is connected to two variables without further qualification, generally direct proportionality can be assumed.



If the product of two variables ($x \cdot y$) is equal to a constant ($k = x \cdot y$), then the two are said to be inversely proportional to each other with the proportionality constant $k$. Equivalently, both variables are directly proportional to the reciprocal of the respective other with proportionality constant $k$ ($x = k \cdot \frac{1}{y}$ and $y = k \cdot \frac{1}{x}$).



If several pairs of variables share the same direct proportionality constant, the equation expressing the equality of these ratios is called a proportion, e.g., $\frac{a}{b} = \frac{x}{y} = \cdots = k$ (for details see Ratio).



\paragraph{Direct Proportionality}





\textbf{The variable $y$ is directly proportional to the variable $x$ with proportionality constant \textasciitilde{}$0.6$.}

% [Image: direct]





Given two variables $x$ and $y$, $y$ is directly proportional to $x$ if there is a non-zero constant $k$ such that



$$

y = kx

$$



The relation is often denoted using the symbols "$\propto$" (not to be confused with the Greek letter alpha) or "$\sim$":



$$

y \propto x \text{ or } y \sim x

$$



For $x \neq 0$, the proportionality constant can be expressed as the ratio



$$

k = \frac{y}{x}

$$



It is also called the \textbf{constant of variation} or \textbf{constant of proportionality}.



A direct proportionality can also be viewed as a linear equation in two variables with a $y$-intercept of $0$ and a slope of $k$. This corresponds to linear growth.



\paragraph{Examples}



\begin{itemize}
    \item If an object travels at a constant speed, then the distance traveled is directly proportional to the time spent traveling, with the speed being the constant of proportionality.
    \item The circumference of a circle is directly proportional to its diameter, with the constant of proportionality equal to $\pi$.
    \item On a map of a sufficiently small geographical area, drawn to scale distances, the distance between any two points on the map is directly proportional to the beeline distance between the two locations represented by those points; the constant of proportionality is the scale of the map.
    \item The force acting on a small object with small mass by a nearby large extended mass due to gravity is directly proportional to the object''s mass; the constant of proportionality between the force and the mass is known as gravitational acceleration.
    \item The net force acting on an object is proportional to the acceleration of that object with respect to an inertial frame of reference. The constant of proportionality in this, Newton''s second law, is the classical mass of the object.
\end{itemize}

\paragraph{Inverse proportionality}



\textbf{Inverse proportionality with a function of y = 1/x}

% [Image: inverse]







The concept of inverse proportionality can be contrasted with direct proportionality. Consider two variables said to be "inversely proportional" to each other. If all other variables are held constant, the magnitude or absolute value of one inversely proportional variable decreases if the other variable increases, while their product (the constant of proportionality $k$) is always the same. As an example, the time taken for a journey is inversely proportional to the speed of travel.



Formally, two variables are inversely proportional (also called varying inversely, in inverse variation, in inverse proportion, in reciprocal proportion) if each of the variables is directly proportional to the multiplicative inverse (reciprocal) of the other, or equivalently if their product is a constant. It follows that the variable $y$ is inversely proportional to the variable $x$ if there exists a non-zero constant $k$ such that



$$

y = \frac{k}{x}

$$



or equivalently,



$$

xy = k

$$



Hence the constant "$k$" is the product of $x$ and $y$.



The graph of two variables varying inversely on the Cartesian coordinate plane is a rectangular hyperbola. The product of the $x$ and $y$ values of each point on the curve equals the constant of proportionality ($k$). Since neither $x$ nor $y$ can equal zero (because $k$ is non-zero), the graph never crosses either axis.



\paragraph{Hyperbolic Coordinates}



The concepts of direct and inverse proportion lead to the location of points in the Cartesian plane by hyperbolic coordinates; the two coordinates correspond to the constant of direct proportionality that specifies a point as being on a particular ray and the constant of inverse proportionality that specifies a point as being on a particular hyperbola.



\paragraph{Linearity}



Linearity is the property of a mathematical relationship (function) that can be graphically represented as a straight line. Linearity is closely related to proportionality. Examples in physics include the linear relationship of voltage and current in an electrical conductor (Ohm''s law), and the relationship of mass and weight. By contrast, more complicated relationships are nonlinear.



Generalized for functions in more than one dimension, linearity means the property of a function of being compatible with addition and scaling, also known as the superposition principle.



The word linear comes from Latin \textit{linearis}, "pertaining to or resembling a line".



\paragraph{In Mathematics}



In mathematics, a linear map or linear function $f(x)$ is a function that satisfies the two properties:



\begin{itemize}
    \item \textbf{Additivity:} $f(x + y) = f(x) + f(y)$.
    \item \textbf{Homogeneity of degree 1:} $f(\alpha x) = \alpha f(x)$ for all $\alpha$.
\end{itemize}

These properties are known as the superposition principle. In this definition, $x$ is not necessarily a real number, but can in general be an element of any vector space. A more special definition of linear function, not coinciding with the definition of linear map, is used in elementary mathematics (see below).



Additivity alone implies homogeneity for rational $\alpha$, since $f(x + x) = f(x) + f(x)$ implies $f(nx) = nf(x)$ for any natural number $n$ by mathematical induction, and then $nf(x) = f(nx) = f\left(m \frac{n}{m} x\right) = mf\left(\frac{n}{m} x\right)$ implies $f\left(\frac{n}{m} x\right) = \frac{n}{m} f(x)$.  The density of the rational numbers in the reals implies that any additive continuous function is homogeneous for any real number $\alpha$, and is therefore linear.



The concept of linearity can be extended to linear operators. Important examples of linear operators include the derivative considered as a differential operator, and other operators constructed from it, such as del and the Laplacian. When a differential equation can be expressed in linear form, it can generally be solved by breaking the equation up into smaller pieces, solving each of those pieces, and summing the solutions.



Linear algebra is the branch of mathematics concerned with the study of vectors, vector spaces (also called ''linear spaces''), linear transformations (also called ''linear maps''), and systems of linear equations.



For a description of linear and nonlinear equations, see linear equation.



\paragraph{Linear Polynomials}



In a different usage to the above definition, a polynomial of degree $1$ is said to be linear, because the graph of a function of that form is a straight line.



Over the reals, a linear equation is one of the forms:



$$

f(x) = mx + b

$$



where $m$ is often called the slope or gradient; $b$ the $y$-intercept, which gives the point of intersection between the graph of the function and the $y$-axis.



Note that this usage of the term linear is not the same as in the section above, because linear polynomials over the real numbers do not in general satisfy either additivity or homogeneity. In fact, they do so if and only if $b = 0$. Hence, if $b \neq 0$, the function is often called an affine function (see in greater generality affine transformation).



\paragraph{Boolean Functions}



In Boolean algebra, a linear function is a function $f$ for which there exist $a\_0, a\_1, \ldots, a\_n \in \{0, 1\}$ such that



$$

f(b\_1, \ldots, b\_n) = a\_0 \oplus (a\_1 \land b\_1) \oplus \cdots \oplus (a\_n \land b\_n)

$$



where $b\_1, \ldots, b\_n \in \{0, 1\}$.



Note that if $a\_0 = 1$, the above function is considered affine in linear algebra (i.e., not linear).



A Boolean function is linear if one of the following holds for the function''s truth table:



\textbf{Hasse diagram of a linear Boolean function}



% [Image: hasse]



\begin{itemize}
    \item In every row in which the truth value of the function is $T$, there are an odd number of $T$s assigned to the arguments, and in every row in which the function is $F$ there is an even number of $T$s assigned to arguments. Specifically, $f(F, F, \ldots, F) = F$, and these functions correspond to linear maps over the Boolean vector space.
    \item In every row in which the value of the function is $T$, there is an even number of $T$s assigned to the arguments of the function; and in every row in which the truth value of the function is $F$, there are an odd number of $T$s assigned to arguments. In this case, $f(F, F, \ldots, F) = T$.
\end{itemize}

Another way to express this is that each variable always makes a difference in the truth value of the operation or it never makes a difference.



Negation, logical biconditional, exclusive or, tautology, and contradiction are linear functions.



\paragraph{Licensing}



Content obtained and/or adapted from:



\begin{itemize}
    \item \href{https://en.wikipedia.org/wiki/Proportionality_(mathematics}{Proportionality (mathematics), Wikipedia}) under a CC BY-SA license
    \item \href{https://en.wikipedia.org/wiki/Linearity}{Linearity, Wikipedia} under a CC BY-SA license
\end{itemize}

\subsection{Learning Objectives}

\begin{enumerate}
    \item 1.\textbf{Understand Proportionality}
    \item Students will be able to identify and apply the concept of direct proportionality by solving equations of the form $y = kx$ and interpreting the proportionality constant $k$. 2.\textbf{Analyze Inverse Proportionality}
    \item Students will be able to recognize and solve problems involving inverse proportionality, using the equation $y = \frac{k}{x}$ and understanding the relationship between variables and the constant $k$. 3.\textbf{Apply Linearity in Functions}
    \item Students will demonstrate the ability to apply the properties of linear functions, including additivity and homogeneity, to solve problems and graph linear equations of the form $f(x) = mx + b$.
\end{enumerate}

\subsection{Assessment Instrument}

\begin{enumerate}
    \item Instructions: For each question, provide a detailed solution, including any necessary steps and interpretations. 1.\textbf{Direct Proportionality Problem}
    \item Given that $y$ is directly proportional to $x$ with a proportionality constant of $0.8$, find the value of $y$ when $x = 5$. Express your answer in the form $y = kx$ and calculate $y$. 2.\textbf{Inverse Proportionality Problem}
    \item If $y$ is inversely proportional to $x$ with a constant of proportionality $k = 12$, determine the value of $y$ when $x = 4$. Solve for $y$ using the equation $y = \frac{k}{x}$. 3.\textbf{Linearity in Functions}
    \item Consider the linear function $f(x) = 3x - 7$. Verify if the function satisfies the properties of linearity by demonstrating additivity ($f(x + y) = f(x) + f(y)$) and homogeneity ($f(\alpha x) = \alpha f(x)$) with appropriate examples. 4.\textbf{Graphing Linear Functions}
    \item Graph the function $f(x) = -2x + 5$ and identify the slope and $y$-intercept. Explain how the graph illustrates the concept of linearity and proportionality in the context of this function.
\end{enumerate}

%%===============================================================
\section{Simple and Compound Interest (Linear and Exponential Models)}
\label{sec:simple-and-compound-interest-linear-and-exponential-models}
%%===============================================================

\paragraph{Simple and Compound Interest (Linear and Exponential Models)}

Interest, in finance and economics, is payment from a borrower or deposit-taking financial institution to a lender or depositor of an amount above repayment of the principal sum (that is, the amount borrowed), at a particular rate. It is distinct from a fee which the borrower may pay the lender or some third party. It is also distinct from dividend which is paid by a company to its shareholders (owners) from its profit or reserve, but not at a particular rate decided beforehand, rather on a pro rata basis as a share in the reward gained by risk-taking entrepreneurs when the revenue earned exceeds the total costs.

For example, a customer would usually pay interest to borrow from a bank, so they pay the bank an amount which is more than the amount they borrowed; or a customer may earn interest on their savings, and so they may withdraw more than they originally deposited. In the case of savings, the customer is the lender, and the bank plays the role of the borrower.

Interest differs from profit, in that interest is received by a lender, whereas profit is received by the owner of an asset, investment or enterprise. (Interest may be part or the whole of the profit on an investment, but the two concepts are distinct from each other from an accounting perspective.)

The rate of interest is equal to the interest amount paid or received over a particular period divided by the principal sum borrowed or lent (usually expressed as a percentage).

Compound interest means that interest is earned on prior interest in addition to the principal. Due to compounding, the total amount of debt grows exponentially, and its mathematical study led to the discovery of the number $e$. In practice, interest is most often calculated on a daily, monthly, or yearly basis, and its impact is influenced greatly by its compounding rate.

\paragraph{Economics}

In economics, the rate of interest is the price of credit, and it plays the role of the cost of capital. In a free market economy, interest rates are subject to the law of supply and demand of the money supply, and one explanation of the tendency of interest rates to be generally greater than zero is the scarcity of loanable funds.

Over centuries, various schools of thought have developed explanations of interest and interest rates. The School of Salamanca justified paying interest in terms of the benefit to the borrower, and interest received by the lender in terms of a premium for the risk of default. In the sixteenth century, Martín de Azpilcueta applied a time preference argument: it is preferable to receive a given good now rather than in the future. Accordingly, interest is compensation for the time the lender forgoes the benefit of spending the money.

On the question of why interest rates are normally greater than zero, in 1770, French economist Anne-Robert-Jacques Turgot, Baron de Laune proposed the theory of fructification. By applying an opportunity cost argument, comparing the loan rate with the rate of return on agricultural land, and a mathematical argument, applying the formula for the value of a perpetuity to a plantation, he argued that the land value would rise without limit, as the interest rate approached zero. For the land value to remain positive and finite keeps the interest rate above zero.

Adam Smith, Carl Menger, and Frédéric Bastiat also propounded theories of interest rates. In the late 19th century, Swedish economist Knut Wicksell in his 1898 Interest and Prices elaborated a comprehensive theory of economic crises based upon a distinction between natural and nominal interest rates. In the 1930s, Wicksell''s approach was refined by Bertil Ohlin and Dennis Robertson and became known as the loanable funds theory. Other notable interest rate theories of the period are those of Irving Fisher and John Maynard Keynes.

\paragraph{Calculation}

\paragraph{Simple Interest}

Simple interest is calculated only on the principal amount, or on that portion of the principal amount that remains. It excludes the effect of compounding. Simple interest can be applied over a time period other than a year, for example, every month.

Simple interest is calculated according to the following formula:

$$
\frac{r \cdot B \cdot m}{n}
$$

where

\begin{itemize}
    \item $r$ is the simple annual interest rate
    \item $B$ is the initial balance
    \item $m$ is the number of time periods elapsed
    \item $n$ is the frequency of applying interest.
\end{itemize}

For example, imagine that a credit card holder has an outstanding balance of `$`2500 and that the simple annual interest rate is 12.99\% per annum, applied monthly, so the frequency of applying interest is 12 per year. Over one month,

$$
\frac{0.1299 \times \$2500}{12} = \$27.06
$$

interest is due (rounded to the nearest cent).

Simple interest applied over 3 months would be

$$
\frac{0.1299 \times \$2500 \times 3}{12} = \$81.19
$$

If the card holder pays off only interest at the end of each of the 3 months, the total amount of interest paid would be

$$
\frac{0.1299 \times \$2500}{12} \times 3 = \$27.06 \text{ per month} \times 3 \text{ months} = \$81.18
$$

which is the simple interest applied over 3 months, as calculated above. (The one cent difference arises due to rounding to the nearest cent.)

\paragraph{Compound Interest}

Compound interest includes interest earned on the interest that was previously accumulated.

Compare, for example, a bond paying 6 percent semiannually (that is, coupons of 3 percent twice a year) with a certificate of deposit (GIC) that pays 6 percent interest once a year. The total interest payment is `$`6 per `$`100 par value in both cases, but the holder of the semiannual bond receives half the `$`6 per year after only 6 months (time preference), and so has the opportunity to reinvest the first `$`3 coupon payment after the first 6 months, and earn additional interest.

For example, suppose an investor buys `$`10,000 par value of a US dollar bond, which pays coupons twice a year, and that the bond''s simple annual coupon rate is 6 percent per year. This means that every 6 months, the issuer pays the holder of the bond a coupon of `$`3 per `$`100 par value. At the end of 6 months, the issuer pays the holder:

$$
\frac{r \cdot B \cdot m}{n} = \frac{6\% \times \$10{,}000 \times 1}{2} = \$300
$$

Assuming the market price of the bond is 100, so it is trading at par value, suppose further that the holder immediately reinvests the coupon by spending it on another `$`300 par value of the bond. In total, the investor therefore now holds:

$$
\$10{,}000 + \$300 = \left(1 + \frac{r}{n}\right) \cdot B = \left(1 + \frac{6\%}{2}\right) \times \$10{,}000
$$

and so earns a coupon at the end of the next 6 months of:

$$
\frac{r \cdot B \cdot m}{n} = \frac{6\% \times (\$10{,}000 + \$300)}{2} = \frac{6\% \times \left(1 + \frac{6\%}{2}\right) \times \$10{,}000}{2} = \$309
$$

Assuming the bond remains priced at par, the investor accumulates at the end of a full 12 months a total value of:

$$
\$10{,}000 + \$300 + \$309 = \$10{,}000 + \frac{6\% \times \$10{,}000}{2} + \frac{6\% \times \left(1 + \frac{6\%}{2}\right) \times \$10{,}000}{2} = \$10{,}000 \times \left(1 + \frac{6\%}{2}\right)^{2}
$$

and the investor earned in total:

$$
\$10{,}000 \times \left(1 + \frac{6\%}{2}\right)^{2} - \$10{,}000 = \$10{,}000 \times \left(\left(1 + \frac{6\%}{2}\right)^{2} - 1\right)
$$

The formula for the annual equivalent compound interest rate is:

$$
\left(1 + \frac{r}{n}\right)^{n} - 1
$$

where

\begin{itemize}
    \item $r$ is the simple annual rate of interest
    \item $n$ is the frequency of applying interest.
\end{itemize}

For example, in the case of a 6\% simple annual rate, the annual equivalent compound rate is:

$$
\left(1 + \frac{6\%}{2}\right)^{2} - 1 = 1.03^{2} - 1 = 6.09\%
$$

\paragraph{Licensing}

Content obtained and/or adapted from:

\begin{itemize}
    \item \href{https://en.wikipedia.org/wiki/Interest}{Interest, Wikipedia} under a CC BY-SA license
\end{itemize}

\subsection{Learning Objectives}

\begin{enumerate}
    \item \textbf{Calculate Simple Interest} Students will be able to compute simple interest using the formula: $$ \frac{r \cdot B \cdot m}{n} $$ where $r$ is the annual interest rate, $B$ is the initial balance, $m$ is the number of periods elapsed, and $n$ is the frequency of applying interest.
    \item \textbf{Determine Compound Interest} Students will be able to determine the total amount and interest earned with compound interest, using the formula: $$ A = B \left(1 + \frac{r}{n}\right)^{nt} $$ where $A$ is the amount, $B$ is the initial principal, $r$ is the annual interest rate, $n$ is the number of compounding periods per year, and $t$ is the number of years.
    \item \textbf{Compare Simple and Compound Interest} Students will be able to compare the impact of simple versus compound interest on financial growth over a given time period and interest rate.
\end{enumerate}

\subsection{Assessment Instrument}

\begin{enumerate}
    \item Questions
    \item \textbf{Simple Interest Calculation} Given a principal amount of `$`5000, an annual interest rate of 5\%, and a time period of 3 years, calculate the simple interest.
    \item \textbf{Compound Interest Calculation} If you invest `$`2000 at an annual interest rate of 4\% compounded quarterly for 2 years, calculate the total amount.
    \item \textbf{Comparison of Interest Types} Compare the total interest earned from a `$`1000 investment at a 6\% annual interest rate over 2 years, compounded annually versus semiannually.
\end{enumerate}

%%===============================================================
\section{Regression}
\label{sec:regression}
%%===============================================================

\paragraph{Regression}



Regression analysis is the process of building a model of the relationship between variables in the form of mathematical equations. The general purpose is to explain how one variable, the dependent variable, is systematically related to the values of one or more independent variables. The independent variable is so called because we imagine its value varying freely across its range while the dependent variable is dependent upon the values taken by the independent. The mathematical function is expressed in terms of a number of parameters that are the coefficients of the equation, and the values of the independent variable. The coefficients are numeric constants by which variable values in the equation are multiplied or which are added to a variable value to determine the unknown. A simple example is the equation for the line:



$$

y = mx + b

$$



Here, by convention, $x$ and $y$ are the variables of interest in our data with $y$ the unknown or dependent variable and $x$ the known or independent variable. The constant $m$ is the slope of the line and $b$ is the $y$-intercept - the value where the line crosses the $y$-axis. So, $m$ and $b$ are the coefficients of the equation.



If we can build a robust regression model of known data, then we can use the equation to predict the values for unobserved cases. Regression also involves the estimation of the strength of the association between the dependent and independent variables, most often through the computation of the correlation coefficient which, as noted above, is itself part of a linear model of the data. The correlation coefficient is squared in reporting a regression analysis and we call it, perhaps rather obviously, $R^2$.



Models are often only approximately like the observed data. Most often, some error in the data will mean that no mathematical function produces exactly the data observed and only that data. Therefore, we are explicitly involved in estimation and our models involve recognition of error. It is for this reason that the equation for the line is often given as:



$$

\hat{y} = mx + b \pm \epsilon

$$



Where $\epsilon$ quantifies the error in our data.



\paragraph{Linear Regression}



In linear regression, the model consists of linear equations. A linear equation is an equation involving only constants or single variable values multiplied by a constant. The variable values in a linear equation must be of the first power, that is, they cannot involve raising a value to a power other than one—they cannot, for example, be squared or cubed. Any value raised to the power of one is just the original value: $x^1 = x$.



To carry out a linear regression, first make a scatter plot of your data.



Let us look again at the following measurements of baby height against age:



$$

\begin{array}{|c|c|}

\hline

\text{Age (months)} & \text{Height (cm)} \\\\ \hline

0 & 53.0 \\\\ \hline

3 & 59.5 \\\\ \hline

6 & 66.0 \\\\ \hline

9 & 71.5 \\\\ \hline

12 & 76.0 \\\\ \hline

18 & 85.0 \\\\ \hline

24 & 90.0 \\\\ \hline

\end{array}

$$





This data can be visualized in this scatter plot:



% [Image: scatter]



\paragraph{Licensing}



Content obtained and/or adapted from:



\begin{itemize}
    \item \href{https://en.wikibooks.org/wiki/Regression}{Regression, Wikibooks} under a CC BY-SA license
\end{itemize}

\subsection{Learning Objectives}

\begin{enumerate}
    \item 1.\textbf{Identify and interpret the components of a linear regression model.}
    \item Students will be able to identify the dependent variable ($y$), independent variable ($x$), slope ($m$), and y-intercept ($b$) in the linear regression equation $y = mx + b$. 2.\textbf{Calculate and interpret the correlation coefficient ($R^2$) and the error term ($\epsilon$) in regression analysis.}
    \item Students will compute the correlation coefficient to assess the strength of the relationship between variables and explain how the error term quantifies discrepancies in predictions. 3.\textbf{Construct and analyze a scatter plot to perform linear regression.}
    \item Students will create scatter plots from given data, fit a linear regression line, and interpret the results to make predictions based on the model.
\end{enumerate}

\subsection{Assessment Instrument}

\begin{enumerate}
    \item Part 1: Multiple Choice 1.What does the slope ($m$) in the linear regression equation $y = mx + b$ represent?
    \item A. The point where the line crosses the y-axis
    \item B. The rate of change of $y$ with respect to $x$
    \item C. The total error in the regression model
    \item D. The correlation between $x$ and $y$ 2.In a regression model, what does $R^2$ indicate?
    \item A. The intercept of the regression line
    \item B. The amount of variance in the dependent variable explained by the independent variable
    \item C. The exact value of the error term
    \item D. The slope of the regression line Part 2: Calculation
    \item Given the linear regression model $y = 2x + 5$, calculate the predicted value of $y$ when $x = 7$.
    \item Given the scatter plot data: $$ \begin{array}{|c|c|} \hline \text{Age (months)} & \text{Height (cm)} \\\\ \hline 0 & 53.0 \\\\ \hline 3 & 59.5 \\\\ \hline 6 & 66.0 \\\\ \hline 9 & 71.5 \\\\ \hline 12 & 76.0 \\\\ \hline 18 & 85.0 \\\\ \hline 24 & 90.0 \\\\ \hline \end{array} $$ Perform a linear regression to determine the equation of the line. Part 3: Interpretation
    \item Explain what the error term ($\epsilon$) represents in the regression equation $\hat{y} = mx + b \pm \epsilon$.
    \item Interpret the slope and y-intercept in the context of the baby height data provided. Part 4: Scatter Plot
    \item Construct a scatter plot using the provided data, fit a linear regression line, and describe the relationship between age and height.
\end{enumerate}

%%===============================================================
\section{Piecewise Linear Function}
\label{sec:piecewise-linear-function}
%%===============================================================

\paragraph{Piecewise Linear Function}



In mathematics and statistics, a piecewise linear, PL, or segmented function is a real-valued function of a real variable, whose graph is composed of straight-line segments.





\paragraph{Definition}



A piecewise linear function is a function defined on a (possibly unbounded) interval of real numbers, such that there is a collection of intervals on each of which the function is an affine function. If the domain of the function is compact, there needs to be a finite collection of such intervals; if the domain is not compact, it may either be required to be finite or to be locally finite in the reals.



\paragraph{Examples}



\paragraph{A continuous piecewise linear function}

The function defined by

$$

f(x) = \begin{cases} 

-x - 3 & \text{if } x \leq -3 \\\\ x + 3 & \text{if } -3 < x < 0 \\\\ -2x + 3 & \text{if } 0 \leq x < 3 \\\\ 0.5x - 4.5 & \text{if } x \geq 3 

\end{cases}

$$

is piecewise linear with four pieces. The graph of this function is shown below. Since the graph of a linear function is a line, the graph of a piecewise linear function consists of line segments and rays. The $x$ values (in the above example $-3$, $0$, and $3$) where the slope changes are typically called breakpoints, changepoints, threshold values, or knots. As in many applications, this function is also continuous. The graph of a continuous piecewise linear function on a compact interval is a polygonal chain.



% [Image: linear]



A piecewise function is a function that is given by different expressions on different intervals. The graph of 

$$

f(x) = \begin{cases}

x^2 & \text{if } x < 2 \\\\ -(x - 3)^2 + 9 & \text{if } x \geq 2

\end{cases}

$$

is shown below. The open circle indicates that the point is not included on the graph and the closed circle indicates that the point is included on the graph. So $f(2) = 8$ not $4$.





% [Image: Upper semi]



Other examples of piecewise linear functions include the absolute value function, the sawtooth function, and the floor function.



\paragraph{Fitting to a Curve}



\textbf{A function (blue) and a piecewise linear approximation to it (red)}



% [Image: fit]





An approximation to a known curve can be found by sampling the curve and interpolating linearly between the points. An algorithm for computing the most significant points subject to a given error tolerance has been published.



\paragraph{Fitting to Data}



If partitions, and then breakpoints, are already known, linear regression can be performed independently on these partitions. However, continuity is not preserved in that case, and also there is no unique reference model underlying the observed data. A stable algorithm with this case has been derived.



If partitions are not known, the residual sum of squares can be used to choose optimal separation points. However, efficient computation and joint estimation of all model parameters (including the breakpoints) may be obtained by an iterative procedure currently implemented in the package `segmented` for the R language.



A variant of decision tree learning called model trees learns piecewise linear functions.



\textbf{A piecewise linear function in two dimensions (top) and the convex polytopes on which it is linear (bottom)}

% [Image: fit]



\paragraph{Notation}





The notion of a piecewise linear function makes sense in several different contexts. Piecewise linear functions may be defined on $n$-dimensional Euclidean space, or more generally any vector space or affine space, as well as on piecewise linear manifolds, simplicial complexes, and so forth. In each case, the function may be real-valued, or it may take values from a vector space, an affine space, a piecewise linear manifold, or a simplicial complex. (In these contexts, the term “linear” does not refer solely to linear transformations, but to more general affine linear functions.)



In dimensions higher than one, it is common to require the domain of each piece to be a polygon or polytope. This guarantees that the graph of the function will be composed of polygonal or polytopal pieces.



Important sub-classes of piecewise linear functions include the continuous piecewise linear functions and the convex piecewise linear functions. In general, for every $n$-dimensional continuous piecewise linear function $f: \mathbb{R}^n \to \mathbb{R}$, there is a $\Pi \in \mathcal{P}(\mathcal{P}(\mathbb{R}^{n+1}))$ such that

$$

f(\vec{x}) = \min\_{\Sigma \in \Pi} \max\_{(\vec{a}, b) \in \Sigma} \vec{a} \cdot \vec{x} + b.

$$

If $f$ is convex and continuous, then there is a $\Sigma \in \mathcal{P}(\mathbb{R}^{n+1})$ such that

$$

f(\vec{x}) = \max\_{(\vec{a}, b) \in \Sigma} \vec{a} \cdot \vec{x} + b.

$$

Splines generalize piecewise linear functions to higher-order polynomials, which are in turn contained in the category of piecewise-differentiable functions, PDIFF.



\paragraph{Applications}



In agriculture, piecewise regression analysis of measured data is used to detect the range over which growth factors affect the yield and the range over which the crop is not sensitive to changes in these factors.



The image below shows that at shallow watertables the yield declines, whereas at deeper ( $> 7 \, \text{dm}$ ) watertables the yield is unaffected. The graph is made using the method of least squares to find the two segments with the best fit.



\textbf{Crop Response to Depth of the Watertable}

% [Image: watertable]



The graph below reveals that crop yields tolerate a soil salinity up to $EC\_e = 8 \, \text{dS/m}$ ($EC\_e$ is the electric conductivity of an extract of a saturated soil sample), while beyond that value the crop production reduces. The graph is made with the method of partial regression to find the longest range of "no effect," i.e., where the line is horizontal. The two segments need not join at the same point. Only for the second segment is the method of least squares used.



\textbf{Example of Crop Response to Soil Salinity}

% [Image: crop]

\paragraph{Licensing}



Content obtained and/or adapted from:



\begin{itemize}
    \item \href{https://en.wikipedia.org/wiki/Piecewise_linear_function}{Piecewise linear function, Wikipedia} under a CC BY-SA license
\end{itemize}

\subsection{Learning Objectives}

\begin{enumerate}
    \item \textbf{Identify and Define Piecewise Linear Functions}:
    \item Understand the definition of a piecewise linear function.
    \item Identify the intervals and affine functions that compose a piecewise linear function.
    \item \textbf{Graph Piecewise Linear Functions}:
    \item Graph piecewise linear functions given their piecewise definitions.
    \item Recognize and label breakpoints, changepoints, threshold values, or knots on the graph.
    \item \textbf{Analyze and Apply Piecewise Linear Functions}:
    \item Use piecewise linear functions to model real-world scenarios, such as crop yield response to environmental factors.
    \item Apply methods such as least squares and partial regression to fit piecewise linear functions to data.
\end{enumerate}

\subsection{Assessment Instrument}

\begin{enumerate}
    \item Question 1 \textbf{Define the piecewise linear function $f(x)$ given below and identify its intervals and corresponding affine functions.} $$ f(x) = \begin{cases} -x - 3 & \text{if } x \leq -3 \\\\ x + 3 & \text{if } -3 < x < 0 \\\\ -2x + 3 & \text{if } 0 \leq x < 3 \\\\ 0.5x - 4.5 & \text{if } x \geq 3 \end{cases} $$ Question 2 \textbf{Graph the piecewise linear function $f(x)$ defined in Question 1. Label the breakpoints on the graph.} Question 3 \textbf{Given the following piecewise linear function, determine $f(2)$ and explain why it is not equal to $4$.} $$ f(x) = \begin{cases} x^2 & \text{if } x < 2 \\\\ -(x - 3)^2 + 9 & \text{if } x \geq 2 \end{cases} $$ Question 4 \textbf{Explain how piecewise regression analysis can be used in agriculture to determine the effect of soil salinity on crop yield.} Question 5 \textbf{Given the data points below, fit a piecewise linear function using the method of least squares to find the best fit.}
    \item Data points: $(1, 2)$, $(2, 4)$, $(3, 6)$, $(4, 5)$, $(5, 3)$
    \item Assume a breakpoint at $x = 3$.
\end{enumerate}

%%===============================================================
\section{Depreciation}
\label{sec:depreciation}
%%===============================================================

\paragraph{Depreciation}





In accountancy, depreciation refers to two aspects of the same concept: first, the actual decrease of fair value of an asset, such as the decrease in value of factory equipment each year as it is used and wears, and second, the allocation in accounting statements of the original cost of the assets to periods in which the assets are used (depreciation with the matching principle).



Depreciation is thus the decrease in the value of assets and the method used to reallocate, or "write down" the cost of a tangible asset (such as equipment) over its useful life span. Businesses depreciate long-term assets for both accounting and tax purposes. The decrease in value of the asset affects the balance sheet of a business or entity, and the method of depreciating the asset, accounting-wise, affects the net income, and thus the income statement that they report. Generally, the cost is allocated as depreciation expense among the periods in which the asset is expected to be used.



Methods of computing depreciation, and the periods over which assets are depreciated, may vary between asset types within the same business and may vary for tax purposes. These may be specified by law or accounting standards, which may vary by country. There are several standard methods of computing depreciation expense, including fixed percentage, straight line, and declining balance methods. Depreciation expense generally begins when the asset is placed in service. For example, a depreciation expense of 100 per year for five years may be recognized for an asset costing 500. Depreciation has been defined as the diminution in the utility or value of an asset and is a non-cash expense. It does not result in any cash outflow; it just means that the asset is not worth as much as it used to be. Causes of depreciation are natural wear and tear.





\paragraph{Accounting Concept}



In determining the net income (profits) from an activity, the receipts from the activity must be reduced by appropriate costs. One such cost is the cost of assets used but not immediately consumed in the activity. Such cost allocated in a given period is equal to the reduction in the value placed on the asset, which is initially equal to the amount paid for the asset and subsequently may or may not be related to the amount expected to be received upon its disposal. Depreciation is any method of allocating such net cost to those periods in which the organization is expected to benefit from the use of the asset. Depreciation is a process of deducting the cost of an asset over its useful life. Assets are sorted into different classes, and each has its own useful life. The asset is referred to as a depreciable asset. Depreciation is technically a method of allocation, not valuation, even though it determines the value placed on the asset in the balance sheet.



Any business or income-producing activity using tangible assets may incur costs related to those assets. If an asset is expected to produce a benefit in future periods, some of these costs must be deferred rather than treated as a current expense. The business then records depreciation expense in its financial reporting as the current period''s allocation of such costs. This is usually done in a rational and systematic manner. Generally, this involves four criteria:



\begin{itemize}
    \item Cost of the asset
    \item Expected salvage value, also known as the residual value of the assets
    \item Estimated useful life of the asset
    \item A method of apportioning the cost over such life
\end{itemize}

\paragraph{Depreciable Basis}



Cost generally is the amount paid for the asset, including all costs related to acquiring and bringing the asset into use. In some countries or for some purposes, salvage value may be ignored. The rules of some countries specify lives and methods to be used for particular types of assets. However, in most countries, the life is based on business experience, and the method may be chosen from one of several acceptable methods.



\paragraph{Impairment}



Accounting rules also require that an impairment charge or expense be recognized if the value of assets declines unexpectedly. Such charges are usually nonrecurring and may relate to any type of asset. Many companies consider write-offs of some of their long-lived assets because some property, plant, and equipment have suffered partial obsolescence. Accountants reduce the asset''s carrying amount by its fair value. For example, if a company continues to incur losses because prices of a particular product or service are higher than the operating costs, companies consider write-offs of the particular asset. These write-offs are referred to as impairments. There are events and changes in circumstances that might lead to impairment. Some examples are:



\begin{itemize}
    \item Large amount of decrease in the fair value of an asset
    \item A change in the manner in which the asset is used
    \item Accumulation of costs that were not originally expected to acquire or construct an asset
    \item A projection of incurring losses associated with the particular asset
\end{itemize}

Events or changes in circumstances indicate that the company may not be able to recover the carrying amount of the asset. In this case, companies use the recoverability test to determine whether impairment has occurred. The steps to determine are:

1. Estimate the future cash flow of the asset (from the use of the asset to disposition)

2. If the sum of the expected cash flow is less than the carrying amount of the asset, the asset is considered impaired



\paragraph{Depletion and Amortization}



Depletion and amortization are similar concepts for natural resources (including oil) and intangible assets, respectively.



\paragraph{Effect on Cash}



Depreciation expense does not require a current outlay of cash. However, since depreciation is an expense to the P\&L account, provided the enterprise is operating in a manner that covers its expenses (e.g., operating at a profit), depreciation is a source of cash in a statement of cash flows, which generally offsets the cash cost of acquiring new assets required to continue operations when existing assets reach the end of their useful lives.



\paragraph{Accumulated Depreciation}



While depreciation expense is recorded on the income statement of a business, its impact is generally recorded in a separate account and disclosed on the balance sheet as accumulated under fixed assets, according to most accounting principles. Accumulated depreciation is known as a contra account because it separately shows a negative amount that is directly associated with an accumulated depreciation account on the balance sheet. Depreciation expense is usually charged against the relevant asset directly. The values of the fixed assets stated on the balance sheet will decline, even if the business has not invested in or disposed of any assets. Theoretically, the amounts will roughly approximate fair value. Otherwise, depreciation expense is charged against accumulated depreciation. Showing accumulated depreciation separately on the balance sheet has the effect of preserving the historical cost of assets on the balance sheet. If there have been no investments or dispositions in fixed assets for the year, then the values of the assets will be the same on the balance sheet for the current and prior year (P/Y).



\paragraph{Methods for Depreciation}



There are several methods for calculating depreciation, generally based on either the passage of time or the level of activity (or use) of the asset.



\paragraph{Straight-Line Depreciation}



Straight-line depreciation is the simplest and most often used method. The straight-line depreciation is calculated by dividing the difference between the asset''s cost and its expected salvage value by the number of years for its expected useful life. (The salvage value may be zero or even negative due to costs required to retire it; however, for depreciation purposes, salvage value is not generally calculated at below zero.) The company will then charge the same amount to depreciation each year over that period until the value shown for the asset has reduced from the original cost to the salvage value.



\textbf{Straight-line method:}

$$

\text{Annual Depreciation Expense} = \frac{\text{Cost of Fixed Asset} - \text{Residual Value}}{\text{Useful Life of Asset (Years)}}

$$



For example, a vehicle that depreciates over 5 years is purchased at a cost of `$`17,000 and will have a salvage value of `$`2,000. Then this vehicle will depreciate at `$`3,000 per year, i.e.:

$$

\text{DE} = \frac{17,000 - 2,000}{5} = 3,000

$$



This table illustrates the straight-line method of depreciation. Book value at the beginning of the first year of depreciation is the original cost of the asset. At any time, book value equals original cost minus accumulated depreciation.



$$

\text{Book Value} = \text{Original Cost} - \text{Accumulated Depreciation}

$$



Book value at the end of the year becomes the book value at the beginning of the next year. The asset is depreciated until the book value equals scrap value.



$$

\begin{array}{|c|c|c|}

\hline

\text{Depreciation Expense} & \text{Accumulated Depreciation at Year-End} & \text{Book Value at Year-End} \\\\ 

\hline

\$3,000 & \$3,000 & \$14,000 \\\\ 

\hline

\$3,000 & \$6,000 & \$11,000 \\\\ 

\hline

\$3,000 & \$9,000 & \$8,000 \\\\ 

\hline

\$3,000 & \$12,000 & \$5,000 \\\\ 

\hline

\$3,000 & \$15,000 & \text{(scrap value) } \$2,000 \\\\ 

\hline

\end{array}

$$





If the vehicle were to be sold and the sales price exceeded the depreciated value (net book value), then the excess would be considered a gain and subject to depreciation recapture. In addition, this gain above the depreciated value would be recognized as ordinary income by the tax office. If the sales price is ever less than the book value, the resulting capital loss is tax-deductible. If the sale price were ever more than the original book value, then the gain above the original book value is recognized as a capital gain.



If a company chooses to depreciate an asset at a different rate from that used by the tax office, then this generates a timing difference in the income statement due to the difference (at a point in time) between the taxation department''s and company''s view of the profit.



\textbf{Diminishing Balance Method}



$$

\begin{array}{|c|c|c|c|}

\hline

\textbf{Depreciation rate} & \textbf{Depreciation expense} & \textbf{Accumulated depreciation} & \textbf{Book value at year-end} \\\\ 

\hline

\text{Original cost: } \$1,000.00 & & & \\\\ 

\hline

40\% & 400.00 & 400.00 & 600.00 \\\\ 

\hline

40\% & 240.00 & 640.00 & 360.00 \\\\ 

\hline

40\% & 144.00 & 784.00 & 216.00 \\\\ 

\hline

40\% & 86.40 & 870.40 & 129.60 \\\\ 

\hline

129.60 - 100.00 & 29.60 & 900.00 & \text{Scrap value: } 100.00 \\\\ 

\hline

\end{array}

$$



\paragraph{Double-Declining-Balance Method}



The double-declining-balance method is used to calculate an asset''s accelerated rate of depreciation against its non-depreciated balance during the earlier years of the asset''s useful life. When using the double-declining-balance method, the salvage value is not considered in determining the annual depreciation, but the book value of the asset being depreciated is never brought below its salvage value, regardless of the method used. Depreciation ceases when either the salvage value or the end of the asset''s useful life is reached.



Since double-declining-balance depreciation does not always depreciate an asset fully by its end of life, some methods also compute a straight-line depreciation each year and apply the greater of the two. This has the effect of converting from declining-balance depreciation to straight-line depreciation at a midpoint in the asset''s life. The double-declining-balance method is also a better representation of how vehicles depreciate and can more accurately match cost with benefit from asset use. The company in the future may want to allocate as little depreciation expenses as possible to help with additional expenses.



With the declining balance method, one can find the depreciation rate that would allow exactly for full depreciation by the end of the period, using the formula:



$$

\text{{depreciation rate}} = 1 - \left(\frac{{\text{{residual value}}}}{{\text{{cost of fixed asset}}}}\right)^{\frac{1}{N}}

$$



where $N$ is the estimated life of the asset (for example, in years).



$$

\text{{DE}} = 2 \times \text{{SLDP}} \times \text{{BV}}

$$



\paragraph{Annuity Depreciation}



Annuity depreciation methods are not based on time, but on a level of annuity. This could be miles driven for a vehicle, or a cycle count for a machine. When the asset is acquired, its life is estimated in terms of this level of activity. Assume the vehicle above is estimated to go 50,000 miles in its lifetime. The per-mile depreciation rate is calculated as:



$$

\frac{{\$17,000 \text{{ cost}} - \$2,000 \text{{ salvage}}}}{50,000 \text{{ miles}}} = \$0.30 \text{{ per mile}}

$$



Each year, the depreciation expense is then calculated by multiplying the number of miles driven by the per-mile depreciation rate.



\paragraph{Sum-of-Years-Digits Method}



Sum-of-years-digits is a depreciation method that results in a more accelerated write-off than the straight-line method, and typically also more accelerated than the declining balance method. Under this method, the annual depreciation is determined by multiplying the depreciable cost by a schedule of fractions.



Sum of the years'' digits method of depreciation is one of the accelerated depreciation techniques which are based on the assumption that assets are generally more productive when they are new and their productivity decreases as they become old. The formula to calculate depreciation under the SYD method is:



$$

\text{{SYD depreciation}} = \text{{depreciable base}} \times \left(\frac{{\text{{remaining useful life}}}}{{\text{{sum of the years'' digits}}}}\right)

$$



where:



$$

\text{{depreciable base}} = \text{{cost}} - \text{{salvage value}}

$$



\textbf{Example}: If an asset has an original cost of \$1000, a useful life of 5 years, and a salvage value of \$100, compute its depreciation schedule.



1. Determine the years'' digits. Since the asset has a useful life of 5 years, the years'' digits are: $5, 4, 3, 2, 1$.

2. Calculate the sum of the digits: $5 + 4 + 3 + 2 + 1 = 15$.

3. The sum of the digits can also be determined by using the formula $\frac{n^2 + n}{2}$, where $n$ is equal to the useful life of the asset in years. The example would be shown as \(\frac{5^2 + 5}{2} = 15\).



Depreciation rates are as follows:



$$

\frac{5}{15} \text{ for the 1st year}, \frac{4}{15} \text{ for the 2nd year}, \frac{3}{15} \text{ for the 3rd year}, \frac{2}{15} \text{ for the 4th year}, \text{ and } \frac{1}{15} \text{ for the 5th year}.

$$



$$

\begin{array}{|c|c|c|c|c|}

\hline

\text{Depreciable base} & \text{Depreciation rate} & \text{Depreciation expense} & \text{Accumulated depreciation} & \text{Book value at end of year} \\\\ 

\hline

\$1,000 \text{ (original cost)} & & & & \\\\ 

\hline

900 & 5/15 & 300 = (900 \times \frac{5}{15}) & 300 & 700 \\\\ 

\hline

900 & 4/15 & 240 = (900 \times \frac{4}{15}) & 540 & 460 \\\\ 

\hline

900 & 3/15 & 180 = (900 \times \frac{3}{15}) & 720 & 280 \\\\ 

\hline

900 & 2/15 & 120 = (900 \times \frac{2}{15}) & 840 & 160 \\\\ 

\hline

900 & 1/15 & 60 = (900 \times \frac{1}{15}) & 900 & 100 \text{ (scrap value)} \\\\ 

\hline

\end{array}

$$





\paragraph{Units-of-production depreciation method}





The units-of-production depreciation method calculates greater deductions for depreciation in years when the asset is heavily used.



$$

\text{{annual depreciation expense}} = \frac{{\text{{cost of fixed asset}} - \text{{residual value}}}}{{\text{{estimated total production}}}} \times \text{{actual production}}

$$



Suppose, an asset has an original cost of \$70,000, salvage value \$10,000, and is expected to produce 6,000 units.



$$

\text{{Depreciation per unit}} = \frac{{\$70,000 - 10,000}}{{6,000}} = \$10

$$



\$10 \times \text{{actual production}} will give the depreciation cost of the current year.



The table below illustrates the units-of-production depreciation schedule of the asset.



$$

\begin{array}{|c|c|c|c|c|}

\hline

\textbf{Units of production} & \textbf{Depreciation cost per unit} & \textbf{Depreciation expense} & \textbf{Accumulated depreciation} & \textbf{Book value at year-end} \\\\ 

\hline

\text{\$70,000 (original cost)} & & & & \\\\ 

\hline

1,000 & 10 & 10,000 & 10,000 & 60,000 \\\\ 

\hline

1,100 & 10 & 11,000 & 21,000 & 49,000 \\\\ 

\hline

1,200 & 10 & 12,000 & 33,000 & 37,000 \\\\ 

\hline

1,300 & 10 & 13,000 & 46,000 & 24,000 \\\\ 

\hline

1,400 & 10 & 14,000 & 60,000 & \text{10,000 (scrap value)} \\\\ 

\hline

\end{array}

$$



Depreciation stops when book value is equal to the scrap value of the asset. In the end, the sum of accumulated depreciation and scrap value equals the original cost.







\paragraph{Group depreciation method}





The group depreciation method is used for depreciating multiple-asset accounts using a similar depreciation method. The assets must be similar in nature and have approximately the same useful lives.







\paragraph{Composite depreciation method}





The composite method is applied to a collection of assets that are not similar, and have different service lives. For example, computers and printers are not similar, but both are part of the office equipment. Depreciation on all assets is determined by using the straight-line-depreciation method.



$$

\begin{array}{|c|c|c|c|c|c|}

\hline

\textbf{Asset} & \textbf{Historical cost} & \textbf{Salvage value} & \textbf{Depreciable cost} & \textbf{Life (years)} & \textbf{Depreciation per year} \\\\ 

\hline

\text{Computers} & \$5,500 & \$500 & \$5,000 & 5 & \$1,000 \\\\ 

\hline

\text{Printers} & \$1,000 & \$100 & \$900 & 3 & \$300 \\\\ 

\hline

\textbf{Total} & \$6,500 & \$600 & \$5,900 & 4.5 & \$1,300 \\\\ 

\hline

\end{array}

$$



\textbf{Composite life} equals the total depreciable cost divided by the total depreciation per year: 



$$

\text{{Composite life}} = \frac{{\$5,900}}{{\$1,300}} = 4.5 \text{{ years}}

$$



\textbf{Composite depreciation rate} equals depreciation per year divided by total historical cost:



$$

\text{{Composite depreciation rate}} = \frac{{\$1,300}}{{\$6,500}} = 0.20 = 20\%

$$



\textbf{Depreciation expense} equals the composite depreciation rate times the balance in the asset account (historical cost):



$$

\text{{Depreciation expense}} = 0.20 \times \$6,500 = \$1,300

$$



When an asset is sold, debit cash for the amount received and credit the asset account for its original cost. Debit the difference between the two to accumulated depreciation. Under the composite method, no gain or loss is recognized on the sale of an asset.



To calculate the composite depreciation rate, divide depreciation per year by total historical cost. To calculate depreciation expense, multiply the result by the same total historical cost. The result will equal the total depreciation per year again.







\paragraph{Tax Depreciation}





Most income tax systems allow a tax deduction for recovery of the cost of assets used in a business or for the production of income. Such deductions are allowed for individuals and companies. Where the assets are consumed currently, the cost may be deducted currently as an expense or treated as part of the cost of goods sold. The cost of assets not currently consumed generally must be deferred and recovered over time, such as through depreciation. Some systems permit the full deduction of the cost, at least in part, in the year the assets are acquired. Other systems allow depreciation expense over some life using some depreciation method or percentage. Rules vary highly by country, and may vary within a country based on the type of asset or type of taxpayer. Many systems that specify depreciation lives and methods for financial reporting require the same lives and methods to be used for tax purposes. Most tax systems provide different rules for real property (buildings, etc.) and personal property (equipment, etc.).







\paragraph{Capital allowances}





A common system is to allow a fixed percentage of the cost of depreciable assets to be deducted each year. This is often referred to as a capital allowance, as it is called in the United Kingdom. Deductions are permitted to individuals and businesses based on assets placed in service during or before the assessment year. Canada''s Capital Cost Allowance are fixed percentages of assets within a class or type of asset. Fixed percentage rates are specified by the type of asset. The fixed percentage is multiplied by the tax basis of assets in service to determine the capital allowance deduction. The tax law or regulations of the country specify these percentages. Capital allowance calculations may be based on the total set of assets, on sets or pools by year (vintage pools), or pools by classes of assets. Depreciation has three methods only.







\paragraph{Tax lives and methods}





Some systems specify lives based on classes of property defined by the tax authority. The Canada Revenue Agency specifies numerous classes based on the type of property and how it is used. Under the United States depreciation system, the Internal Revenue Service publishes a detailed guide that includes a table of asset lives and the applicable conventions. The table also incorporates specified lives for certain commonly used assets (e.g., office furniture, computers, automobiles) which override the business use lives. U.S. tax depreciation is computed under the double-declining balance method switching to straight line or the straight-line method, at the option of the taxpayer. IRS tables specify percentages to apply to the basis of an asset for each year in which it is in service. Depreciation first becomes deductible when an asset is placed in service.







\paragraph{Additional depreciation}





Many systems allow an additional deduction for a portion of the cost of depreciable assets acquired in the current tax year. The UK system provides a first-year capital allowance of £50,000. In the United States, two such deductions are available. A deduction for the full cost of depreciable tangible personal property is allowed up to \$500,000 through 2013. This deduction is fully phased out for businesses acquiring over \$2,000,000 of such property during the year. In addition, additional first-year depreciation of 50\% of the cost of most other depreciable tangible personal property is allowed as a deduction. Some other systems have similar first-year or accelerated allowances.







\paragraph{Real property}





Many tax systems prescribe longer depreciable lives for buildings and land improvements. Such lives may vary by type of use. Many such systems, including the United States and Canada, permit depreciation for real property using only the straight-line method or a small fixed percentage of the cost. Generally, no depreciation tax deduction is allowed for bare land. In the United States, residential rental buildings are depreciable over a 27.5-year or 40-year life, other buildings over a 39-year or 40-year life, and land improvements over a 15-year or 20-year life, all using the straight-line method.







\paragraph{Averaging conventions}





Depreciation calculations require a lot of record-keeping if done for each asset a business owns, especially if assets are added after they are acquired, or partially disposed of. However, many tax systems permit all assets of a similar type acquired in the same year to be combined in a "pool". Depreciation is then computed for all assets in the pool as a single calculation. These calculations must make assumptions about the date of acquisition. The United States system allows a taxpayer to use a half-year convention for personal property or a mid-month convention for real property. Under such a convention, all property of a particular type is considered to have been acquired at the midpoint of the acquisition period. One-half of a full period''s depreciation is allowed in the acquisition period (and also in the final depreciation period if the life of the assets is a whole number of years). United States rules require a mid-quarter convention for personal property if more than 40\% of the acquisitions for the year are in the final quarter.







\paragraph{Licensing}





Content obtained and/or adapted from:



\href{https://en.wikipedia.org/wiki/Depreciation}{Depreciation, Wikipedia} under a CC BY-SA license

\subsection{Learning Objectives}

\begin{enumerate}
    \item SLO 1: Straight-Line Depreciation
    \item \textbf{Objective:} Students will be able to calculate annual depreciation expense using the straight-line depreciation method.
    \item \textbf{Performance Criteria:} Given the cost, salvage value, and useful life of an asset, students will determine the annual depreciation expense and prepare a table showing depreciation expense, accumulated depreciation, and book value for each year. SLO 2: Declining Balance Depreciation
    \item \textbf{Objective:} Students will be able to compute depreciation expense using the declining balance method.
    \item \textbf{Performance Criteria:} Given the cost of an asset, a specified depreciation rate, and its useful life, students will calculate the annual depreciation expense for each year and prepare a table showing depreciation expense, accumulated depreciation, and book value. SLO 3: Composite Depreciation Method
    \item \textbf{Objective:} Students will be able to calculate composite depreciation using the composite method.
    \item \textbf{Performance Criteria:} Given historical costs, salvage values, and useful lives of multiple assets, students will compute the composite depreciation rate and total depreciation expense for a collection of assets.
\end{enumerate}

\subsection{Assessment Instrument}

\begin{enumerate}
    \item Part 1: Straight-Line Depreciation
    \item \textbf{Problem:} An asset costs `$`12,000, has a salvage value of `$`2,000, and a useful life of 8 years. Calculate the annual depreciation expense. Prepare a table showing the annual depreciation expense, accumulated depreciation, and book value at the end of each year.
    \item \textbf{Problem:} Given an asset purchased for `$`25,000 with a salvage value of `$`5,000 and a useful life of 10 years, prepare a table showing the depreciation expense, accumulated depreciation, and book value for each year using the straight-line method. Part 2: Declining Balance Depreciation
    \item \textbf{Problem:} An asset costs `$`8,000, has a useful life of 5 years, and a depreciation rate of 40\%. Calculate the depreciation expense for each year using the declining balance method. Prepare a table showing the annual depreciation expense, accumulated depreciation, and book value.
    \item \textbf{Problem:} Given an asset with an initial cost of `$`15,000, a useful life of 6 years, and a declining balance rate of 30\%, calculate the annual depreciation expense for each year. Prepare a table to show the depreciation expense, accumulated depreciation, and book value at the end of each year. Part 3: Composite Depreciation Method
    \item \textbf{Problem:} A company has two types of assets: Computers costing `$`7,000 with a salvage value of `$`500 and a useful life of 4 years, and Printers costing `$`1,800 with a salvage value of `$`200 and a useful life of 3 years. Calculate the composite depreciation rate and total annual depreciation expense for the combined assets.
    \item \textbf{Problem:} Given a set of assets with total historical cost of `$`10,000, total salvage value of `$`1,000, and total annual depreciation expense of `$`2,000, calculate the composite depreciation rate and prepare a table showing the total depreciation expense and the effective life of the assets.
\end{enumerate}

%%===============================================================
\section{Geometric interpretation of interpolation}
\label{sec:geometric-interpretation-of-interpolation}
%%===============================================================

\paragraph{Geometric Interpretation of Interpolation}



In the mathematical field of numerical analysis, interpolation is a type of estimation, a method of constructing (finding) new data points based on the range of a discrete set of known data points.



In engineering and science, one often has a number of data points, obtained by sampling or experimentation, which represent the values of a function for a limited number of values of the independent variable. It is often required to interpolate; that is, estimate the value of that function for an intermediate value of the independent variable.



A closely related problem is the approximation of a complicated function by a simple function. Suppose the formula for some given function is known, but too complicated to evaluate efficiently. A few data points from the original function can be interpolated to produce a simpler function which is still fairly close to the original. The resulting gain in simplicity may outweigh the loss from interpolation error and give better performance in calculation process.



% [Image: display]

\textbf{An interpolation of a finite set of points on an epitrochoid. The points in red are connected by blue interpolated spline curves deduced only from the red points. The interpolated curves have polynomial formulas much simpler than that of the original epitrochoid curve.}



\paragraph{Example}



This array gives some values of an unknown function $f(x)$.



$$

\begin{array}{cc}

x & f(x) \\\\ 

0 & 0 \\\\ 

1 & 0.8415 \\\\ 

2 & 0.9093 \\\\ 

3 & 0.1411 \\\\ 

4 & -0.7568 \\\\ 

5 & -0.9589 \\\\ 

6 & -0.2794 \\\\ 

\end{array}

$$



% [Image: plot]



\textbf{Plot of the data points as given in the array.}



Interpolation provides a means of estimating the function at intermediate points, such as $x=2.5$.



We describe some methods of interpolation, differing in such properties as: accuracy, cost, number of data points needed, and smoothness of the resulting interpolant function.



\paragraph{Piecewise Constant Interpolation}



% [Image: piece]



\textbf{Piecewise constant interpolation, or nearest-neighbor interpolation.}



The simplest interpolation method is to locate the nearest data value, and assign the same value. In simple problems, this method is unlikely to be used, as linear interpolation (see below) is almost as easy, but in higher-dimensional multivariate interpolation, this could be a favourable choice for its speed and simplicity.



\paragraph{Linear Interpolation}



% [Image: super]



\textbf{Plot of the data with linear interpolation superimposed.}



One of the simplest methods is linear interpolation (sometimes known as lerp). Consider the above example of estimating $f(2.5)$. Since 2.5 is midway between 2 and 3, it is reasonable to take $f(2.5)$ midway between $f(2) = 0.9093$ and $f(3) = 0.1411$, which yields 0.5252.



Generally, linear interpolation takes two data points, say $(x\_a, y\_a)$ and $(x\_b, y\_b)$, and the interpolant is given by:



$$y = y\_a + \left(y\_b - y\_a\right) \frac{x - x\_a}{x\_b - x\_a} \text{ at the point } (x, y)$$



$$\frac{y - y\_a}{y\_b - y\_a} = \frac{x - x\_a}{x\_b - x\_a}$$



$$\frac{y - y\_a}{x - x\_a} = \frac{y\_b - y\_a}{x\_b - x\_a}$$



This previous equation states that the slope of the new line between $(x\_a, y\_a)$ and $(x, y)$ is the same as the slope of the line between $(x\_a, y\_a)$ and $(x\_b, y\_b)$.



Linear interpolation is quick and easy, but it is not very precise. Another disadvantage is that the interpolant is not differentiable at the point $x\_k$.



The following error estimate shows that linear interpolation is not very precise. Denote the function which we want to interpolate by $g$, and suppose that $x$ lies between $x\_a$ and $x\_b$ and that $g$ is twice continuously differentiable. Then the linear interpolation error is



$$|f(x) - g(x)| \leq C (x\_b - x\_a)^2 \text{ where }C = \frac{1}{8} \max\_{r \in [x\_a, x\_b]} |g''''(r)|.$$



In words, the error is proportional to the square of the distance between the data points. The error in some other methods, including polynomial interpolation and spline interpolation (described below), is proportional to higher powers of the distance between the data points. These methods also produce smoother interpolants.



\paragraph{Polynomial Interpolation}



% [Image: poly]



\textbf{Plot of the data with polynomial interpolation applied.}



Polynomial interpolation is a generalization of linear interpolation. Note that the linear interpolant is a linear function. We now replace this interpolant with a polynomial of higher degree.



Consider again the problem given above. The following sixth degree polynomial goes through all the seven points:



$$f(x) = -0.0001521 x^6 - 0.003130 x^5 + 0.07321 x^4 - 0.3577 x^3 + 0.2255 x^2 + 0.9038 x.$$



Substituting $x = 2.5$, we find that $f(2.5) \approx 0.59678$.



Generally, if we have $n$ data points, there is exactly one polynomial of degree at most $n - 1$ going through all the data points. The interpolation error is proportional to the distance between the data points to the power $n$. Furthermore, the interpolant is a polynomial and thus infinitely differentiable. So, we see that polynomial interpolation overcomes most of the problems of linear interpolation.



However, polynomial interpolation also has some disadvantages. Calculating the interpolating polynomial is computationally expensive (see computational complexity) compared to linear interpolation. Furthermore, polynomial interpolation may exhibit oscillatory artifacts, especially at the end points (see Runge''s phenomenon).



Polynomial interpolation can estimate local maxima and minima that are outside the range of the samples, unlike linear interpolation. For example, the interpolant above has a local maximum at $x \approx 1.566$, $f(x) \approx 1.003$ and a local minimum at $x \approx 4.708$, $f(x) \approx -1.003$. However, these maxima and minima may exceed the theoretical range of the function; for example, a function that is always positive may have an interpolant with negative values, and whose inverse therefore contains false vertical asymptotes.



More generally, the shape of the resulting curve, especially for very high or low values of the independent variable, may be contrary to commonsense; that is, to what is known about the experimental system which has generated the data points. These disadvantages can be reduced by using spline interpolation or restricting attention to Chebyshev polynomials.



\paragraph{Spline Interpolation}



% [Image: spline]



\textbf{Plot of the data with spline interpolation applied.}



Remember that linear interpolation uses a linear function for each of intervals $[x\_k, x\_{k+1}]$. Spline interpolation uses low-degree polynomials in each of the intervals, and chooses the polynomial pieces such that they fit smoothly together. The resulting function is called a spline.



For instance, the natural cubic spline is piecewise cubic and twice continuously differentiable. Furthermore, its second derivative is zero at the end points. The natural cubic spline interpolating the points in the array above is given by



$f(x) =

\begin{cases}

-0.1522 x^3 + 0.9937 x, \& \text{if } x \in [0, 1], \\\\ 

-0.01258 x^3 - 0.4189 x^2 + 1.4126 x - 0.1396, \& \text{if } x \in [1, 2], \\\\ 

0.1403 x^3 - 1.3359 x^2 + 3.2467 x - 1.3623, \& \text{if } x \in [2, 3], \\\\ 

0.1579 x^3 - 1.4945 x^2 + 3.7225 x - 1.8381, \& \text{if } x \in [3, 4], \\\\ 

0.05375 x^3 - 0.2450 x^2 - 1.2756 x + 4.8259, \& \text{if } x \in [4, 5], \\\\ 

-0.1871 x^3 + 3.3673 x^2 - 19.3370 x + 34.9282, \& \text{if } x \in [5, 6].

\end{cases}$



In this case we get $f(2.5) = 0.5972$.



Like polynomial interpolation, spline interpolation incurs a smaller error than linear interpolation, while the interpolant is smoother and easier to evaluate than the high-degree polynomials used in polynomial interpolation. However, the global nature of the basis functions leads to ill-conditioning. This is completely mitigated by using splines of compact support, such as are implemented in Boost.Math and discussed in Kress.



\paragraph{Function Approximation}



Interpolation is a common way to approximate functions. Given a function $f: [a, b] \to \mathbb{R}$ with a set of points $x\_1, x\_2, \dots, x\_n \in [a, b]$, one can form a function $s: [a, b] \to \mathbb{R}$ such that $f(x\_i) = s(x\_i)$ for $i = 1, 2, \dots, n$ (that is, that $s$ interpolates $f$ at these points). In general, an interpolant need not be a good approximation, but there are well known and often reasonable conditions where it will. For example, if $f \in C^4([a, b])$ (four times continuously differentiable) then cubic spline interpolation has an error bound given by



$\|f - s\|\_{\infty} \leq C \|f^{(4)}\|\_{\infty} h^4$



where $h = \max\_{i=1, 2, \dots, n-1} |x\_{i+1} - x\_i|$ and $C$ is a constant.



\paragraph{Via Gaussian Processes}



Gaussian process is a powerful non-linear interpolation tool. Many popular interpolation tools are actually equivalent to particular Gaussian processes. Gaussian processes can be used not only for fitting an interpolant that passes exactly through the given data points but also for regression; that is, for fitting a curve through noisy data. In the geostatistics community, Gaussian process regression is also known as Kriging.



\paragraph{Other Forms}



Other forms of interpolation can be constructed by picking a different class of interpolants. For instance, rational interpolation is interpolation by rational functions using Padé approximant, and trigonometric interpolation is interpolation by trigonometric polynomials using Fourier series. Another possibility is to use wavelets.



The Whittaker–Shannon interpolation formula can be used if the number of data points is infinite or if the function to be interpolated has compact support.



Sometimes, we know not only the value of the function that we want to interpolate at some points but also its derivative. This leads to Hermite interpolation problems.



When each data point is itself a function, it can be useful to see the interpolation problem as a partial advection problem between each data point. This idea leads to the displacement interpolation problem used in transportation theory.



\paragraph{In Higher Dimensions}



Multivariate interpolation is the interpolation of functions of more than one variable. Methods include bilinear interpolation and bicubic interpolation in two dimensions, and trilinear interpolation in three dimensions. They can be applied to gridded or scattered data.



$\phantom{--}$% [Image: near]$\phantom{---}$  % [Image: bilinear] $\phantom{---}$ % [Image: bicubic]



$\phantom{-}$\textbf{Nearest Neighbor $\phantom{-----}$ Bilinear $\phantom{-------}$ Bicubic}



\paragraph{In Digital Signal Processing}



In the domain of digital signal processing, the term interpolation refers to the process of converting a sampled digital signal (such as a sampled audio signal) to that of a higher sampling rate (Upsampling) using various digital filtering techniques (for example, convolution with a frequency-limited impulse signal). In this application, there is a specific requirement that the harmonic content of the original signal be preserved without creating aliased harmonic content of the original signal above the original Nyquist limit of the signal (that is, above $f\_s/2$ of the original signal sample rate). An early and fairly elementary discussion on this subject can be found in Rabiner and Crochiere''s book \textit{Multirate Digital Signal Processing}.



\paragraph{Related Concepts}



The term extrapolation is used to find data points outside the range of known data points.



In curve fitting problems, the constraint that the interpolant has to go exactly through the data points is relaxed. It is only required to approach the data points as closely as possible (within some other constraints). This requires parameterizing the potential interpolants and having some way of measuring the error. In the simplest case, this leads to least squares approximation.



Approximation theory studies how to find the best approximation to a given function by another function from some predetermined class, and how good this approximation is. This clearly yields a bound on how well the interpolant can approximate the unknown function.



\paragraph{Generalization}



If we consider $x$ as a variable in a topological space, and the function $f(x)$ mapping to a Banach space, then the problem is treated as "interpolation of operators". The classical results about interpolation of operators are the Riesz–Thorin theorem and the Marcinkiewicz theorem. There are also many other subsequent results.



\paragraph{Licensing}



Content obtained and/or adapted from:



\href{https://en.wikipedia.org/wiki/Interpolation}{Interpolation, Wikipedia} under a CC BY-SA license

\subsection{Learning Objectives}

\begin{enumerate}
    \item \textbf{SLO 1:} Given two data points $(x\_a, f(x\_a))$ and $(x\_b, f(x\_b))$, students will be able to derive and use the formula for linear interpolation to estimate the value of the function at an intermediate point, $x$. \textbf{SLO 2} Students will compare the accuracy and computational complexity of linear interpolation, polynomial interpolation, and spline interpolation using provided datasets, and will be able to explain the advantages and disadvantages of each method. \textbf{SLO 3:} Given a set of data points, students will be able to construct a natural cubic spline interpolation and use it to estimate values at intermediate points, demonstrating an understanding of the smoothness and differentiability of spline interpolants.
\end{enumerate}

\subsection{Assessment Instrument}

\begin{enumerate}
    \item Part A: Linear Interpolation
    \item \textbf{Problem:}
    \item Given the data points $(1, 0.8415)$ and $(2, 0.9093)$, estimate the value of the function at $x = 1.5$ using linear interpolation. Part B: Method Comparison
    \item \textbf{Problem:}
    \item Given a set of data points, briefly describe the differences in accuracy, computational complexity, and smoothness between linear interpolation, polynomial interpolation, and spline interpolation. Use an example to illustrate your points. Part C: Spline Interpolation
    \item \textbf{Problem:}
    \item Construct a natural cubic spline interpolation for the following data points: $(0, 0)$, $(1, 0.8415)$, $(2, 0.9093)$, $(3, 0.1411)$. Estimate the value at $x = 1.5$ using the spline function.
\end{enumerate}

%%===============================================================
\section{Univariate, Bivariate, Correlation and Causation}
\label{sec:univariate-bivariate-correlation-and-causation}
%%===============================================================

\paragraph{Univariate, Bivariate, Correlation and Causation}



\paragraph{Univariate Data}



In mathematics, a univariate object is an expression, equation, function, or polynomial involving only one variable. Objects involving more than one variable are multivariate. For example, the fundamental theorem of algebra and Euclid''s algorithm for polynomials are fundamental properties of univariate polynomials that cannot be generalized to multivariate polynomials.



In statistics, a univariate distribution characterizes one variable. For instance, univariate data are composed of a single scalar component. In time series analysis, the entire time series is the "variable": a univariate time series represents the series of values over time of a single quantity. A "multivariate time series" characterizes the changing values over time of several quantities. Sometimes, univariate time series may be treated using multivariate statistical analyses and represented using multivariate distributions.



\paragraph{Measures of Location and Variation}



1. \textbf{Measures of Location} (e.g., mode, median, arithmetic mean) describe where the data is centrally arranged.

2. \textbf{Measures of Variation} (e.g., span, interquartile distance, standard deviation) describe how similar or different the data are scattered.



\paragraph{Bivariate Data}



Bivariate data involves two quantitative variables. This is analogous to univariate data, which involves one quantitative variable.



Bivariate data is used in statistics to study measures of central tendencies (e.g., mean, median, mode, midrange), variability, and spread. More than one variable is collected, such as height and weight in medical studies. Bivariate data compares two quantitative variables of an individual and can be shown graphically through histograms, scatterplots, or dotplots.



\paragraph{Example}



Given the following data:



\textbf{Men''s ages when they were married:}



$$

25, 26, 27, 29, 30, 31, 33, 36, 38, 40

$$



\textbf{Women''s ages when they were married:}



$$

19, 20, 21, 22, 23, 25, 26, 28, 29, 30

$$



The mean age of marriage for men is:



$$

\frac{25 + 26 + 27 + 29 + 30 + 31 + 33 + 36 + 38 + 40}{10} = 31.5

$$



The mean age of marriage for women is:



$$

\frac{19 + 20 + 21 + 22 + 23 + 25 + 26 + 28 + 29 + 30}{10} = 24.3

$$



Both distributions are slightly skewed to the right, indicating that more data values occur to the left of the mean. Based on this data, women generally marry at a younger age than men.



\paragraph{Statistical and Deterministic Relationships}



A deterministic relationship implies an exact mathematical relationship or dependence between variables. For example, Newton''s law of gravity:



$$

F = k \left(\frac{m\_1 m\_2}{r^2}\right)

$$



where $F$ is the force, $k$ is a constant, $m\_1$ and $m\_2$ are masses, and $r$ is the distance.



A random or stochastic relationship allows variation due to factors like measurement errors, reporting errors, computing errors, and other influences. For instance, crop yield relative to rainfall may have variability due to these factors.



\paragraph{Regression versus Causation}



Regression deals with dependence among variables within a model but does not imply causation. For example, while rainfall affects crop yield, crop yield does not affect rainfall. A statistical relationship does not necessarily imply causation.



\paragraph{Regression versus Correlation}



Correlation analysis measures the degree of linear association between two variables. For example, a correlation coefficient of 0.69 between crop yield and rainfall suggests a mathematical relationship.



In statistics, both variables are treated as having random error. In regression analysis, crop yield is the dependent variable, and rainfall is the explanatory variable. The independent variable is assumed to have no random component.



This distinction will be important in the section on measurement.



\paragraph{Correlation}



% [Image: pearson]



\textbf{Several sets of $(x, y)$ points, with the Pearson correlation coefficient of $x$ and $y$ for each set. The correlation reflects the noisiness and direction of a linear relationship (top row), but not the slope of that relationship (middle), nor many aspects of nonlinear relationships (bottom). N.B.: the figure in the center has a slope of 0 but in that case the correlation coefficient is undefined because the variance of $Y$ is zero.}



In statistics, correlation or dependence is any statistical relationship, whether causal or not, between two random variables or bivariate data. In the broadest sense, correlation is any statistical association, though it commonly refers to the degree to which a pair of variables are linearly related. Familiar examples of dependent phenomena include the correlation between the height of parents and their offspring, and the correlation between the price of a good and the quantity the consumers are willing to purchase, as depicted in the so-called demand curve.



Correlations are useful because they can indicate a predictive relationship that can be exploited in practice. For example, an electrical utility may produce less power on a mild day based on the correlation between electricity demand and weather. In this example, there is a causal relationship, because extreme weather causes people to use more electricity for heating or cooling. However, in general, the presence of a correlation is not sufficient to infer the presence of a causal relationship (i.e., correlation does not imply causation).



Formally, random variables are dependent if they do not satisfy a mathematical property of probabilistic independence. In informal parlance, correlation is synonymous with dependence. However, when used in a technical sense, correlation refers to any of several specific types of mathematical operations between the tested variables and their respective expected values. Essentially, correlation is the measure of how two or more variables are related to one another. There are several correlation coefficients, often denoted $\rho$ or $r$, measuring the degree of correlation. The most common of these is the Pearson correlation coefficient, which is sensitive only to a linear relationship between two variables (which may be present even when one variable is a nonlinear function of the other). Other correlation coefficients – such as Spearman''s rank correlation – have been developed to be more robust than Pearson''s, i.e., more sensitive to nonlinear relationships. Mutual information can also be applied to measure dependence between two variables.



\paragraph{Pearson''s Product-Moment Coefficient}



% [Image: scatter]



\textbf{Example Scatterplots of various datasets with various correlation coefficients}



\paragraph{Definition}



The most familiar measure of dependence between two quantities is the Pearson product-moment correlation coefficient (PPMCC), or "Pearson''s correlation coefficient", commonly called simply "the correlation coefficient". Mathematically, it is defined as the quality of least squares fitting to the original data. It is obtained by taking the ratio of the covariance of the two variables in question of our numerical dataset, normalized to the square root of their variances. Mathematically, one simply divides the covariance of the two variables by the product of their standard deviations. Karl Pearson developed the coefficient from a similar but slightly different idea by Francis Galton.



A Pearson product-moment correlation coefficient attempts to establish a line of best fit through a dataset of two variables by essentially laying out the expected values and the resulting Pearson''s correlation coefficient indicates how far away the actual dataset is from the expected values. Depending on the sign of our Pearson''s correlation coefficient, we can end up with either a negative or positive correlation if there is any sort of relationship between the variables of our dataset.



The population correlation coefficient $\rho\_{X,Y}$ between two random variables $X$ and $Y$ with expected values $\mu\_X$ and $\mu\_Y$ and standard deviations $\sigma\_X$ and $\sigma\_Y$ is defined as:



$$

\rho\_{X,Y} = \operatorname{corr}(X, Y) = \frac{\operatorname{cov}(X, Y)}{\sigma\_X \sigma\_Y} = \frac{\operatorname{E}[(X - \mu\_X)(Y - \mu\_Y)]}{\sigma\_X \sigma\_Y}

$$



where $\operatorname{E}$ is the expected value operator, $\operatorname{cov}$ means covariance, and $\operatorname{corr}$ is a widely used alternative notation for the correlation coefficient. The Pearson correlation is defined only if both standard deviations are finite and positive. An alternative formula purely in terms of moments is:



$$

\rho\_{X,Y} = \frac{\operatorname{E}(XY) - \operatorname{E}(X) \operatorname{E}(Y)}{\sqrt{\operatorname{E}(X^2) - \operatorname{E}(X)^2} \cdot \sqrt{\operatorname{E}(Y^2) - \operatorname{E}(Y)^2}}

$$



\paragraph{Symmetry Property}



The correlation coefficient is symmetric:



$$

\operatorname{corr}(X, Y) = \operatorname{corr}(Y, X)

$$



This is verified by the commutative property of multiplication.



\paragraph{Correlation and Independence}



It is a corollary of the Cauchy–Schwarz inequality that the absolute value of the Pearson correlation coefficient is not greater than 1. Therefore, the value of a correlation coefficient ranges between -1 and +1. The correlation coefficient is +1 in the case of a perfect direct (increasing) linear relationship (correlation), -1 in the case of a perfect inverse (decreasing) linear relationship (anti-correlation), and some value in the open interval $(-1, 1)$ in all other cases, indicating the degree of linear dependence between the variables. As it approaches zero, there is less of a relationship (closer to uncorrelated). The closer the coefficient is to either -1 or 1, the stronger the correlation between the variables.



If the variables are independent, Pearson''s correlation coefficient is 0, but the converse is not true because the correlation coefficient detects only linear dependencies between two variables.



$$

\begin{aligned}

X, Y \text{ independent} &\Rightarrow \rho\_{X, Y} = 0 \text{ (X, Y uncorrelated)} \\rho\_{X, Y} = 0 \text{ (X, Y uncorrelated)} &\nRightarrow X, Y \text{ independent}

\end{aligned}

$$



For example, suppose the random variable $X$ is symmetrically distributed about zero, and $Y = X^2$. Then $Y$ is completely determined by $X$, so that $X$ and $Y$ are perfectly dependent, but their correlation is zero; they are uncorrelated. However, in the special case when $X$ and $Y$ are jointly normal, uncorrelatedness is equivalent to independence.



Even though uncorrelated data does not necessarily imply independence, one can check if random variables are independent if their mutual information is 0.



\paragraph{Sample Correlation Coefficient}



Given a series of $n$ measurements of the pair $(X\_i, Y\_i)$ indexed by $i = 1, \ldots, n$, the sample correlation coefficient can be used to estimate the population Pearson correlation $\rho\_{X, Y}$ between $X$ and $Y$. The sample correlation coefficient is defined as:



% [Image: formula]



where $\overline{x}$ and $\overline{y}$ are the sample means of $X$ and $Y$, and $s\_x$ and $s\_y$ are the corrected sample standard deviations of $X$ and $Y$.



Equivalent expressions for $r\_{xy}$ are



% [Image: formula]



where $s''\_{x}$ and $s''\_{y}$ are the uncorrected sample standard deviations of $X$ and $Y$.



If $x$ and $y$ are results of measurements that contain measurement error, the realistic limits on the correlation coefficient are not $-1$ to $+1$ but a smaller range. For the case of a linear model with a single independent variable, the coefficient of determination ($R^2$) is the square of $r\_{xy}$, Pearson''s product-moment coefficient.



Example

Consider the joint probability distribution of $X$ and $Y$ given in the table below.



$$

\begin{array}{|c|c|c|c|}

\hline

\operatorname{P}(X=x, Y=y) & y = -1 & y = 0 & y = 1 \\\\ 

\hline

x = 0 & 0 & \frac{1}{3} & 0 \\\\ 

\hline

x = 1 & \frac{1}{3} & 0 & \frac{1}{3} \\\\ 

\hline

\end{array}

$$



For this joint distribution, the marginal distributions are:



$$

\operatorname{P}(X=x) =

\begin{cases}

\frac{1}{3} & \text{for } x=0 \\\\ 

\frac{2}{3} & \text{for } x=1

\end{cases}

$$



$$

\operatorname{P}(Y=y) =

\begin{cases}

\frac{1}{3} & \text{for } y=-1 \\\\ 

\frac{1}{3} & \text{for } y=0 \\\\ 

\frac{1}{3} & \text{for } y=1

\end{cases}

$$



This yields the following expectations and variances:



$$

\mu\_X = \frac{2}{3}

$$



$$

\mu\_Y = 0

$$



$$

\sigma\_X^2 = \frac{2}{9}

$$



$$

\sigma\_Y^2 = \frac{2}{3}

$$



Therefore:



$$

\rho\_{X,Y} = \frac{1}{\sigma\_X \sigma\_Y} \operatorname{E}[(X - \mu\_X)(Y - \mu\_Y)]

$$

$$

= \frac{1}{\sigma\_X \sigma\_Y} \sum\_{x,y} (x - \mu\_X)(y - \mu\_Y) \operatorname{P}(X=x, Y=y)

$$

$$

= \frac{1}{\sigma\_X \sigma\_Y} \left[(1 - \frac{2}{3})(-1 - 0)\frac{1}{3} + (0 - \frac{2}{3})(0 - 0)\frac{1}{3} + (1 - \frac{2}{3})(1 - 0)\frac{1}{3}\right]

= 0

$$



\paragraph{Rank correlation coefficients}



Rank correlation coefficients, such as Spearman''s rank correlation coefficient and Kendall''s rank correlation coefficient ($\tau$), measure the extent to which, as one variable increases, the other variable tends to increase, without requiring that increase to be represented by a linear relationship. If, as one variable increases, the other decreases, the rank correlation coefficients will be negative. It is common to regard these rank correlation coefficients as alternatives to Pearson''s coefficient, used either to reduce the amount of calculation or to make the coefficient less sensitive to non-normality in distributions. However, this view has little mathematical basis, as rank correlation coefficients measure a different type of relationship than the Pearson product-moment correlation coefficient, and are best seen as measures of a different type of association, rather than as an alternative measure of the population correlation coefficient.



To illustrate the nature of rank correlation, and its difference from linear correlation, consider the following four pairs of numbers $(x, y)$:



$(0, 1), (10, 100), (101, 500), (102, 2000)$.



As we go from each pair to the next pair $x$ increases, and so does $y$. This relationship is perfect, in the sense that an increase in $x$ is always accompanied by an increase in $y$. This means that we have a perfect rank correlation, and both Spearman''s and Kendall''s correlation coefficients are 1, whereas in this example Pearson product-moment correlation coefficient is 0.7544, indicating that the points are far from lying on a straight line. In the same way, if $y$ always decreases when $x$ increases, the rank correlation coefficients will be $-1$, while the Pearson product-moment correlation coefficient may or may not be close to $-1$, depending on how close the points are to a straight line. Although in the extreme cases of perfect rank correlation the two coefficients are both equal (being both +1 or both -1), this is not generally the case, and so values of the two coefficients cannot meaningfully be compared. For example, for the three pairs $(1, 1), (2, 3), (3, 2)$ Spearman''s coefficient is $\frac{1}{2}$, while Kendall''s coefficient is $\frac{1}{3}$.



\paragraph{Other measures of dependence among random variables}



The information given by a correlation coefficient is not enough to define the dependence structure between random variables. The correlation coefficient completely defines the dependence structure only in very particular cases, for example when the distribution is a multivariate normal distribution. (See diagram above.) In the case of elliptical distributions, it characterizes the (hyper-)ellipses of equal density; however, it does not completely characterize the dependence structure (for example, a multivariate t-distribution''s degrees of freedom determine the level of tail dependence).



Distance correlation was introduced to address the deficiency of Pearson''s correlation that it can be zero for dependent random variables; zero distance correlation implies independence.



The Randomized Dependence Coefficient is a computationally efficient, copula-based measure of dependence between multivariate random variables. RDC is invariant with respect to non-linear scalings of random variables, is capable of discovering a wide range of functional association patterns and takes value zero at independence.



For two binary variables, the odds ratio measures their dependence, and takes non-negative numbers, possibly infinity: $[0, +\infty]$. Related statistics such as Yule''s $Y$ and Yule''s $Q$ normalize this to the correlation-like range $[-1, 1]$. The odds ratio is generalized by the logistic model to model cases where the dependent variables are discrete and there may be one or more independent variables.



The correlation ratio, entropy-based mutual information, total correlation, dual total correlation, and polychoric correlation are all also capable of detecting more general dependencies, as is consideration of the copula between them, while the coefficient of determination generalizes the correlation coefficient to multiple regression.



\paragraph{Sensitivity to the data distribution}



The degree of dependence between variables $X$ and $Y$ does not depend on the scale on which the variables are expressed. That is, if we are analyzing the relationship between $X$ and $Y$, most correlation measures are unaffected by transforming $X$ to $a + bX$ and $Y$ to $c + dY$, where $a$, $b$, $c$, and $d$ are constants ($b$ and $d$ being positive). This is true of some correlation statistics as well as their population analogues. Some correlation statistics, such as the rank correlation coefficient, are also invariant to monotone transformations of the marginal distributions of $X$ and/or $Y$.



Most correlation measures are sensitive to the manner in which $X$ and $Y$ are sampled. Dependencies tend to be stronger if viewed over a wider range of values. Thus, if we consider the correlation coefficient between the heights of fathers and their sons over all adult males, and compare it to the same correlation coefficient calculated when the fathers are selected to be between $165$ cm and $170$ cm in height, the correlation will be weaker in the latter case. Several techniques have been developed that attempt to correct for range restriction in one or both variables, and are commonly used in meta-analysis; the most common are Thorndike''s case II and case III equations.



Various correlation measures in use may be undefined for certain joint distributions of $X$ and $Y$. For example, the Pearson correlation coefficient is defined in terms of moments, and hence will be undefined if the moments are undefined. Measures of dependence based on quantiles are always defined. Sample-based statistics intended to estimate population measures of dependence may or may not have desirable statistical properties such as being unbiased, or asymptotically consistent, based on the spatial structure of the population from which the data were sampled.



Sensitivity to the data distribution can be used to an advantage. For example, scaled correlation is designed to use the sensitivity to the range in order to pick out correlations between fast components of time series. By reducing the range of values in a controlled manner, the correlations on long time scale are filtered out and only the correlations on short time scales are revealed.



% [Image: pearson]



\textbf{Pearson/Spearman correlation coefficients between $X$ and $Y$ are shown when the two variables'' ranges are unrestricted, and when the range of $X$ is restricted to the interval $(0,1)$.}



\paragraph{Correlation matrices}



The correlation matrix of $n$ random variables $X\_1, \ldots, X\_n$ is the $n \times n$ matrix whose $(i, j)$ entry is $\operatorname{corr}(X\_i, X\_j)$. Thus the diagonal entries are all identically unity. If the measures of correlation used are product-moment coefficients, the correlation matrix is the same as the covariance matrix of the standardized random variables $X\_i / \sigma(X\_i)$ for $i = 1, \ldots, n$. This applies both to the matrix of population correlations (in which case $\sigma$ is the population standard deviation), and to the matrix of sample correlations (in which case $\sigma$ denotes the sample standard deviation). Consequently, each is necessarily a positive-semidefinite matrix. Moreover, the correlation matrix is strictly positive definite if no variable can have all its values exactly generated as a linear function of the values of the others.



The correlation matrix is symmetric because the correlation between $X\_i$ and $X\_j$ is the same as the correlation between $X\_j$ and $X\_i$.



A correlation matrix appears, for example, in one formula for the coefficient of multiple determination, a measure of goodness of fit in multiple regression.



In statistical modelling, correlation matrices representing the relationships between variables are categorized into different correlation structures, which are distinguished by factors such as the number of parameters required to estimate them. For example, in an exchangeable correlation matrix, all pairs of variables are modeled as having the same correlation, so all non-diagonal elements of the matrix are equal to each other. On the other hand, an autoregressive matrix is often used when variables represent a time series, since correlations are likely to be greater when measurements are closer in time. Other examples include independent, unstructured, M-dependent, and Toeplitz.



In exploratory data analysis, the iconography of correlations consists in replacing a correlation matrix by a diagram where the “remarkable” correlations are represented by a solid line (positive correlation), or a dotted line (negative correlation).



\paragraph{Nearest valid correlation matrix}



In some applications (e.g., building data models from only partially observed data) one wants to find the "nearest" correlation matrix to an "approximate" correlation matrix (e.g., a matrix which typically lacks semi-definite positiveness due to the way it has been computed).



In $2002$, Higham formalized the notion of nearness using the Frobenius norm and provided a method for computing the nearest correlation matrix using the Dykstra''s projection algorithm, of which an implementation is available as an online Web API.



This sparked interest in the subject, with new theoretical (e.g., computing the nearest correlation matrix with factor structure) and numerical (e.g., usage of Newton''s method for computing the nearest correlation matrix) results obtained in the subsequent years.



\paragraph{Uncorrelatedness and independence of stochastic processes}



Similarly for two stochastic processes $({X\_t})\_{t \in \mathcal{T}} \text{ and } ({Y\_t})\_{t \in \mathcal{T}}$: If they are independent, then they are uncorrelated. The opposite of this statement might not be true. Even if two variables are uncorrelated, they might not be independent of each other.



\paragraph{Common misconceptions}



\paragraph{Correlation and causality}



The conventional dictum that "correlation does not imply causation" means that correlation cannot be used by itself to infer a causal relationship between the variables. This dictum should not be taken to mean that correlations cannot indicate the potential existence of causal relations. However, the causes underlying the correlation, if any, may be indirect and unknown, and high correlations also overlap with identity relations (tautologies), where no causal process exists. Consequently, a correlation between two variables is not a sufficient condition to establish a causal relationship (in either direction).



A correlation between age and height in children is fairly causally transparent, but a correlation between mood and health in people is less so. Does improved mood lead to improved health, or does good health lead to good mood, or both? Or does some other factor underlie both? In other words, a correlation can be taken as evidence for a possible causal relationship, but cannot indicate what the causal relationship, if any, might be.



\paragraph{Simple linear correlations}



% [Image: data]



\textbf{Four sets of data with the same correlation of $0.816$}



The Pearson correlation coefficient indicates the strength of a linear relationship between two variables, but its value generally does not completely characterize their relationship. In particular, if the conditional mean of $Y$ given $X$, denoted $\operatorname{E}(Y \mid X)$, is not linear in $X$, the correlation coefficient will not fully determine the form of $\operatorname{E}(Y \mid X)$.



The adjacent image shows scatter plots of Anscombe''s quartet, a set of four different pairs of variables created by Francis Anscombe. The four $y$ variables have the same mean ($7.5$), variance ($4.12$), correlation ($0.816$) and regression line ($y = 3 + 0.5x$). However, as can be seen on the plots, the distribution of the variables is very different. The first one (top left) seems to be distributed normally, and corresponds to what one would expect when considering two variables correlated and following the assumption of normality. The second one (top right) is not distributed normally; while an obvious relationship between the two variables can be observed, it is not linear. In this case, the Pearson correlation coefficient does not indicate that there is an exact functional relationship: only the extent to which that relationship can be approximated by a linear relationship. In the third case (bottom left), the linear relationship is perfect, except for one outlier which exerts enough influence to lower the correlation coefficient from $1$ to $0.816$. Finally, the fourth example (bottom right) shows another example when one outlier is enough to produce a high correlation coefficient, even though the relationship between the two variables is not linear.



These examples indicate that the correlation coefficient, as a summary statistic, cannot replace visual examination of the data. The examples are sometimes said to demonstrate that the Pearson correlation assumes that the data follow a normal distribution, but this is only partially correct. The Pearson correlation can be accurately calculated for any distribution that has a finite covariance matrix, which includes most distributions encountered in practice. However, the Pearson correlation coefficient (taken together with the sample mean and variance) is only a sufficient statistic if the data is drawn from a multivariate normal distribution. As a result, the Pearson correlation coefficient fully characterizes the relationship between variables if and only if the data are drawn from a multivariate normal distribution.



\paragraph{Bivariate normal distribution}



If a pair $(X, Y)$ of random variables follows a bivariate normal distribution, the conditional mean $\operatorname{E}(X \mid Y)$ is a linear function of $Y$, and the conditional mean $\operatorname{E}(Y \mid X)$ is a linear function of $X$. The correlation coefficient $\rho\\_{X,Y}$ between $X$ and $Y$, along with the marginal means and variances of $X$ and $Y$, determines this linear relationship:



$$\operatorname{E}(Y \mid X) = \operatorname{E}(Y) + \rho\_{X,Y} \cdot \frac{\sigma\_Y}{\sigma\_X} \cdot (X - \operatorname{E}(X)),$$



where $\operatorname{E}(X)$ and $\operatorname{E}(Y)$ are the expected values of $X$ and $Y$, respectively, and $\sigma\_X$ and $\sigma\_Y$ are the standard deviations of $X$ and $Y$, respectively.



\paragraph{Standard error}



If $x$ and $y$ are random variables, a standard error is associated with the correlation that is:



$$SE\_{r} = \frac{1 - r^2}{\sqrt{n - 2}},$$



where $r$ is the correlation and $n$ the number of samples.



\paragraph{Licensing}



Content obtained and/or adapted from:



\begin{itemize}
    \item \href{https://en.wikipedia.org/wiki/Univariate}{Univariate, Wikipedia} under a CC BY-SA license
    \item \href{https://en.wikibooks.org/wiki/Applicable_Mathematics/Bivariate_Data}{Bivariate Data, Wikibooks} under a CC BY-SA license
    \item \href{https://en.wikibooks.org/wiki/Econometric_Theory/Regression_versus_Causation_and_Correlation}{Regression versus Causation and Correlation, Wikibooks} under a CC BY-SA license
    \item \href{https://en.wikipedia.org/wiki/Correlation}{Correlation, Wikipedia} under a CC BY-SA license
\end{itemize}

\subsection{Learning Objectives}

\begin{enumerate}
    \item SLO 1: Measures of Location and Variation Students will be able to compute and interpret measures of location (mean, median, mode) and measures of variation (range, interquartile range, standard deviation) for univariate data. They will understand how these measures describe the central tendency and dispersion of a dataset. SLO 2: Correlation and Regression Analysis Students will be able to calculate and interpret the Pearson correlation coefficient between two variables, and distinguish between correlation and causation. They will apply linear regression techniques to analyze the relationship between bivariate data and understand how regression models differ from correlation. SLO 3: Visualizing and Analyzing Bivariate Data Students will be able to create and interpret scatterplots to visualize the relationship between two quantitative variables. They will be able to identify patterns, trends, and potential outliers in bivariate data and understand how these visualizations relate to correlation coefficients and regression lines.
\end{enumerate}

\subsection{Assessment Instrument}

\begin{enumerate}
    \item Assessment Instrument Part 1: Univariate Data Analysis \textbf{1.1 Compute the following for the given dataset:} $[12, 15, 18, 20, 22, 25, 30]$ a) Mean b) Median c) Mode d) Range e) Standard Deviation \textbf{1.2 Explain what the computed measures reveal about the dataset in terms of central tendency and dispersion.} Part 2: Correlation and Regression Analysis \textbf{2.1 Given the following bivariate data:} $$ \begin{array}{|c|c|} \hline X & Y \\\\ \hline 2 & 3 \\\\ 4 & 5 \\\\ 6 & 8 \\\\ 8 & 10 \\\\ 10 & 12 \\\\ \hline \end{array} $$ a) Compute the Pearson correlation coefficient between $X$ and $Y$. b) Perform a simple linear regression analysis to find the line of best fit for the data. Provide the equation of the regression line. \textbf{2.2 Interpret the correlation coefficient and the regression line. Discuss the relationship between $X$ and $Y$ and whether or not you can infer causation from this analysis.} Part 3: Visualizing and Analyzing Bivariate Data \textbf{3.1 Create a scatterplot for the following dataset:} $$ \begin{array}{|c|c|} \hline X & Y \\\\ \hline 1 & 2 \\\\ 2 & 4 \\\\ 3 & 6 \\\\ 4 & 8 \\\\ 5 & 10 \\\\ \hline \end{array} $$ \textbf{3.2 Analyze the scatterplot to determine the type of relationship between $X$ and $Y$. Identify any trends or patterns and explain how these relate to the computed correlation coefficient and regression analysis.} \textbf{3.3 Discuss how changes in the scale or range of the data might affect the correlation coefficient and the interpretation of the scatterplot.}
\end{enumerate}

%%===============================================================
\section{Time series model of exponential growth}
\label{sec:time-series-model-of-exponential-growth}
%%===============================================================

\paragraph{Time Series Model of Exponential Growth}



% [Image: time]



\textbf{Time Series: Random Data Plus Trend, with Best-Fit Line and Different Applied Filters}



In mathematics, a time series is a series of data points indexed (or listed or graphed) in time order. Most commonly, a time series is a sequence taken at successive equally spaced points in time. Thus it is a sequence of discrete-time data. Examples of time series are heights of ocean tides, counts of sunspots, and the daily closing value of the Dow Jones Industrial Average.



A time series is very frequently plotted via a run chart (which is a temporal line chart). Time series are used in statistics, signal processing, pattern recognition, econometrics, mathematical finance, weather forecasting, earthquake prediction, electroencephalography, control engineering, astronomy, communications engineering, and largely in any domain of applied science and engineering which involves temporal measurements.



Time series analysis comprises methods for analyzing time series data in order to extract meaningful statistics and other characteristics of the data. Time series forecasting is the use of a model to predict future values based on previously observed values. While regression analysis is often employed in such a way as to test relationships between one or more different time series, this type of analysis is not usually called "time series analysis", which refers in particular to relationships between different points in time within a single series. Interrupted time series analysis is used to detect changes in the evolution of a time series from before to after some intervention which may affect the underlying variable.



Time series data have a natural temporal ordering. This makes time series analysis distinct from cross-sectional studies, in which there is no natural ordering of the observations (e.g., explaining people's wages by reference to their respective education levels, where the individuals' data could be entered in any order). Time series analysis is also distinct from spatial data analysis where the observations typically relate to geographical locations (e.g., accounting for house prices by the location as well as the intrinsic characteristics of the houses). A stochastic model for a time series will generally reflect the fact that observations close together in time will be more closely related than observations further apart. In addition, time series models will often make use of the natural one-way ordering of time so that values for a given period will be expressed as deriving in some way from past values, rather than from future values (see time reversibility).



Time series analysis can be applied to real-valued, continuous data, discrete numeric data, or discrete symbolic data (i.e., sequences of characters, such as letters and words in the English language).







\paragraph{Methods for Analysis}



Methods for time series analysis may be divided into two classes: frequency-domain methods and time-domain methods. The former include spectral analysis and wavelet analysis; the latter include auto-correlation and cross-correlation analysis. In the time domain, correlation and analysis can be made in a filter-like manner using scaled correlation, thereby mitigating the need to operate in the frequency domain.



Additionally, time series analysis techniques may be divided into parametric and non-parametric methods. The parametric approaches assume that the underlying stationary stochastic process has a certain structure which can be described using a small number of parameters (for example, using an autoregressive or moving average model). In these approaches, the task is to estimate the parameters of the model that describes the stochastic process. By contrast, non-parametric approaches explicitly estimate the covariance or the spectrum of the process without assuming that the process has any particular structure.



Methods of time series analysis may also be divided into linear and non-linear, and univariate and multivariate.



\paragraph{Panel Data}



A time series is one type of panel data. Panel data is the general class, a multidimensional data set, whereas a time series data set is a one-dimensional panel (as is a cross-sectional dataset). A data set may exhibit characteristics of both panel data and time series data. One way to tell is to ask what makes one data record unique from the other records. If the answer is the time data field, then this is a time series data set candidate. If determining a unique record requires a time data field and an additional identifier which is unrelated to time (e.g., student ID, stock symbol, country code), then it is a panel data candidate. If the differentiation lies on the non-time identifier, then the data set is a cross-sectional data set candidate.



\paragraph{Analysis}



There are several types of motivation and data analysis available for time series which are appropriate for different purposes.



\paragraph{Motivation}



In the context of statistics, econometrics, quantitative finance, seismology, meteorology, and geophysics the primary goal of time series analysis is forecasting. In the context of signal processing, control engineering and communication engineering it is used for signal detection. Other applications are in data mining, pattern recognition and machine learning, where time series analysis can be used for clustering, classification, query by content, anomaly detection as well as forecasting.



\paragraph{Exploratory Analysis}



% [Image: TB]



\textbf{Tuberculosis incidence US 1953-2009}



A straightforward way to examine a regular time series is manually with a line chart. An example chart is shown on the right for tuberculosis incidence in the United States, made with a spreadsheet program. The number of cases was standardized to a rate per 100,000 and the percent change per year in this rate was calculated. The nearly steadily dropping line shows that the TB incidence was decreasing in most years, but the percent change in this rate varied by as much as +/- 10\%, with 'surges' in 1975 and around the early 1990s. The use of both vertical axes allows the comparison of two time series in one graphic.



Other techniques include:



\begin{itemize}
    \item Autocorrelation analysis to examine serial dependence
    \item Spectral analysis to examine cyclic behavior which need not be related to seasonality. For example, sunspot activity varies over 11 year cycles. Other common examples include celestial phenomena, weather patterns, neural activity, commodity prices, and economic activity.
    \item Separation into components representing trend, seasonality, slow and fast variation, and cyclical irregularity: see trend estimation and decomposition of time series
\end{itemize}

\paragraph{Curve Fitting}



Curve fitting is the process of constructing a curve, or mathematical function, that has the best fit to a series of data points, possibly subject to constraints. Curve fitting can involve either interpolation, where an exact fit to the data is required, or smoothing, in which a "smooth" function is constructed that approximately fits the data. A related topic is regression analysis, which focuses more on questions of statistical inference such as how much uncertainty is present in a curve that is fit to data observed with random errors. Fitted curves can be used as an aid for data visualization, to infer values of a function where no data are available, and to summarize the relationships among two or more variables. Extrapolation refers to the use of a fitted curve beyond the range of the observed data, and is subject to a degree of uncertainty since it may reflect the method used to construct the curve as much as it reflects the observed data.



The construction of economic time series involves the estimation of some components for some dates by interpolation between values ("benchmarks") for earlier and later dates. Interpolation is estimation of an unknown quantity between two known quantities (historical data), or drawing conclusions about missing information from the available information ("reading between the lines"). Interpolation is useful where the data surrounding the missing data is available and its trend, seasonality, and longer-term cycles are known. This is often done by using a related series known for all relevant dates. Alternatively polynomial interpolation or spline interpolation is used where piecewise polynomial functions are fit into time intervals such that they fit smoothly together. A different problem which is closely related to interpolation is the approximation of a complicated function by a simple function (also called regression). The main difference between regression and interpolation is that polynomial regression gives a single polynomial that models the entire data set. Spline interpolation, however, yield a piecewise continuous function composed of many polynomials to model the data set.



Extrapolation is the process of estimating, beyond the original observation range, the value of a variable on the basis of its relationship with another variable. It is similar to interpolation, which produces estimates between known observations, but extrapolation is subject to greater uncertainty and a higher risk of producing meaningless results.



\paragraph{Function Approximation}



In general, a function approximation problem asks us to select a function among a well-defined class that closely matches ("approximates") a target function in a task-specific way. One can distinguish two major classes of function approximation problems:



1. For known target functions, approximation theory is the branch of numerical analysis that investigates how certain known functions (for example, special functions) can be approximated by a specific class of functions (for example, polynomials or rational functions) that often have desirable properties (inexpensive computation, continuity, integral and limit values, etc.).



2. The target function, call it $g$, may be unknown; instead of an explicit formula, only a set of points (a time series) of the form $(x, g(x))$ is provided. Depending on the structure of the domain and codomain of $g$, several techniques for approximating $g$ may be applicable. For example, if $g$ is an operation on the real numbers, techniques of interpolation, extrapolation, regression analysis, and curve fitting can be used. If the codomain (range or target set) of $g$ is a finite set, one is dealing with a classification problem instead. A related problem of online time series approximation is to summarize the data in one-pass and construct an approximate representation that can support a variety of time series queries with bounds on worst-case error.



To some extent, the different problems (regression, classification, fitness approximation) have received a unified treatment in statistical learning theory, where they are viewed as supervised learning problems.



\paragraph{Prediction and Forecasting}



In statistics, prediction is a part of statistical inference. One particular approach to such inference is known as predictive inference, but the prediction can be undertaken within any of the several approaches to statistical inference. Indeed, one description of statistics is that it provides a means of transferring knowledge about a sample of a population to the whole population, and to other related populations, which is not necessarily the same as prediction over time. When information is transferred across time, often to specific points in time, the process is known as forecasting.



\begin{itemize}
    \item Fully formed statistical models for stochastic simulation purposes, so as to generate alternative versions of the time series, representing what might happen over non-specific time-periods in the future
    \item Simple or fully formed statistical models to describe the likely outcome of the time series in the immediate future, given knowledge of the most recent outcomes (forecasting).
    \item Forecasting on time series is usually done using automated statistical software packages and programming languages, such as Julia, Python, R, SAS, SPSS and many others.
    \item Forecasting on large scale data can be done with Apache Spark using the Spark-TS library, a third-party package.
\end{itemize}

\paragraph{Classification}



Assigning time series pattern to a specific category, for example identify a word based on series of hand movements in sign language.



\paragraph{Signal Estimation}



This approach is based on harmonic analysis and filtering of signals in the frequency domain using the Fourier transform, and spectral density estimation, the development of which was significantly accelerated during World War II by mathematician Norbert Wiener, electrical engineers Rudolf E. Kálmán, Dennis Gabor and others for filtering signals from noise and predicting signal values at a certain point in time.



\paragraph{Segmentation}



Splitting a time-series into a sequence of segments. It is often the case that a time-series can be represented as a sequence of individual segments, each with its own characteristic properties. For example, the audio signal from a conference call can be partitioned into pieces corresponding to the times during which each person was speaking. In time-series segmentation, the goal is to identify the segment boundary points in the time-series, and to characterize the dynamical properties associated with each segment. One can approach this problem using change-point detection, or by modeling the time-series as a more sophisticated system, such as a Markov jump linear system.



\paragraph{Models}



Models for time series data can have many forms and represent different stochastic processes. When modeling variations in the level of a process, three broad classes of practical importance are the autoregressive (AR) models, the integrated (I) models, and the moving average (MA) models. These three classes depend linearly on previous data points. Combinations of these ideas produce autoregressive moving average (ARMA) and autoregressive integrated moving average (ARIMA) models. The autoregressive fractionally integrated moving average (ARFIMA) model generalizes the former three. Extensions of these classes to deal with vector-valued data are available under the heading of multivariate time-series models and sometimes the preceding acronyms are extended by including an initial "V" for "vector", as in VAR for vector autoregression. An additional set of extensions of these models is available for use where the observed time-series is driven by some "forcing" time-series (which may not have a causal effect on the observed series): the distinction from the multivariate case is that the forcing series may be deterministic or under the experimenter's control. For these models, the acronyms are extended with a final "X" for "exogenous".



Non-linear dependence of the level of a series on previous data points is of interest, partly because of the possibility of producing a chaotic time series. However, more importantly, empirical investigations can indicate the advantage of using predictions derived from non-linear models, over those from linear models, as for example in nonlinear autoregressive exogenous models. Further references on nonlinear time series analysis: (Kantz and Schreiber), and (Abarbanel).



Among other types of non-linear time series models, there are models to represent the changes of variance over time (heteroskedasticity). These models represent autoregressive conditional heteroskedasticity (ARCH) and the collection comprises a wide variety of representation (GARCH, TARCH, EGARCH, FIGARCH, CGARCH, etc.). Here changes in variability are related to, or predicted by, recent past values of the observed series. This is in contrast to other possible representations of locally varying variability, where the variability might be modelled as being driven by a separate time-varying process, as in a doubly stochastic model.



In recent work on model-free analyses, wavelet transform based methods (for example locally stationary wavelets and wavelet decomposed neural networks) have gained favor. Multiscale (often referred to as multiresolution) techniques decompose a given time series, attempting to illustrate time dependence at multiple scales. See also Markov switching multifractal (MSMF) techniques for modeling volatility evolution.



A Hidden Markov model (HMM) is a statistical Markov model in which the system being modeled is assumed to be a Markov process with unobserved (hidden) states. An HMM can be considered as the simplest dynamic Bayesian network. HMM models are widely used in speech recognition, for translating a time series of spoken words into text.



\paragraph{Notation}



A number of different notations are in use for time-series analysis. A common notation specifying a time series $X$ that is indexed by the natural numbers is written



$$X = (X\_1, X\_2, \ldots).$$



Another common notation is



$$Y = (Y\_t: t \in T),$$



where $T$ is the index set.



\paragraph{Conditions}



There are two sets of conditions under which much of the theory is built:



\begin{itemize}
    \item Stationary process
    \item Ergodic process
\end{itemize}

However, ideas of stationarity must be expanded to consider two important ideas: strict stationarity and second-order stationarity. Both models and applications can be developed under each of these conditions, although the models in the latter case might be considered as only partly specified.



In addition, time-series analysis can be applied where the series are seasonally stationary or non-stationary. Situations where the amplitudes of frequency components change with time can be dealt with in time-frequency analysis which makes use of a time–frequency representation of a time-series or signal.



\paragraph{Visualization}



Time series can be visualized with two categories of chart: Overlapping Charts and Separated Charts. Overlapping Charts display all-time series on the same layout while Separated Charts presents them on different layouts (but aligned for comparison purposes).



\paragraph{Overlapping Charts}



\begin{itemize}
    \item Braided graphs
    \item Line charts
    \item Slope graphs
    \item GapChart
\end{itemize}

\paragraph{Separated Charts}



\begin{itemize}
    \item Horizon graphs
    \item Reduced line chart (small multiples)
    \item Silhouette graph
    \item Circular silhouette graph
\end{itemize}

\paragraph{Licensing}



Content obtained and/or adapted from:



\href{https://en.wikipedia.org/wiki/Time_series}{Time series, Wikipedia} under a CC BY-SA license

\subsection{Learning Objectives}

\begin{enumerate}
    \item SLO 1: Understanding Time Series and Exponential Growth Define a time series and describe its application in various fields such as economics, meteorology, and finance. Explain the concept of exponential growth and how it can be represented and analyzed within a time series context. SLO 2: Analyzing Trends and Filtering in Time Series Analyze a time series data set to identify underlying trends, seasonal variations, and noise. Apply appropriate filters to the data to isolate the trend component and evaluate its impact on predictions and forecasting. SLO 3: Curve Fitting and Exponential Models Use curve fitting techniques to model a time series with exponential growth. Compare different fitting methods (e.g., polynomial vs. exponential) and assess their accuracy and applicability to forecasting future values.
\end{enumerate}

\subsection{Assessment Instrument}

\begin{enumerate}
    \item Assessment for SLO 1 \textbf{Example Data:}
    \item Time: [1, 2, 3, 4, 5, 6, 7, 8, 9, 10]
    \item Values: [10, 15, 22, 32, 47, 68, 98, 142, 205, 290] \textbf{Questions:}
    \item What is a time series, and how is it used in forecasting?
    \item Identify and describe the exponential growth pattern in the provided data. Assessment for SLO 2 \textbf{Example Data:}
    \item Time: [Jan, Feb, Mar, Apr, May, Jun, Jul, Aug, Sep, Oct]
    \item Values: [120, 130, 145, 160, 175, 190, 210, 225, 240, 260] \textbf{Questions:}
    \item Describe the trend observed in the time series data.
    \item Apply a moving average filter with a window size of 3 to the data. What is the smoothed trend? Assessment for SLO 3 \textbf{Example Data:}
    \item Time: [1, 2, 3, 4, 5, 6, 7, 8, 9, 10]
    \item Values: [12, 18, 27, 38, 52, 70, 92, 118, 148, 182] \textbf{Questions:}
    \item Fit both a polynomial and an exponential model to the provided time series data. Provide the equations of the models.
    \item Compare the models based on their goodness of fit (e.g., R-squared value). Which model provides a better fit for the data? Justify your answer.
    \item Use the chosen model to forecast the next 3 time points and provide the predicted values.
\end{enumerate}

%%===============================================================
\section{Linear and Exponential Models}
\label{sec:linear-and-exponential-models}
%%===============================================================

\paragraph{Linear and Exponential Models}



\paragraph{\href{https://www.khanacademy.org/math/algebra/x2f8bb11595b61c86:exponential-growth-decay/x2f8bb11595b61c86:exponential-vs-linear-growth/v/linear-exponential-models}{Exponential vs. Linear Models: Verbal}, Khan Academy}







\paragraph{\href{https://www.khanacademy.org/math/algebra/x2f8bb11595b61c86:exponential-growth-decay/x2f8bb11595b61c86:exponential-vs-linear-growth/v/exponential-vs-linear-models}{Exponential vs. Linear Models: Table}, Khan Academy}







\paragraph{\href{https://www.khanacademy.org/math/algebra/x2f8bb11595b61c86:exponential-growth-decay/x2f8bb11595b61c86:exponential-vs-linear-models/v/linear-and-exponential-growth-from-data}{Linear vs. Exponential Growth: From Data}, Khan Academy}

\subsection{Learning Objectives}

\begin{enumerate}
    \item Explore the differences between linear and exponential models using verbal descriptions. Learn how to identify whether a situation can be modeled by a linear or exponential function based on the context.
    \item Compare linear and exponential growth by analyzing data in table form. Understand how linear growth adds a constant amount while exponential growth multiplies by a constant factor.
    \item Analyze real-world data to distinguish between linear and exponential growth. Learn how to create and interpret graphs that represent linear and exponential relationships.
\end{enumerate}

\subsection{Assessment Instrument}

\begin{enumerate}
    \item Question 1 Identify whether the following scenarios describe a linear or exponential relationship:
    \item A bank account earns `$`5 in interest every month.
    \item A population of bacteria doubles in size every hour.
    \item A car rental company charges a fixed fee of `$`20 per day plus `$`0.10 per mile driven.
    \item The value of a stock increases by 2\% each day. Question 2 Given the following situation, determine if it is best modeled by a linear or exponential function, and explain why: "A company’s sales increase by 15\% each quarter." Question 3 Given the following data, identify whether the growth is linear or exponential: $$ \begin{array}{|c|c|} \hline \text{Time (days)} & \text{Amount} \\\\ \hline 0 & 100 \\\\ \hline 1 & 150 \\\\ \hline 2 & 200 \\\\ \hline 3 & 250 \\\\ \hline 4 & 300 \\\\ \hline \end{array} $$ Explain your reasoning. Question 4 Analyze the following data and determine whether the data represents linear or exponential growth: $$ \begin{array}{|c|c|} \hline \text{Time (years)} & \text{Population} \\\\ \hline 0 & 1000 \\\\ \hline 1 & 1200 \\\\ \hline 2 & 1440 \\\\ \hline 3 & 1728 \\\\ \hline 4 & 2074 \\\\ \hline \end{array} $$ Explain your conclusion based on the pattern observed. Question 5 Plot the following data points on a graph and determine if the relationship is linear or exponential: $$ \begin{array}{|c|c|} \hline \text{Year} & \text{Value} \\\\ \hline 2010 & 500 \\\\ \hline 2011 & 550 \\\\ \hline 2012 & 605 \\\\ \hline 2013 & 665.5 \\\\ \hline 2014 & 732.05 \\\\ \hline \end{array} $$ Explain your reasoning based on the shape of the graph. Question 6 Given the real-world data below, create a graph and explain whether it shows linear or exponential growth: $$ \begin{array}{|c|c|} \hline \text{Month} & \text{Subscribers} \\\\ \hline \text{Jan} & 200 \\\\ \hline \text{Feb} & 400 \\\\ \hline \text{Mar} & 800 \\\\ \hline \text{Apr} & 1600 \\\\ \hline \text{May} & 3200 \\\\ \hline \end{array} $$ Describe the characteristics that led you to your conclusion.
\end{enumerate}

%%===============================================================
\section{Continuous Growth}
\label{sec:continuous-growth}
%%===============================================================

\paragraph{Continuous Growth}



% [Image: growth]



\textbf{The graph illustrates how exponential growth (green) surpasses both linear (red) and cubic (blue) growth. Exponential growth is a process that increases quantity over time. It occurs when the instantaneous rate of change (that is, the derivative) of a quantity with respect to time is proportional to the quantity itself. Described as a function, a quantity undergoing exponential growth is an exponential function of time, that is, the variable representing time is the exponent (in contrast to other types of growth, such as quadratic growth).}



If the constant of proportionality is negative, then the quantity decreases over time, and is said to be undergoing exponential decay instead. In the case of a discrete domain of definition with equal intervals, it is also called geometric growth or geometric decay since the function values form a geometric progression.



The formula for exponential growth of a variable $x$ at the growth rate $r$, as time $t$ goes on in discrete intervals (that is, at integer times 0, 1, 2, 3, ...), is:



$$

x\_t = x\_0 (1 + r)^t

$$



where $x\_0$ is the value of $x$ at time $0$. The growth of a colony is often used to illustrate this. One bacterium splits itself into two, each of which splits itself resulting in four, then eight, 16, 32, and so on. The rate of increase keeps increasing because it is proportional to the ever-increasing number of bacteria. Growth like this is observed in real-life activity or phenomena, such as the spread of virus infection, the growth of debt due to compound interest, and the spread of viral videos. In real cases, initial exponential growth often does not last forever, instead slowing down eventually due to upper limits caused by external factors and turning into logistic growth.



Terms like ''''exponential growth'''' are sometimes incorrectly interpreted as ''''rapid growth''''. Indeed, something that is growing exponentially can in fact be growing slowly.





\paragraph{Examples}



\paragraph{Biology}



% [Image: bacteria]



\textbf{Bacteria exhibit exponential growth under optimal conditions.}



\begin{itemize}
    \item The number of microorganisms in a culture will increase exponentially until an essential nutrient is exhausted. Typically, the first organism splits into two daughter organisms, who then each split to form four, who split to form eight, and so on. Because exponential growth indicates constant growth rate, it is frequently assumed that exponentially growing cells are at a steady-state. However, cells can grow exponentially at a constant rate while remodeling their metabolism and gene expression.
\end{itemize}

\begin{itemize}
    \item A virus (for example, COVID-19, or smallpox) typically will spread exponentially at first if no artificial immunization is available. Each infected person can infect multiple new people.
\end{itemize}

\paragraph{Physics}



\begin{itemize}
    \item Avalanche breakdown within a dielectric material. A free electron becomes sufficiently accelerated by an externally applied electrical field that it frees up additional electrons as it collides with atoms or molecules of the dielectric media. These secondary electrons also are accelerated, creating larger numbers of free electrons. The resulting exponential growth of electrons and ions may rapidly lead to complete dielectric breakdown of the material.
\end{itemize}

\begin{itemize}
    \item Nuclear chain reaction (the concept behind nuclear reactors and nuclear weapons). Each uranium nucleus that undergoes fission produces multiple neutrons, each of which can be absorbed by adjacent uranium atoms, causing them to fission in turn. If the probability of neutron absorption exceeds the probability of neutron escape (a function of the shape and mass of the uranium), the production rate of neutrons and induced uranium fissions increases exponentially, in an uncontrolled reaction. "Due to the exponential rate of increase, at any point in the chain reaction 99\% of the energy will have been released in the last 4.6 generations. It is a reasonable approximation to think of the first 53 generations as a latency period leading up to the actual explosion, which only takes 3–4 generations."
\end{itemize}

\begin{itemize}
    \item Positive feedback within the linear range of electrical or electroacoustic amplification can result in the exponential growth of the amplified signal, although resonance effects may favor some component frequencies of the signal over others.
\end{itemize}

\paragraph{Economics}



\begin{itemize}
    \item Economic growth is expressed in percentage terms, implying exponential growth.
\end{itemize}

\paragraph{Finance}



\begin{itemize}
    \item Compound interest at a constant interest rate provides exponential growth of the capital.
\end{itemize}

\begin{itemize}
    \item Pyramid schemes or Ponzi schemes also show this type of growth, resulting in high profits for a few initial investors and losses among great numbers of investors.
\end{itemize}

\paragraph{Computer Science}



\begin{itemize}
    \item Processing power of computers. See also Moore''''s law and technological singularity. (Under exponential growth, there are no singularities. The singularity here is a metaphor, meant to convey an unimaginable future. The link of this hypothetical concept with exponential growth is most vocally made by futurist Ray Kurzweil.)
\end{itemize}

\begin{itemize}
    \item In computational complexity theory, computer algorithms of exponential complexity require an exponentially increasing amount of resources (e.g., time, computer memory) for only a constant increase in problem size. So, for an algorithm of time complexity $2^x$, if a problem of size $x = 10$ requires 10 seconds to complete, and a problem of size $x = 11$ requires 20 seconds, then a problem of size $x = 12$ will require 40 seconds. This kind of algorithm typically becomes unusable at very small problem sizes, often between 30 and 100 items (most computer algorithms need to be able to solve much larger problems, up to tens of thousands or even millions of items in reasonable times, something that would be physically impossible with an exponential algorithm). Also, the effects of Moore''''s Law do not help the situation much because doubling processor speed merely allows you to increase the problem size by a constant. E.g., if a slow processor can solve problems of size $x$ in time $t$, then a processor twice as fast could only solve problems of size $x + \text{constant}$ in the same time $t$. So, exponentially complex algorithms are most often impractical, and the search for more efficient algorithms is one of the central goals of computer science today.
\end{itemize}

\paragraph{Internet Phenomena}



\begin{itemize}
    \item Internet contents, such as internet memes or videos, can spread in an exponential manner, often said to "go viral" as an analogy to the spread of viruses. With media such as social networks, one person can forward the same content to many people simultaneously, who then spread it to even more people, and so on, causing rapid spread. For example, the video Gangnam Style was uploaded to YouTube on July 15, 2012, reaching hundreds of thousands of viewers on the first day, millions on the twentieth day, and was cumulatively viewed by hundreds of millions in less than two months.
\end{itemize}

\paragraph{Basic Formula}



% [Image: graph1]



The graph above shows Exponential Growth with the following values:



$$

\begin{aligned}

a &= 3 \\\\ 

b &= 2 \\\\ 

r &= 5

\end{aligned}

$$





% [Image: graph2]



The graph above shows Exponential Growth with the following values:



$$

\begin{aligned}

a &= 24 \\\\ 

b &= \frac{1}{2} \\\\ 

r &= 5

\end{aligned}

$$



A quantity $ x $ depends exponentially on time $ t $ if:



$$

x(t) = a \cdot b^{t/\tau}

$$



where the constant $ a $ is the initial value of $ x $,



$$

x(0) = a

$$



the constant $ b $ is a positive growth factor and $ \tau $ is the time constant—the time required for $ x $ to increase by one factor of $ b $:



$$

x(t + \tau) = a \cdot b^{\frac{t + \tau}{\tau}} = a \cdot b^{\frac{t}{\tau}} \cdot b^{\frac{\tau}{\tau}} = x(t) \cdot b

$$



If $ \tau > 0 $ and $ b > 1 $, then $ x $ has exponential growth. If $ \tau < 0 $ and $ b > 1 $, or $ \tau > 0 $ and $ 0 < b < 1 $, then $ x $ has exponential decay.



\textbf{Example:} If a species of bacteria doubles every ten minutes, starting out with only one bacterium, how many bacteria would be present after one hour? The question implies $ a = 1 $, $ b = 2 $, and $ \tau = 10 $ min.



$$

x(t) = a \cdot b^{t/\tau} = 1 \cdot 2^{(60\ \text{min}) / (10\ \text{min})}

$$



$$

x(1\ \text{hr}) = 1 \cdot 2^6 = 64

$$



After one hour, or six ten-minute intervals, there would be sixty-four bacteria.



Many pairs $ (b, \tau) $ of a dimensionless non-negative number $ b $ and an amount of time $ \tau $ (a physical quantity which can be expressed as the product of a number of units and a unit of time) represent the same growth rate, with $ \tau $ proportional to $ \log b $. For any fixed $ b $ not equal to 1 (e.g., $ e $ or 2), the growth rate is given by the non-zero time $ \tau $. For any non-zero time $ \tau $, the growth rate is given by the dimensionless positive number $ b $.



Thus, the law of exponential growth can be written in different but mathematically equivalent forms, by using a different base. The most common forms are the following:



$$

x(t) = x\_0 \cdot e^{kt} = x\_0 \cdot e^{t/\tau} = x\_0 \cdot 2^{t/T} = x\_0 \cdot \left(1 + \frac{r}{100}\right)^{t/p}

$$



where $ x\_0 $ expresses the initial quantity $ x(0) $.



\textbf{Parameters} (negative in the case of exponential decay):



\begin{itemize}
    \item The growth constant $ k $ is the frequency (number of times per unit time) of growing by a factor $ e $; in finance, it is also called the logarithmic return, continuously compounded return, or force of interest.
    \item The $ e $-folding time $ \tau $ is the time it takes to grow by a factor $ e $.
    \item The doubling time $ T $ is the time it takes to double.
    \item The percent increase $ r $ (a dimensionless number) in a period $ p $.
\end{itemize}

The quantities $ k $, $ \tau $, and $ T $, and for a given $ p $ also $ r $, have a one-to-one connection given by the following equation (which can be derived by taking the natural logarithm of the above):



$$

k = \frac{1}{\tau} = \frac{\ln 2}{T} = \frac{\ln \left(1 + \frac{r}{100}\right)}{p}

$$



where $ k = 0 $ corresponds to $ r = 0 $ and to $ \tau $ and $ T $ being infinite.



If $ p $ is the unit of time, the quotient $ t/p $ is simply the number of units of time. Using the notation $ t $ for the (dimensionless) number of units of time rather than the time itself, $ t/p $ can be replaced by $ t $, but for uniformity, this has been avoided here. In this case, the division by $ p $ in the last formula is not a numerical division either, but converts a dimensionless number to the correct quantity including the unit.



A popular approximated method for calculating the doubling time from the growth rate is the rule of 70, that is, 



$$

T \simeq \frac{70}{r}

$$



\paragraph{Reformulation as Log-Linear Growth}



% [Image: graph3]



\textbf{Graphs comparing doubling times and half lives of exponential growths (bold lines) and decay (faint lines), and their 70/t and 72/t approximations. In the SVG version, hover over a graph to highlight it and its complement.}



If a variable $ x $ exhibits exponential growth according to 



$$

x(t) = x\_0 (1 + r)^t

$$



then the logarithm (to any base) of $ x $ grows linearly over time, as can be seen by taking logarithms of both sides of the exponential growth equation:



$$

\log x(t) = \log x\_0 + t \cdot \log(1 + r)

$$



This allows an exponentially growing variable to be modeled with a log-linear model. For example, if one wishes to empirically estimate the growth rate from intertemporal data on $ x $, one can linearly regress $ \log x $ on $ t $.



\paragraph{Differential Equation}



The exponential function 



$$

x(t) = x(0) e^{kt}

$$



satisfies the linear differential equation:



$$

\frac{dx}{dt} = kx

$$



saying that the change per instant of time of $ x $ at time $ t $ is proportional to the value of $ x(t) $, and $ x(t) $ has the initial value $ x(0) $.



The differential equation is solved by direct integration:



$$

\frac{dx}{dt} = kx

$$



$$

\frac{dx}{x} = k\,dt

$$



$$

\int\_{x(0)}^{x(t)} \frac{dx}{x} = k \int\_{0}^{t}\,dt

$$



$$

\ln \frac{x(t)}{x(0)} = kt

$$



so that



$$

x(t) = x(0) e^{kt}

$$



In the above differential equation, if $ k < 0 $, then the quantity experiences exponential decay.



For a nonlinear variation of this growth model, see the logistic function.





\paragraph{Other Growth Rates}



In the long run, exponential growth of any kind will overtake linear growth of any kind (that is the basis of the Malthusian catastrophe) as well as any polynomial growth, that is, for all \( \alpha \):



$$

\lim\_{t \rightarrow \infty} \frac{t^{\alpha}}{ae^{t}} = 0.

$$



There is a whole hierarchy of conceivable growth rates that are slower than exponential and faster than linear (in the long run).



Growth rates may also be faster than exponential. In the most extreme case, when growth increases without bound in finite time, it is called hyperbolic growth. In between exponential and hyperbolic growth lie more classes of growth behavior, like the hyperoperations beginning at tetration, and \( A(n,n) \), the diagonal of the Ackermann function.





\paragraph{Logistic Growth}



% [Image: graph4]



\textbf{The J-shaped exponential growth (left, blue) and the S-shaped logistic growth (right, red).}



In reality, initial exponential growth is often not sustained forever. After some period, it will be slowed by external or environmental factors. For example, population growth may reach an upper limit due to resource limitations. In 1845, the Belgian mathematician Pierre François Verhulst first proposed a mathematical model of growth like this, called the "logistic growth."





\paragraph{Limitations of Models}



Exponential growth models of physical phenomena only apply within limited regions, as unbounded growth is not physically realistic. Although growth may initially be exponential, the modeled phenomena will eventually enter a region in which previously ignored negative feedback factors become significant (leading to a logistic growth model) or other underlying assumptions of the exponential growth model, such as continuity or instantaneous feedback, break down.





\paragraph{Exponential Growth Bias}



Studies show that human beings have difficulty understanding exponential growth. Exponential growth bias is the tendency to underestimate compound growth processes. This bias can have financial implications as well. Below are some stories that emphasize this bias.





\paragraph{Rice on a Chessboard}



According to an old legend, vizier Sissa Ben Dahir presented an Indian King Sharim with a beautiful handmade chessboard. The king asked what he would like in return for his gift, and the courtier surprised the king by asking for one grain of rice on the first square, two grains on the second, four grains on the third, etc. The king readily agreed and asked for the rice to be brought. All went well at first, but the requirement for \(2^{n-1}\) grains on the \(n\)th square demanded over a million grains on the 21st square, more than a million million (aka trillion) on the 41st, and there simply was not enough rice in the whole world for the final squares. (From Swirski, 2006)



The second half of the chessboard is the time when an exponentially growing influence is having a significant economic impact on an organization''s overall business strategy.





\paragraph{Water Lily}



French children are offered a riddle, which appears to be an aspect of exponential growth: "the apparent suddenness with which an exponentially growing quantity approaches a fixed limit." The riddle imagines a water lily plant growing in a pond. The plant doubles in size every day and, if left alone, it would smother the pond in 30 days, killing all the other living things in the water. Day after day, the plant''s growth is small, so it is decided that it won''t be a concern until it covers half of the pond. Which day will that be? The 29th day, leaving only one day to save the pond.





\paragraph{Licensing}



Content obtained and/or adapted from:



\href{https://en.wikipedia.org/wiki/Exponential_growth}{Exponential growth, Wikipedia} under a CC BY-SA license.

\subsection{Learning Objectives}

\begin{enumerate}
    \item \textbf{Identify and Explain Exponential Growth} Students will be able to identify exponential growth in various contexts and explain its characteristics, including how the quantity increases over time with a rate proportional to its current value.
    \item \textbf{Apply the Exponential Growth Formula} Students will be able to apply the exponential growth formula to solve problems related to population growth, finance, and other areas where quantities increase exponentially.
    \item \textbf{Differentiate Between Exponential Growth and Decay} Students will be able to differentiate between exponential growth and decay, and apply the appropriate formula to model each type of growth.
\end{enumerate}

\subsection{Assessment Instrument}

\begin{enumerate}
    \item \textbf{Problem 1: Exponential Growth Identification} Provide a scenario where a quantity exhibits exponential growth. For instance, a population of bacteria doubles every hour. Identify and explain the exponential growth pattern in this scenario. \textbf{Problem 2: Applying the Exponential Growth Formula} A virus infects a population such that the number of infected individuals triples every 5 days. If there are initially 100 infected individuals, calculate the number of infected individuals after 15 days using the formula $x\_t = x\_0 (1 + r)^t$. \textbf{Problem 3: Exponential Decay Calculation} A certain radioactive substance has a decay rate of 10\% per day. If the initial amount of the substance is 500 grams, calculate the remaining amount of the substance after 10 days using the formula $x(t) = a \cdot b^{t/\tau}$. Assume the decay constant $b = 0.9$ and $\tau = 1$ day. \textbf{Problem 4: Growth and Decay Comparison} Compare exponential growth and decay by providing an example of each. Use the formula $x(t) = x\_0 e^{kt}$ for both scenarios and describe how the parameters $k$ influence growth or decay. \textbf{Problem 5: Log-Linear Model Transformation} Given a dataset showing exponential growth, transform it to a log-linear model by taking the logarithm of the data and plotting it against time. Explain how the log-linear model helps in analyzing exponential growth. \textbf{Problem 6: Differential Equation Verification} Verify that the function $x(t) = x(0) e^{kt}$ satisfies the differential equation $\frac{dx}{dt} = kx$. Show your work through integration and differentiation.
\end{enumerate}

%%===============================================================
\section{Carrying Capacity and Logistic Growth Rate}
\label{sec:carrying-capacity-and-logistic-growth-rate}
%%===============================================================

\paragraph{Carrying Capacity and Logistic Growth Rate}



% [Image: log]



\textbf{Standard logistic sigmoid function where $L = 1$, $k = 1$, $x\_0 = 0$}  





A logistic function or logistic curve is a common S-shaped curve (sigmoid curve) with equation



$$

f(x) = \frac{L}{1 + e^{-k(x - x\_0)}}

$$



where:



\begin{itemize}
    \item $x\_0$, the $x$ value of the sigmoid's midpoint;
    \item $L$, the curve's maximum value;
    \item $k$, the logistic growth rate or steepness of the curve.
\end{itemize}

For values of $x$ in the domain of real numbers from $-\infty$ to $+\infty$, the S-curve shown on the right is obtained, with the graph of $f$ approaching $L$ as $x$ approaches $+\infty$ and approaching zero as $x$ approaches $-\infty$.





\paragraph{Modeling population growth}



% [Image: Pierre]



\textbf{Pierre-François Verhulst (1804–1849)}

 

A typical application of the logistic equation is a common model of population growth, originally due to Pierre-François Verhulst in 1838, where the rate of reproduction is proportional to both the existing population and the amount of available resources, all else being equal. The Verhulst equation was published after Verhulst had read Thomas Malthus' \textit{An Essay on the Principle of Population}, which describes the Malthusian growth model of simple (unconstrained) exponential growth. Verhulst derived his logistic equation to describe the self-limiting growth of a biological population. The equation was rediscovered in 1911 by A. G. McKendrick for the growth of bacteria in broth and experimentally tested using a technique for nonlinear parameter estimation. The equation is also sometimes called the Verhulst-Pearl equation following its rediscovery in 1920 by Raymond Pearl (1879–1940) and Lowell Reed (1888–1966) of the Johns Hopkins University. Another scientist, Alfred J. Lotka derived the equation again in 1925, calling it the law of population growth.



Letting $P$ represent population size ($N$ is often used in ecology instead) and $t$ represent time, this model is formalized by the differential equation:



$$

\frac{dP}{dt} = rP \left(1 - \frac{P}{K}\right)

$$



where the constant $r$ defines the growth rate and $K$ is the carrying capacity.



In the equation, the early, unimpeded growth rate is modeled by the first term $+rP$. The value of the rate $r$ represents the proportional increase of the population $P$ in one unit of time. Later, as the population grows, the modulus of the second term (which multiplied out is $-rP^2 / K$) becomes almost as large as the first, as some members of the population $P$ interfere with each other by competing for some critical resource, such as food or living space. This antagonistic effect is called the bottleneck, and is modeled by the value of the parameter $K$. The competition diminishes the combined growth rate, until the value of $P$ ceases to grow (this is called maturity of the population). The solution to the equation (with $P\_0$ being the initial population) is



$$

P(t) = \frac{KP\_0 e^{rt}}{K + P\_0 \left(e^{rt} - 1\right)} = \frac{K}{1 + \left(\frac{K - P\_0}{P\_0}\right) e^{-rt}}

$$



where



$$

\lim\_{t \to \infty} P(t) = K

$$



Which is to say that $K$ is the limiting value of $P$: the highest value that the population can reach given infinite time (or come close to reaching in finite time). It is important to stress that the carrying capacity is asymptotically reached independently of the initial value $P(0) > 0$, and also in the case that $P(0) > K$.



In ecology, species are sometimes referred to as $r$-strategist or $K$-strategist depending upon the selective processes that have shaped their life history strategies. Choosing the variable dimensions so that $n$ measures the population in units of carrying capacity, and $\tau$ measures time in units of $1 / r$, gives the dimensionless differential equation



$$

\frac{dn}{d\tau} = n(1 - n)

$$



\paragraph{Time-varying carrying capacity}



Since the environmental conditions influence the carrying capacity, as a consequence it can be time-varying, with $K(t) > 0$, leading to the following mathematical model:



$$

\frac{dP}{dt} = rP \cdot \left(1 - \frac{P}{K(t)}\right)

$$



A particularly important case is that of carrying capacity that varies periodically with period $T$:



$$

K(t + T) = K(t)

$$



It can be shown that in such a case, independently from the initial value $P(0) > 0$, $P(t)$ will tend to a unique periodic solution $P\_*(t)$, whose period is $T$.



A typical value of $T$ is one year: In such case $K(t)$ may reflect periodical variations of weather conditions.



Another interesting generalization is to consider that the carrying capacity $K(t)$ is a function of the population at an earlier time, capturing a delay in the way population modifies its environment. This leads to a logistic delay equation, which has a very rich behavior, with bistability in some parameter range, as well as a monotonic decay to zero, smooth exponential growth, punctuated unlimited growth (i.e., multiple S-shapes), punctuated growth or alternation to a stationary level, oscillatory approach to a stationary level, sustainable oscillations, finite-time singularities as well as finite-time death.



\paragraph{Resources}

\begin{itemize}
    \item \href{https://math.libretexts.org/Bookshelves/Calculus/Book\%3A_Calculus_(OpenStax}{The Logistic Equation}/08\%3A_Introduction_to_Differential_Equations/8.4\%3A_The_Logistic_Equation), Mathematics LibreTexts
    \item \href{https://courses.lumenlearning.com/boundless-biology/chapter/environmental-limits-to-population-growth/}{Environmental Limits to Population Growth}, Lumen Learning
\end{itemize}

\paragraph{Licensing}

Content obtained and/or adapted from:



\begin{itemize}
    \item \href{https://en.wikipedia.org/wiki/Logistic_function}{Logistic function, Wikipedia} under a CC BY-SA license
\end{itemize}

\subsection{Learning Objectives}

\begin{enumerate}
    \item \textbf{Identify and Explain the Logistic Function} Students will be able to identify the logistic function $f(x) = \frac{L}{1 + e^{-k(x - x\_0)}}$, explain its components ($x\_0$, $L$, and $k$), and describe how the function models growth with a carrying capacity.
    \item \textbf{Apply the Logistic Growth Model} Students will be able to apply the logistic growth differential equation $\frac{dP}{dt} = rP \left(1 - \frac{P}{K}\right)$ to model population growth, calculate the population at time $t$, and interpret the effects of carrying capacity $K$ and growth rate $r$.
    \item \textbf{Analyze Time-Varying Carrying Capacity} Students will be able to analyze and solve the logistic growth equation with a time-varying carrying capacity $K(t)$, including cases where $K(t)$ varies periodically with period $T$, and discuss the implications for population dynamics.
\end{enumerate}

\subsection{Assessment Instrument}

\begin{enumerate}
    \item Problem 1: Logistic Function Identification
    \item \textbf{Given the logistic function} $f(x) = \frac{1}{1 + e^{-2(x - 3)}}$, identify the values of $L$, $k$, and $x\_0$.
    \item \textbf{Explain how these parameters affect the shape and position of the logistic curve}. Problem 2: Application of Logistic Growth Model
    \item \textbf{For a population with an initial size of $P\_0 = 50$, growth rate $r = 0.1$, and carrying capacity $K = 1000$, calculate the population size $P(t)$ at $t = 10$} using the equation $$ P(t) = \frac{KP\_0 e^{rt}}{K + P\_0 \left(e^{rt} - 1\right)} $$
    \item \textbf{Discuss the behavior of the population as $t \to \infty$}. Problem 3: Time-Varying Carrying Capacity
    \item \textbf{Solve the logistic growth equation} with a time-varying carrying capacity $K(t)$ where $K(t) = 1000 + 200 \sin\left(\frac{2\pi t}{12}\right)$, and the growth rate $r = 0.05$.
    \item \textbf{Describe the effect of the periodic variation in $K(t)$ on the population dynamics} and determine if $P(t)$ reaches a unique periodic solution.
\end{enumerate}

%%===============================================================
\section{Logistic Growth Model}
\label{sec:logistic-growth-model}
%%===============================================================

\paragraph{Logistic Growth Model}



A logistic function or logistic curve is a common S-shaped curve (sigmoid curve) with the equation



$$

f(x) = \frac{L}{1 + e^{-k(x - x\_0)}}

$$



where:



\begin{itemize}
    \item $x\_0$, the $x$ value of the sigmoid''s midpoint;
    \item $L$, the curve''s maximum value;
    \item $k$, the logistic growth rate or steepness of the curve.
\end{itemize}



% [Image: curve]



For values of $x$ in the domain of real numbers from $-\infty$ to $+\infty$, the S-curve shown above is obtained, with the graph of $f$ approaching $L$ as $x$ approaches $+\infty$ and approaching zero as $x$ approaches $-\infty$.



\paragraph{Mathematical Properties}



The standard logistic function is the logistic function with parameters $k = 1$, $x\_0 = 0$, $L = 1$, which yields



$$

f(x) = \frac{1}{1 + e^{-x}} = \frac{e^x}{e^x + 1} = \frac{1}{2} + \frac{1}{2} \tanh \left( \frac{x}{2} \right)

$$



In practice, due to the nature of the exponential function $e^{-x}$, it is often sufficient to compute the standard logistic function for $x$ over a small range of real numbers, such as a range contained in $[-6, +6]$, as it quickly converges very close to its saturation values of 0 and 1.



The logistic function has the symmetry property that



$$

1 - f(x) = f(-x)

$$



Thus, $x \mapsto f(x) - \frac{1}{2}$ is an odd function.



The logistic function is an offset and scaled hyperbolic tangent function:



$$

f(x) = \frac{1}{2} + \frac{1}{2} \tanh \left( \frac{x}{2} \right)

$$



or



$$

\tanh(x) = 2f(2x) - 1

$$



This follows from



$$

\tanh(x) = \frac{e^x - e^{-x}}{e^x + e^{-x}} = \frac{e^x \cdot (1 - e^{-2x})}{e^x \cdot (1 + e^{-2x})} = f(2x) - \frac{e^{-2x}}{1 + e^{-2x}} = f(2x) - \frac{e^{-2x} + 1 - 1}{1 + e^{-2x}} = 2f(2x) - 1

$$



\paragraph{Derivative}



The standard logistic function has an easily calculated derivative. The derivative is known as the logistic distribution:



$$

f(x) = \frac{1}{1 + e^{-x}} = \frac{e^x}{1 + e^x}

$$



$$

\frac{d}{dx} f(x) = \frac{e^x \cdot (1 + e^x) - e^x \cdot e^x}{(1 + e^x)^2} = \frac{e^x}{(1 + e^x)^2} = f(x) \left(1 - f(x)\right)

$$



\paragraph{Integral}



Conversely, its antiderivative can be computed by the substitution $u = 1 + e^x$, since



$$

f(x) = \frac{e^x}{1 + e^x} = \frac{u''}{u}

$$



so (dropping the constant of integration)



$$

\int \frac{e^x}{1 + e^x} \, dx = \int \frac{1}{u} \, du = \ln u = \ln(1 + e^x)

$$



In artificial neural networks, this is known as the softplus function and (with scaling) is a smooth approximation of the ramp function, just as the logistic function (with scaling) is a smooth approximation of the Heaviside step function.



\paragraph{Logistic Differential Equation}



The standard logistic function is the solution of the simple first-order non-linear ordinary differential equation



$$

\frac{d}{dx} f(x) = f(x) \left(1 - f(x)\right)

$$



with boundary condition $f(0) = \frac{1}{2}$. This equation is the continuous version of the logistic map. Note that the reciprocal logistic function is a solution to a simple first-order linear ordinary differential equation.



The qualitative behavior is easily understood in terms of the phase line: the derivative is 0 when the function is 1; and the derivative is positive for $f$ between 0 and 1, and negative for $f$ above 1 or less than 0 (though negative populations do not generally accord with a physical model). This yields an unstable equilibrium at 0 and a stable equilibrium at 1, and thus for any function value greater than 0 and less than 1, it grows to 1.



The logistic equation is a special case of the Bernoulli differential equation and has the following solution:



$$

f(x) = \frac{e^x}{e^x + C}

$$



Choosing the constant of integration $C = 1$ gives the other well-known form of the definition of the logistic curve:



$$

f(x) = \frac{e^x}{e^x + 1} = \frac{1}{1 + e^{-x}}

$$



More quantitatively, as can be seen from the analytical solution, the logistic curve shows early exponential growth for negative argument, which reaches linear growth of slope $\frac{1}{4}$ for an argument near 0, then approaches 1 with an exponentially decaying gap.



The logistic function is the inverse of the natural logit function and so can be used to convert the logarithm of odds into a probability. In mathematical notation, the logistic function is sometimes written as $\text{expit}$, in the same form as $\text{logit}$. The conversion from the log-likelihood ratio of two alternatives also takes the form of a logistic curve.



The differential equation derived above is a special case of a general differential equation that only models the sigmoid function for $x > 0$. In many modeling applications, the more general form



$$

\frac{df(x)}{dx} = \frac{k}{a}f(x)(a - f(x)), \quad f(0) = \frac{a}{1 + e^{kr}}

$$



can be desirable. Its solution is the shifted and scaled sigmoid $a S(k(x - r))$.



The hyperbolic-tangent relationship leads to another form for the logistic function''s derivative:



$$

\frac{d}{dx} f(x) = \frac{1}{4} \, \text{sech}^2 \left( \frac{x}{2} \right)

$$



which ties the logistic function into the logistic distribution.



\paragraph{Rotational Symmetry About $(0, \frac{1}{2})$}



The sum of the logistic function and its reflection about the vertical axis, $f(-x)$, is



$$

\frac{1}{1 + e^{-x}} + \frac{1}{1 + e^{-(-x)}} = \frac{e^x}{e^x + 1} + \frac{1}{e^x + 1} = 1

$$



The logistic function is thus rotationally symmetrical about the point $(0, \frac{1}{2})$.



\paragraph{Applications}



Link created an extension of Wald''s theory of sequential analysis to a distribution-free accumulation of random variables until either a positive or negative bound is first equaled or exceeded. Link derives the probability of first equaling or exceeding the positive boundary as $(1+e^{-\theta A})$, the logistic function. This is the first proof that the logistic function may have a stochastic process as its basis. Link provides a century of examples of "logistic" experimental results and a newly derived relation between this probability and the time of absorption at the boundaries.



\paragraph{In ecology: modeling population growth}



\textbf{Pierre-François Verhulst (1804–1849)} 



% [Image: pierre]



A typical application of the logistic equation is a common model of population growth (see also population dynamics), originally due to Pierre-François Verhulst in 1838, where the rate of reproduction is proportional to both the existing population and the amount of available resources, all else being equal. The Verhulst equation was published after Verhulst had read Thomas Malthus'' \textit{An Essay on the Principle of Population}, which describes the Malthusian growth model of simple (unconstrained) exponential growth. Verhulst derived his logistic equation to describe the self-limiting growth of a biological population. The equation was rediscovered in 1911 by A. G. McKendrick for the growth of bacteria in broth and experimentally tested using a technique for nonlinear parameter estimation. The equation is also sometimes called the Verhulst-Pearl equation following its rediscovery in 1920 by Raymond Pearl (1879–1940) and Lowell Reed (1888–1966) of the Johns Hopkins University. Another scientist, Alfred J. Lotka derived the equation again in 1925, calling it the law of population growth.



Letting $P$ represent population size ($N$ is often used in ecology instead) and $t$ represent time, this model is formalized by the differential equation:



$$

\frac{dP}{dt} = rP \left(1-\frac{P}{K}\right),

$$



where the constant $r$ defines the growth rate and $K$ is the carrying capacity.



In the equation, the early, unimpeded growth rate is modeled by the first term $+rP$. The value of the rate $r$ represents the proportional increase of the population $P$ in one unit of time. Later, as the population grows, the modulus of the second term (which multiplied out is $-rP^2/K$) becomes almost as large as the first, as some members of the population $P$ interfere with each other by competing for some critical resource, such as food or living space. This antagonistic effect is called the bottleneck, and is modeled by the value of the parameter $K$. The competition diminishes the combined growth rate, until the value of $P$ ceases to grow (this is called maturity of the population). The solution to the equation (with $P\_0$ being the initial population) is



$$

P(t) = \frac{KP\_0 e^{rt}}{K + P\_0 \left(e^{rt} - 1\right)} = \frac{K}{1 + \left(\frac{K - P\_0}{P\_0}\right) e^{-rt}},

$$



where



$$

\lim\_{t \to \infty} P(t) = K.

$$



Which is to say that $K$ is the limiting value of $P$: the highest value that the population can reach given infinite time (or come close to reaching in finite time). It is important to stress that the carrying capacity is asymptotically reached independently of the initial value $P(0) > 0$, and also in the case that $P(0) > K$.



In ecology, species are sometimes referred to as $r$-strategist or $K$-strategist depending upon the selective processes that have shaped their life history strategies. Choosing the variable dimensions so that $n$ measures the population in units of carrying capacity, and $\tau$ measures time in units of $1/r$, gives the dimensionless differential equation



$$

\frac{dn}{d\tau} = n(1 - n).

$$



\paragraph{Time-varying carrying capacity}



Since the environmental conditions influence the carrying capacity, as a consequence it can be time-varying, with $K(t) > 0$, leading to the following mathematical model:



$$

\frac{dP}{dt} = rP \cdot \left(1-\frac{P}{K(t)}\right).

$$



A particularly important case is that of carrying capacity that varies periodically with period $T$:



$$

K(t + T) = K(t).

$$



It can be shown that in such a case, independently from the initial value $P(0) > 0$, $P(t)$ will tend to a unique periodic solution $P\_*(t)$, whose period is $T$.



A typical value of $T$ is one year: In such case $K(t)$ may reflect periodical variations of weather conditions.



Another interesting generalization is to consider that the carrying capacity $K(t)$ is a function of the population at an earlier time, capturing a delay in the way population modifies its environment. This leads to a logistic delay equation, which has a very rich behavior, with bistability in some parameter range, as well as a monotonic decay to zero, smooth exponential growth, punctuated unlimited growth (i.e., multiple S-shapes), punctuated growth or alternation to a stationary level, oscillatory approach to a stationary level, sustainable oscillations, finite-time singularities as well as finite-time death.



\paragraph{In statistics and machine learning}



Logistic functions are used in several roles in statistics. For example, they are the cumulative distribution function of the logistic family of distributions, and they are, a bit simplified, used to model the chance a chess player has to beat his opponent in the Elo rating system. More specific examples now follow.



\paragraph{Logistic regression}



Logistic functions are used in logistic regression to model how the probability $p$ of an event may be affected by one or more explanatory variables: an example would be to have the model



$$

p = f(a + bx),

$$



where $x$ is the explanatory variable, $a$ and $b$ are model parameters to be fitted, and $f$ is the standard logistic function.



Logistic regression and other log-linear models are also commonly used in machine learning. A generalisation of the logistic function to multiple inputs is the softmax activation function, used in multinomial logistic regression.



Another application of the logistic function is in the Rasch model, used in item response theory. In particular, the Rasch model forms a basis for maximum likelihood estimation of the locations of objects or persons on a continuum, based on collections of categorical data, for example the abilities of persons on a continuum based on responses that have been categorized as correct and incorrect.



\paragraph{Neural networks}



Logistic functions are often used in neural networks to introduce nonlinearity in the model or to clamp signals to within a specified interval. A popular neural net element computes a linear combination of its input signals, and applies a bounded logistic function as the activation function to the result; this model can be seen as a "smoothed" variant of the classical threshold neuron.



A common choice for the activation or "squashing" functions, used to clip for large magnitudes to keep the response of the neural network bounded is



$$

g(h) = \frac{1}{1 + e^{-2\beta h}},

$$



which is a logistic function.



These relationships result in simplified implementations of artificial neural networks with artificial neurons. Practitioners caution that sigmoidal functions which are antisymmetric about the origin (e.g., the hyperbolic tangent) lead to faster convergence when training networks with backpropagation.



The logistic function is itself the derivative of another proposed activation function, the softplus.



\paragraph{In medicine: modeling of growth of tumors}



Another application of logistic curve is in medicine, where the logistic differential equation is used to model the growth of tumors. This application can be considered an extension of the above-mentioned use in the framework of ecology (see also the Generalized logistic curve, allowing for more parameters). Denoting with $X(t)$ the size of the tumor at time $t$, its dynamics are governed by



$$

X'' = r \left(1 - \frac{X}{K}\right)X,

$$



which is of the type



$$

X'' = F(X)X, \quad F''(X) \leq 0,

$$



where $F(X)$ is the proliferation rate of the tumor.



If a chemotherapy is started with a log-kill effect, the equation may be revised to be



$$

X'' = r \left(1 - \frac{X}{K}\right)X - c(t)X,

$$



where $c(t)$ is the therapy-induced death rate. In the idealized case of very long therapy, $c(t)$ can be modeled as a periodic function (of period $T$) or (in case of continuous infusion therapy) as a constant function, and one has that



$$

\frac{1}{T} \int\_0^T c(t) \, dt > r \rightarrow \lim\_{t \to +\infty} x(t) = 0,

$$



i.e., if the average therapy-induced death rate is greater than the baseline proliferation rate, then there is the eradication of the disease. Of course, this is an oversimplified model of both the growth and the therapy (e.g., it does not take into account the phenomenon of clonal resistance).



\paragraph{In medicine: modeling of a pandemic}



A novel infectious pathogen to which a population has no immunity will generally spread exponentially in the early stages, while the supply of susceptible individuals is plentiful. The SARS-CoV-2 virus that causes COVID-19 exhibited exponential growth early in the course of infection in several countries in early 2020. Many factors, ranging from lack of susceptibles (either through the continued spread of infection until it passes the threshold for herd immunity or reduction in the accessibility of susceptibles through physical distancing measures), exponential-looking epidemic curves may first linearize (replicating the "logarithmic" to "logistic" transition first noted by Pierre-François Verhulst, as noted above) and then reach a maximal limit.



A logistic function, or related functions (e.g., the Gompertz function) are usually used in a descriptive or phenomenological manner because they fit well not only to the early exponential rise, but to the eventual leveling off of the pandemic as the population develops a herd immunity. This is in contrast to actual models of pandemics which attempt to formulate a description based on the dynamics of the pandemic (e.g., contact rates, incubation times, social distancing, etc.). Some simple models have been developed, however, which yield a logistic solution.



\paragraph{Modeling COVID-19 infection trajectory}



\textbf{Generalized logistic function (Richards growth curve) in epidemiological modeling}



% [Image: generalized]





A generalized logistic function, also called the Richards growth curve, is widely used in modeling COVID-19 infection trajectories. Infection trajectory is a daily time series data for the cumulative number of infected cases for a subject such as a country, city, state, etc. There are variant re-parameterizations in the literature; one frequently used form is



$$

f(t; \theta\_1, \theta\_2, \theta\_3, \xi) = \frac{\theta\_1}{\left[1 + \xi \exp(-\theta\_2 \cdot (t - \theta\_3))\right]^{1/\xi}}

$$



where $ \theta\_1, \theta\_2, \theta\_3 $ are real numbers, and $ \xi $ is a positive real number. The flexibility of the curve $ f $ is due to the parameter $ \xi $: (i) if $ \xi = 1 $ then the curve reduces to the logistic function, and (ii) if $ \xi $ converges to zero, then the curve converges to the Gompertz function. In epidemiological modeling, $ \theta\_1 $, $ \theta\_2 $, and $ \theta\_3 $ represent the final epidemic size, infection rate, and lag phase, respectively.



\textbf{Extrapolated infection trajectories of 40 countries severely affected by COVID-19 and grand (population) average through May 14th}



% [Image: extrapolated]



One of the benefits of using a growth function such as the generalized logistic function in epidemiological modeling is its relatively easy expansion to the multilevel model framework by using the growth function to describe infection trajectories from multiple subjects (countries, cities, states, etc.). Such a modeling framework can also be widely called the nonlinear mixed effect model or hierarchical nonlinear model. An example of using the generalized logistic function in a Bayesian multilevel model is the Bayesian hierarchical Richards model.



\paragraph{In chemistry: reaction models}



The concentration of reactants and products in autocatalytic reactions follows the logistic function. The degradation of Platinum group metal-free (PGM-free) oxygen reduction reaction (ORR) catalysts in fuel cell cathodes follows the logistic decay function, suggesting an autocatalytic degradation mechanism.



\paragraph{In physics: Fermi–Dirac distribution}



The logistic function determines the statistical distribution of fermions over the energy states of a system in thermal equilibrium. In particular, it is the distribution of the probabilities that each possible energy level is occupied by a fermion, according to Fermi–Dirac statistics.



\paragraph{In linguistics: language change}



In linguistics, the logistic function can be used to model language change: an innovation that is at first marginal begins to spread more quickly with time, and then more slowly as it becomes more universally adopted.



\paragraph{In agriculture: modeling crop response}



The logistic S-curve can be used for modeling the crop response to changes in growth factors. There are two types of response functions: positive and negative growth curves. For example, crop yield may increase with increasing value of the growth factor up to a certain level (positive function), or it may decrease with increasing growth factor values (negative function owing to a negative growth factor), which situation requires an inverted S-curve.



$\phantom{--------}$\textbf{S-curve model for crop yield vs. depth of water table.} 



% [Image: depth]  



$\phantom{--------}$ \textbf{Inverted S-curve model for crop yield vs. soil salinity.}



% [Image: soil]



\paragraph{In economics and sociology: diffusion of innovations}



The logistic function can be used to illustrate the progress of the diffusion of an innovation through its life cycle.



In \textit{The Laws of Imitation} (1890), Gabriel Tarde describes the rise and spread of new ideas through imitative chains. In particular, Tarde identifies three main stages through which innovations spread: the first one corresponds to the difficult beginnings, during which the idea has to struggle within a hostile environment full of opposing habits and beliefs; the second one corresponds to the properly exponential take-off of the idea, with $f(x) = 2^x$; finally, the third stage is logarithmic, with $f(x) = \log(x)$, and corresponds to the time when the impulse of the idea gradually slows down while, simultaneously new opponent ideas appear. The ensuing situation halts or stabilizes the progress of the innovation, which approaches an asymptote.



In a Sovereign state, the subnational units (Constituent states or cities) may use loans to finance their projects. However, this funding source is usually subject to strict legal rules as well as to economy scarcity constraints, specially the resources the banks can lend (due to their equity or Basel limits). These restrictions, which represent a saturation level, along with an exponential rush in an economic competition for money, create a public finance diffusion of credit pleas and the aggregate national response is a sigmoid curve.



In the history of economy, when new products are introduced there is an intense amount of research and development which leads to dramatic improvements in quality and reductions in cost. This leads to a period of rapid industry growth. Some of the more famous examples are: railroads, incandescent light bulbs, electrification, cars and air travel. Eventually, dramatic improvement and cost reduction opportunities are exhausted, the product or process are in widespread use with few remaining potential new customers, and markets become saturated.



Logistic analysis was used in papers by several researchers at the International Institute of Applied Systems Analysis (IIASA). These papers deal with the diffusion of various innovations, infrastructures and energy source substitutions and the role of work in the economy as well as with the long economic cycle. Long economic cycles were investigated by Robert Ayres (1989). Cesare Marchetti published on long economic cycles and on diffusion of innovations. Arnulf Grübler''s book (1990) gives a detailed account of the diffusion of infrastructures including canals, railroads, highways and airlines, showing that their diffusion followed logistic shaped curves.



Carlota Perez used a logistic curve to illustrate the long (Kondratiev) business cycle with the following labels: beginning of a technological era as irruption, the ascent as frenzy, the rapid build out as synergy and the completion as maturity.



\paragraph{Resources}



\begin{itemize}
    \item \href{https://www.khanacademy.org/science/ap-biology/ecology-ap/population-ecology-ap/a/exponential-logistic-growth}{Exponential and Logistic Growth}, Khan Academy
\end{itemize}

\paragraph{Licensing}



Content obtained and/or adapted from:



\begin{itemize}
    \item \href{https://en.wikipedia.org/wiki/Logistic_function}{Logistic function, Wikipedia} under a CC BY-SA license
\end{itemize}

\subsection{Learning Objectives}

\begin{enumerate}
    \item \textbf{1. Understand and Analyze the Logistic Growth Model}
    \item Identify the key components of the logistic growth model, including the parameters $L$, $k$, and $x\_0$.
    \item Explain the significance of the S-shaped curve and how it models growth in various applications. \textbf{2. Derive and Interpret the Derivatives of the Logistic Function}
    \item Derive the first and second derivatives of the logistic function.
    \item Interpret the meaning of these derivatives in the context of population growth and other applications. \textbf{3. Apply the Logistic Growth Model to Real-World Scenarios}
    \item Use the logistic growth model to analyze and predict population growth, tumor growth, or the spread of diseases.
    \item Apply the model to determine the carrying capacity and growth rate in given scenarios.
\end{enumerate}

\subsection{Assessment Instrument}

\begin{enumerate}
    \item Problem 1: Identifying Components of the Logistic Growth Model Given the logistic growth function: $$ f(x) = \frac{L}{1 + e^{-k(x - x\_0)}} $$ a. Identify and describe the roles of the parameters $L$, $k$, and $x\_0$. b. Sketch the curve of this function and label the key points, including the carrying capacity and the inflection point. Problem 2: Derivative of the Logistic Function Given the standard logistic function: $$ f(x) = \frac{1}{1 + e^{-x}} $$ a. Derive the first derivative of the function and interpret its meaning in the context of population growth. b. Calculate the derivative at $x = 0$ and describe its significance. Problem 3: Application to Real-World Scenario Consider a population that follows the logistic growth model: $$ \frac{dP}{dt} = rP \left(1-\frac{P}{K}\right) $$ a. If the initial population is $P\_0 = 100$, the carrying capacity $K = 1000$, and the growth rate $r = 0.2$, find the population $P(t)$ at time $t = 10$. b. Explain how changes in the carrying capacity $K$ and growth rate $r$ would affect the population growth over time. Problem 4: Logistic Growth in Epidemiology Using the generalized logistic function: $$ f(t; \theta\_1, \theta\_2, \theta\_3, \xi) = \frac{\theta\_1}{\left[1 + \xi \exp(-\theta\_2 \cdot (t - \theta\_3))\right]^{1/\xi}} $$ a. Describe how this model can be used to predict the spread of a disease like COVID-19. b. Discuss the significance of the parameter $\xi$ and how it alters the shape of the growth curve.
\end{enumerate}

%%===============================================================
\section{Complex Population Growth and Decay Models}
\label{sec:complex-population-growth-and-decay-models}
%%===============================================================

\paragraph{Complex Population Growth and Decay Models}





\paragraph{Propotionality}



\textbf{The variable $y$ is directly proportional to the variable $x$ with proportionality constant $0.6$.}



% [Image: direct]



\textbf{The variable $y$ is inversely proportional to the variable $x$ with proportionality constant $1$.}



% [Image: inverse]



In mathematics, two varying quantities are said to be in a relation of proportionality, multiplicatively connected to a constant; that is, when either their ratio or their product yields a constant. The value of this constant is called the coefficient of proportionality or proportionality constant.



If the ratio $\frac{y}{x}$ of two variables ($x$ and $y$) is equal to a constant ($k = \frac{y}{x}$), then the variable in the numerator of the ratio ($y$) can be product of the other variable and the constant ($y = k \cdot x$). In this case $y$ is said to be directly proportional to $x$ with proportionality constant $k$. Equivalently one may write $\frac{1}{k} \cdot y$; that is, $x$ is directly proportional to $y$ with proportionality constant $\frac{1}{k} = (\frac{x}{y})$. If the term proportional is connected to two variables without further qualification, generally direct proportionality can be assumed.

If the product of two variables ($x \cdot y$) is equal to a constant ($k = x \cdot y$), then the two are said to be inversely proportional to each other with the proportionality constant $k$. Equivalently, both variables are directly proportional to the reciprocal of the respective other with proportionality constant $k$ ($x = k \cdot \frac{1}{y}$ and $y = k \cdot \frac{1}{x}$).

If several pairs of variables share the same direct proportionality constant, the equation expressing the equality of these ratios is called a proportion, e.g., $\frac{a}{b} = \frac{x}{y}$ = $\cdots$ = $k$ (for details see Ratio).



\paragraph{Direct proportionality}



Given two variables $x$ and $y$, $y$ is directly proportional to $x$ if there is a non-zero constant $k$ such that



$y = kx.$



The relation is often denoted using the symbols "$\propto$" (not to be confused with the Greek letter alpha) or "$\sim$":



$y \propto x,$ or $y \sim x.$



For $x \ne 0$ the proportionality constant can be expressed as the ratio



$k = \frac{y}{x}.$



It is also called the constant of variation or constant of proportionality.



A direct proportionality can also be viewed as a linear equation in two variables with a $y$-intercept of $0$ and a slope of $k$. This corresponds to linear growth.



\paragraph{Examples}



\begin{itemize}
    \item If an object travels at a constant speed, then the distance traveled is directly proportional to the time spent traveling, with the speed being the constant of proportionality.
    \item The circumference of a circle is directly proportional to its diameter, with the constant of proportionality equal to $\pi$.
    \item On a map of a sufficiently small geographical area, drawn to scale distances, the distance between any two points on the map is directly proportional to the beeline distance between the two locations represented by those points; the constant of proportionality is the scale of the map.
    \item The force, acting on a small object with small mass by a nearby large extended mass due to gravity, is directly proportional to the object's mass; the constant of proportionality between the force and the mass is known as gravitational acceleration.
    \item The net force acting on an object is proportional to the acceleration of that object with respect to an inertial frame of reference. The constant of proportionality in this, Newton's second law, is the classical mass of the object.
\end{itemize}

\paragraph{Solve problems involving joint variation}



Many situations are more complicated than a basic direct variation or inverse variation model. One variable often depends on multiple other variables. When a variable is dependent on the product or quotient of two or more variables, this is called joint variation. For example, the cost of busing students for each school trip varies with the number of students attending and the distance from the school. The variable $c$, cost, varies jointly with the number of students, $n$, and the distance, $d$.



\paragraph{A GENERAL NOTE: JOINT VARIATION}



Joint variation occurs when a variable varies directly or inversely with multiple variables.

For instance, if $x$ varies directly with both $y$ and $z$, we have $x = kyz$. If $x$ varies directly with $y$ and inversely with $z$, we have $x = \frac{ky}{z}$.

Notice that we only use one constant in a joint variation equation.



\paragraph{Example}



A quantity $x$ varies directly with the square of $y$ and inversely with the cube root of $z$. If $x = 6$ when $y = 2$ and $z = 8$, find $x$ when $y = 1$ and $z = 27$.



\textbf{Solution}



Begin by writing an equation to show the relationship between the variables.



$x = \frac{k y^2}{\sqrt[3]{z}}$



Substitute $x = 6$, $y = 2$, and $z = 8$ to find the value of the constant $k$.



$6 = \frac{k 2^2}{\sqrt[3]{8}}$



$6 = \frac{4k}{2}$



$3 = k$



Now we can substitute the value of the constant into the equation for the relationship.



$x = \frac{3 y^2}{\sqrt[3]{z}}$



To find $x$ when $y = 1$ and $z = 27$, we will substitute values for $y$ and $z$ into our equation.



$x = \frac{3 (1)^2}{\sqrt[3]{27}}$



$ = 1$



\paragraph{Licensing}



Content obtained and/or adapted from:



\begin{itemize}
    \item \href{https://en.wikipedia.org/wiki/Proportionality_(mathematics}{Proportionality (mathematics)}), Wikipedia under a CC BY-SA license
    \item \href{https://courses.lumenlearning.com/ivytech-collegealgebra/chapter/solve-problems-involving-joint-variation/}{Joint variation}, Lumen Learning under a CC BY license
\end{itemize}

\subsection{Learning Objectives}

\begin{enumerate}
    \item \textbf{Identify and apply direct proportionality:} Given a relationship between two variables, $x$ and $y$, determine if $y$ is directly proportional to $x$ and calculate the proportionality constant $k$ where $y = kx$.
    \item \textbf{Identify and apply inverse proportionality:} Given a relationship where the product of two variables, $x$ and $y$, is constant, determine if $y$ is inversely proportional to $x$ and calculate the proportionality constant $k$ where $x \cdot y = k$.
    \item \textbf{Solve joint variation problems:} Given a relationship where a variable $x$ varies jointly with two variables $y$ and $z$, solve for $x$ using the formula $x = \frac{k y^2}{\sqrt[3]{z}}$ and find the value of the constant $k$ from provided values.
\end{enumerate}

\subsection{Assessment Instrument}

\begin{enumerate}
    \item \textbf{1. Direct Proportionality} \textit{Question:} Given the equation $y = 0.8x$, identify the proportionality constant and verify if $y$ is directly proportional to $x$. \textbf{2. Inverse Proportionality} \textit{Question:} If $x \cdot y = 12$ and $x = 3$, find $y$ and determine the proportionality constant $k$. \textbf{3. Joint Variation} \textit{Question:} A quantity $x$ varies directly with the square of $y$ and inversely with the cube root of $z$. If $x = 10$ when $y = 3$ and $z = 27$, find $x$ when $y = 2$ and $z = 8$.
\end{enumerate}

%%===============================================================
\section{Analyzing Complex Population Growth and Decay Models}
\label{sec:analyzing-complex-population-growth-and-decay-models}
%%===============================================================

Analyzing Complex Population Growth and Decay Models

%%===============================================================
\section{Periodic Function}
\label{sec:periodic-function}
%%===============================================================

\paragraph{Periodic Function}



A periodic function is a function that repeats its values at regular intervals, for example, the trigonometric functions, which repeat at intervals of $2\pi$ radians. Periodic functions are used throughout science to describe oscillations, waves, and other phenomena that exhibit periodicity. Any function that is not periodic is called aperiodic.



\textbf{An illustration of a periodic function with period $P$.}



% [Image: periodic]





\paragraph{Definition}



A function $f$ is said to be periodic if, for some nonzero constant $P$, it is the case that



$$f(x+P) = f(x)$$



for all values of $x$ in the domain. A nonzero constant $P$ for which this is the case is called a period of the function. If there exists a least positive constant $P$ with this property, it is called the fundamental period (also primitive period, basic period, or prime period). Often, "the" period of a function is used to mean its fundamental period. A function with period $P$ will repeat on intervals of length $P$, and these intervals are sometimes also referred to as periods of the function.



Geometrically, a periodic function can be defined as a function whose graph exhibits translational symmetry, i.e., a function $f$ is periodic with period $P$ if the graph of $f$ is invariant under translation in the $x$-direction by a distance of $P$. This definition of periodicity can be extended to other geometric shapes and patterns, as well as be generalized to higher dimensions, such as periodic tessellations of the plane. A sequence can also be viewed as a function defined on the natural numbers, and for a periodic sequence, these notions are defined accordingly.



\paragraph{Examples}



\textbf{A graph of the sine function, showing two complete periods.}



% [Image: sine]



\paragraph{Real number examples}



The sine function is periodic with period $2\pi$, since



$$\sin(x + 2\pi) = \sin x$$



for all values of $x$. This function repeats on intervals of length $2\pi$ (see the graph to the right).



Everyday examples are seen when the variable is time; for instance, the hands of a clock or the phases of the moon show periodic behaviour. Periodic motion is motion in which the position(s) of the system are expressible as periodic functions, all with the same period.



For a function on the real numbers or on the integers, that means that the entire graph can be formed from copies of one particular portion, repeated at regular intervals.



A simple example of a periodic function is the function $f$ that gives the "fractional part" of its argument. Its period is $1$. In particular,



$$f(0.5) = f(1.5) = f(2.5) = \cdots = 0.5$$



The graph of the function $f$ is the sawtooth wave.



The trigonometric functions sine and cosine are common periodic functions, with period $2\pi$ (see the figure on the below). The subject of Fourier series investigates the idea that an 'arbitrary' periodic function is a sum of trigonometric functions with matching periods.



% [Image: sincos]



According to the definition above, some exotic functions, for example, the Dirichlet function, are also periodic; in the case of the Dirichlet function, any nonzero rational number is a period.



\paragraph{Complex number examples}



Using complex analysis, we have the common period function:



$$e^{ikx} = \cos(kx) + i \sin(kx).$$



Since the cosine and sine functions are both periodic with period $2\pi$, the complex exponential is made up of cosine and sine waves. This means that Euler's formula (above) has the property such that if $L$ is the period of the function, then



$$L = \frac{2\pi}{k}.$$



\paragraph{Double-periodic functions}



A function whose domain is the complex numbers can have two incommensurate periods without being constant. The elliptic functions are such functions. ("Incommensurate" in this context means not real multiples of each other.)



\paragraph{Properties}



Periodic functions can take on values many times. More specifically, if a function $f$ is periodic with period $P$, then for all $x$ in the domain of $f$ and all positive integers $n$,



$$f(x + nP) = f(x).$$



If $f(x)$ is a function with period $P$, then $f(ax)$, where $a$ is a non-zero real number such that $ax$ is within the domain of $f$, is periodic with period $\frac{P}{a}$. For example, $f(x) = \sin(x)$ has period $2\pi$; therefore, $\sin(5x)$ will have period $\frac{2\pi}{5}$.



Some periodic functions can be described by Fourier series. For instance, for $L^p$ space functions, Carleson's theorem states that they have a pointwise almost everywhere convergent Fourier series. Fourier series can only be used for periodic functions, or for functions on a bounded (compact) interval. If $f$ is a periodic function with period $P$ that can be described by a Fourier series, the coefficients of the series can be described by an integral over an interval of length $P$.



Any function that consists only of periodic functions with the same period is also periodic (with period equal or smaller), including:



\begin{itemize}
    \item addition, subtraction, multiplication and division of periodic functions, and
    \item taking a power or a root of a periodic function (provided it is defined for all $x$).
\end{itemize}

\paragraph{Generalizations}



\paragraph{Antiperiodic functions}



One common subset of periodic functions is that of antiperiodic functions. This is a function $f$ such that $f(x + P) = -f(x)$ for all $x$. (Thus, a $P$-antiperiodic function is a $2P$-periodic function.) For example, the sine and cosine functions are $\pi$-antiperiodic and $2\pi$-periodic. While a $P$-antiperiodic function is a $2P$-periodic function, the converse is not necessarily true.



\paragraph{Bloch-periodic functions}



A further generalization appears in the context of Bloch's theorems and Floquet theory, which govern the solution of various periodic differential equations. In this context, the solution (in one dimension) is typically a function of the form:



$$f(x + P) = e^{ikP} f(x),$$



where $k$ is a real or complex number (the Bloch wavevector or Floquet exponent). Functions of this form are sometimes called Bloch-periodic in this context. A periodic function is the special case $k = 0$, and an antiperiodic function is the special case $k = \frac{\pi}{P}$.



\paragraph{Quotient spaces as domain}



In signal processing, you encounter the problem that Fourier series represent periodic functions and that Fourier series satisfy convolution theorems (i.e., convolution of Fourier series corresponds to multiplication of represented periodic function and vice versa), but periodic functions cannot be convolved with the usual definition, since the involved integrals diverge. A possible way out is to define a periodic function on a bounded but periodic domain. To this end, you can use the notion of a quotient space:



$$\mathbb{R} / \mathbb{Z} = \{x + \mathbb{Z} : x \in \mathbb{R} \} = \{\{y : y \in \mathbb{R} \land y - x \in \mathbb{Z} \} : x \in \mathbb{R} \}.$$



That is, each element in $\mathbb{R} / \mathbb{Z}$ is an equivalence class of real numbers that share the same fractional part. Thus, a function like $f: \mathbb{R} / \mathbb{Z} \to \mathbb{R}$ is a representation of a 1-periodic function.



\paragraph{Calculating period}



Consider a real waveform consisting of superimposed frequencies, expressed in a set as ratios to a fundamental frequency, $f$: $F = \frac{1}{f} [f\_1 \, f\_2 \, f\_3 \, \cdots \, f\_N]$ where all non-zero elements $\ge 1$ and at least one of the elements of the set is $1$. To find the period, $T$, first find the least common denominator of all the elements in the set. The period can be found as $T = \frac{LCD}{f}.$ 



Consider that for a simple sinusoid, $T = \frac{1}{f}$. Therefore, the LCD can be seen as a periodicity multiplier.



For the set representing all notes of the Western major scale: $[1 \, \frac{9}{8} \, \frac{5}{4} \, \frac{4}{3} \, \frac{3}{2} \, \frac{5}{3} \, \frac{15}{8}]$ the LCD is $24$ therefore $T = \frac{24}{f}.$ 



For the set representing all notes of a major triad: $[1 \, \frac{5}{4} \, \frac{3}{2}]$ the LCD is $4$ therefore $T = \frac{4}{f}.$ 



For the set representing all notes of a minor triad: $[1 \, \frac{6}{5} \, \frac{3}{2}]$ the LCD is $10$ therefore $T = \frac{10}{f}.$ 



If no least common denominator exists, for instance, if one of the above elements were irrational, then the wave would not be periodic.



\paragraph{Licensing}



Content obtained and/or adapted from:



\href{https://en.wikipedia.org/wiki/Periodic_function}{Periodic function}, Wikipedia under a CC BY-SA license

\subsection{Learning Objectives}

\begin{enumerate}
    \item \textbf{1. Understand and Define Periodic Functions:}
    \item Students will be able to define a periodic function and identify its period. They will demonstrate this by explaining how the sine function, with a period of $2\pi$, repeats its values. \textbf{2. Identify Periods in Various Functions:}
    \item Students will calculate the period of different periodic functions, including trigonometric functions and complex exponentials. For example, they will determine the period of $e^{ikx}$ and explain the role of $k$ in defining the period. \textbf{3. Apply Periodicity to Real-World Examples:}
    \item Students will apply their understanding of periodic functions to real-world scenarios such as waveforms and musical scales. They will find the period of waveforms given as sets of frequencies and determine the least common denominator (LCD) to calculate the period.
\end{enumerate}

\subsection{Assessment Instrument}

\begin{enumerate}
    \item \textbf{Part 1: Short Answer Questions}
    \item Define a periodic function. How is the period of a periodic function related to its graph? Provide an example.
    \item Given the function $f(x) = \sin(x)$, determine its period and explain why this period is significant.
    \item Explain how to calculate the period of a waveform composed of superimposed frequencies. Use the set $[1, \frac{5}{4}, \frac{3}{2}]$ as an example and find the period. \textbf{Part 2: Problem Solving}
    \item Given the function $f(x) = \cos(kx)$, find the period of the function and describe how changing the value of $k$ affects the period.
    \item Consider the waveform with frequencies represented as ratios to a fundamental frequency: $F = \frac{1}{f} [1, \frac{9}{8}, \frac{5}{4}, \frac{4}{3}, \frac{3}{2}, \frac{5}{3}, \frac{15}{8}]$. Calculate the period of this waveform.
    \item For a periodic function $f(x) = \sin(x) + \sin(2x)$, determine the period of the function and explain the process of finding it. \textbf{Part 3: True or False}
    \item A function that repeats every $2\pi$ radians is always periodic. (True/False)
    \item The period of the function $e^{ikx}$ is $2\pi$. (True/False)
    \item If a function is periodic with period $P$, then $f(x + 2P) = f(x)$. (True/False)
\end{enumerate}

%%===============================================================
\section{The Sine Function}
\label{sec:the-sine-function}
%%===============================================================

\paragraph{The Sine Function}







\paragraph{Sine and Cosine Graphs}



And here in this graph, the colors we've chosen to draw the waves in have no relation to the frequency. Both waves have the same frequency, but we've used different colors.



% [Image: Sine cosine plot]



Notice that in the graph of sine and cosine above that the two graphs have the same shape. If we slide the sine graph slightly to the left, it coincides exactly with the cosine graph. Using the terminology used to describe sinusoidal waves, they have the same amplitude, the same frequency and different phases.



\paragraph{Definitions}



A sinusoidal wave is characterized by three parameters: \textbf{amplitude}, \textbf{frequency} and \textbf{phase}.



\begin{itemize}
    \item The amplitude is the amount the function varies, positively or negatively, from zero in the $y$ direction.
    \item The frequency is how many complete cycles there are of the wave in unit distance on the $x$ axis (which often measures time).
    \item The phase is relevant when comparing two waves of the same frequency. It is how much (measured in degrees or radians) one is shifted relative to the other on the $x$ axis.
\end{itemize}

This terminology comes from sound engineering where higher frequency sounds have higher pitch and waves of greater amplitude are louder.



As an alternative to specifying the frequency, the number of cycles in unit distance, we can instead specify the wavelength, the length of one cycle. The higher the frequency, the shorter the wavelength. The lower the frequency the longer the wavelength.



You can check this on the diagram of waves of different frequencies at the top of the page.



\paragraph{In More Detail}



To fully understand how each of the following definitions relate, look at the graph of the function $f(t) = 4 \sin \left(\frac{\pi}{6} t\right) + 40$ below. For the figure representing (a)-(b), read the corresponding definitions.



% [Image: graph]





\paragraph{Amplitude}



The amplitude is the maximum amount that the wave differs from the sinusoidal axis value, $c$ – the value by which the function is shifted by the average of the maximum to the minimum range of the function, 



$$\frac{y\_{max} + y\_{min}}{2}.$$ 



The amplitude comes from two different locations: the maximum value minus $c$ or the minimum value minus $c$. The positive magnitude of the value is taken (absolute value). Look at the data below for an example.



$$

\begin{array}{|c|c|}

\hline

x & g(x) \\\\ 

\hline

0 & 20 \\\\ 

\hline

2 & 0 \\\\ 

\hline

4 & -20 \\\\ 

\hline

6 & 0 \\\\ 

\hline

8 & 20 \\\\ 

\hline

\vdots & \vdots \\\\ 

\end{array}

$$



Let us not worry about finding the function for this set of data for now. Instead, focus on finding the amplitude. The maximum $y$-value $y\_{max}=20$. The minimum $y$-value $y\_{min}=-20$. Therefore, $\frac{20 + (-20)}{2} = \frac{0}{2} = 0$.  The maximum difference comes from two different locations: the maximum value minus the sinusoidal axis value or the minimum value minus the sinusoidal axis value. In this example, we will choose to use the difference between the maximum value to the sinusoidal axis value, which is $20 - 0 = 20$.  Therefore, the amplitude of the above data is $20$.



In general, for $\sin(b\theta) + cMATH2\cos(b\theta) + c$ (where $b > 0$) the amplitude is $1$. For $a\sin(b\theta) + c$ and $a\cos(b\theta) + c$ (where $a > 0$ and $b > 0$) the amplitude is $a$. The difference between the greatest and least value is called the double amplitude or peak to peak amplitude. Knowing this, the way to find the amplitude given the greatest and least value of the range is to take the difference between the two and then divide by 2:



$$a = \frac{y\_{max} - y\_{min}}{2}$$



\paragraph{Frequency}

The frequency (sometimes $f$) is the number of cycles a wave goes through in a standard distance or time. If this standard distance is $360^{\circ}$ or $2\pi$ radians, the frequency of $\sin(\theta)MATH4\cos(\theta)$ is $1$, and that of $\sin(b\theta)$ and $\cos(b\theta)$ is $b$. A graph of a high frequency wave shows more complete cycles in the same horizontal span as a low frequency wave. Notice $b$ is the frequency value, $f$, so $b = f$.



The wavelength (or period), $\lambda$, decreases as the frequency increases. Mathematically, this relationship is written as $\lambda \propto \frac{1}{f}$ (the wavelength is inversely proportional to the frequency). It is the distance over which a complete cycle occurs, that is, the wave goes from $c$ to its maximum positive value, back to $c$, down to the lowest negative value, then back up to $c$.



The functions 

$\sin(\theta)$ and $\cos(\theta)$ both go through a complete cycle as the angle increases by $360^{\circ}$ or $2\pi$ radians, hence this is the wavelength. The functions $\sin(b\theta)$ and $\cos(b\theta)$ (where $b$ is any positive number) go through one cycle as the angle $\theta$ increases by $\frac{360^{\circ}}{b}$ or $\frac{2\pi}{b}$, hence their wavelength $\lambda = \frac{360^{\circ}}{|b|}$ or $\lambda = \frac{2\pi}{|b|}$. This makes sense, because if $b = f$ and $\lambda = \frac{2\pi}{|b|}$, then $\lambda \propto \frac{1}{b}$.



Finally, using wavelength, one can directly find the frequency. If $\lambda = \frac{2\pi}{|b|}$, then



$$\Rightarrow \lambda \times \frac{1}{\lambda} \times b = \frac{2\pi}{b} \times b \times \frac{1}{\lambda}$$



$$\Rightarrow b = \frac{2\pi}{\lambda}$$



\paragraph{Phase Shift}



The phase is the fraction of the cycle that you have reached at any given point. If the point is not specified, it may be assumed to be $0^{\circ}$. The start of the cycle is the point where the wave is $0$ and going from negative to positive. (If the wave is not a sine curve, there may be more than one such point in a complete cycle so the starting point may be arbitrary.) For $\sin(\theta)$ the phase at $0^{\circ}$ is $0$; for $\cos(\theta)$ it is $\frac{1}{4}$. If the cycle has wavelength $360^{\circ}$ or $2\pi$ radians, the phase is often expressed as an angle corresponding to the given fraction of the wavelength, so the phase of $\cos(\theta)$ at $0^{\circ}$ is $90^{\circ}$ or $\frac{\pi}{2}$ radians.



It is assumed that the average value of the wave is $0$. If the wave is say $\sin(\theta) + 4$ or $2\pi$, it is never $0$ so to define the phase we would need the point where it reaches its average value ($4$ in this case).



\paragraph{Examples}



\paragraph{Larger Amplitude}



$10\cos(\theta)$ has ten times the amplitude of $\cos(\theta)$



The graph oscillates between $10$ and $-10$ on the $y$ axis.



\paragraph{Larger Frequency}



$\cos(100\theta)$ has a higher frequency than $\cos(\theta)$



There will be $100$ times as many cycles in the same distance along the $x$ axis. The amplitude is unchanged.



The wavelength is also decreased. The distance along the $x$ axis for a complete cycle is one hundredth of what it was before.



\paragraph{Different Phase}



$\sin(\theta)$ has a different phase from $\cos(\theta)$. The phase difference happens to be $90^{\circ}$, because $\sin(\theta + 90^{\circ}) = \cos(\theta)$. If two waves are $180^{\circ}$ out of phase then where one is positive the other has the same value only negative. $-\sin(\theta)$ is $180^{\circ}$ out of phase with $\sin(\theta)$



\paragraph{Offset}



Two sinusoidal waves may also have a different offset to each other, i.e. their average values are different, but they still have the same shape. Their graphs are $y$ shifted relative to each other. In sound engineering this is called DC Bias.



\paragraph{Different Offset}



$1 + \cos(\theta)$ has a different offset to $\cos(\theta)$. Whereas $\cos(\theta)$ takes both positive and negative values, $1 + \cos(\theta)$ is never below $0$.



This offset is based off the vertical shift, also known as the sinusoidal axis value. Keep in mind, unlike other functions, this vertical shift does not denote the $y$-intercept. The only way to determine the $y$-intercept correctly is to substitute zero into a function. Quick quiz: what is $\cos(0)$? In case you forgot, it is $1$. Therefore, the $y$-intercept for $f(\theta) = 1 + \cos(\theta)$ is $f(0) = 1 + \cos(0) = 1 + 1 = 2$.



\paragraph{The Graphs and Their Differences}



Comparing and contrasting two different items may seem like a pointless task to the student, at first. However, the purpose of such an activity is to make students analyze information and help cement differences between the two ideas. Here, it is imperative that the student understands the differences between cosine and sine because otherwise, the student is not able to model situations correctly! While this Wikibooks will not make the student analyze the two functions themselves, it may be a good idea to remember the differences.



\paragraph{The Parent Functions}



\textbf{The sine and cosine curves are both graphed over the period $ \lambda = 2\pi $ with a frequency $ f = 1 $.} 



% [Image: graphs]



The sine and cosine functions are shown above, where $ 0 \leq x \leq 2\pi $. Notice that the frequency and period of the functions are both identical. Given that these are also parent functions, the vertical shift is zero. However, the similarities start to end. 

 First, look at the pattern of the function. Over $ \lambda = 2\pi $, the periodic function $ \sin(x) $ follows the basic pattern of "middle," "top," "middle," "bottom," "middle." That is, the value of $ f(x) = \sin(x) $ starts at $ f(0) = 0 $. However, $ f\left(\frac{\pi}{2}\right) = 1 $. This is the highest possible value of sine when looking at the positive values, because the unit circle is peaked when $ \theta = \frac{\pi}{2} $. At $ x = \pi $, $ f(\pi) = 0 $. This has again to do with the unit circle. If you remembered the patterns you identified in the circle (or even the unit circle itself), you can infer how $ f(x) = \sin(x) $ will look. The minimum value of the range is at $ -1 $ precisely for this reason. As such, the pattern we identified is always true.



By contrast, the cosine function follows a more simplistic pattern. For $ g(x) = \cos(x) $, you already start at the maximum value because $ g(0) = \cos(0) = 1 $ is the largest $ x $-value on the circle, while $ g(\pi) = \cos(\pi) = -1 $ is the lowest $ x $-value on the circle. Because of the unit circle, the simplistic pattern results: "maximum," "medium," "minimum," "medium," "maximum."



The patterns mentioned above are things that you should keep in the back of your mind. These patterns are the main ways in which a person can easily identify when to use one particular function over another. When a phase shift is applied, choosing to use a cosine function or a sine function is best left to student judgment.



From the information we have so far, we get the following information:



The vertical shift $ c = \frac{y\_{max} + y\_{min}}{2} = \frac{1 + (-1)}{2} = \frac{0}{2} = 0 $.



The amplitude $ a = \frac{y\_{max} - y\_{min}}{2} = \frac{1 - (-1)}{2} = \frac{2}{2} = 1 $.



The frequency $ b = \frac{2\pi}{\lambda} = \frac{2\pi}{2\pi} = 1 $.



The general form of the sinusoidal function is $ A\sin(B(x - D)) + C $.



The sine can be substituted for cosine, and it would work just the same.



$ A $ is the amplitude.



$ B $ is the frequency.



$ C $ is the vertical shift (sinusoidal axis value).



$ D $ is the phase (horizontal) shift.



We have described the general pattern for the cosine and sine functions. However, it is important to step back and appreciate the underlying math for which the cosine and sine curves rely on, the unit circle. While these functions are beautiful, the underlying math behind why these functions behave the way they do is the reason for this beauty, and in a way, makes it more beautiful than the end result.



Draw two circles with radii $ r $ and $ R $ centered at the origin, where $ r < R $. If a line extending from the origin reaches an endpoint of each circle at the angle $ \theta $, then $ \cos \theta = \frac{x\_1}{r} $, and $ \cos \theta = \frac{x\_2}{R} $. As such,



% [Image: circles]



$$ x\_1 = r\cos\theta $$

$$ x\_2 = R\cos\theta $$



Because $ \cos\theta $ is multiplied directly to $ r $ or $ R $, for $ r < R $, $ r\cos\theta < R\cos\theta $.



Let $ f(\theta) = r\cos\theta $, $ g(\theta) = R\cos\theta $, and $ r = 1 $. Because we know how to graph $ f(\theta) = \cos\theta $, $ R $ is multiplied directly to $ \cos\theta $, and $ g(\theta) = Rf(\theta) $, the resulting function must be "expanded" over the $ y $-axis. This expansion must be related to the amplitude of the function, because $ Ry\_{max} $ and $ Ry\_{min} $ gives $ \frac{Ry\_{max} - Ry\_{min}}{2} = \frac{R(y\_{max} - y\_{min})}{2} = R \cdot \frac{y\_{max} - y\_{min}}{2} = Ra = R $.  (Recall $ a = 1 $).



\begin{itemize}
    \item The amplitude of the function is given by the radius of the circle.
\end{itemize}

\begin{itemize}
    \item Changing the radius of the circle does not change the frequency, the vertical shift, or the phase shift. (\textbf{Exercise}: show that this is true for at least the frequency and the vertical shift.)
\end{itemize}

This result will be useful when working with problems later.



As of now, increasing the radius is all that we could connect these sinusoidal functions to the circle. We need to understand more about additions and such before one can address the beauty of this math. Nevertheless, always keep this result locked in your heart.



\paragraph{Modeling Situations}



The sinusoidal functions are very useful in real-life applications, where these periodic (repeating) functions model situations that follow a pattern over a set amount of time, such as in modeling day length, planetary motion, sound waves, energy waves, etcetera.



\paragraph{Fish Length}



\begin{itemize}
    \item Marine biologists determined the average length of $ N = 50,000 $ fish in a pond in millimeters (mm) is dependent on the time, in years. The table is shown below, where $ t $ is the time, in years, and $ l(t) $ is the average length of the fish, in millimeters. Find the sinusoidal function that best models the information provided.
\end{itemize}

$$

\begin{array}{|c|c|}

t & l(t) \\\\ 

\hline

0 & 44 \\\\ 

\hline

3 & 40 \\\\ 

\hline

6 & 36 \\\\ 

\hline

9 & 40 \\\\ 

\hline

12 & 44 \\\\ 

\hline

\vdots & \vdots \\\\ 

\end{array}

$$



\begin{itemize}
    \item The situation is given to be periodic and sinusoidal. Therefore, it is best to use the sinusoidal function and its properties: $ a \cdot \text{trig}(b(x + d)) + c $.
\end{itemize}

\begin{itemize}
    \item Find the vertical shift, $ c $. Recall that $ c = \frac{y\_{max} + y\_{min}}{2} $. According to the table, the maximum y-value is 44, while the minimum value is 36. Therefore, the vertical shift is $ \frac{44 + 36}{2} = 40 $.
\end{itemize}

\begin{itemize}
    \item Find the amplitude, $ a $. Recall that $ a = \frac{y\_{max} - y\_{min}}{2} $. According to the table, the maximum y-value is 44, while the minimum value is 36. Therefore, the amplitude is $ \frac{44 - 36}{2} = 4 $.
\end{itemize}

\begin{itemize}
    \item Find the frequency, $ b $. Recall that $ b = \frac{2\pi}{\lambda} $. If the wavelength is the distance whereby a complete cycle occurs, then from $ t = 0 \text{ sec} $ to $ t = 12 \text{ sec} $, the cycle is complete. As such, $ \lambda = 12 - 0 = 12 $. Given this information, $ b = \frac{2\pi}{12} = \frac{\pi}{6} $.
\end{itemize}

\begin{itemize}
    \item For now, let us ignore the phase shift, $ d $. Substitute the information we already know into the function above: $ 4 \cdot \text{trig}\left(\frac{\pi}{6}t\right) + 40 $. Recall that a sinusoidal function that has a maximum value reappear again over a wavelength (from maximum to medium to minimum to medium to maximum) is cosine. As such $ 4\cos\left(\frac{\pi}{6}t\right) + 40 $ models the situation sufficiently.
\end{itemize}

\begin{itemize}
    \item Equivalently, assuming we want to keep it in terms of sine, we know the following is true:
\end{itemize}

$$ \cos\left(\theta\right) = \sin\left(\frac{\pi}{2} - \theta\right) = \sin\left(\frac{\pi}{2} + \theta\right) = -\sin\left(\theta - \frac{\pi}{2}\right) $$



\begin{itemize}
    \item Therefore, $ 4\cos\left(\frac{\pi}{6}t\right) + 40 = 4\sin\left(\frac{\pi}{6}\left(\frac{\pi}{2} - t\right)\right) + 40 = 4\sin\left(\frac{\pi}{6}\left(\frac{\pi}{2} + t\right)\right) + 40 = -4\sin\left(\frac{\pi}{6}\left(t - \frac{\pi}{2}\right)\right) + 40 $, where the information we already know is ignored.
\end{itemize}

The above question is comparatively easier than other similar problems like this. In the real world, you are less likely to see a steady increase/decrease in fish length because of two factors:



1. \textbf{Evolution}: Bigger fish are usually seen as predators, which are often avoided. Smaller fish are more likely to be seen as weaker.

2. \textbf{Internal Competition}: Size among females may be seen as more desirable in mates, and therefore, those with these traits are more likely to pass on offspring. As such, competition may exist among females for desirable mates.



There are many questions a person may be asking: where is the above circle in the question? Why use radians as opposed to degrees? There are a few reasons, but it requires us to delve more deeply into mathematics that is usually considered advanced (polar coordinates). Unfortunately, Book 1 will not delve into this topic. However, for those that are curious, Book 3 will have polar coordinates that fully explain this behavior.



\paragraph{Licensing}



Content obtained and/or adapted from:



\href{https://en.wikibooks.org/wiki/Trigonometry/Phase_and_Frequency}{Phase and Frequency}, Wikibooks under a CC BY-SA license.

\subsection{Learning Objectives}

\begin{enumerate}
    \item \textbf{1. Graph the Sine and Cosine Functions:}
    \item Students will be able to graph the sine and cosine functions over one complete cycle and identify key characteristics such as amplitude, frequency, wavelength, and phase shift. \textbf{2. Analyze the Effects of Parameters on the Sine Function:}
    \item Students will be able to describe how changes in amplitude, frequency, phase shift, and vertical shift affect the graph of the sine function and explain the real-world implications of these changes. \textbf{3. Model Periodic Phenomena Using the Sine Function:}
    \item Students will be able to apply the sine function to model and solve problems related to periodic phenomena such as sound waves, seasonal variations, and other real-life scenarios.
\end{enumerate}

\subsection{Assessment Instrument}

\begin{enumerate}
    \item Part 1: Graphing the Sine and Cosine Functions \textbf{1. Graph the Functions:}
    \item Graph $f(x) = \sin(x)$ and $g(x) = \cos(x)$ on the same set of axes for $0 \leq x \leq 2\pi$. Label key points, including maximum, minimum, and intercepts. \textbf{2. Identify Characteristics:}
    \item For each function, identify the amplitude, frequency, period, and phase shift. Part 2: Analyzing the Effects of Parameters \textbf{3. Parameter Variations:}
    \item Given the function $h(x) = 3\sin(2x - \frac{\pi}{4}) + 2$, describe the effect of each parameter on the graph of the sine function. Specifically, explain how the amplitude, frequency, phase shift, and vertical shift are affected by the values in the function. Part 3: Modeling Real-World Scenarios \textbf{4. Modeling with Sine:}
    \item A city experiences temperature variations throughout the year that can be modeled using the function $T(t) = 10\sin\left(\frac{\pi}{6}(t - 3)\right) + 25$, where $T(t)$ is the temperature in degrees Celsius, and $t$ is the time in months.
    \item (a) Determine the amplitude, period, phase shift, and vertical shift of the temperature function.
    \item (b) Interpret the meaning of these parameters in the context of the city''s climate.
    \item (c) Predict the temperature at $t = 6$ months using the function. Part 4: Application \textbf{5. Compare and Contrast:}
    \item Compare the graphs of $y = \sin(x)$ and $y = \sin(2x)$. Describe how the frequency and wavelength differ and provide a real-life scenario where a higher frequency sine wave might be observed. \textbf{6. Phase Shift Exploration:}
    \item Explain what it means for two sine waves to be "out of phase" and provide an example of how phase shift could be observed in a real-world application, such as sound waves or electrical circuits. \textbf{7. Amplitude and Real-Life Implications:}
    \item Discuss how changes in the amplitude of a sine wave could impact a physical system, such as sound intensity in a concert hall or the height of ocean waves during a storm.
\end{enumerate}
